% Generated mechanically from frozen input p12/ja/index.tex.
% Do not edit this generated file; edit the bound locale source instead.
\chapter*{自動生成索引}
\addcontentsline{toc}{chapter}{自動生成索引}

\phantomsection
\label{index-section-phantom}


\frenchspacing


\begin{multicols}{2}[\section{英語原語のアルファベット順による定義}\label{index-section-alphabetized}]
\noindent
{\it $(-1)$-指数}
\ref{models-definition-type-minus-one} を参照

\noindent
{\it $(-2)$-指数}
\ref{models-definition-type-minus-two} を参照

\noindent
{\it $(2, 1)$-圏}
\ref{categories-definition-2-1-category} を参照

\noindent
{\it $(2, 1)$-周期複体}
\ref{chow-definition-periodic-complex} を参照

\noindent
{\it $(\mathcal{F}_n)$ は標準的に $X$ へ延長される}
\ref{algebraization-definition-canonically-algebraizable} を参照

\noindent
{\it $(\mathcal{F}_n)$ は $X$ へ延長される}
\ref{algebraization-definition-algebraizable} を参照

\noindent
{\it $(\mathcal{F}_n)$ は $(a, b)$-不等式を満たす}
\ref{algebraization-definition-s-d-inequalities} を参照

\noindent
{\it $(\mathcal{F}_n)$ は狭義 $(a, b)$-不等式を満たす}
\ref{algebraization-definition-s-d-inequalities} を参照

\noindent
{\it $(\textit{Spaces}/S)_\etale$}
\ref{spaces-topologies-definition-big-etale-site} を参照

\noindent
{\it $(\textit{Spaces}/S)_{fppf}$}
\ref{spaces-topologies-definition-big-fppf-site} を参照

\noindent
{\it $(\textit{Spaces}/S)_{ph}$}
\ref{spaces-topologies-definition-big-ph-site} を参照

\noindent
{\it $(\textit{Spaces}/X)_\etale$}
\ref{spaces-topologies-definition-big-small-etale} を参照

\noindent
{\it $(\textit{Spaces}/X)_{fppf}$}
\ref{spaces-topologies-definition-big-small-fppf} を参照

\noindent
{\it $(\textit{Spaces}/X)_{ph}$}
\ref{spaces-topologies-definition-big-small-ph} を参照

\noindent
{\it $(A, B)$-両側加群}
\ref{algebra-definition-bimodule} を参照

\noindent
{\it $(A, B)$-両側加群}
\ref{dga-definition-bimodule} を参照

\noindent
{\it $(R_k)$}
\ref{algebra-definition-conditions} を参照

\noindent
{\it $(R_k)$}
\ref{properties-definition-Rk} を参照

\noindent
{\it $(S_k)$}
\ref{algebra-definition-conditions} を参照

\noindent
{\it $(S_k)$}
\ref{properties-definition-Rk} を参照

\noindent
{\it $(S_k)$}
\ref{coherent-definition-depth} を参照

\noindent
{\it $(S_k)$}
\ref{coherent-definition-depth} を参照

\noindent
{\it $(U', R', s', t', c')$ は $(U, R, s, t, c)$ 上カルテジアンである}
\ref{groupoids-definition-cartesian-morphism} を参照

\noindent
{\it $1$-射}
\ref{categories-definition-2-category} を参照

\noindent
{\it $S$ 上の代数スタックの $2$-圏}
\ref{algebraic-definition-morphism-algebraic-stacks} を参照

\noindent
{\it $\mathcal{C}$ 上グルーポイドに ファイバー化された圏の $2$-圏}
\ref{categories-definition-categories-fibred-in-groupoids-over-C} を参照

\noindent
{\it $\mathcal{C}$ 上セトイドにファイバー化された圏の $2$-圏}
\ref{categories-definition-categories-fibred-in-setoids-over-C} を参照

\noindent
{\it $\mathcal{C}$ 上集合にファイバー化された圏の $2$-圏}
\ref{categories-definition-categories-fibred-in-sets-over-C} を参照

\noindent
{\it $\mathcal{C}$ 上の圏の $2$-圏}
\ref{categories-definition-categories-over-C} を参照

\noindent
{\it $\mathcal{C}$ 上のファイバー圏の $2$-圏}
\ref{categories-definition-fibred-categories-over-C} を参照

\noindent
{\it $\mathcal{C}$ 上のグルーポイドのスタックの $2$-圏}
\ref{stacks-definition-stacks-in-groupoids-over-C} を参照

\noindent
{\it $\mathcal{C}$ 上のセトイドのスタックの $2$-圏}
\ref{stacks-definition-stacks-in-setoids-over-C} を参照

\noindent
{\it $\mathcal{C}$ 上のスタックの $2$-圏}
\ref{stacks-definition-stacks-over-C} を参照

\noindent
{\it $2$-圏}
\ref{categories-definition-2-category} を参照

\noindent
{\it $2$-射}
\ref{categories-definition-2-category} を参照

\noindent
{\it $2$-周期複体}
\ref{chow-definition-periodic-complex} を参照

\noindent
{\it $I$ に関して $\alpha$-小}
\ref{injectives-definition-small} を参照

\noindent
{\it $\delta$ は $\gamma$ と両立する}
\ref{crystalline-definition-compatible} を参照

\noindent
{\it $T$ の $\delta$-次元}
\ref{spaces-chow-definition-delta-dimension} を参照

\noindent
{\it $Z$ の $\delta$-次元}
\ref{chow-definition-delta-dimension} を参照

\noindent
{\it $\mathcal{A}$ から $\mathcal{D}$ への $\delta$-関手}
\ref{derived-definition-delta-functor} を参照

\noindent
{\it $\delta$-関手}
\ref{homology-definition-cohomological-delta-functor} を参照

\noindent
{\it $A$ の $\delta$-不変量}
\ref{varieties-definition-delta-invariant-algebra} を参照

\noindent
{\it $x$ における $X$ の $\delta$-不変量}
\ref{varieties-definition-delta-invariant} を参照

\noindent
{\it $\delta(\tau)$}
\ref{exercises-definition-delta} を参照

\noindent
{\it $\delta^n_j : [n-1] \to [n]$}
\ref{simplicial-definition-face-degeneracy} を参照

\noindent
{\it $\ell$-進コホモロジー}
\ref{trace-definition-cohomology-l-adic} を参照

\noindent
{\it $\ell$-進層}
\ref{trace-definition-l-adic-sheaf} を参照

\noindent
{\it $\epsilon$-不変量}
\ref{obsolete-definition-epsilon} を参照

\noindent
{\it $\Ext$-群}
\ref{homology-definition-ext-group} を参照

\noindent
{\it $\Hom(U, V)$}
\ref{simplicial-definition-hom-deltak-cosimplicial} を参照

\noindent
{\it $\Hom(U, V)$}
\ref{simplicial-definition-hom-deltak-simplicial} を参照

\noindent
{\it $\Hom(U, V)$}
\ref{simplicial-definition-hom-from-simplicial-set} を参照

\noindent
{\it $\kappa$-生成}
\ref{properties-definition-kappa-generated} を参照

\noindent
{\it $\mathbf{Z}_\ell$-層}
\ref{trace-definition-l-adic-sheaf} を参照

\noindent
{\it $\mathcal{A}^0$}
\ref{dga-definition-H0-of-graded-category} を参照

\noindent
{\it $\mathcal{C}_\Lambda$}
\ref{formal-defos-definition-CLambda} を参照

\noindent
{\it $x$ における $\mathcal{F}$ の長さは $d$ である}
\ref{spaces-chow-definition-length-at-x} を参照

\noindent
{\it $\mathcal{F}$ は次元 $\geq n$ で $S$ 上平坦である}
\ref{flat-definition-flat-dimension-n} を参照

\noindent
{\it $\mathcal{F}$ は次元 $\geq n$ において $Y$ 上平坦}
\ref{spaces-flat-definition-flat-dimension-n} を参照

\noindent
{\it $\mathcal{G}$-トーサー}
\ref{cohomology-definition-torsor} を参照

\noindent
{\it $\mathcal{G}$-トーサー}
\ref{sites-cohomology-definition-torsor} を参照

\noindent
{\it $\mathcal{I}$ は $\mathcal{J}$ において共終}
\ref{categories-definition-cofinal} を参照

\noindent
{\it $\mathcal{I}$ は $\mathcal{J}$ において始的}
\ref{categories-definition-initial} を参照

\noindent
{\it $\mathcal{K}_X$}
\ref{divisors-definition-sheaf-meromorphic-functions} を参照

\noindent
{\it $\mathcal{K}_X$}
\ref{spaces-divisors-definition-sheaf-meromorphic-functions} を参照

\noindent
{\it $\mathcal{O}^*$}
\ref{sites-modules-definition-invertible-sheaf} を参照

\noindent
{\it $\mathcal{O}_1$-導分}
\ref{modules-definition-derivation} を参照

\noindent
{\it $\mathcal{O}_1$-導分}
\ref{sites-modules-definition-derivation} を参照

\noindent
{\it $\mathcal{O}_\mathcal{X}$-加群}
\ref{stacks-sheaves-definition-modules} を参照

\noindent
{\it $\mathcal{S}$ は $\mathcal{C}$ から誘導される位相を備える}
\ref{stacks-definition-topology-inherited} を参照

\noindent
{\it $\mathcal{S}_F$}
\ref{categories-definition-split-fibred-category} を参照

\noindent
{\it $\mathcal{S}_F$}
\ref{categories-definition-split-category-fibred-in-groupoids} を参照

\noindent
{\it $\mathcal{X}$ が $\mathcal{Y}$ 上で相対的に表現可能}
\ref{categories-definition-representable-map-categories-fibred-in-groupoids} を参照

\noindent
{\it $\mathfrak g^r_d$}
\ref{curves-definition-linear-series} を参照

\noindent
{\it $\mathit{QC}(\mathcal{A}, \text{d})$}
\ref{sdga-definition-cartesian} を参照

\noindent
{\it $\mathit{QC}(\mathcal{O})$}
\ref{sites-cohomology-definition-cartesian} を参照

\noindent
{\it $\phi$ は $f$ 上にある}
\ref{categories-definition-fibre-category} を参照

\noindent
{\it $\Sh(\mathcal{C})$}
\ref{sites-definition-category-sheaves-sets} を参照

\noindent
{\it $\sigma^n_j : [n + 1]  \to [n]$}
\ref{simplicial-definition-face-degeneracy} を参照

\noindent
{\it $\tau$ $G$-トーサー}
\ref{groupoids-definition-principal-homogeneous-space} を参照

\noindent
{\it $\tau$ $G$-トーサー}
\ref{spaces-groupoids-definition-principal-homogeneous-space} を参照

\noindent
{\it 基底上 $\tau$-局所的}
\ref{descent-definition-property-morphisms-local} を参照

\noindent
{\it 基底上 $\tau$-局所的}
\ref{spaces-descent-definition-property-morphisms-local} を参照

\noindent
{\it 始域上 $\tau$-局所的}
\ref{descent-definition-property-morphisms-local-source} を参照

\noindent
{\it 始域上 $\tau$-局所的}
\ref{spaces-descent-definition-property-morphisms-local-source} を参照

\noindent
{\it 標的上 $\tau$-局所的}
\ref{descent-definition-property-morphisms-local} を参照

\noindent
{\it 標的上 $\tau$-局所的}
\ref{spaces-descent-definition-property-morphisms-local} を参照

\noindent
{\it $\tau$ トーサー}
\ref{groupoids-definition-principal-homogeneous-space} を参照

\noindent
{\it $\tau$ トーサー}
\ref{spaces-groupoids-definition-principal-homogeneous-space} を参照

\noindent
{\it $\tau$-被覆}
\ref{etale-cohomology-definition-tau-covering} を参照

\noindent
{\it $\textit{Adeq}((\Sch/S)_\tau, \mathcal{O})$}
\ref{adequate-definition-category-adequate-modules} を参照

\noindent
{\it $\textit{Adeq}(\mathcal{O})$}
\ref{adequate-definition-category-adequate-modules} を参照

\noindent
{\it $\textit{Adeq}(S)$}
\ref{adequate-definition-category-adequate-modules} を参照

\noindent
{\it $\text{Fil}^f(\mathcal{A})$}
\ref{exercises-definition-finite-filtration-category} を参照

\noindent
{\it $\underline{U}$}
\ref{sites-definition-representable-sheaf} を参照

\noindent
{\it $\varphi$-導分}
\ref{modules-definition-derivation} を参照

\noindent
{\it $\varphi$-導分}
\ref{sites-modules-definition-derivation} を参照

\noindent
{\it $\widehat{\mathcal{C}}_\Lambda$}
\ref{formal-defos-definition-completion-CLambda} を参照

\noindent
{\it $A$ は $A_1, \ldots, A_n$ のウェッジである}
\ref{varieties-definition-wedge} を参照

\noindent
{\it $A$-双導分}
\ref{cotangent-definition-biderivation} を参照

\noindent
{\it $X$ から $Y$ への $B$-有理写像}
\ref{spaces-morphisms-definition-rational-map} を参照

\noindent
{\it $c$-進}
\ref{formal-spaces-definition-weakly-adic} を参照

\noindent
{\it $C$}
\ref{descent-definition-C} を参照

\noindent
{\it $C_r$}
\ref{etale-cohomology-definition-Cr} を参照

\noindent
{\it $C_{S/R}$}
\ref{algebra-definition-universal-thickening} を参照

\noindent
{\it $d(M)$}
\ref{algebra-definition-d} を参照

\noindent
{\it $D_c(X_\etale, \Lambda)$}
\ref{etale-cohomology-definition-c} を参照

\noindent
{\it $D_{ctf}(X_\etale, \Lambda)$}
\ref{etale-cohomology-definition-ctf} を参照

\noindent
{\it $f s^{-1}$}
\ref{categories-definition-right-localization-as-fraction} を参照

\noindent
{\it $f$ は $x$ において相対次元 $d$ をもつ}
\ref{spaces-morphisms-definition-dimension-fibre} を参照

\noindent
{\it $F$ は $G$ 上相対的に表現可能}
\ref{categories-definition-representable-map-presheaves} を参照

\noindent
{\it $f$-豊富}
\ref{morphisms-definition-relatively-ample} を参照

\noindent
{\it $f$-豊富}
\ref{spaces-divisors-definition-relatively-ample} を参照

\noindent
{\it $X/S$ 上の（$\sigma$ に関する）$F$-クリスタル}
\ref{crystalline-definition-F-crystal} を参照

\noindent
{\it $f$-写像 $\varphi : \mathcal{G} \to \mathcal{F}$}
\ref{spaces-properties-definition-f-map} を参照

\noindent
{\it $f$-写像 $\xi : \mathcal{G} \to \mathcal{F}$}
\ref{sheaves-definition-f-map} を参照

\noindent
{\it $f$-相対的に豊富}
\ref{morphisms-definition-relatively-ample} を参照

\noindent
{\it $f$-相対的に豊富}
\ref{spaces-divisors-definition-relatively-ample} を参照

\noindent
{\it $f$-相対的に非常に豊富}
\ref{morphisms-definition-very-ample} を参照

\noindent
{\it $f$-非常に豊富}
\ref{morphisms-definition-very-ample} を参照

\noindent
{\it $f^{-1}\mathcal{S}$}
\ref{stacks-definition-pullback-stack} を参照

\noindent
{\it $f_*$}
\ref{descent-definition-pushforward} を参照

\noindent
{\it $f_*\mathcal{S}$}
\ref{stacks-definition-pushforward-stack} を参照

\noindent
{\it $G$-同変準連接 $\mathcal{O}_X$-加群}
\ref{groupoids-definition-equivariant-module} を参照

\noindent
{\it $G$-同変準連接 $\mathcal{O}_X$-加群}
\ref{spaces-groupoids-definition-equivariant-module} を参照

\noindent
{\it $G$-同変}
\ref{groupoids-definition-action-group-scheme} を参照

\noindent
{\it $G$-同変}
\ref{spaces-groupoids-definition-action-group-space} を参照

\noindent
{\it $G$-不変}
\ref{groupoids-quotients-definition-invariant} を参照

\noindent
{\it $G$-加群}
\ref{etale-cohomology-definition-G-module-continuous} を参照

\noindent
{\it $G$-集合}
\ref{pione-definition-G-set-continuous} を参照

\noindent
{\it $\tau$ 位相における $G$-トーサー}
\ref{groupoids-definition-principal-homogeneous-space} を参照

\noindent
{\it $\tau$ 位相における $G$-トーサー}
\ref{spaces-groupoids-definition-principal-homogeneous-space} を参照

\noindent
{\it $G$-トーサー}
\ref{groupoids-definition-principal-homogeneous-space} を参照

\noindent
{\it $P$ 上の $f$ の $G$-トレース}
\ref{trace-definition-trace-G} を参照

\noindent
{\it $G\textit{-Sets}$}
\ref{pione-definition-G-set-continuous} を参照

\noindent
{\it $g_!\mathcal{F} = (g_{p!}\mathcal{F})^\#$}
\ref{sites-modules-definition-g-shriek} を参照

\noindent
{\it $g_{p!}\mathcal{F}$}
\ref{sites-modules-definition-g-shriek} を参照

\noindent
{\it $H^{i + k}(A^\bullet) \longrightarrow H^i(A[k]^\bullet)$}
\ref{homology-definition-cohomology-shift} を参照

\noindent
{\it $H_1$-正則イデアル}
\ref{more-algebra-definition-regular-ideal} を参照

\noindent
{\it $H_1$-正則埋め込み}
\ref{divisors-definition-regular-immersion} を参照

\noindent
{\it $H_1$-正則埋め込み}
\ref{spaces-more-morphisms-definition-regular-immersion} を参照

\noindent
{\it $H_1$-正則}
\ref{more-algebra-definition-koszul-regular-sequence} を参照

\noindent
{\it $H_1$-正則}
\ref{divisors-definition-regular-ideal-sheaf} を参照

\noindent
{\it $H_{i + k}(A_\bullet) \rightarrow H_i(A[k]_\bullet)$}
\ref{homology-definition-homology-shift} を参照

\noindent
{\it $I$-進完備}
\ref{algebra-definition-complete} を参照

\noindent
{\it $I$-進完備}
\ref{algebra-definition-complete} を参照

\noindent
{\it $I$-深さ}
\ref{algebra-definition-depth} を参照

\noindent
{\it $I$-深さ}
\ref{local-cohomology-definition-depth-complex} を参照

\noindent
{\it $I$-冪ねじれ加群}
\ref{more-algebra-definition-f-power-torsion} を参照

\noindent
{\it $I$-射影的}
\ref{more-algebra-definition-near-projective} を参照

\noindent
{\it $\mathcal{E}$ の第 $i$ Chern 類}
\ref{spaces-chow-definition-chern-classes} を参照

\noindent
{\it 第 $i$ Chern 類}
\ref{chow-definition-chern-classes-final} を参照

\noindent
{\it $M$ の第 $i$ Chow 群}
\ref{weil-definition-chow-group-motives} を参照

\noindent
{\it 第 $i$ 拡大群}
\ref{derived-definition-ext} を参照

\noindent
{\it $F$ の第 $i$ 右導来関手 $R^iF$}
\ref{derived-definition-higher-derived-functors} を参照

\noindent
{\it $\mathcal{F}$ に付随する $k$-サイクル}
\ref{chow-definition-cycle-associated-to-coherent-sheaf} を参照

\noindent
{\it $\mathcal{F}$ に付随する $k$-サイクル}
\ref{spaces-chow-definition-cycle-associated-to-coherent-sheaf} を参照

\noindent
{\it $Y$ に付随する $k$-サイクル}
\ref{spaces-chow-definition-cycle-associated-to-closed-subscheme} を参照

\noindent
{\it $Z$ に付随する $k$-サイクル}
\ref{chow-definition-cycle-associated-to-closed-subscheme} を参照

\noindent
{\it $k$-サイクル}
\ref{chow-definition-cycles} を参照

\noindent
{\it $k$-サイクル}
\ref{spaces-chow-definition-cycles} を参照

\noindent
{\it $k$-シフト鎖複体 $A[k]_\bullet$}
\ref{homology-definition-shift} を参照

\noindent
{\it $k$-シフト余鎖複体 $A[k]^\bullet$}
\ref{homology-definition-shift-cochain} を参照

\noindent
{\it $k$-シフト加群}
\ref{dga-definition-shift} を参照

\noindent
{\it 第 $k$ Fitting イデアル}
\ref{more-algebra-definition-fitting-ideal} を参照

\noindent
{\it 第 $k$ シフト $A$-加群}
\ref{dga-definition-shift-graded-module} を参照

\noindent
{\it 第 $k$ シフト $A$-加群}
\ref{dga-definition-shift-graded-module} を参照

\noindent
{\it $\mathcal{F}$ の $L$-関数}
\ref{trace-definition-L-function-finite-ring} を参照

\noindent
{\it $\mathcal{F}$ の $L$-関数}
\ref{trace-definition-L-function-l-adic} を参照

\noindent
{\it $M \mapsto M^\vee$}
\ref{more-algebra-definition-simple-functors} を参照

\noindent
{\it $M$-$H_1$-正則}
\ref{more-algebra-definition-koszul-regular-sequence} を参照

\noindent
{\it $M$-Koszul-正則}
\ref{more-algebra-definition-koszul-regular-sequence} を参照

\noindent
{\it $R$ に関して $m$-擬連接}
\ref{more-algebra-definition-relatively-pseudo-coherent} を参照

\noindent
{\it $R$ に関して $m$-擬連接}
\ref{more-algebra-definition-relatively-pseudo-coherent} を参照

\noindent
{\it $S$ に相対的に $m$-擬連接}
\ref{more-morphisms-definition-relative-pseudo-coherence} を参照

\noindent
{\it $S$ に相対的に $m$-擬連接}
\ref{more-morphisms-definition-relative-pseudo-coherence} を参照

\noindent
{\it $Y$ に関して $m$-擬連接}
\ref{spaces-more-morphisms-definition-relative-pseudo-coherence} を参照

\noindent
{\it $Y$ に関して $m$-擬連接}
\ref{spaces-more-morphisms-definition-relative-pseudo-coherence} を参照

\noindent
{\it $m$-擬連接}
\ref{more-algebra-definition-pseudo-coherent} を参照

\noindent
{\it $m$-擬連接}
\ref{more-algebra-definition-pseudo-coherent} を参照

\noindent
{\it $m$-擬連接}
\ref{cohomology-definition-pseudo-coherent} を参照

\noindent
{\it $m$-擬連接}
\ref{cohomology-definition-pseudo-coherent} を参照

\noindent
{\it $m$-擬連接}
\ref{sites-cohomology-definition-pseudo-coherent} を参照

\noindent
{\it $m$-擬連接}
\ref{sites-cohomology-definition-pseudo-coherent} を参照

\noindent
{\it $M$-準正則}
\ref{algebra-definition-quasi-regular-sequence} を参照

\noindent
{\it $I$ 内の $M$-正則列}
\ref{algebra-definition-regular-sequence} を参照

\noindent
{\it $M$-正則列}
\ref{algebra-definition-regular-sequence} を参照

\noindent
{\it $m$-正則}
\ref{varieties-definition-regularity} を参照

\noindent
{\it $U$ の $n$-単体}
\ref{simplicial-definition-terminology-simplicial-sets} を参照

\noindent
{\it $\mathcal{C}$ の $n$-切断単体対象}
\ref{simplicial-definition-truncated-simplicial-object} を参照

\noindent
{\it $n$ 次無限小近傍}
\ref{more-morphisms-definition-first-order-infinitesimal-neighbourhood} を参照

\noindent
{\it $n$ 次無限小近傍}
\ref{spaces-more-morphisms-definition-first-order-infinitesimal-neighbourhood} を参照

\noindent
{\it $n$ 次肥厚化}
\ref{more-morphisms-definition-thickening} を参照

\noindent
{\it $n$ 次肥厚化}
\ref{spaces-more-morphisms-definition-thickening} を参照

\noindent
{\it $R$-双線形}
\ref{algebra-definition-bilinear} を参照

\noindent
{\it $R$-導分}
\ref{algebra-definition-derivation} を参照

\noindent
{\it $R$-同値}
\ref{groupoids-quotients-definition-geometric-orbits} を参照

\noindent
{\it $R$-不変}
\ref{groupoids-definition-invariant-open} を参照

\noindent
{\it $R$-不変}
\ref{groupoids-definition-invariant-open} を参照

\noindent
{\it $R$-不変}
\ref{groupoids-definition-invariant-open} を参照

\noindent
{\it $R$-不変}
\ref{spaces-groupoids-definition-invariant-open} を参照

\noindent
{\it $R$-不変}
\ref{spaces-groupoids-definition-invariant-open} を参照

\noindent
{\it $R$-不変}
\ref{spaces-groupoids-definition-invariant-open} を参照

\noindent
{\it $R$-不変}
\ref{groupoids-quotients-definition-invariant} を参照

\noindent
{\it $R$-線形圏 $\mathcal{A}$}
\ref{dga-definition-linear-category} を参照

\noindent
{\it $R$-線形関手}
\ref{dga-definition-functor-linear-categories} を参照

\noindent
{\it $R$-線形}
\ref{formal-defos-definition-linear} を参照

\noindent
{\it 有限表示の $R$-加群}
\ref{algebra-definition-module-finite-type} を参照

\noindent
{\it $R$-軌道}
\ref{groupoids-quotients-definition-orbit} を参照

\noindent
{\it $R$-軌道}
\ref{groupoids-quotients-definition-geometric-orbits} を参照

\noindent
{\it $R$-完全}
\ref{more-algebra-definition-relatively-perfect} を参照

\noindent
{\it $R\text{-}G$-加群}
\ref{etale-cohomology-definition-G-module-continuous} を参照

\noindent
{\it $R_{(f)}$}
\ref{exercises-definition-Dplus-Vplus} を参照

\noindent
{\it $S$ は有限型 $R$-代数である}
\ref{algebra-definition-finite-type} を参照

\noindent
{\it $S$-双有理}
\ref{morphisms-definition-birational-schemes} を参照

\noindent
{\it $S$-導分 $D : \mathcal{O}_{X/S} \to \mathcal{F}$}
\ref{crystalline-definition-global-derivation} を参照

\noindent
{\it $S$-導分}
\ref{modules-definition-differentials} を参照

\noindent
{\it $S$-完全}
\ref{perfect-definition-relatively-perfect} を参照

\noindent
{\it $S$-純粋}
\ref{flat-definition-pure} を参照

\noindent
{\it $S$-純粋}
\ref{flat-definition-pure} を参照

\noindent
{\it $S$-有理写像（$X$ から $Y$ への）}
\ref{morphisms-definition-rational-map} を参照

\noindent
{\it $s^{-1} f$}
\ref{categories-definition-left-localization-as-fraction} を参照

\noindent
{\it $T$ は $Y$ 上固有}
\ref{spaces-perfect-definition-proper-over-base} を参照

\noindent
{\it $U$-許容ブローアップ}
\ref{divisors-definition-admissible-blowup} を参照

\noindent
{\it $U$-許容ブローアップ}
\ref{spaces-divisors-definition-admissible-blowup} を参照

\noindent
{\it $x$ は $X$ 上の余次元 $d$ の点である}
\ref{spaces-properties-definition-dimension-local-ring} を参照

\noindent
{\it $x$ は $X$ の付随点}
\ref{spaces-divisors-definition-weakly-associated} を参照

\noindent
{\it $x$ は $\mathcal{F}$ に付随する}
\ref{spaces-divisors-definition-weakly-associated} を参照

\noindent
{\it $X$ は $x$ において正則}
\ref{spaces-properties-definition-regular-at-point} を参照

\noindent
{\it $x$ は $U$ 上にある}
\ref{categories-definition-fibre-category} を参照

\noindent
{\it $X_{affine, \etale}$}
\ref{spaces-properties-definition-affine-etale-site} を参照

\noindent
{\it $X_{spaces, \etale}$}
\ref{spaces-properties-definition-spaces-etale-site} を参照

\noindent
{\it $Y$ は $X$ 上カルテジアンである}
\ref{spaces-simplicial-definition-cartesian-morphism} を参照

\noindent
{\it $Y$-導分}
\ref{sites-modules-definition-sheaf-differentials} を参照

\noindent
{\it $Y$-完全}
\ref{spaces-more-morphisms-definition-relatively-perfect} を参照

\noindent
{\it $Y$-純粋}
\ref{spaces-flat-definition-pure} を参照

\noindent
{\it $Y$-純粋}
\ref{spaces-flat-definition-pure} を参照

\noindent
{\it $Z$ は $S$ 上固有}
\ref{coherent-definition-proper-over-base} を参照

\noindent
{\it （加法的）Herbrand 商}
\ref{chow-definition-periodic-length} を参照

\noindent
{\it $f$ と $g$ の 2-ファイバー積}
\ref{categories-definition-2-fibre-products} を参照

\noindent
{\it $f$ から $g$ への 2-射}
\ref{sites-definition-2-morphism-topoi} を参照

\noindent
{\it $f$ から $g$ への 2-射}
\ref{sites-modules-definition-2-morphism-ringed-topoi} を参照

\noindent
{\it $\mathfrak q$ においてエタール}
\ref{algebra-definition-etale} を参照

\noindent
{\it $x \in X$ においてエタール}
\ref{morphisms-definition-etale} を参照

\noindent
{\it $x \in X$ においてエタール}
\ref{etale-definition-etale-schemes-1} を参照

\noindent
{\it $x$ においてエタール}
\ref{spaces-morphisms-definition-etale} を参照

\noindent
{\it $T$ のエタール被覆}
\ref{topologies-definition-etale-covering} を参照

\noindent
{\it $X$ のエタール被覆}
\ref{spaces-topologies-definition-etale-covering} を参照

\noindent
{\it エタール被覆}
\ref{etale-cohomology-definition-etale-covering-initial} を参照

\noindent
{\it エタール被覆}
\ref{etale-cohomology-definition-etale-covering} を参照

\noindent
{\it エタール同値関係}
\ref{spaces-definition-etale-equivalence-relation} を参照

\noindent
{\it 局所環のエタール準同型}
\ref{etale-definition-etale-ring} を参照

\noindent
{\it 始域と標的上エタール局所的}
\ref{descent-definition-local-source-target} を参照

\noindent
{\it 始域と標的上エタール局所的}
\ref{descent-definition-local-source-target-at-point} を参照

\noindent
{\it $\overline{s}$ における $S$ のエタール局所環}
\ref{etale-cohomology-definition-etale-local-rings} を参照

\noindent
{\it $\overline{x}$ における $X$ のエタール局所環}
\ref{spaces-properties-definition-etale-local-rings} を参照

\noindent
{\it エタール局所構成可能}
\ref{spaces-properties-definition-locally-constructible} を参照

\noindent
{\it エタール局所的}
\ref{descent-definition-local-at-point} を参照

\noindent
{\it エタール近傍}
\ref{etale-cohomology-definition-geometric-point} を参照

\noindent
{\it エタール近傍}
\ref{spaces-properties-definition-etale-neighbourhood} を参照

\noindent
{\it $(S, s)$ のエタール近傍}
\ref{more-morphisms-definition-etale-neighbourhood} を参照

\noindent
{\it エタール層}
\ref{stacks-sheaves-definition-sheaves} を参照

\noindent
{\it エタールトポス}
\ref{etale-cohomology-definition-etale-topos} を参照

\noindent
{\it エタールトポス}
\ref{spaces-properties-definition-etale-topos} を参照

\noindent
{\it エタール平滑は始域と終域について局所的}
\ref{spaces-descent-definition-etale-smooth-local-source-target} を参照

\noindent
{\it エタール}
\ref{algebra-definition-etale} を参照

\noindent
{\it エタール}
\ref{morphisms-definition-etale} を参照

\noindent
{\it エタール}
\ref{descent-definition-etale-morphism-germs} を参照

\noindent
{\it エタール}
\ref{etale-definition-etale-schemes-1} を参照

\noindent
{\it エタール}
\ref{etale-cohomology-definition-etale-morphism} を参照

\noindent
{\it エタール}
\ref{spaces-properties-definition-etale} を参照

\noindent
{\it エタール}
\ref{stacks-morphisms-definition-etale} を参照

\noindent
{\it Serre 関手が存在する}
\ref{equiv-definition-Serre-functor} を参照

\noindent
{\it $X$ 上のアーベル前層}
\ref{sheaves-definition-abelian-presheaves} を参照

\noindent
{\it アーベル群の前層}
\ref{etale-cohomology-definition-presheaf} を参照

\noindent
{\it $X$ 上のアーベル層}
\ref{sheaves-definition-abelian-sheaf} を参照

\noindent
{\it アーベル群の層}
\ref{etale-cohomology-definition-category-sheaves} を参照

\noindent
{\it アーベル多様体}
\ref{groupoids-definition-abelian-variety} を参照

\noindent
{\it アーベル}
\ref{homology-definition-abelian-category} を参照

\noindent
{\it $X$ の絶対 Frobenius}
\ref{varieties-definition-absolute-frobenius} を参照

\noindent
{\it 絶対ガロア群}
\ref{etale-cohomology-definition-algebraic-geometric-point} を参照

\noindent
{\it 絶対分岐指数}
\ref{more-algebra-definition-mixed} を参照

\noindent
{\it 絶対弱正規化}
\ref{morphisms-definition-seminormalization} を参照

\noindent
{\it 絶対平坦}
\ref{more-algebra-definition-weakly-etale} を参照

\noindent
{\it 絶対平坦}
\ref{more-algebra-definition-weakly-etale} を参照

\noindent
{\it 絶対平坦}
\ref{more-morphisms-definition-weakly-etale} を参照

\noindent
{\it 絶対整閉}
\ref{more-algebra-definition-absolutely-integrally-closed} を参照

\noindent
{\it 絶対弱正規}
\ref{morphisms-definition-seminormal-ring} を参照

\noindent
{\it 絶対弱正規}
\ref{morphisms-definition-seminormal} を参照

\noindent
{\it $H(K)$ に収束する}
\ref{homology-definition-filtered-differential-ss-converges} を参照

\noindent
{\it $H^*(K^\bullet)$ に収束する}
\ref{homology-definition-filtered-complex-ss-converges} を参照

\noindent
{\it $H^n(\text{Tot}(K^{\bullet, \bullet}))$ に収束する}
\ref{homology-definition-ss-double-complex-converge} を参照

\noindent
{\it $H^n(\text{Tot}(K^{\bullet, \bullet}))$ に収束する}
\ref{homology-definition-ss-double-complex-converge} を参照

\noindent
{\it 代数空間 $X/B$ への $G$ の作用}
\ref{spaces-groupoids-definition-action-group-space} を参照

\noindent
{\it $G$ のスキーム $X/S$ への作用}
\ref{groupoids-definition-action-group-scheme} を参照

\noindent
{\it 自由に作用する}
\ref{spaces-definition-quotient} を参照

\noindent
{\it $LF$ に関して非輪状}
\ref{derived-definition-derived-functor} を参照

\noindent
{\it $RF$ に関して非輪状}
\ref{derived-definition-derived-functor} を参照

\noindent
{\it 非輪状}
\ref{homology-definition-quasi-isomorphism} を参照

\noindent
{\it 非輪状}
\ref{homology-definition-quasi-isomorphism-cochain} を参照

\noindent
{\it 加法的モノイド圏}
\ref{homology-definition-additive-monoidal} を参照

\noindent
{\it 加法的}
\ref{homology-definition-preadditive} を参照

\noindent
{\it 加法的}
\ref{homology-definition-additive-category} を参照

\noindent
{\it 適切}
\ref{adequate-definition-adequate-functor} を参照

\noindent
{\it 適切}
\ref{adequate-definition-adequate} を参照

\noindent
{\it アディック構成可能}
\ref{proetale-definition-adic} を参照

\noindent
{\it アディック構成可能}
\ref{proetale-definition-adic-constructible} を参照

\noindent
{\it アディック滑らか}
\ref{proetale-definition-adic} を参照

\noindent
{\it アディック滑らか}
\ref{proetale-definition-adic-constructible} を参照

\noindent
{\it アディック射}
\ref{formal-spaces-definition-adic-morphism} を参照

\noindent
{\it アディック*}
\ref{formal-spaces-definition-types-affine-formal-algebraic-space} を参照

\noindent
{\it アディック}
\ref{more-algebra-definition-topological-ring} を参照

\noindent
{\it アディック}
\ref{formal-spaces-definition-adic-homomorphism} を参照

\noindent
{\it アディック}
\ref{formal-spaces-definition-types-affine-formal-algebraic-space} を参照

\noindent
{\it 許容全射}
\ref{dga-definition-admissible-ses} を参照

\noindent
{\it 許容単射}
\ref{dga-definition-admissible-ses} を参照

\noindent
{\it 許容関係}
\ref{chow-definition-determinant} を参照

\noindent
{\it 許容短完全列}
\ref{dga-definition-admissible-ses} を参照

\noindent
{\it 許容}
\ref{more-algebra-definition-topological-ring} を参照

\noindent
{\it 許容}
\ref{chow-definition-determinant} を参照

\noindent
{\it アフィン $n$-空間（$R$ 上）}
\ref{constructions-definition-affine-n-space} を参照

\noindent
{\it アフィン $n$-空間（$S$ 上）}
\ref{constructions-definition-affine-n-space} を参照

\noindent
{\it アフィン・ブローアップ代数}
\ref{algebra-definition-blow-up} を参照

\noindent
{\it $\mathcal{A}$ に付随するアフィン錐}
\ref{constructions-definition-cone} を参照

\noindent
{\it アフィン形式代数空間}
\ref{formal-spaces-definition-affine-formal-algebraic-space} を参照

\noindent
{\it アフィンスキーム}
\ref{schemes-definition-affine-scheme} を参照

\noindent
{\it アフィン層別数}
\ref{more-morphisms-definition-affine-stratification-number} を参照

\noindent
{\it アフィン層別}
\ref{more-morphisms-definition-affine-stratification} を参照

\noindent
{\it アフィン多様体}
\ref{varieties-definition-variety-type} を参照

\noindent
{\it アフィン}
\ref{morphisms-definition-affine} を参照

\noindent
{\it アフィン}
\ref{spaces-morphisms-definition-affine} を参照

\noindent
{\it アフィン}
\ref{stacks-morphisms-definition-affine} を参照

\noindent
{\it 代数的 $k$-スキーム}
\ref{varieties-definition-algebraic-scheme} を参照

\noindent
{\it $K$ における $k$ の代数閉包}
\ref{fields-definition-algebraically-closed-in} を参照

\noindent
{\it 代数閉包}
\ref{fields-definition-algebraic-closure} を参照

\noindent
{\it 代数拡大}
\ref{fields-definition-algebraic} を参照

\noindent
{\it $S$ 上の代数空間}
\ref{spaces-definition-algebraic-space} を参照

\noindent
{\it $Z$ 上の代数空間構造}
\ref{spaces-properties-definition-reduced-induced-space} を参照

\noindent
{\it $S$ 上の代数スタック}
\ref{algebraic-definition-algebraic-stack} を参照

\noindent
{\it $Z$ 上の代数スタック構造}
\ref{stacks-properties-definition-reduced-induced-stack} を参照

\noindent
{\it 代数スタック}
\ref{stacks-introduction-definition-algebraic-stack} を参照

\noindent
{\it $K$ において代数的に閉じている}
\ref{fields-definition-algebraically-closed-in} を参照

\noindent
{\it 代数閉}
\ref{fields-definition-algebraically-closed} を参照

\noindent
{\it 代数的独立}
\ref{fields-definition-transcendence} を参照

\noindent
{\it 代数的}
\ref{fields-definition-algebraic} を参照

\noindent
{\it 代数的}
\ref{fields-definition-separable-algebraic} を参照

\noindent
{\it 代数的}
\ref{etale-cohomology-definition-algebraic-geometric-point} を参照

\noindent
{\it 代数的}
\ref{criteria-definition-algebraic} を参照

\noindent
{\it ほとんど余連続}
\ref{sites-definition-almost-cocontinuous} を参照

\noindent
{\it $R$ 上概整}
\ref{algebra-definition-almost-integral} を参照

\noindent
{\it $X$ のオルタレーション}
\ref{morphisms-definition-alteration} を参照

\noindent
{\it $X$ のオルタレーション}
\ref{spaces-over-fields-definition-alteration} を参照

\noindent
{\it 交代 {\v C}ech 複体}
\ref{cohomology-definition-alternating-cech-complex} を参照

\noindent
{\it 交代 {\v C}ech 複体}
\ref{spaces-cohomology-definition-alternating-cech-complex} を参照

\noindent
{\it 融合和}
\ref{categories-definition-coproducts} を参照

\noindent
{\it $X$ 上の豊富な可逆層の族}
\ref{morphisms-definition-family-ample-invertible-modules} を参照

\noindent
{\it $X/S$ 上豊富}
\ref{morphisms-definition-relatively-ample} を参照

\noindent
{\it $X/Y$ 上豊富}
\ref{spaces-divisors-definition-relatively-ample} を参照

\noindent
{\it 豊富}
\ref{properties-definition-ample} を参照

\noindent
{\it $f$-冪ねじれ加群}
\ref{more-algebra-definition-f-power-torsion} を参照

\noindent
{\it $R$ の定義イデアル}
\ref{algebra-definition-ideal-definition} を参照

\noindent
{\it 解析的不分岐}
\ref{algebra-definition-analytically-unramified} を参照

\noindent
{\it 解析的不分岐}
\ref{algebra-definition-analytically-unramified} を参照

\noindent
{\it $m$ の零化イデアル}
\ref{algebra-definition-annihilator} を参照

\noindent
{\it $M$ の零化イデアル}
\ref{algebra-definition-annihilator} を参照

\noindent
{\it 零化イデアル}
\ref{modules-definition-annihilator-sheaf} を参照

\noindent
{\it 完全複体による近似が 成立する}
\ref{perfect-definition-approximation} を参照

\noindent
{\it 完全複体による近似が 成立する}
\ref{spaces-perfect-definition-approximation} を参照

\noindent
{\it 三つ組に対して近似が成立する}
\ref{perfect-definition-approximation-holds} を参照

\noindent
{\it 三つ組に対して近似が成立する}
\ref{spaces-perfect-definition-approximation-holds} を参照

\noindent
{\it 算術的フロベニウス}
\ref{trace-definition-arithmetic-frobenius} を参照

\noindent
{\it アルティン的}
\ref{algebra-definition-artinian} を参照

\noindent
{\it アルティン的}
\ref{homology-definition-Artinian} を参照

\noindent
{\it アルティン的}
\ref{homology-definition-Artinian} を参照

\noindent
{\it アルティン的}
\ref{exercises-definition-Noetherian-space} を参照

\noindent
{\it 随伴エタールサイト}
\ref{stacks-sheaves-definition-inherited-topologies} を参照

\noindent
{\it 付随アフィンエタール・サイト}
\ref{stacks-sheaves-definition-alternative-inherited-topologies} を参照

\noindent
{\it 付随するアフィン fppf サイト}
\ref{stacks-sheaves-definition-alternative-inherited-topologies} を参照

\noindent
{\it 付随するアフィンサイト}
\ref{stacks-sheaves-definition-alternative} を参照

\noindent
{\it 付随アフィン滑らかサイト}
\ref{stacks-sheaves-definition-alternative-inherited-topologies} を参照

\noindent
{\it 付随アフィンシントミック・サイト}
\ref{stacks-sheaves-definition-alternative-inherited-topologies} を参照

\noindent
{\it 付随するアフィン Zariski サイト}
\ref{stacks-sheaves-definition-alternative-inherited-topologies} を参照

\noindent
{\it 随伴 fppf サイト}
\ref{stacks-sheaves-definition-inherited-topologies} を参照

\noindent
{\it 付随次数付き環}
\ref{modules-definition-gamma-star} を参照

\noindent
{\it 随伴する fppf トポスの射}
\ref{stacks-sheaves-definition-morphism} を参照

\noindent
{\it $X$ の随伴点}
\ref{divisors-definition-associated} を参照

\noindent
{\it 付随単複体}
\ref{homology-definition-associated-simple-complex} を参照

\noindent
{\it 随伴滑らかなサイト}
\ref{stacks-sheaves-definition-inherited-topologies} を参照

\noindent
{\it 付随シントミック・サイト}
\ref{stacks-sheaves-definition-inherited-topologies} を参照

\noindent
{\it 付随全複体}
\ref{homology-definition-associated-simple-complex} を参照

\noindent
{\it 随伴 Zariski サイト}
\ref{stacks-sheaves-definition-inherited-topologies} を参照

\noindent
{\it 随伴する}
\ref{algebra-definition-associated} を参照

\noindent
{\it 随伴する}
\ref{divisors-definition-associated} を参照

\noindent
{\it 同伴}
\ref{algebra-definition-irreducible-prime-element} を参照

\noindent
{\it 相対次元 $1$ の高々節点的}
\ref{curves-definition-nodal-family} を参照

\noindent
{\it 相対次元 $1$ の高々節点的}
\ref{spaces-more-morphisms-definition-nodal-family} を参照

\noindent
{\it Atiyah 類}
\ref{cotangent-definition-atiyah-class} を参照

\noindent
{\it Atiyah 類}
\ref{cotangent-definition-atiyah-class-general} を参照

\noindent
{\it $\mathcal{C}$ の対象 $X$ へ向かう $U$ の増大写像 $\epsilon : U \to X$}
\ref{simplicial-definition-augmentation} を参照

\noindent
{\it 自己随伴}
\ref{more-algebra-definition-auto-ass} を参照

\noindent
{\it $x$ の自己同型関手}
\ref{formal-defos-definition-automorphism-functor} を参照

\noindent
{\it $F$ 上の $E$ の自己同型}
\ref{fields-definition-automorphisms} を参照

\noindent
{\it $E/F$ の自己同型}
\ref{fields-definition-automorphisms} を参照

\noindent
{\it B\'ezout 整域}
\ref{more-algebra-definition-bezout} を参照

\noindent
{\it $F'$ の $S$ への基底変換}
\ref{spaces-definition-base-change} を参照

\noindent
{\it 基底変換}
\ref{algebra-definition-base-change} を参照

\noindent
{\it 基底変換}
\ref{algebra-definition-base-change} を参照

\noindent
{\it 基底変換}
\ref{schemes-definition-base-change} を参照

\noindent
{\it 基底変換}
\ref{schemes-definition-base-change} を参照

\noindent
{\it 基底変換}
\ref{schemes-definition-base-change} を参照

\noindent
{\it 基底変換}
\ref{groupoids-quotients-definition-base-change} を参照

\noindent
{\it $f$ に沿う基底拡大}
\ref{descent-definition-effective-descent} を参照

\noindent
{\it $X$ の位相の開基}
\ref{topology-definition-base} を参照

\noindent
{\it 基点}
\ref{pione-definition-fundamental-group} を参照

\noindent
{\it $X$ の位相の基底}
\ref{topology-definition-base} を参照

\noindent
{\it $S$ の大 $\tau$-サイト}
\ref{etale-cohomology-definition-tau-site} を参照

\noindent
{\it 大 $\tau$-トポス}
\ref{etale-cohomology-definition-etale-topos} を参照

\noindent
{\it $S$ の大エタール・サイト}
\ref{topologies-definition-big-small-etale} を参照

\noindent
{\it $S$ 上の大エタール・サイト}
\ref{etale-cohomology-definition-big-etale-site} を参照

\noindent
{\it 大エタールサイト}
\ref{topologies-definition-big-etale-site} を参照

\noindent
{\it $S$ の大アフィンエタール・サイト}
\ref{topologies-definition-big-small-etale} を参照

\noindent
{\it $S$ の大アフィン fppf サイト}
\ref{topologies-definition-big-small-fppf} を参照

\noindent
{\it $S$ の大アフィン h サイト}
\ref{flat-definition-big-small-h} を参照

\noindent
{\it $S$ の大アフィン ph サイト}
\ref{topologies-definition-big-small-ph} を参照

\noindent
{\it $S$ の大アフィン pro-エタール・サイト}
\ref{proetale-definition-big-small-proetale} を参照

\noindent
{\it $S$ の大アフィン滑らかサイト}
\ref{topologies-definition-big-small-smooth} を参照

\noindent
{\it $S$ の大アフィンシントミック・サイト}
\ref{topologies-definition-big-small-syntomic} を参照

\noindent
{\it $S$ の大アフィン Zariski サイト}
\ref{topologies-definition-big-small-Zariski} を参照

\noindent
{\it 大結晶サイト}
\ref{crystalline-definition-big-crystalline-site} を参照

\noindent
{\it $S$ の大 fppf サイト}
\ref{topologies-definition-big-small-fppf} を参照

\noindent
{\it 大 fppf サイト}
\ref{topologies-definition-big-fppf-site} を参照

\noindent
{\it $S$ の大 h サイト}
\ref{flat-definition-big-small-h} を参照

\noindent
{\it 大 h サイト}
\ref{flat-definition-big-h-site} を参照

\noindent
{\it $S$ の大 ph サイト}
\ref{topologies-definition-big-small-ph} を参照

\noindent
{\it 大 ph サイト}
\ref{topologies-definition-big-ph-site} を参照

\noindent
{\it $S$ の大 pro-エタール・サイト}
\ref{proetale-definition-big-small-proetale} を参照

\noindent
{\it 大 pro-エタール・サイト}
\ref{proetale-definition-big-proetale-site} を参照

\noindent
{\it $S$ の大滑らかサイト}
\ref{topologies-definition-big-small-smooth} を参照

\noindent
{\it 大平滑サイト}
\ref{topologies-definition-big-smooth-site} を参照

\noindent
{\it $S$ の大シントミック・サイト}
\ref{topologies-definition-big-small-syntomic} を参照

\noindent
{\it 大シントミック・サイト}
\ref{topologies-definition-big-syntomic-site} を参照

\noindent
{\it $S$ の大 Zariski サイト}
\ref{topologies-definition-big-small-Zariski} を参照

\noindent
{\it 大 Zariski サイト}
\ref{topologies-definition-big-zariski-site} を参照

\noindent
{\it 大}
\ref{etale-cohomology-definition-big-etale-site} を参照

\noindent
{\it 双有理}
\ref{morphisms-definition-birational-schemes} を参照

\noindent
{\it 双有理}
\ref{morphisms-definition-birational} を参照

\noindent
{\it 双有理}
\ref{spaces-morphisms-definition-birational-spaces} を参照

\noindent
{\it 双有理}
\ref{decent-spaces-definition-birational} を参照

\noindent
{\it $f$ に対する次数 $p$ の双変類 $c$}
\ref{chow-definition-bivariant-class} を参照

\noindent
{\it $f$ に対する次数 $p$ の双変類 $c$}
\ref{spaces-chow-definition-bivariant-class} を参照

\noindent
{\it $x$ における $X$ のブローアップ $X' \to X$}
\ref{spaces-resolve-definition-blowup-at-point} を参照

\noindent
{\it $X$ の $Z$ に沿ったブローアップ}
\ref{divisors-definition-blow-up} を参照

\noindent
{\it $X$ の $Z$ に沿ったブローアップ}
\ref{spaces-divisors-definition-blow-up} を参照

\noindent
{\it $X$ のイデアル層 $\mathcal{I}$ におけるブローアップ}
\ref{divisors-definition-blow-up} を参照

\noindent
{\it $X$ のイデアル層 $\mathcal{I}$ におけるブローアップ}
\ref{spaces-divisors-definition-blow-up} を参照

\noindent
{\it ブローアップ代数}
\ref{algebra-definition-blow-up} を参照

\noindent
{\it 上に有界}
\ref{homology-definition-bounded-ss} を参照

\noindent
{\it 上に有界}
\ref{derived-definition-complexes-notation} を参照

\noindent
{\it 下に有界}
\ref{homology-definition-bounded-ss} を参照

\noindent
{\it 下に有界}
\ref{derived-definition-complexes-notation} を参照

\noindent
{\it 有界導来圏}
\ref{derived-definition-unbounded-derived-category} を参照

\noindent
{\it 有界フィルトレーション付き導来圏}
\ref{derived-definition-filtered-derived-bounded} を参照

\noindent
{\it 有界}
\ref{homology-definition-bounded-ss} を参照

\noindent
{\it 有界}
\ref{derived-definition-complexes-notation} を参照

\noindent
{\it $f$ のファイバーの次数を有界にする}
\ref{morphisms-definition-universally-bounded} を参照

\noindent
{\it Bourbaki 固有}
\ref{topology-definition-proper-map} を参照

\noindent
{\it ブラウアー群}
\ref{brauer-definition-brauer-group} を参照

\noindent
{\it ブラウアー群}
\ref{etale-cohomology-definition-brauer-group} を参照

\noindent
{\it 標準降下データ}
\ref{stacks-definition-effective-descent-datum} を参照

\noindent
{\it 標準降下データ}
\ref{descent-definition-descent-datum-effective-quasi-coherent} を参照

\noindent
{\it 標準降下データ}
\ref{descent-definition-effective} を参照

\noindent
{\it 標準降下データ}
\ref{descent-definition-effective-family} を参照

\noindent
{\it 標準降下データ}
\ref{spaces-descent-definition-descent-datum-effective-quasi-coherent} を参照

\noindent
{\it 標準降下データ}
\ref{spaces-descent-definition-effective} を参照

\noindent
{\it 標準降下データ}
\ref{spaces-descent-definition-effective-family} を参照

\noindent
{\it $T$ 上の標準的なスキーム構造}
\ref{morphisms-definition-scheme-structure-connected-component} を参照

\noindent
{\it 標準切断}
\ref{divisors-definition-invertible-sheaf-effective-Cartier-divisor} を参照

\noindent
{\it 標準位相}
\ref{sites-definition-canonical-topology} を参照

\noindent
{\it Cartan--Eilenberg 分解}
\ref{derived-definition-cartan-eilenberg} を参照

\noindent
{\it カルテジアン}
\ref{categories-definition-cartesian} を参照

\noindent
{\it カルテジアン}
\ref{groupoids-definition-cartesian-morphism} を参照

\noindent
{\it カルテジアン}
\ref{spaces-simplicial-definition-cartesian-sheaf} を参照

\noindent
{\it カルテジアン}
\ref{spaces-simplicial-definition-cartesian-sheaf} を参照

\noindent
{\it カルテジアン}
\ref{spaces-simplicial-definition-cartesian-sheaf} を参照

\noindent
{\it カルテジアン}
\ref{spaces-simplicial-definition-cartesian-sheaf} を参照

\noindent
{\it カルテジアン}
\ref{spaces-simplicial-definition-cartesian-derived} を参照

\noindent
{\it カルテジアン}
\ref{spaces-simplicial-definition-cartesian-derived-modules} を参照

\noindent
{\it カルテジアン}
\ref{spaces-simplicial-definition-cartesian-morphism} を参照

\noindent
{\it Cartier 因子}
\ref{exercises-definition-divisor} を参照

\noindent
{\it $\mathcal{C}$ における圏論的モジュライ空間}
\ref{stacks-more-morphisms-definition-categorical-quotient} を参照

\noindent
{\it 圏論的モジュライ空間}
\ref{stacks-more-morphisms-definition-categorical-quotient} を参照

\noindent
{\it $\mathcal{C}$ における圏論的商}
\ref{groupoids-quotients-definition-categorical} を参照

\noindent
{\it スキームにおける圏論的商}
\ref{groupoids-quotients-definition-categorical} を参照

\noindent
{\it スキームの圏における圏論的商}
\ref{groupoids-quotients-definition-categorical} を参照

\noindent
{\it 圏論的商}
\ref{groupoids-quotients-definition-categorical} を参照

\noindent
{\it 圏論的コンパクト}
\ref{categories-definition-compact-object} を参照

\noindent
{\it $\mathcal{F}$ の形式対象の圏 $\widehat{\mathcal{F}}$}
\ref{formal-defos-definition-formal-objects} を参照

\noindent
{\it $\mathcal{C}$ 上の亜群余ファイバー圏}
\ref{formal-defos-definition-category-cofibred-groupoids} を参照

\noindent
{\it 離散圏にファイバー化された圏}
\ref{categories-definition-category-fibred-sets} を参照

\noindent
{\it セトイドにファイバー化された圏}
\ref{categories-definition-category-fibred-setoids} を参照

\noindent
{\it 集合にファイバー化された圏}
\ref{categories-definition-category-fibred-sets} を参照

\noindent
{\it （余鎖）複体の圏}
\ref{derived-definition-complexes-notation} を参照

\noindent
{\it $\mathcal{A}$ の複体の圏}
\ref{dga-definition-homotopy-category-of-dga-category} を参照

\noindent
{\it $\mathcal{C}$ の余単体的対象の圏}
\ref{simplicial-definition-cosimplicial-object} を参照

\noindent
{\it $\mathcal{A}$ の有限フィルトレーション付き 対象の圏}
\ref{derived-definition-finite-filtered} を参照

\noindent
{\it $\mathcal{A}$ の次数付き対象の圏}
\ref{homology-definition-graded} を参照

\noindent
{\it $\mathcal{C}$ 上の関手における亜群の圏}
\ref{formal-defos-definition-groupoid-in-functors} を参照

\noindent
{\it 集合の層の圏}
\ref{etale-cohomology-definition-category-sheaves} を参照

\noindent
{\it $\mathcal{C}$ の単体的対象の圏}
\ref{simplicial-definition-simplicial-object} を参照

\noindent
{\it 圏}
\ref{categories-definition-category} を参照

\noindent
{\it 鎖状}
\ref{topology-definition-catenary} を参照

\noindent
{\it 鎖状}
\ref{algebra-definition-catenary} を参照

\noindent
{\it 鎖状}
\ref{properties-definition-catenary} を参照

\noindent
{\it 鎖状}
\ref{decent-spaces-definition-catenary} を参照

\noindent
{\it 鎖状}
\ref{exercises-definition-catenary} を参照

\noindent
{\it 鎖状}
\ref{exercises-definition-catenary} を参照

\noindent
{\it 中心をもつ}
\ref{algebra-definition-valuation-ring} を参照

\noindent
{\it 中心}
\ref{divisors-definition-blow-up} を参照

\noindent
{\it 中心}
\ref{spaces-divisors-definition-blow-up} を参照

\noindent
{\it 中心的}
\ref{brauer-definition-central} を参照

\noindent
{\it 既約閉部分集合の鎖}
\ref{topology-definition-Krull} を参照

\noindent
{\it 素イデアルの鎖}
\ref{algebra-definition-chain-primes} を参照

\noindent
{\it $\mathcal{X}'$ の基底変換}
\ref{algebraic-definition-change-of-base} を参照

\noindent
{\it 標数}
\ref{fields-definition-characteristic} を参照

\noindent
{\it $X$ 上の $\mathcal{E}$ の Chern 類}
\ref{chow-definition-chern-classes} を参照

\noindent
{\it $E$ の Chern 類は定義される}
\ref{chow-definition-defined-on-perfect} を参照

\noindent
{\it 引き戻しの選択}
\ref{categories-definition-pullback-functor-fibred-category} を参照

\noindent
{\it Chow コホモロジー}
\ref{chow-definition-chow-cohomology} を参照

\noindent
{\it Chow コホモロジー}
\ref{spaces-chow-definition-chow-cohomology} を参照

\noindent
{\it $X$ 上の有理同値を法とする $k$-サイクルの Chow 群}
\ref{chow-definition-rational-equivalence} を参照

\noindent
{\it $X$ 上の有理同値を法とする $k$-サイクルの Chow 群}
\ref{spaces-chow-definition-rational-equivalence} を参照

\noindent
{\it $X$ 上の $k$-サイクルの Chow 群}
\ref{chow-definition-rational-equivalence} を参照

\noindent
{\it $X$ 上の $k$-サイクルの Chow 群}
\ref{spaces-chow-definition-rational-equivalence} を参照

\noindent
{\it $A$ の類群}
\ref{exercises-definition-class-group} を参照

\noindent
{\it 古典的場合}
\ref{formal-defos-definition-CLambda} を参照

\noindent
{\it 古典的生成元}
\ref{derived-definition-generators} を参照

\noindent
{\it 古典的 Weil コホモロジー理論}
\ref{weil-definition-weil-cohomology-theory-classical} を参照

\noindent
{\it 古典的}
\ref{formal-spaces-definition-types-affine-formal-algebraic-space} を参照

\noindent
{\it 環付き空間の閉埋め込み}
\ref{modules-definition-closed-immersion} を参照

\noindent
{\it 閉埋め込み}
\ref{sites-definition-immersion-topoi} を参照

\noindent
{\it 閉埋め込み}
\ref{schemes-definition-closed-immersion-locally-ringed-spaces} を参照

\noindent
{\it 閉埋め込み}
\ref{schemes-definition-immersion} を参照

\noindent
{\it 閉埋め込み}
\ref{spaces-definition-immersion} を参照

\noindent
{\it 閉埋め込み}
\ref{formal-spaces-definition-closed-immersion} を参照

\noindent
{\it 閉埋め込み}
\ref{stacks-properties-definition-immersion} を参照

\noindent
{\it 閉部分群スキーム}
\ref{groupoids-definition-closed-subgroup-scheme} を参照

\noindent
{\it 閉部分スキーム}
\ref{schemes-definition-immersion} を参照

\noindent
{\it $X$ の、イデアルの層 $\mathcal{I}$ に付随する 閉部分空間}
\ref{schemes-definition-closed-subspace} を参照

\noindent
{\it 閉部分空間}
\ref{spaces-definition-immersion} を参照

\noindent
{\it 閉部分スタック}
\ref{stacks-properties-definition-substacks} を参照

\noindent
{\it 閉部分トポス}
\ref{sites-definition-closed-subtopos} を参照

\noindent
{\it 閉}
\ref{topology-definition-proper-map} を参照

\noindent
{\it 閉}
\ref{spaces-morphisms-definition-closed} を参照

\noindent
{\it 閉}
\ref{stacks-morphisms-definition-closed} を参照

\noindent
{\it 閉}
\ref{exercises-definition-closed} を参照

\noindent
{\it スキームにおける粗商}
\ref{groupoids-quotients-definition-coarse} を参照

\noindent
{\it 粗商}
\ref{groupoids-quotients-definition-coarse} を参照

\noindent
{\it より粗い}
\ref{sites-definition-finer} を参照

\noindent
{\it 余カルテジアン}
\ref{categories-definition-cocartesian} を参照

\noindent
{\it 余連続}
\ref{sites-definition-cocontinuous} を参照

\noindent
{\it コサイクル条件}
\ref{stacks-definition-descent-data} を参照

\noindent
{\it コサイクル条件}
\ref{descent-definition-descent-datum-quasi-coherent} を参照

\noindent
{\it コサイクル条件}
\ref{descent-definition-descent-datum-modules} を参照

\noindent
{\it コサイクル条件}
\ref{descent-definition-descent-datum} を参照

\noindent
{\it コサイクル条件}
\ref{etale-cohomology-definition-descent-datum} を参照

\noindent
{\it コサイクル条件}
\ref{spaces-descent-definition-descent-datum-quasi-coherent} を参照

\noindent
{\it コサイクル条件}
\ref{spaces-descent-definition-descent-datum} を参照

\noindent
{\it 余次元}
\ref{topology-definition-codimension} を参照

\noindent
{\it 余有向}
\ref{categories-definition-codirected} を参照

\noindent
{\it 余有向}
\ref{categories-definition-codirected} を参照

\noindent
{\it 係数環}
\ref{algebra-definition-coefficient-ring} を参照

\noindent
{\it 余等化子}
\ref{categories-definition-coequalizers} を参照

\noindent
{\it 余フィルター付き}
\ref{categories-definition-codirected} を参照

\noindent
{\it 余フィルター付き}
\ref{categories-definition-codirected} を参照

\noindent
{\it 共終}
\ref{categories-definition-cofinal} を参照

\noindent
{\it Cohen 環}
\ref{algebra-definition-cohen-ring} を参照

\noindent
{\it 点 $x$ で Cohen--Macaulay}
\ref{more-morphisms-definition-CM} を参照

\noindent
{\it 点 $x$ で Cohen--Macaulay}
\ref{spaces-more-morphisms-definition-CM} を参照

\noindent
{\it Cohen--Macaulay 射}
\ref{more-morphisms-definition-CM} を参照

\noindent
{\it Cohen--Macaulay 射}
\ref{spaces-more-morphisms-definition-CM} を参照

\noindent
{\it Cohen--Macaulay}
\ref{algebra-definition-CM} を参照

\noindent
{\it Cohen--Macaulay}
\ref{algebra-definition-module-CM} を参照

\noindent
{\it Cohen--Macaulay}
\ref{algebra-definition-local-ring-CM} を参照

\noindent
{\it Cohen--Macaulay}
\ref{algebra-definition-ring-CM} を参照

\noindent
{\it Cohen--Macaulay}
\ref{properties-definition-Cohen-Macaulay} を参照

\noindent
{\it Cohen--Macaulay}
\ref{coherent-definition-Cohen-Macaulay} を参照

\noindent
{\it Cohen--Macaulay}
\ref{exercises-definition-CM} を参照

\noindent
{\it 連接 $\mathcal{O}_X$-加群の層}
\ref{modules-definition-coherent} を参照

\noindent
{\it コヒーレント加群}
\ref{algebra-definition-coherent} を参照

\noindent
{\it コヒーレント環}
\ref{algebra-definition-coherent} を参照

\noindent
{\it 連接}
\ref{sites-modules-definition-site-local} を参照

\noindent
{\it 連接}
\ref{spaces-cohomology-definition-coherent} を参照

\noindent
{\it 連接}
\ref{stacks-cohomology-definition-coherent} を参照

\noindent
{\it 連接}
\ref{exercises-definition-coherent} を参照

\noindent
{\it コホモロジー的 $\delta$ 関手}
\ref{homology-definition-cohomological-delta-functor} を参照

\noindent
{\it $f$ のコホモロジー次元}
\ref{etale-cohomology-definition-cd-f} を参照

\noindent
{\it $I$ の $A$ におけるコホモロジー次元}
\ref{local-cohomology-definition-cd} を参照

\noindent
{\it $X$ のコホモロジー次元}
\ref{etale-cohomology-definition-cd} を参照

\noindent
{\it コホモロジー的}
\ref{derived-definition-homological} を参照

\noindent
{\it コホモロジー加群}
\ref{chow-definition-periodic-complex} を参照

\noindent
{\it コホモロジー加群}
\ref{chow-definition-periodic-complex} を参照

\noindent
{\it コンパクト台をもつ $K$ のコホモロジー}
\ref{more-etale-definition-cohomology-compact-support} を参照

\noindent
{\it $f$ の余像}
\ref{homology-definition-kernel} を参照

\noindent
{\it 余核}
\ref{homology-definition-kernel} を参照

\noindent
{\it 余極限}
\ref{categories-definition-colimit} を参照

\noindent
{\it 余極限}
\ref{exercises-definition-colimit} を参照

\noindent
{\it 組合せ的に同値}
\ref{sites-definition-combinatorial-tautological} を参照

\noindent
{\it 可換}
\ref{dga-definition-cdga} を参照

\noindent
{\it コンパクト対象}
\ref{derived-definition-compact-object} を参照

\noindent
{\it コンパクト生成}
\ref{derived-definition-compactly-generated} を参照

\noindent
{\it $K$ のコンパクト台付きコホモロジー}
\ref{more-etale-definition-cohomology-compact-support} を参照

\noindent
{\it 微分次数付き構造と両立する}
\ref{dpa-definition-divided-powers-dga} を参照

\noindent
{\it 三角構造と両立する}
\ref{derived-definition-localization} を参照

\noindent
{\it $x$ における $\mathcal{F}/X/S$ の完全デヴィサージュ}
\ref{flat-definition-complete-devissage-at-x} を参照

\noindent
{\it $s$ 上の $\mathcal{F}/X/S$ の完全デヴィサージュ}
\ref{flat-definition-complete-devissage} を参照

\noindent
{\it $\mathfrak q$ における $N/S/R$ の完全デヴィサージュ}
\ref{flat-definition-complete-devissage-at-x-algebra} を参照

\noindent
{\it $\mathfrak r$ 上の $N/S/R$ の完全デヴィサージュ}
\ref{flat-definition-complete-devissage-algebra} を参照

\noindent
{\it $k$ 上完全交叉}
\ref{algebra-definition-lci-local-ring} を参照

\noindent
{\it 完全交叉}
\ref{dpa-definition-lci} を参照

\noindent
{\it 完備局所環}
\ref{algebra-definition-complete-local-ring} を参照

\noindent
{\it 完備化された主局所化}
\ref{restricted-definition-completed-principal-localization} を参照

\noindent
{\it 完備化テンソル積}
\ref{formal-spaces-definition-toplogy-tensor-product} を参照

\noindent
{\it 完全分解}
\ref{more-morphisms-definition-cd-morphism} を参照

\noindent
{\it 完全分解}
\ref{more-morphisms-definition-cd-morphism} を参照

\noindent
{\it 完整閉}
\ref{algebra-definition-almost-integral} を参照

\noindent
{\it $(U, R, s, t, c)$ の完備化 $(U, R, s, t, c)^{\wedge}$}
\ref{formal-defos-definition-completion-groupoid-in-functors} を参照

\noindent
{\it $\mathcal{F}$ の完備化}
\ref{formal-defos-definition-completion} を参照

\noindent
{\it $X$ の $T$ に沿う完備化}
\ref{formal-spaces-definition-completion} を参照

\noindent
{\it $X$ の $T$ に沿う完備化}
\ref{formal-spaces-definition-completion-formal-algebraic-space} を参照

\noindent
{\it $Z$ に沿う $X$ の完備化}
\ref{formal-spaces-definition-completion-subspace} を参照

\noindent
{\it 複体}
\ref{homology-definition-exact} を参照

\noindent
{\it 合成 $f\circ g$}
\ref{sites-definition-topos} を参照

\noindent
{\it $\varphi$ と $\psi$ の合成}
\ref{sheaves-definition-composition-f-maps} を参照

\noindent
{\it 芽の射の合成}
\ref{descent-definition-germs} を参照

\noindent
{\it 環付きサイトの射の合成}
\ref{sites-modules-definition-ringed-site} を参照

\noindent
{\it 環付き空間の射の合成}
\ref{sheaves-definition-composition-maps-ringed-spaces} を参照

\noindent
{\it 環付きトポスの射の合成}
\ref{sites-modules-definition-ringed-topos} を参照

\noindent
{\it 合成}
\ref{categories-definition-2-category} を参照

\noindent
{\it 合成}
\ref{sites-definition-composition-morphisms-sites} を参照

\noindent
{\it $\Omega$ における $K$ と $L$ の合成体}
\ref{fields-definition-compositum} を参照

\noindent
{\it 計算する}
\ref{derived-definition-computes} を参照

\noindent
{\it 計算する}
\ref{derived-definition-computes} を参照

\noindent
{\it 条件 (RS)}
\ref{formal-defos-definition-RS} を参照

\noindent
{\it 条件 (RS)}
\ref{artin-definition-RS} を参照

\noindent
{\it 条件 (RS*)}
\ref{artin-definition-RS-star} を参照

\noindent
{\it 条件 (S1) および (S2)}
\ref{formal-defos-definition-S1-S2} を参照

\noindent
{\it 錐 $\pi : C \to S$（$S$ 上）}
\ref{constructions-definition-abstract-cone} を参照

\noindent
{\it $\mathcal{A}$ に付随する錐}
\ref{constructions-definition-cone} を参照

\noindent
{\it 錐}
\ref{derived-definition-cone} を参照

\noindent
{\it 錐}
\ref{dga-definition-cone} を参照

\noindent
{\it 錐}
\ref{sdga-definition-cone} を参照

\noindent
{\it 連結成分}
\ref{topology-definition-connected-components} を参照

\noindent
{\it 連結成分}
\ref{exercises-definition-connected-component} を参照

\noindent
{\it 連結}
\ref{categories-definition-category-connected} を参照

\noindent
{\it 連結}
\ref{topology-definition-connected-components} を参照

\noindent
{\it 連結}
\ref{exercises-definition-connected-component} を参照

\noindent
{\it 余法代数 $\mathcal{C}_{Z/X, *}$、すなわち $Z$ の $X$ における余法代数}
\ref{divisors-definition-conormal-sheaf} を参照

\noindent
{\it 余法代数 $\mathcal{C}_{Z/X, *}$、すなわち $Z$ の $X$ における余法代数}
\ref{spaces-more-morphisms-definition-conormal-algebra} を参照

\noindent
{\it $f$ の余法代数}
\ref{divisors-definition-conormal-sheaf} を参照

\noindent
{\it $i$ の余法代数}
\ref{spaces-more-morphisms-definition-conormal-algebra} を参照

\noindent
{\it 余法加群}
\ref{algebra-definition-universal-thickening} を参照

\noindent
{\it $Z$ の $X$ における余法層 $\mathcal{C}_{Z/X}$}
\ref{morphisms-definition-conormal-sheaf} を参照

\noindent
{\it $Z$ の $X$ における余法層 $\mathcal{C}_{Z/X}$}
\ref{spaces-more-morphisms-definition-conormal-sheaf} を参照

\noindent
{\it $i$ の余法層}
\ref{morphisms-definition-conormal-sheaf} を参照

\noindent
{\it $i$ の余法層}
\ref{spaces-more-morphisms-definition-conormal-sheaf} を参照

\noindent
{\it $Z$ の $X$ 上の余法層}
\ref{more-morphisms-definition-universal-thickening} を参照

\noindent
{\it $Z$ の $X$ 上の余法層}
\ref{spaces-more-morphisms-definition-universal-thickening} を参照

\noindent
{\it 保守的}
\ref{sites-definition-enough-points} を参照

\noindent
{\it 値 $A$ の定数前層}
\ref{sheaves-definition-constant-presheaf} を参照

\noindent
{\it 値 $A$ の定数層}
\ref{sheaves-definition-constant-sheaf} を参照

\noindent
{\it 値 $A$ の定数層}
\ref{etale-cohomology-definition-finite-locally-constant} を参照

\noindent
{\it 値 $E$ をもつ定数層}
\ref{etale-cohomology-definition-finite-locally-constant} を参照

\noindent
{\it 値 $M$ をもつ定数層}
\ref{etale-cohomology-definition-finite-locally-constant} を参照

\noindent
{\it 定数層}
\ref{sites-modules-definition-locally-constant} を参照

\noindent
{\it 定数層}
\ref{etale-cohomology-definition-additive-sheaf} を参照

\noindent
{\it 定数層}
\ref{etale-cohomology-definition-finite-locally-constant} を参照

\noindent
{\it 定数層}
\ref{etale-cohomology-definition-finite-locally-constant} を参照

\noindent
{\it 定数層}
\ref{etale-cohomology-definition-finite-locally-constant} を参照

\noindent
{\it 構成可能 $\Lambda$-層}
\ref{proetale-definition-adic} を参照

\noindent
{\it 構成可能}
\ref{topology-definition-constructible} を参照

\noindent
{\it 構成可能}
\ref{etale-cohomology-definition-constructible} を参照

\noindent
{\it 構成可能}
\ref{etale-cohomology-definition-constructible} を参照

\noindent
{\it 構成可能}
\ref{etale-cohomology-definition-constructible} を参照

\noindent
{\it 構成可能}
\ref{proetale-definition-constructible} を参照

\noindent
{\it 構成可能}
\ref{proetale-definition-Dbc} を参照

\noindent
{\it $x$ の内容イデアル}
\ref{more-algebra-definition-content-ideal} を参照

\noindent
{\it 連続群コホモロジー群}
\ref{etale-cohomology-definition-galois-cohomology} を参照

\noindent
{\it 連続}
\ref{sites-definition-continuous} を参照

\noindent
{\it 反変}
\ref{categories-definition-contravariant} を参照

\noindent
{\it $H^*(K^\bullet)$ に収束する}
\ref{homology-definition-filtered-complex-ss-converges} を参照

\noindent
{\it $H^n(\text{Tot}(K^{\bullet, \bullet}))$ に収束する}
\ref{homology-definition-ss-double-complex-converge} を参照

\noindent
{\it $H^n(\text{Tot}(K^{\bullet, \bullet}))$ に収束する}
\ref{homology-definition-ss-double-complex-converge} を参照

\noindent
{\it 余積}
\ref{categories-definition-coproducts} を参照

\noindent
{\it 余積}
\ref{categories-definition-coproduct} を参照

\noindent
{\it 余正則}
\ref{homology-definition-bounded-ss} を参照

\noindent
{\it 余単体アーベル群}
\ref{simplicial-definition-cosimplicial-object} を参照

\noindent
{\it $\mathcal{C}$ の余単体的対象 $U$}
\ref{simplicial-definition-cosimplicial-object} を参照

\noindent
{\it 余単体集合}
\ref{simplicial-definition-cosimplicial-object} を参照

\noindent
{\it $Y$ 上の $X$ の余接複体 $L_{X/Y}$}
\ref{cotangent-definition-cotangent-morphism-schemes} を参照

\noindent
{\it $Y$ 上の $X$ の余接複体 $L_{X/Y}$}
\ref{cotangent-definition-cotangent-morphism-spaces} を参照

\noindent
{\it 余接複体}
\ref{cotangent-definition-cotangent-complex-ring-map} を参照

\noindent
{\it 余接複体}
\ref{cotangent-definition-cotangent-complex-morphism-sheaves-rings} を参照

\noindent
{\it 余接複体}
\ref{cotangent-definition-cotangent-complex-morphism-ringed-spaces} を参照

\noindent
{\it 余接複体}
\ref{cotangent-definition-cotangent-complex-morphism-ringed-topoi} を参照

\noindent
{\it 可算添字付き}
\ref{formal-spaces-definition-countable} を参照

\noindent
{\it $\mathcal{C}$ の被覆}
\ref{sites-definition-site} を参照

\noindent
{\it 被覆族}
\ref{etale-cohomology-definition-site} を参照

\noindent
{\it 被覆}
\ref{hypercovering-definition-covering-SR} を参照

\noindent
{\it 被覆}
\ref{hypercovering-definition-covering-SR} を参照

\noindent
{\it $F$ を被覆する}
\ref{schemes-definition-representable-by-open-immersions} を参照

\noindent
{\it $\mathcal{O}_{X/S}$-加群のクリスタル}
\ref{crystalline-definition-modules} を参照

\noindent
{\it 有限局所自由加群のクリスタル}
\ref{crystalline-definition-crystal-quasi-coherent-modules} を参照

\noindent
{\it 準連接加群のクリスタル}
\ref{crystalline-definition-crystal-quasi-coherent-modules} を参照

\noindent
{\it 結晶サイト}
\ref{crystalline-definition-crystalline-site} を参照

\noindent
{\it 曲線}
\ref{varieties-definition-curve} を参照

\noindent
{\it 曲線}
\ref{etale-cohomology-definition-variety} を参照

\noindent
{\it $X$ 上のサイクル}
\ref{chow-definition-cycles} を参照

\noindent
{\it $X$ 上のサイクル}
\ref{spaces-chow-definition-cycles} を参照

\noindent
{\it $\mathcal{A}$ 上の $\mathcal{B}$ の de Rham 複体}
\ref{modules-definition-de-rham-complex} を参照

\noindent
{\it $Y \subset X$ に対する $S$ 上の対数極ド・ラーム複体}
\ref{derham-definition-log-complex} を参照

\noindent
{\it $Y \subset X$ に対する $S$ 上の対数極ド・ラーム複体が定義される}
\ref{derham-definition-local-product} を参照

\noindent
{\it de Rham 複体}
\ref{modules-definition-de-rham-complex-morphism-ringed-spaces} を参照

\noindent
{\it 良好}
\ref{decent-spaces-definition-very-reasonable} を参照

\noindent
{\it 良好}
\ref{decent-spaces-definition-relative-conditions} を参照

\noindent
{\it 良好}
\ref{stacks-morphisms-definition-decent} を参照

\noindent
{\it $\mathfrak m$ の分解群}
\ref{more-algebra-definition-decomposition-inertia} を参照

\noindent
{\it 減少フィルトレーション}
\ref{homology-definition-filtered} を参照

\noindent
{\it Dedekind 整域}
\ref{algebra-definition-dedekind-domain} を参照

\noindent
{\it 点 $x \in X$ で定義される}
\ref{morphisms-definition-domain-of-definition} を参照

\noindent
{\it 点 $x \in |X|$ において定義される}
\ref{spaces-morphisms-definition-domain-of-definition} を参照

\noindent
{\it 節点特異点を定める}
\ref{curves-definition-multicross} を参照

\noindent
{\it 節点特異点を定める}
\ref{curves-definition-nodal} を参照

\noindent
{\it 有理特異点を定める}
\ref{resolve-definition-reduce-to-rational} を参照

\noindent
{\it 変形圏}
\ref{formal-defos-definition-deformation-category} を参照

\noindent
{\it $x$ の退化}
\ref{simplicial-definition-terminology-simplicial-sets} を参照

\noindent
{\it $E_r$ で退化する}
\ref{homology-definition-limit-spectral-sequence} を参照

\noindent
{\it 退化している}
\ref{simplicial-definition-terminology-simplicial-sets} を参照

\noindent
{\it $\mathcal{Y}$ 上の $\mathcal{X}$ の次数 $d$ の有限 Hilbert スタック}
\ref{examples-stacks-definition-hilbert-d-stack} を参照

\noindent
{\it $X$ の $Y$ 上の次数}
\ref{morphisms-definition-degree} を参照

\noindent
{\it $X$ の $Y$ 上の次数}
\ref{spaces-over-fields-definition-degree} を参照

\noindent
{\it $\mathcal{L}$ に関する $Z$ の次数}
\ref{varieties-definition-degree} を参照

\noindent
{\it 零サイクルの次数}
\ref{chow-definition-degree-zero-cycle} を参照

\noindent
{\it 零サイクルの次数}
\ref{spaces-chow-definition-degree-zero-cycle} を参照

\noindent
{\it 非分離度}
\ref{fields-definition-insep-degree} を参照

\noindent
{\it 次数}
\ref{fields-definition-degree} を参照

\noindent
{\it 次数}
\ref{morphisms-definition-finite-locally-free} を参照

\noindent
{\it 次数}
\ref{varieties-definition-degree-invertible-sheaf} を参照

\noindent
{\it 次数}
\ref{varieties-definition-degree-invertible-sheaf} を参照

\noindent
{\it 次数}
\ref{spaces-morphisms-definition-finite-locally-free} を参照

\noindent
{\it Deligne--Mumford スタック}
\ref{algebraic-definition-deligne-mumford} を参照

\noindent
{\it 点における深さ $k$}
\ref{coherent-definition-depth} を参照

\noindent
{\it 点における深さ $k$}
\ref{coherent-definition-depth} を参照

\noindent
{\it 深さ}
\ref{algebra-definition-depth} を参照

\noindent
{\it 深さ}
\ref{local-cohomology-definition-depth-complex} を参照

\noindent
{\it 導分}
\ref{algebra-definition-derivation} を参照

\noindent
{\it $(\mathcal{A}, \text{d})$ の導来圏}
\ref{sdga-definition-derived-category} を参照

\noindent
{\it $(A, \text{d})$ の導来圏}
\ref{dga-definition-unbounded-derived-category} を参照

\noindent
{\it $\mathcal{A}$ の導来圏}
\ref{derived-definition-unbounded-derived-category} を参照

\noindent
{\it 準連接なコホモロジー層をもつ $\mathcal{O}_\mathcal{X}$-加群の導来圏}
\ref{stacks-perfect-definition-derived} を参照

\noindent
{\it 準連接なコホモロジー層をもつ $\mathcal{O}_X$-加群の導来圏}
\ref{spaces-perfect-definition-derived-quasi-coherent} を参照

\noindent
{\it 導来余極限}
\ref{derived-definition-derived-colimit} を参照

\noindent
{\it $\mathcal{I}$ に関して導来完備}
\ref{algebraization-definition-derived-complete} を参照

\noindent
{\it $I$ に関して導来完備}
\ref{more-algebra-definition-derived-complete} を参照

\noindent
{\it $I$ に関して導来完備}
\ref{more-algebra-definition-derived-complete} を参照

\noindent
{\it 導来同値}
\ref{equiv-definition-derived-equivalent} を参照

\noindent
{\it 導来内部 Hom}
\ref{sdga-definition-pushforward} を参照

\noindent
{\it 導来極限}
\ref{derived-definition-derived-limit} を参照

\noindent
{\it 導来引き戻し}
\ref{sdga-definition-pullback} を参照

\noindent
{\it 導来順像}
\ref{sdga-definition-pushforward} を参照

\noindent
{\it 導来テンソル積}
\ref{more-algebra-definition-derived-tor} を参照

\noindent
{\it 導来テンソル積}
\ref{cohomology-definition-derived-tor} を参照

\noindent
{\it 導来テンソル積}
\ref{sites-cohomology-definition-derived-tor} を参照

\noindent
{\it 導来テンソル積}
\ref{sdga-definition-pullback} を参照

\noindent
{\it 準連接層の降下データ $(\mathcal{F}_i, \varphi_{ij})$}
\ref{descent-definition-descent-datum-quasi-coherent} を参照

\noindent
{\it 準連接層の降下データ $(\mathcal{F}_i, \varphi_{ij})$}
\ref{spaces-descent-definition-descent-datum-quasi-coherent} を参照

\noindent
{\it $R \to A$ に関する加群の降下データ $(N, \varphi)$}
\ref{descent-definition-descent-datum-modules} を参照

\noindent
{\it 族 $\{X_i \to S\}$ に関する降下データ $(V_i, \varphi_{ij})$}
\ref{descent-definition-descent-datum-for-family-of-morphisms} を参照

\noindent
{\it 族 $\{X_i \to X\}$ に関する降下データ $(V_i, \varphi_{ij})$}
\ref{spaces-descent-definition-descent-datum-for-family-of-morphisms} を参照

\noindent
{\it 族 $\{f_i : U_i \to U\}$ に関する $\mathcal{S}$ の降下データ $(X_i, \varphi_{ij})$}
\ref{stacks-definition-descent-data} を参照

\noindent
{\it $V/X/S$ の降下データ}
\ref{descent-definition-descent-datum} を参照

\noindent
{\it $V/Y/X$ に対する降下データ}
\ref{spaces-descent-definition-descent-datum} を参照

\noindent
{\it $X \to S$ に関する降下データ}
\ref{descent-definition-descent-datum} を参照

\noindent
{\it $Y \to X$ に関する降下データ}
\ref{spaces-descent-definition-descent-datum} を参照

\noindent
{\it 降下データ}
\ref{etale-cohomology-definition-descent-datum} を参照

\noindent
{\it 降下データ}
\ref{etale-cohomology-definition-descent-datum-modules} を参照

\noindent
{\it 加群の降下射}
\ref{descent-definition-effective-descent} を参照

\noindent
{\it $(M, \varphi, \psi)$ の行列式}
\ref{chow-definition-periodic-determinant} を参照

\noindent
{\it 有限長 $R$-加群 $M$ の行列式}
\ref{chow-definition-determinant} を参照

\noindent
{\it $\varphi$ の微分 $d \varphi : T \mathcal{F} \to T \mathcal{G}$}
\ref{formal-defos-definition-differential} を参照

\noindent
{\it 微分次数付き $(\mathcal{A}, \mathcal{B})$-双加群}
\ref{sdga-definition-dg-bimodule} を参照

\noindent
{\it 微分次数付き $(A, B)$-両側加群}
\ref{dga-definition-bimodule} を参照

\noindent
{\it 微分次数付き $\mathcal{A}$-加群}
\ref{sdga-definition-dgm} を参照

\noindent
{\it $R$ 上の微分次数付き代数}
\ref{dga-definition-dga} を参照

\noindent
{\it $R$ 上の微分次数付き圏 $\mathcal{A}$}
\ref{dga-definition-dga-category} を参照

\noindent
{\it 微分次数付き直和}
\ref{dga-definition-dg-direct-sum} を参照

\noindent
{\it 微分次数付き加群}
\ref{dga-definition-dgm} を参照

\noindent
{\it 微分次数付き加群}
\ref{sdga-definition-dgm} を参照

\noindent
{\it 微分対象}
\ref{homology-definition-differential-object} を参照

\noindent
{\it $k$ 階微分作用素 $D : \mathcal{F} \to \mathcal{G}$}
\ref{modules-definition-differential-operators} を参照

\noindent
{\it $k$ 階微分作用素 $D : \mathcal{F} \to \mathcal{G}$}
\ref{sites-modules-definition-differential-operators} を参照

\noindent
{\it $k$ 階の微分作用素 $D : M \to N$}
\ref{algebra-definition-differential-operators} を参照

\noindent
{\it $X/S$ 上の $k$ 階微分作用素}
\ref{modules-definition-relative-differential-operators} を参照

\noindent
{\it 差イデアル}
\ref{discriminant-definition-different} を参照

\noindent
{\it 次元関数}
\ref{topology-definition-dimension-function} を参照

\noindent
{\it $x$ における $\mathcal{X}$ の次元}
\ref{stacks-properties-definition-dimension-at-point} を参照

\noindent
{\it $x$ における $X$ の次元}
\ref{properties-definition-dimension} を参照

\noindent
{\it $x$ における $X$ の次元}
\ref{spaces-properties-definition-dimension-at-point} を参照

\noindent
{\it $x$ における $\mathcal{X}$ の局所環の次元}
\ref{stacks-geometry-definition-dimension-local-ring} を参照

\noindent
{\it $x$ における $X$ の局所環の次元}
\ref{spaces-properties-definition-dimension-local-ring} を参照

\noindent
{\it $x$ における $f$ のファイバーの局所環の次元}
\ref{spaces-morphisms-definition-dimension-fibre} を参照

\noindent
{\it 次元}
\ref{topology-definition-Krull} を参照

\noindent
{\it 次元}
\ref{properties-definition-dimension} を参照

\noindent
{\it 次元}
\ref{spaces-properties-definition-dimension} を参照

\noindent
{\it 次元}
\ref{stacks-properties-definition-dimension} を参照

\noindent
{\it 順像関手}
\ref{sites-definition-localize} を参照

\noindent
{\it 順像関手}
\ref{sites-modules-definition-localize-ringed-site} を参照

\noindent
{\it コンパクト台付き直像}
\ref{more-etale-definition-f-shriek-separated} を参照

\noindent
{\it コンパクト台付き直像}
\ref{more-etale-definition-f-shriek-lqf} を参照

\noindent
{\it 順像}
\ref{etale-cohomology-definition-direct-image-presheaf} を参照

\noindent
{\it 順像}
\ref{etale-cohomology-definition-direct-image-sheaf} を参照

\noindent
{\it 直和デヴィサージュ}
\ref{algebra-definition-devissage} を参照

\noindent
{\it 直和}
\ref{homology-definition-direct-sum} を参照

\noindent
{\it 有向逆系}
\ref{categories-definition-directed-system} を参照

\noindent
{\it 有向半順序集合}
\ref{categories-definition-directed-set} を参照

\noindent
{\it 有向集合}
\ref{categories-definition-directed-set} を参照

\noindent
{\it 有向集合}
\ref{exercises-definition-directed-poset} を参照

\noindent
{\it 有向系}
\ref{categories-definition-directed-system} を参照

\noindent
{\it 有向系}
\ref{algebra-definition-directed-system} を参照

\noindent
{\it 有向}
\ref{categories-definition-directed} を参照

\noindent
{\it 有向}
\ref{categories-definition-directed} を参照

\noindent
{\it 離散 $G$-加群}
\ref{etale-cohomology-definition-G-module-continuous} を参照

\noindent
{\it 離散 $G$-集合}
\ref{pione-definition-G-set-continuous} を参照

\noindent
{\it 離散付値環}
\ref{algebra-definition-value-group} を参照

\noindent
{\it 離散的}
\ref{categories-definition-discrete} を参照

\noindent
{\it $L/K$ の判別式}
\ref{fields-definition-discriminant} を参照

\noindent
{\it $M$ と $M'$ の距離}
\ref{algebra-definition-distance} を参照

\noindent
{\it $K(\mathcal{A})$ の完全三角}
\ref{derived-definition-distinguished-triangle} を参照

\noindent
{\it 完全三角}
\ref{derived-definition-triangulated-category} を参照

\noindent
{\it 識別三角形}
\ref{dga-definition-distinguished-triangle} を参照

\noindent
{\it 分割冪 $A$-導分}
\ref{crystalline-definition-derivation} を参照

\noindent
{\it $(A, I, \gamma)$ に相対的な $B$ における $J$ の分割冪包絡}
\ref{crystalline-definition-divided-power-envelope} を参照

\noindent
{\it 分割冪環}
\ref{dpa-definition-divided-power-ring} を参照

\noindent
{\it 分割冪スキーム}
\ref{crystalline-definition-divided-power-scheme} を参照

\noindent
{\it 分割冪構造 $\gamma$}
\ref{crystalline-definition-divided-power-structure} を参照

\noindent
{\it 分割冪構造}
\ref{dpa-definition-divided-powers} を参照

\noindent
{\it 分割冪構造}
\ref{dpa-definition-divided-powers-graded} を参照

\noindent
{\it $(S, \mathcal{I}, \gamma)$ に相対的な $X$ の分割冪肥厚化}
\ref{crystalline-definition-divided-power-thickening-X} を参照

\noindent
{\it 分割冪肥厚化}
\ref{crystalline-definition-affine-thickening} を参照

\noindent
{\it 分割冪肥厚化}
\ref{crystalline-definition-divided-power-thickening} を参照

\noindent
{\it $S$ 上 DM}
\ref{stacks-morphisms-definition-absolute-separated} を参照

\noindent
{\it DM}
\ref{stacks-morphisms-definition-separated} を参照

\noindent
{\it DM}
\ref{stacks-morphisms-definition-absolute-separated} を参照

\noindent
{\it 定義域}
\ref{morphisms-definition-domain-of-definition} を参照

\noindent
{\it 定義域}
\ref{spaces-morphisms-definition-domain-of-definition} を参照

\noindent
{\it 整域}
\ref{fields-definition-domain} を参照

\noindent
{\it 支配的}
\ref{morphisms-definition-dominant} を参照

\noindent
{\it 支配的}
\ref{morphisms-definition-dominant-rational} を参照

\noindent
{\it 支配的}
\ref{spaces-morphisms-definition-dominant} を参照

\noindent
{\it 支配的}
\ref{spaces-morphisms-definition-dominant-rational} を参照

\noindent
{\it 支配する}
\ref{algebra-definition-valuation-ring} を参照

\noindent
{\it 支配する}
\ref{algebra-definition-domination} を参照

\noindent
{\it 点線矢印}
\ref{stacks-morphisms-definition-fill-in-diagram} を参照

\noindent
{\it 二重複体}
\ref{homology-definition-double-complex} を参照

\noindent
{\it 双対数}
\ref{varieties-definition-dual-numbers} を参照

\noindent
{\it 双対数}
\ref{exercises-definition-dual-numbers} を参照

\noindent
{\it $\omega_S^\bullet$ に関して正規化された双対化複体}
\ref{duality-definition-good-dualizing} を参照

\noindent
{\it 双対化複体}
\ref{dualizing-definition-dualizing} を参照

\noindent
{\it 双対化複体}
\ref{duality-definition-dualizing-scheme} を参照

\noindent
{\it 双対化複体}
\ref{spaces-duality-definition-dualizing-scheme} を参照

\noindent
{\it 有効 Cartier 因子}
\ref{divisors-definition-effective-Cartier-divisor} を参照

\noindent
{\it 有効 Cartier 因子}
\ref{spaces-divisors-definition-effective-Cartier-divisor} を参照

\noindent
{\it 有効 Cartier 因子}
\ref{exercises-definition-divisor} を参照

\noindent
{\it 加群の有効降下射}
\ref{descent-definition-effective-descent} を参照

\noindent
{\it 有効エピ族}
\ref{sites-definition-universal-effective-epimorphisms} を参照

\noindent
{\it 有効}
\ref{stacks-definition-effective-descent-datum} を参照

\noindent
{\it 有効}
\ref{descent-definition-descent-datum-effective-quasi-coherent} を参照

\noindent
{\it 有効}
\ref{descent-definition-descent-datum-effective-module} を参照

\noindent
{\it 有効}
\ref{descent-definition-effective} を参照

\noindent
{\it 有効}
\ref{descent-definition-effective-family} を参照

\noindent
{\it 有効}
\ref{chow-definition-effective-cycle} を参照

\noindent
{\it 有効}
\ref{etale-cohomology-definition-descent-datum} を参照

\noindent
{\it 有効}
\ref{etale-cohomology-definition-effective-modules} を参照

\noindent
{\it 有効}
\ref{relative-cycles-definition-effective} を参照

\noindent
{\it 有効}
\ref{spaces-descent-definition-descent-datum-effective-quasi-coherent} を参照

\noindent
{\it 有効}
\ref{spaces-descent-definition-effective} を参照

\noindent
{\it 有効}
\ref{spaces-descent-definition-effective-family} を参照

\noindent
{\it 有効}
\ref{artin-definition-effective} を参照

\noindent
{\it Eilenberg--Mac Lane 対象 $K(A, k)$}
\ref{simplicial-definition-eilenberg-maclane} を参照

\noindent
{\it $\mathfrak q$ における環準同型 $R \to S$ の基本エタール局所化}
\ref{flat-definition-elementary-etale-neighbourhood} を参照

\noindent
{\it 初等エタール近傍}
\ref{more-morphisms-definition-etale-neighbourhood} を参照

\noindent
{\it 初等エタール近傍}
\ref{decent-spaces-definition-elemenary-etale-neighbourhood} を参照

\noindent
{\it 初等識別正方形}
\ref{spaces-perfect-definition-elementary-distinguished-square} を参照

\noindent
{\it 単因子整域}
\ref{more-algebra-definition-bezout} を参照

\noindent
{\it $R$ 上の $A$ における初等標準元}
\ref{smoothing-definition-strictly-standard} を参照

\noindent
{\it 埋め込まれた随伴点}
\ref{divisors-definition-embedded} を参照

\noindent
{\it 埋没随伴素イデアル}
\ref{algebra-definition-embedded-primes} を参照

\noindent
{\it 埋め込み成分}
\ref{divisors-definition-embedded} を参照

\noindent
{\it 埋め込み点}
\ref{divisors-definition-embedded} を参照

\noindent
{\it $R$ の埋没素イデアル}
\ref{algebra-definition-embedded-primes} を参照

\noindent
{\it $x$ における $X$ の埋め込み次元}
\ref{varieties-definition-embed-dim} を参照

\noindent
{\it $x$ における $X/k$ の埋め込み次元}
\ref{varieties-definition-embedding-dimension} を参照

\noindent
{\it 埋め込み}
\ref{sites-definition-embedding} を参照

\noindent
{\it 十分多くの $P$ 対象をもつ}
\ref{sites-definition-w-contractible} を参照

\noindent
{\it 十分な入射対象をもつ}
\ref{homology-definition-enough-injectives} を参照

\noindent
{\it 十分な射影対象をもつ}
\ref{homology-definition-enough-projectives} を参照

\noindent
{\it 十分多くの弱可縮対象をもつ}
\ref{sites-definition-w-contractible} を参照

\noindent
{\it 包絡射}
\ref{chow-definition-envelope} を参照

\noindent
{\it エピ射}
\ref{categories-definition-mono-epi} を参照

\noindent
{\it 等化子}
\ref{categories-definition-equalizers} を参照

\noindent
{\it 等次元}
\ref{topology-definition-equidimensional} を参照

\noindent
{\it 等次元}
\ref{relative-cycles-definition-equidimensional} を参照

\noindent
{\it 圏同値}
\ref{categories-definition-equivalence-categories} を参照

\noindent
{\it $B$ 上の $U$ の同値関係}
\ref{spaces-groupoids-definition-equivalence-relation} を参照

\noindent
{\it $S$ 上の $U$ の同値関係}
\ref{groupoids-definition-equivalence-relation} を参照

\noindent
{\it 同値な型}
\ref{models-definition-type-equivalent} を参照

\noindent
{\it 同値}
\ref{categories-definition-equivalence} を参照

\noindent
{\it 同値}
\ref{derived-definition-yoneda-extension} を参照

\noindent
{\it 同値}
\ref{morphisms-definition-rational-map} を参照

\noindent
{\it 同値}
\ref{etale-cohomology-definition-brauer-equivalent} を参照

\noindent
{\it 同値}
\ref{spaces-morphisms-definition-rational-map} を参照

\noindent
{\it 同変準連接 $\mathcal{O}_X$-加群}
\ref{groupoids-definition-equivariant-module} を参照

\noindent
{\it 同変準連接 $\mathcal{O}_X$-加群}
\ref{spaces-groupoids-definition-equivariant-module} を参照

\noindent
{\it 同変}
\ref{groupoids-definition-action-group-scheme} を参照

\noindent
{\it 同変}
\ref{spaces-groupoids-definition-action-group-space} を参照

\noindent
{\it 本質拡大}
\ref{dualizing-definition-essential} を参照

\noindent
{\it 本質的全射}
\ref{formal-defos-definition-essential-surjection} を参照

\noindent
{\it 本質的に定値な逆系}
\ref{categories-definition-essentially-constant-system} を参照

\noindent
{\it 本質的に定値な系}
\ref{categories-definition-essentially-constant-system} を参照

\noindent
{\it 値}
\ref{categories-definition-essentially-constant-diagram} を参照

\noindent
{\it 値}
\ref{categories-definition-essentially-constant-diagram} を参照

\noindent
{\it 本質的に有限表示}
\ref{algebra-definition-essentially-finite-p-t} を参照

\noindent
{\it 本質的に有限型}
\ref{algebra-definition-essentially-finite-p-t} を参照

\noindent
{\it 本質的全射}
\ref{categories-definition-faithful} を参照

\noindent
{\it 本質的}
\ref{dualizing-definition-essential} を参照

\noindent
{\it 本質的}
\ref{dualizing-definition-essential} を参照

\noindent
{\it $\mathcal{F}$ の Euler 標数}
\ref{varieties-definition-euler-characteristic} を参照

\noindent
{\it $\mathcal{F}$ の Euler 標数}
\ref{spaces-over-fields-definition-euler-characteristic} を参照

\noindent
{\it Euler--Poincar\'e 函数}
\ref{exercises-definition-graded-module} を参照

\noindent
{\it 至る所で定義される}
\ref{derived-definition-everywhere-defined} を参照

\noindent
{\it 至る所で定義される}
\ref{derived-definition-everywhere-defined} を参照

\noindent
{\it $x_i$ で完全}
\ref{homology-definition-exact} を参照

\noindent
{\it $y$ で完全}
\ref{homology-definition-exact} を参照

\noindent
{\it 完全複体}
\ref{homology-definition-exact} を参照

\noindent
{\it 完全対}
\ref{homology-definition-exact-couple} を参照

\noindent
{\it 完全関手}
\ref{derived-definition-exact-functor-triangulated-categories} を参照

\noindent
{\it 次数付き加群の完全列}
\ref{exercises-definition-graded-algebra} を参照

\noindent
{\it 完全列}
\ref{homology-definition-exact} を参照

\noindent
{\it 完全}
\ref{categories-definition-exact} を参照

\noindent
{\it 完全}
\ref{homology-definition-exact} を参照

\noindent
{\it 完全}
\ref{chow-definition-periodic-complex} を参照

\noindent
{\it 優秀}
\ref{more-algebra-definition-excellent} を参照

\noindent
{\it 優秀}
\ref{morphisms-definition-excellent} を参照

\noindent
{\it 例外因子}
\ref{divisors-definition-blow-up} を参照

\noindent
{\it 例外因子}
\ref{spaces-divisors-definition-blow-up} を参照

\noindent
{\it 網羅的}
\ref{homology-definition-filtered} を参照

\noindent
{\it 付値判定法の存在性部分}
\ref{stacks-morphisms-definition-existence} を参照

\noindent
{\it 延長される}
\ref{dpa-definition-extends} を参照

\noindent
{\it $A$ による $B$ の拡大 $E$}
\ref{homology-definition-extension} を参照

\noindent
{\it $\mathcal{F}$ の $0$ による延長 $j_!\mathcal{F}$}
\ref{sheaves-definition-j-shriek-structures} を参照

\noindent
{\it $\mathcal{F}$ の $e$ による延長 $j_!\mathcal{F}$}
\ref{sheaves-definition-j-shriek-structures} を参照

\noindent
{\it $\mathcal{F}$ の $0$ による延長 $j_{p!}\mathcal{F}$}
\ref{sheaves-definition-j-shriek-structures} を参照

\noindent
{\it $\mathcal{F}$ の $e$ による延長 $j_{p!}\mathcal{F}$}
\ref{sheaves-definition-j-shriek-structures} を参照

\noindent
{\it $0$ による延長}
\ref{sheaves-definition-j-shriek-structures} を参照

\noindent
{\it $0$ による延長}
\ref{sheaves-definition-j-shriek-structures} を参照

\noindent
{\it 零による延長}
\ref{sites-modules-definition-localize-ringed-site} を参照

\noindent
{\it 零による延長}
\ref{etale-cohomology-definition-extension-zero} を参照

\noindent
{\it 零による延長}
\ref{etale-cohomology-definition-extension-zero} を参照

\noindent
{\it 零による延長}
\ref{proetale-definition-extension-zero} を参照

\noindent
{\it 零による延長}
\ref{proetale-definition-extension-zero} を参照

\noindent
{\it $\mathcal{F}$ の空集合による延長 $j_!\mathcal{F}$}
\ref{sheaves-definition-j-shriek} を参照

\noindent
{\it $\mathcal{F}$ の空集合による延長 $j_{p!}\mathcal{F}$}
\ref{sheaves-definition-j-shriek} を参照

\noindent
{\it 空集合による $\mathcal{G}$ の延長}
\ref{sites-definition-localize} を参照

\noindent
{\it 離散付値環の拡大}
\ref{more-algebra-definition-extension-discrete-valuation-rings} を参照

\noindent
{\it 付値環の拡大}
\ref{more-algebra-definition-extension-valuation-rings} を参照

\noindent
{\it 極値非連結}
\ref{topology-definition-extremally-disconnected} を参照

\noindent
{\it $x$ の面}
\ref{simplicial-definition-terminology-simplicial-sets} を参照

\noindent
{\it 忠実平坦}
\ref{algebra-definition-flat} を参照

\noindent
{\it 忠実平坦}
\ref{algebra-definition-flat} を参照

\noindent
{\it 忠実平坦}
\ref{etale-definition-flat-rings} を参照

\noindent
{\it 忠実平坦}
\ref{etale-definition-flat-schemes} を参照

\noindent
{\it 忠実}
\ref{categories-definition-faithful} を参照

\noindent
{\it 終域を固定した射の族}
\ref{sites-definition-family-morphisms-fixed-target} を参照

\noindent
{\it 終域を固定した射の族}
\ref{etale-cohomology-definition-family-morphisms-fixed-target} を参照

\noindent
{\it ファイバー圏}
\ref{categories-definition-fibre-category} を参照

\noindent
{\it ファイバー（$f$ の $s$ 上のもの）}
\ref{schemes-definition-fibre} を参照

\noindent
{\it $U$ 上の $V$ と $W$ のファイバー積}
\ref{simplicial-definition-fibre-product} を参照

\noindent
{\it $U$ 上の $V$ と $W$ のファイバー積}
\ref{simplicial-definition-fibre-product-cosimplicial-objects} を参照

\noindent
{\it ファイバー積}
\ref{categories-definition-fibre-products} を参照

\noindent
{\it ファイバー積}
\ref{schemes-definition-fibre-product} を参照

\noindent
{\it $\mathcal{C}$ 上のファイバー圏}
\ref{categories-definition-fibred-category} を参照

\noindent
{\it グルーポイドにファイバー化されている}
\ref{categories-definition-fibred-groupoids} を参照

\noindent
{\it $f$ のファイバーは普遍的に有界である}
\ref{morphisms-definition-universally-bounded} を参照

\noindent
{\it $f$ のファイバーは普遍的に有界である}
\ref{decent-spaces-definition-universally-bounded} を参照

\noindent
{\it 拡大体}
\ref{fields-definition-extension} を参照

\noindent
{\it 有理関数体}
\ref{morphisms-definition-function-field} を参照

\noindent
{\it 有理関数体}
\ref{spaces-over-fields-definition-function-field} を参照

\noindent
{\it 体}
\ref{fields-definition-field} を参照

\noindent
{\it フィルター付き非輪状}
\ref{derived-definition-filtered-acyclic} を参照

\noindent
{\it フィルター付き非輪状}
\ref{exercises-definition-filtered-acyclic} を参照

\noindent
{\it $\mathcal{A}$ のフィルター付き複体 $K^\bullet$}
\ref{homology-definition-filtered-complex} を参照

\noindent
{\it $\mathcal{A}$ のフィルトレーション付き導来圏}
\ref{derived-definition-filtered-derived} を参照

\noindent
{\it フィルター付き導来関手}
\ref{trace-definition-filtered-derived-functors} を参照

\noindent
{\it フィルトレーション付き微分対象}
\ref{homology-definition-filtered-differential} を参照

\noindent
{\it フィルター付き入射的}
\ref{derived-definition-filtered-complexes-notation} を参照

\noindent
{\it フィルター付き入射的}
\ref{trace-definition-filtered} を参照

\noindent
{\it フィルター付き入射的}
\ref{exercises-definition-injective-filtered} を参照

\noindent
{\it $\mathcal{A}$ のフィルター付き対象}
\ref{homology-definition-filtered} を参照

\noindent
{\it フィルター付き擬同型}
\ref{derived-definition-filtered-acyclic} を参照

\noindent
{\it フィルター付き擬同型}
\ref{trace-definition-filtered} を参照

\noindent
{\it フィルター付き擬同型}
\ref{exercises-definition-filtered-quasi-isomorphism} を参照

\noindent
{\it フィルター付き}
\ref{categories-definition-directed} を参照

\noindent
{\it フィルター付き}
\ref{categories-definition-directed} を参照

\noindent
{\it 終対象}
\ref{categories-definition-final-object-2-category} を参照

\noindent
{\it 終対象}
\ref{categories-definition-initial-final} を参照

\noindent
{\it より細かい}
\ref{sites-definition-finer} を参照

\noindent
{\it 有限 $\text{Tor}$ 次元}
\ref{trace-definition-finite-tor-dimension} を参照

\noindent
{\it 有限 $R$-加群}
\ref{algebra-definition-module-finite-type} を参照

\noindent
{\it 有限自由}
\ref{sites-modules-definition-global} を参照

\noindent
{\it 有限大域次元}
\ref{algebra-definition-finite-gl-dim} を参照

\noindent
{\it 有限入射次元}
\ref{more-algebra-definition-injective-dimension} を参照

\noindent
{\it 有限局所定数}
\ref{sites-modules-definition-locally-constant} を参照

\noindent
{\it 有限局所定数}
\ref{etale-cohomology-definition-finite-locally-constant} を参照

\noindent
{\it 有限局所定数}
\ref{etale-cohomology-definition-finite-locally-constant} を参照

\noindent
{\it 階数 $r$ の有限局所自由}
\ref{algebra-definition-locally-free} を参照

\noindent
{\it 階数 $r$ の有限局所自由}
\ref{modules-definition-locally-free} を参照

\noindent
{\it 有限局所自由}
\ref{algebra-definition-locally-free} を参照

\noindent
{\it 有限局所自由}
\ref{modules-definition-locally-free} を参照

\noindent
{\it 有限局所自由}
\ref{sites-modules-definition-site-local} を参照

\noindent
{\it 有限局所自由}
\ref{morphisms-definition-finite-locally-free} を参照

\noindent
{\it 有限局所自由}
\ref{spaces-morphisms-definition-finite-locally-free} を参照

\noindent
{\it 有限局所自由}
\ref{exercises-definition-finite-locally-free} を参照

\noindent
{\it 有限次肥厚化}
\ref{more-morphisms-definition-thickening} を参照

\noindent
{\it 有限次肥厚化}
\ref{spaces-more-morphisms-definition-thickening} を参照

\noindent
{\it $x \in X$ において有限表示}
\ref{morphisms-definition-finite-presentation} を参照

\noindent
{\it $x$ において有限表示}
\ref{spaces-morphisms-definition-locally-finite-presentation} を参照

\noindent
{\it 有限表示}
\ref{algebra-definition-finite-type} を参照

\noindent
{\it 有限表示}
\ref{modules-definition-finite-presentation} を参照

\noindent
{\it 有限表示}
\ref{morphisms-definition-finite-presentation} を参照

\noindent
{\it 有限表示}
\ref{exercises-definition-finite-presentation} を参照

\noindent
{\it 有限射影次元}
\ref{algebra-definition-finite-proj-dim} を参照

\noindent
{\it 有限射影次元}
\ref{more-algebra-definition-projective-dimension} を参照

\noindent
{\it 有限 Tor 次元をもつ}
\ref{more-algebra-definition-tor-amplitude} を参照

\noindent
{\it 有限 Tor 次元をもつ}
\ref{more-algebra-definition-tor-amplitude} を参照

\noindent
{\it 有限 Tor 次元をもつ}
\ref{cohomology-definition-tor-amplitude} を参照

\noindent
{\it 有限 Tor 次元をもつ}
\ref{sites-cohomology-definition-tor-amplitude} を参照

\noindent
{\it $x \in X$ において有限型}
\ref{morphisms-definition-finite-type} を参照

\noindent
{\it $x$ において有限型}
\ref{spaces-morphisms-definition-locally-finite-type} を参照

\noindent
{\it 有限型点}
\ref{morphisms-definition-finite-type-point} を参照

\noindent
{\it 有限型点}
\ref{spaces-morphisms-definition-finite-type-point} を参照

\noindent
{\it 有限型点}
\ref{stacks-morphisms-definition-finite-type-point} を参照

\noindent
{\it 有限型}
\ref{algebra-definition-finite-type} を参照

\noindent
{\it 有限型}
\ref{modules-definition-finite-type} を参照

\noindent
{\it 有限型}
\ref{morphisms-definition-finite-type} を参照

\noindent
{\it 有限型}
\ref{formal-spaces-definition-finite-type} を参照

\noindent
{\it 有限生成 $R$-加群}
\ref{algebra-definition-module-finite-type} を参照

\noindent
{\it 有限生成体拡大}
\ref{fields-definition-generated-by} を参照

\noindent
{\it 有限表示 $R$-加群}
\ref{algebra-definition-module-finite-type} を参照

\noindent
{\it $R$ に関して有限表示}
\ref{more-algebra-definition-relatively-finitely-presented} を参照

\noindent
{\it $S$ に相対的に有限表示される}
\ref{more-morphisms-definition-relatively-finitely-presented-sheaf} を参照

\noindent
{\it 有限}
\ref{fields-definition-degree} を参照

\noindent
{\it 有限}
\ref{algebra-definition-finite-ring-map} を参照

\noindent
{\it 有限}
\ref{brauer-definition-finite} を参照

\noindent
{\it 有限}
\ref{homology-definition-filtered} を参照

\noindent
{\it 有限}
\ref{morphisms-definition-integral} を参照

\noindent
{\it 有限}
\ref{spaces-morphisms-definition-integral} を参照

\noindent
{\it 有限}
\ref{stacks-morphisms-definition-integral} を参照

\noindent
{\it 第一 Chern 類}
\ref{chow-definition-first-chern-class} を参照

\noindent
{\it 一次無限小近傍}
\ref{more-morphisms-definition-first-order-infinitesimal-neighbourhood} を参照

\noindent
{\it 一次無限小近傍}
\ref{spaces-more-morphisms-definition-first-order-infinitesimal-neighbourhood} を参照

\noindent
{\it 一次肥厚化}
\ref{more-morphisms-definition-thickening} を参照

\noindent
{\it 一次肥厚化}
\ref{spaces-more-morphisms-definition-thickening} を参照

\noindent
{\it 一次肥厚化}
\ref{stacks-more-morphisms-definition-first-order-thickening} を参照

\noindent
{\it フラビー}
\ref{cohomology-definition-flasque} を参照

\noindent
{\it フラスク}
\ref{cohomology-definition-flasque} を参照

\noindent
{\it 平坦（それぞれ忠実平坦）}
\ref{etale-definition-flat-rings} を参照

\noindent
{\it $x \in X$ において平坦}
\ref{etale-definition-flat-schemes} を参照

\noindent
{\it $x$ において $Y$ 上平坦}
\ref{spaces-morphisms-definition-flat-module} を参照

\noindent
{\it $x$ で平坦}
\ref{modules-definition-flat-at-point} を参照

\noindent
{\it $x$ で平坦}
\ref{modules-definition-flat-morphism} を参照

\noindent
{\it $x$ で平坦}
\ref{spaces-morphisms-definition-flat} を参照

\noindent
{\it $x$ で平坦}
\ref{stacks-morphisms-definition-flat-at-point} を参照

\noindent
{\it 点 $x \in X$ において平坦}
\ref{morphisms-definition-flat} を参照

\noindent
{\it 平坦基底変換性}
\ref{stacks-cohomology-definition-flat-base-change} を参照

\noindent
{\it 平坦基底変換}
\ref{groupoids-quotients-definition-base-change} を参照

\noindent
{\it 平坦な群スキーム}
\ref{groupoids-definition-smooth-group-scheme} を参照

\noindent
{\it $R$ 上の平坦な局所完全交叉}
\ref{algebra-definition-lci} を参照

\noindent
{\it $(\Sh(\mathcal{D}), \mathcal{O}')$ 上平坦}
\ref{sites-modules-definition-flat-module} を参照

\noindent
{\it $S$ 上、点 $x \in X$ において平坦}
\ref{morphisms-definition-flat} を参照

\noindent
{\it $S$ 上平坦}
\ref{morphisms-definition-flat} を参照

\noindent
{\it $x \in X$ において $Y$ 上平坦}
\ref{etale-definition-flat-schemes} を参照

\noindent
{\it 点 $x \in X$ で $Y$ 上平坦}
\ref{modules-definition-flat-module} を参照

\noindent
{\it $Y$ 上平坦}
\ref{modules-definition-flat-module} を参照

\noindent
{\it $Y$ 上平坦}
\ref{spaces-morphisms-definition-flat-module} を参照

\noindent
{\it $f$ による $\alpha$ の平坦引き戻し}
\ref{chow-definition-flat-pullback} を参照

\noindent
{\it $f$ による $\alpha$ の平坦引き戻し}
\ref{spaces-chow-definition-flat-pullback} を参照

\noindent
{\it 平坦 fppf サイト}
\ref{stacks-cohomology-definition-lisse-etale} を参照

\noindent
{\it 平坦化層別}
\ref{flat-definition-flattening-stratification} を参照

\noindent
{\it 平坦化層別}
\ref{flat-definition-flattening-stratification} を参照

\noindent
{\it 平坦}
\ref{algebra-definition-flat} を参照

\noindent
{\it 平坦}
\ref{algebra-definition-flat} を参照

\noindent
{\it 平坦}
\ref{modules-definition-flat} を参照

\noindent
{\it 平坦}
\ref{modules-definition-flat-morphism} を参照

\noindent
{\it 平坦}
\ref{sites-modules-definition-flat} を参照

\noindent
{\it 平坦}
\ref{sites-modules-definition-flat} を参照

\noindent
{\it 平坦}
\ref{sites-modules-definition-flat} を参照

\noindent
{\it 平坦}
\ref{sites-modules-definition-flat} を参照

\noindent
{\it 平坦}
\ref{sites-modules-definition-flat-morphism} を参照

\noindent
{\it 平坦}
\ref{sites-modules-definition-flat-morphism} を参照

\noindent
{\it 平坦}
\ref{morphisms-definition-flat} を参照

\noindent
{\it 平坦}
\ref{etale-definition-flat-rings} を参照

\noindent
{\it 平坦}
\ref{etale-definition-flat-schemes} を参照

\noindent
{\it 平坦}
\ref{spaces-morphisms-definition-flat} を参照

\noindent
{\it 平坦}
\ref{restricted-definition-flat} を参照

\noindent
{\it 平坦}
\ref{stacks-morphisms-definition-flat} を参照

\noindent
{\it 形式代数空間}
\ref{formal-spaces-definition-formal-algebraic-space} を参照

\noindent
{\it $\mathcal{X}$ の $x_0$ を通る形式的分枝}
\ref{stacks-geometry-definition-formal-branches} を参照

\noindent
{\it 形式的修正}
\ref{restricted-definition-formal-modification} を参照

\noindent
{\it $\mathcal{F}$ の形式対象 $\xi = (R, \xi_n, f_n)$}
\ref{formal-defos-definition-formal-objects} を参照

\noindent
{\it 形式対象}
\ref{artin-definition-formal-objects} を参照

\noindent
{\it 形式スペクトル}
\ref{formal-spaces-definition-affine-formal-spectrum} を参照

\noindent
{\it $R$ 上形式的エタール}
\ref{algebra-definition-formally-etale} を参照

\noindent
{\it 形式的エタール}
\ref{more-morphisms-definition-formally-etale} を参照

\noindent
{\it 形式的エタール}
\ref{spaces-more-morphisms-definition-formally-smooth-etale-unramified} を参照

\noindent
{\it 形式的エタール}
\ref{spaces-more-morphisms-definition-formally-etale} を参照

\noindent
{\it 形式鎖状}
\ref{more-algebra-definition-formally-catenary} を参照

\noindent
{\it $G$ の下で形式的に主等質}
\ref{groupoids-definition-pseudo-torsor} を参照

\noindent
{\it $G$ の下で形式的に主等質}
\ref{spaces-groupoids-definition-pseudo-torsor} を参照

\noindent
{\it $\mathfrak n$-進位相に関して形式的に 滑らか}
\ref{more-algebra-definition-formally-smooth-adic} を参照

\noindent
{\it $R$ 上形式的に滑らか}
\ref{algebra-definition-formally-smooth} を参照

\noindent
{\it $R$ 上形式的に滑らか}
\ref{more-algebra-definition-formally-smooth} を参照

\noindent
{\it 形式的に滑らか}
\ref{more-morphisms-definition-formally-smooth} を参照

\noindent
{\it 形式的に滑らか}
\ref{spaces-more-morphisms-definition-formally-smooth-etale-unramified} を参照

\noindent
{\it 形式的に滑らか}
\ref{spaces-more-morphisms-definition-formally-smooth} を参照

\noindent
{\it 形式的に滑らか}
\ref{stacks-more-morphisms-definition-formally-smooth} を参照

\noindent
{\it $R$ 上形式的不分岐}
\ref{algebra-definition-formally-unramified} を参照

\noindent
{\it 形式的不分岐}
\ref{more-morphisms-definition-formally-unramified} を参照

\noindent
{\it 形式的不分岐}
\ref{spaces-more-morphisms-definition-formally-smooth-etale-unramified} を参照

\noindent
{\it 形式的不分岐}
\ref{spaces-more-morphisms-definition-formally-unramified} を参照

\noindent
{\it Fourier--Mukai 関手}
\ref{equiv-definition-fourier-mukai-functor} を参照

\noindent
{\it Fourier--Mukai 核}
\ref{equiv-definition-fourier-mukai-functor} を参照

\noindent
{\it $T$ の fppf 被覆}
\ref{topologies-definition-fppf-covering} を参照

\noindent
{\it $X$ の fppf 被覆}
\ref{spaces-topologies-definition-fppf-covering} を参照

\noindent
{\it fppf 層}
\ref{stacks-sheaves-definition-sheaves} を参照

\noindent
{\it $T$ の fpqc 被覆}
\ref{topologies-definition-fpqc-covering} を参照

\noindent
{\it $X$ の fpqc 被覆}
\ref{spaces-topologies-definition-fpqc-covering} を参照

\noindent
{\it fpqc 被覆}
\ref{etale-cohomology-definition-fpqc-covering} を参照

\noindent
{\it 自由 $\mathcal{O}$-加群}
\ref{sites-modules-definition-global} を参照

\noindent
{\it $\mathcal{G}$ 上の自由アーベル前層}
\ref{etale-cohomology-definition-free-abelian-presheaf} を参照

\noindent
{\it 自由アーベル群の前層}
\ref{sites-modules-definition-free-abelian-presheaf-on} を参照

\noindent
{\it 自由アーベル群の層}
\ref{sites-modules-definition-free-abelian-sheaf-on} を参照

\noindent
{\it 自由加群}
\ref{more-algebra-definition-simple-functors} を参照

\noindent
{\it 自由}
\ref{groupoids-definition-free-action} を参照

\noindent
{\it 自由}
\ref{spaces-groupoids-definition-free-action} を参照

\noindent
{\it 充満部分圏}
\ref{categories-definition-subcategory} を参照

\noindent
{\it 充満忠実}
\ref{categories-definition-faithful} を参照

\noindent
{\it 関数体}
\ref{morphisms-definition-function-field} を参照

\noindent
{\it 関数体}
\ref{spaces-over-fields-definition-function-field} を参照

\noindent
{\it $R$-線形圏の間の関手}
\ref{dga-definition-functor-linear-categories} を参照

\noindent
{\it $R$ 上の微分次数付き圏の関手}
\ref{dga-definition-functor-dga-categories} を参照

\noindent
{\it $R$ 上の次数付き圏の間の関手}
\ref{dga-definition-functor-graded-categories} を参照

\noindent
{\it モノイド圏の関手}
\ref{categories-definition-functor-monoidal-categories} を参照

\noindent
{\it 対称モノイド圏の関手}
\ref{categories-definition-functor-symmetric-monoidal-categories} を参照

\noindent
{\it 関手的入射埋め込みをもつ}
\ref{homology-definition-functorial-injective-embedding} を参照

\noindent
{\it 関手的射影全射をもつ}
\ref{homology-definition-functorial-projective-surjections} を参照

\noindent
{\it 関手}
\ref{categories-definition-functor} を参照

\noindent
{\it 関手}
\ref{categories-definition-functor-into-2-category} を参照

\noindent
{\it 基本群}
\ref{pione-definition-fundamental-group} を参照

\noindent
{\it G-環}
\ref{more-algebra-definition-G-ring} を参照

\noindent
{\it G-スキーム}
\ref{properties-definition-G-scheme} を参照

\noindent
{\it $\mathfrak q$ において G-不分岐}
\ref{algebra-definition-unramified} を参照

\noindent
{\it $x \in X$ において G-不分岐}
\ref{morphisms-definition-unramified} を参照

\noindent
{\it $x$ において G-非分岐}
\ref{spaces-morphisms-definition-unramified} を参照

\noindent
{\it G-不分岐}
\ref{algebra-definition-unramified} を参照

\noindent
{\it G-不分岐}
\ref{morphisms-definition-unramified} を参照

\noindent
{\it G-不分岐}
\ref{spaces-morphisms-definition-unramified} を参照

\noindent
{\it ガロア圏}
\ref{pione-definition-galois-category} を参照

\noindent
{\it 係数 $M$ の $K$ のガロアコホモロジー群}
\ref{etale-cohomology-definition-galois-cohomology} を参照

\noindent
{\it ガロアコホモロジー群}
\ref{etale-cohomology-definition-galois-cohomology} を参照

\noindent
{\it ガロア群}
\ref{fields-definition-galois-group} を参照

\noindent
{\it ガロア}
\ref{fields-definition-galois} を参照

\noindent
{\it ガロア}
\ref{fields-definition-separable-algebraic} を参照

\noindent
{\it 一般化は $f$ に沿って 持ち上がる}
\ref{topology-definition-lift-specializations} を参照

\noindent
{\it 一般化}
\ref{topology-definition-specialization} を参照

\noindent
{\it 一般化}
\ref{exercises-definition-closed} を参照

\noindent
{\it 一般化的}
\ref{topology-definition-lift-specializations} を参照

\noindent
{\it $r$ 個の大域切断で生成される}
\ref{sites-modules-definition-global} を参照

\noindent
{\it 有限個の大域切断で生成される}
\ref{sites-modules-definition-global} を参照

\noindent
{\it 大域切断で生成される}
\ref{modules-definition-globally-generated} を参照

\noindent
{\it 大域切断で生成される}
\ref{sites-modules-definition-global} を参照

\noindent
{\it 生成する}
\ref{fields-definition-generated-by} を参照

\noindent
{\it 生成する}
\ref{modules-definition-globally-generated} を参照

\noindent
{\it 生成対象}
\ref{derived-definition-generators} を参照

\noindent
{\it 生成対象}
\ref{injectives-definition-grothendieck-conditions} を参照

\noindent
{\it 生成点}
\ref{topology-definition-generic-point} を参照

\noindent
{\it 生成点}
\ref{exercises-definition-generic-point} を参照

\noindent
{\it 種数}
\ref{pic-definition-genus} を参照

\noindent
{\it 種数}
\ref{curves-definition-genus} を参照

\noindent
{\it 幾何学的フロベニウス}
\ref{trace-definition-geometric-frobenius} を参照

\noindent
{\it 幾何学的フロベニウス}
\ref{trace-definition-geometric-frobenius-on-stalk} を参照

\noindent
{\it 幾何種数}
\ref{curves-definition-geometric-genus} を参照

\noindent
{\it $x$ の上にある幾何学的点}
\ref{spaces-properties-definition-geometric-point} を参照

\noindent
{\it 幾何学的点}
\ref{etale-cohomology-definition-geometric-point} を参照

\noindent
{\it 幾何学的点}
\ref{spaces-properties-definition-geometric-point} を参照

\noindent
{\it 幾何学的商}
\ref{groupoids-quotients-definition-geometric} を参照

\noindent
{\it $k$ 上幾何学的連結}
\ref{algebra-definition-geometrically-connected} を参照

\noindent
{\it 幾何学的連結}
\ref{varieties-definition-geometrically-connected} を参照

\noindent
{\it 幾何学的連結}
\ref{spaces-over-fields-definition-geometrically-connected} を参照

\noindent
{\it $k$ 上幾何学的整}
\ref{algebra-definition-geometrically-integral} を参照

\noindent
{\it 幾何学的整}
\ref{varieties-definition-geometrically-integral} を参照

\noindent
{\it 幾何学的整}
\ref{spaces-over-fields-definition-geometrically-integral} を参照

\noindent
{\it $k$ 上幾何学的既約}
\ref{algebra-definition-geometrically-irreducible} を参照

\noindent
{\it 幾何学的既約}
\ref{varieties-definition-geometrically-irreducible} を参照

\noindent
{\it 幾何学的既約}
\ref{spaces-over-fields-definition-geometrically-irreducible} を参照

\noindent
{\it $x$ において幾何学的に正規}
\ref{varieties-definition-geometrically-normal} を参照

\noindent
{\it 幾何学的に正規}
\ref{algebra-definition-geometrically-normal} を参照

\noindent
{\it 幾何学的に正規}
\ref{varieties-definition-geometrically-normal} を参照

\noindent
{\it $x$ において幾何学的に点ごとに整}
\ref{varieties-definition-geometrically-integral} を参照

\noindent
{\it 幾何学的に点ごとに整}
\ref{varieties-definition-geometrically-integral} を参照

\noindent
{\it $x$ で幾何学的被約}
\ref{varieties-definition-geometrically-reduced} を参照

\noindent
{\it $x$ で幾何学的被約}
\ref{spaces-over-fields-definition-geometrically-reduced} を参照

\noindent
{\it $k$ 上幾何学的に被約}
\ref{algebra-definition-geometrically-reduced} を参照

\noindent
{\it 幾何学的被約}
\ref{varieties-definition-geometrically-reduced} を参照

\noindent
{\it 幾何学的被約}
\ref{spaces-over-fields-definition-geometrically-reduced} を参照

\noindent
{\it $x$ において幾何学的に正則}
\ref{varieties-definition-geometrically-regular} を参照

\noindent
{\it $k$ 上幾何学的に正則}
\ref{varieties-definition-geometrically-regular} を参照

\noindent
{\it 幾何学的に正則}
\ref{algebra-definition-geometrically-regular} を参照

\noindent
{\it $x$ において幾何的単枝}
\ref{properties-definition-unibranch} を参照

\noindent
{\it $x$ において幾何的単枝}
\ref{spaces-properties-definition-unibranch} を参照

\noindent
{\it $x$ において幾何的単枝}
\ref{stacks-properties-definition-number-of-geometric-branches} を参照

\noindent
{\it 幾何的単枝}
\ref{more-algebra-definition-unibranch} を参照

\noindent
{\it 幾何的単枝}
\ref{properties-definition-unibranch} を参照

\noindent
{\it 幾何的単枝}
\ref{spaces-properties-definition-unibranch} を参照

\noindent
{\it 上のガーベ}
\ref{stacks-definition-gerbe-over-stack-in-groupoids} を参照

\noindent
{\it 上のガーベ}
\ref{stacks-morphisms-definition-gerbe} を参照

\noindent
{\it ガーベ}
\ref{stacks-definition-gerbe} を参照

\noindent
{\it ガーベ}
\ref{stacks-morphisms-definition-gerbe} を参照

\noindent
{\it $x$ における $X$ の芽}
\ref{descent-definition-germs} を参照

\noindent
{\it $k$ 上大域的完全交叉}
\ref{algebra-definition-lci-field} を参照

\noindent
{\it 大域次元}
\ref{algebra-definition-finite-gl-dim} を参照

\noindent
{\it 大域有限表示}
\ref{sites-modules-definition-global} を参照

\noindent
{\it 大域 Lefschetz 数}
\ref{trace-definition-global-lefschetz-number} を参照

\noindent
{\it 大域表示}
\ref{sites-modules-definition-global} を参照

\noindent
{\it 大域切断}
\ref{sites-definition-global-sections} を参照

\noindent
{\it 下降性}
\ref{algebra-definition-going-up-down} を参照

\noindent
{\it 上昇性}
\ref{algebra-definition-going-up-down} を参照

\noindent
{\it 下降定理}
\ref{exercises-definition-GU-GD} を参照

\noindent
{\it 上昇定理}
\ref{exercises-definition-GU-GD} を参照

\noindent
{\it 良い商}
\ref{groupoids-quotients-definition-good} を参照

\noindent
{\it 良い還元}
\ref{models-definition-good} を参照

\noindent
{\it 良い層化}
\ref{topology-definition-good-stratification} を参照

\noindent
{\it $x$ において Gorenstein}
\ref{duality-definition-gorenstein-morphism} を参照

\noindent
{\it $x$ において Gorenstein}
\ref{spaces-more-morphisms-definition-gorenstein} を参照

\noindent
{\it Gorenstein 射}
\ref{duality-definition-gorenstein-morphism} を参照

\noindent
{\it Gorenstein 射}
\ref{spaces-more-morphisms-definition-gorenstein} を参照

\noindent
{\it Gorenstein}
\ref{dualizing-definition-gorenstein} を参照

\noindent
{\it Gorenstein}
\ref{dualizing-definition-gorenstein} を参照

\noindent
{\it Gorenstein}
\ref{duality-definition-gorenstein} を参照

\noindent
{\it 次数付き $(\mathcal{A}, \mathcal{B})$-両側加群}
\ref{sdga-definition-bimodule} を参照

\noindent
{\it 次数付き $(A, B)$-両側加群}
\ref{dga-definition-bimodule} を参照

\noindent
{\it 次数付き $\mathcal{A}$-加群}
\ref{sdga-definition-gm} を参照

\noindent
{\it 次数付き $A$-代数}
\ref{exercises-definition-graded-algebra} を参照

\noindent
{\it $R$ 上の次数付き圏 $\mathcal{A}$}
\ref{dga-definition-graded-category} を参照

\noindent
{\it 次数付き直和}
\ref{dga-definition-graded-direct-sum} を参照

\noindent
{\it 次数付き関手}
\ref{dga-definition-functor-graded-categories} を参照

\noindent
{\it 次数付きイデアル}
\ref{exercises-definition-graded-algebra} を参照

\noindent
{\it 次数付き入射的}
\ref{sdga-definition-graded-injective} を参照

\noindent
{\it 次数付き $A$-代数 $B$ 上の次数付き加群 $M$}
\ref{exercises-definition-graded-algebra} を参照

\noindent
{\it 次数付き加群}
\ref{sdga-definition-gm} を参照

\noindent
{\it 次数付き加群}
\ref{exercises-definition-graded-module} を参照

\noindent
{\it 次数付き部分加群}
\ref{exercises-definition-graded-algebra} を参照

\noindent
{\it $\mathbf{Z}$ 上のグラスマン多様体}
\ref{constructions-definition-grassmannian} を参照

\noindent
{\it $R$ 上のグラスマン多様体}
\ref{constructions-definition-grassmannian} を参照

\noindent
{\it $S$ 上のグラスマン多様体}
\ref{constructions-definition-grassmannian} を参照

\noindent
{\it Grothendieck アーベル圏}
\ref{injectives-definition-grothendieck-conditions} を参照

\noindent
{\it $X$ のグロタンディーク群}
\ref{perfect-definition-K-group} を参照

\noindent
{\it $X$ 上のコヒーレント層のグロタンディーク群}
\ref{perfect-definition-K-group} を参照

\noindent
{\it $B$ 上の群代数空間}
\ref{spaces-groupoids-definition-group-space} を参照

\noindent
{\it 群コホモロジー群}
\ref{etale-cohomology-definition-galois-cohomology} を参照

\noindent
{\it $x$ 上の $x'$ の無限小自己同型群}
\ref{formal-defos-definition-relative-infinitesimal-auts} を参照

\noindent
{\it $x_0$ の無限小自己同型群}
\ref{formal-defos-definition-infinitesimal-auts} を参照

\noindent
{\it $S$ 上の群スキーム}
\ref{groupoids-definition-group-scheme} を参照

\noindent
{\it $B$ 上の代数空間における亜群}
\ref{spaces-groupoids-definition-groupoid} を参照

\noindent
{\it $\mathcal{C}$ 上の関手における亜群}
\ref{formal-defos-definition-groupoid-in-functors} を参照

\noindent
{\it $S$ 上の亜群}
\ref{groupoids-definition-groupoid} を参照

\noindent
{\it $S$ 上の亜群スキーム}
\ref{groupoids-definition-groupoid} を参照

\noindent
{\it グルーポイド}
\ref{categories-definition-groupoid} を参照

\noindent
{\it Gysin 準同型}
\ref{chow-definition-gysin-homomorphism} を参照

\noindent
{\it Gysin 準同型}
\ref{spaces-chow-definition-gysin-homomorphism} を参照

\noindent
{\it Gysin 写像}
\ref{chow-definition-lci-gysin} を参照

\noindent
{\it $T$ の h 被覆}
\ref{flat-definition-h-covering} を参照

\noindent
{\it H-射影的}
\ref{morphisms-definition-projective} を参照

\noindent
{\it H-準射影的}
\ref{morphisms-definition-quasi-projective} を参照

\noindent
{\it 二対象の余積をもつ}
\ref{categories-definition-has-coproducts-of-pairs} を参照

\noindent
{\it 十分多くの点をもつ}
\ref{sites-definition-enough-points} を参照

\noindent
{\it ファイバー積をもつ}
\ref{categories-definition-has-fibre-products} を参照

\noindent
{\it 二対象の積をもつ}
\ref{categories-definition-has-products-of-pairs} を参照

\noindent
{\it 性質 $(\beta)$ をもつ}
\ref{decent-spaces-definition-relative-conditions} を参照

\noindent
{\it 性質 $(\beta)$ をもつ}
\ref{decent-spaces-definition-relative-conditions} を参照

\noindent
{\it $x$ において性質 $\mathcal{P}$ をもつ}
\ref{spaces-properties-definition-property-at-point} を参照

\noindent
{\it $x$ において性質 $\mathcal{P}$ をもつ}
\ref{stacks-properties-definition-property-at-point} を参照

\noindent
{\it 性質 $\mathcal{P}$ をもつ}
\ref{spaces-properties-definition-type-property} を参照

\noindent
{\it 性質 $\mathcal{P}$ をもつ}
\ref{spaces-morphisms-definition-P} を参照

\noindent
{\it 性質 $\mathcal{P}$ をもつ}
\ref{stacks-properties-definition-type-property} を参照

\noindent
{\it 性質 $\mathcal{P}$ をもつ}
\ref{stacks-morphisms-definition-P} を参照

\noindent
{\it 性質 $\mathcal{P}$ をもつ}
\ref{stacks-morphisms-definition-etale-smooth-P} を参照

\noindent
{\it $x$ において性質 $\mathcal{Q}$ をもつ}
\ref{spaces-morphisms-definition-P-at-point} を参照

\noindent
{\it Hausdorff}
\ref{exercises-definition-Hausdorff} を参照

\noindent
{\it 高さ}
\ref{algebra-definition-height} を参照

\noindent
{\it $x$ における $X$ のヘンゼル局所環}
\ref{decent-spaces-definition-henselian-local-ring} を参照

\noindent
{\it ヘンゼル対}
\ref{more-algebra-definition-henselian-pair} を参照

\noindent
{\it ヘンゼル的}
\ref{algebra-definition-henselian} を参照

\noindent
{\it ヘンゼル的}
\ref{etale-cohomology-definition-henselian} を参照

\noindent
{\it $\mathcal{O}_{S, s}$ のヘンゼル化}
\ref{etale-cohomology-definition-etale-local-rings} を参照

\noindent
{\it $s$ における $S$ のヘンゼル化}
\ref{etale-cohomology-definition-etale-local-rings} を参照

\noindent
{\it ヘンゼル化}
\ref{algebra-definition-henselization} を参照

\noindent
{\it 高次順像}
\ref{etale-cohomology-definition-higher-direct-images} を参照

\noindent
{\it ヒルベルト関数}
\ref{exercises-definition-graded-module} を参照

\noindent
{\it ヒルベルト多項式}
\ref{algebra-definition-hilbert-polynomial} を参照

\noindent
{\it ヒルベルト多項式}
\ref{varieties-definition-hilbert-polynomial} を参照

\noindent
{\it ヒルベルト多項式}
\ref{exercises-definition-graded-module} を参照

\noindent
{\it ホッジフィルトレーション}
\ref{derham-definition-hodge-filtration} を参照

\noindent
{\it 斉次スペクトル $\text{Proj}(R)$}
\ref{exercises-definition-Proj-R} を参照

\noindent
{\it $\mathcal{A}$ の $S$ 上の斉次スペクトル}
\ref{constructions-definition-relative-proj} を参照

\noindent
{\it $X$ 上の $\mathcal{A}$ の斉次スペクトル}
\ref{spaces-divisors-definition-relative-proj} を参照

\noindent
{\it 斉次スペクトル}
\ref{algebra-definition-proj} を参照

\noindent
{\it 斉次スペクトル}
\ref{constructions-definition-structure-sheaf} を参照

\noindent
{\it 斉次}
\ref{exercises-definition-homogeneous-ideal} を参照

\noindent
{\it ホモロジー的}
\ref{derived-definition-homological} を参照

\noindent
{\it $K$ のホモロジー}
\ref{hypercovering-definition-homology} を参照

\noindent
{\it ホモロジー}
\ref{homology-definition-differential-object-homology} を参照

\noindent
{\it 微分次数付き $(\mathcal{A}, \mathcal{B})$-両側加群の準同型}
\ref{sdga-definition-dg-bimodule} を参照

\noindent
{\it 微分次数付き $\mathcal{A}$-加群の準同型}
\ref{sdga-definition-dgm} を参照

\noindent
{\it 微分次数付き $\mathcal{O}$-代数の準同型}
\ref{sdga-definition-dga} を参照

\noindent
{\it 微分次数付き代数の準同型}
\ref{dga-definition-homomorphism-dga} を参照

\noindent
{\it 微分次数付き加群の準同型}
\ref{dga-definition-dgm} を参照

\noindent
{\it 分割冪環の準同型}
\ref{dpa-definition-divided-power-ring} を参照

\noindent
{\it 分割冪肥厚化の準同型}
\ref{crystalline-definition-affine-thickening} を参照

\noindent
{\it 次数付き $(\mathcal{A}, \mathcal{B})$-両側加群の準同型}
\ref{sdga-definition-bimodule} を参照

\noindent
{\it 次数付き $\mathcal{A}$-加群の準同型}
\ref{sdga-definition-gm} を参照

\noindent
{\it 次数付き $\mathcal{O}$-代数の準同型}
\ref{sdga-definition-ga} を参照

\noindent
{\it 系の準同型}
\ref{algebra-definition-homomorphism-directed-systems} を参照

\noindent
{\it 位相群の準同型}
\ref{topology-definition-topological-group} を参照

\noindent
{\it 位相加群の準同型}
\ref{topology-definition-topological-module} を参照

\noindent
{\it 位相加群の準同型}
\ref{more-algebra-definition-topological-ring} を参照

\noindent
{\it 位相環の準同型}
\ref{topology-definition-topological-ring} を参照

\noindent
{\it 位相環の準同型}
\ref{more-algebra-definition-topological-ring} を参照

\noindent
{\it 次数付き加群・環の準同型}
\ref{exercises-definition-graded-algebra} を参照

\noindent
{\it ホモトピック}
\ref{simplicial-definition-homotopy} を参照

\noindent
{\it ホモトピック}
\ref{simplicial-definition-homotopy-cosimplicial} を参照

\noindent
{\it ホモトピック}
\ref{dga-definition-homotopy} を参照

\noindent
{\it ホモトピック}
\ref{sdga-definition-homotopy} を参照

\noindent
{\it $f$ と $g$ の間のホモトピー}
\ref{dga-definition-homotopy} を参照

\noindent
{\it $f$ と $g$ の間のホモトピー}
\ref{sdga-definition-homotopy} を参照

\noindent
{\it $\mathcal{A}$ のホモトピー圏}
\ref{dga-definition-homotopy-category-of-dga-category} を参照

\noindent
{\it ホモトピー圏}
\ref{dga-definition-complexes-notation} を参照

\noindent
{\it ホモトピー圏}
\ref{sdga-definition-complexes-notation} を参照

\noindent
{\it ホモトピー余極限}
\ref{derived-definition-derived-colimit} を参照

\noindent
{\it ホモトピー同値}
\ref{homology-definition-homotopy-equivalent} を参照

\noindent
{\it ホモトピー同値}
\ref{homology-definition-homotopy-equivalent-cochain} を参照

\noindent
{\it ホモトピー同値}
\ref{simplicial-definition-homotopy-equivalent} を参照

\noindent
{\it ホモトピー同値}
\ref{homology-definition-homotopy-equivalent} を参照

\noindent
{\it ホモトピー同値}
\ref{homology-definition-homotopy-equivalent-cochain} を参照

\noindent
{\it ホモトピー同値}
\ref{simplicial-definition-homotopy-equivalent} を参照

\noindent
{\it $a$ から $b$ へのホモトピー}
\ref{simplicial-definition-homotopy} を参照

\noindent
{\it $a$ から $b$ へのホモトピー}
\ref{simplicial-definition-homotopy-cosimplicial} を参照

\noindent
{\it ホモトピー極限}
\ref{derived-definition-derived-limit} を参照

\noindent
{\it 水平}
\ref{categories-definition-horizontal-composition} を参照

\noindent
{\it 水平}
\ref{categories-definition-2-category} を参照

\noindent
{\it $\mathcal{G}$ の超被覆}
\ref{hypercovering-definition-hypercovering-variant} を参照

\noindent
{\it $X$ の超被覆}
\ref{hypercovering-definition-hypercovering} を参照

\noindent
{\it 超被覆}
\ref{hypercovering-definition-hypercovering-variant} を参照

\noindent
{\it 定義イデアル}
\ref{more-algebra-definition-topological-ring} を参照

\noindent
{\it $s$ の分母イデアル層}
\ref{divisors-definition-regular-meromorphic-ideal-denominators} を参照

\noindent
{\it 局所環を同一視する}
\ref{proetale-definition-local-isomorphism} を参照

\noindent
{\it $\varphi$ の像}
\ref{sites-definition-image} を参照

\noindent
{\it $f$ の像}
\ref{homology-definition-kernel} を参照

\noindent
{\it 与えられた $\delta$-関手による短完全列の像}
\ref{derived-definition-delta-functor} を参照

\noindent
{\it 直近特殊化}
\ref{topology-definition-dimension-function} を参照

\noindent
{\it 埋め込み}
\ref{schemes-definition-immersion} を参照

\noindent
{\it 埋め込み}
\ref{spaces-definition-immersion} を参照

\noindent
{\it 埋め込み}
\ref{stacks-properties-definition-immersion} を参照

\noindent
{\it $s$ 上の $\mathcal{F}$ の非純性}
\ref{flat-definition-impurity} を参照

\noindent
{\it $\mathcal{F}$ の $y$ 上の不純性}
\ref{spaces-flat-definition-impurity} を参照

\noindent
{\it 同じホモトピー類に属する}
\ref{simplicial-definition-homotopy} を参照

\noindent
{\it 同じホモトピー類に属する}
\ref{simplicial-definition-homotopy-cosimplicial} を参照

\noindent
{\it ind-エタール}
\ref{proetale-definition-ind-etale} を参照

\noindent
{\it ind-準アフィン}
\ref{more-morphisms-definition-ind-quasi-affine} を参照

\noindent
{\it ind-準アフィン}
\ref{more-morphisms-definition-ind-quasi-affine} を参照

\noindent
{\it ind-Zariski}
\ref{proetale-definition-ind-zariski} を参照

\noindent
{\it 直既約}
\ref{dualizing-definition-indecomposable} を参照

\noindent
{\it 誘導フィルトレーション}
\ref{homology-definition-filtered} を参照

\noindent
{\it 誘導フィルトレーション}
\ref{homology-definition-filtration-cohomology-filtered-differential} を参照

\noindent
{\it 誘導フィルトレーション}
\ref{homology-definition-filtration-cohomology-filtered-complex} を参照

\noindent
{\it $\mathcal{C}$ における $I$ 上の帰納系}
\ref{categories-definition-system-over-poset} を参照

\noindent
{\it $\mathcal{S}$ の惰性ファイバー圏 $\mathcal{I}_\mathcal{S}$}
\ref{categories-definition-inertia-fibred-category} を参照

\noindent
{\it $\mathfrak m$ の惰性群}
\ref{more-algebra-definition-decomposition-inertia} を参照

\noindent
{\it 始対象}
\ref{categories-definition-initial-final} を参照

\noindent
{\it 始対象}
\ref{categories-definition-initial} を参照

\noindent
{\it 入射包絡}
\ref{dualizing-definition-injective-hull} を参照

\noindent
{\it $A$ の入射分解}
\ref{derived-definition-injective-resolution} を参照

\noindent
{\it $K^\bullet$ の入射分解}
\ref{derived-definition-injective-resolution} を参照

\noindent
{\it $[a, b]$ における入射振幅}
\ref{more-algebra-definition-injective-dimension} を参照

\noindent
{\it 入射的}
\ref{sheaves-definition-injective-surjective} を参照

\noindent
{\it 入射的}
\ref{sheaves-definition-injective-surjective} を参照

\noindent
{\it 入射的}
\ref{sites-definition-presheaves-injective-surjective} を参照

\noindent
{\it 入射的}
\ref{sites-definition-sheaves-injective-surjective} を参照

\noindent
{\it 入射的}
\ref{homology-definition-injective-surjective} を参照

\noindent
{\it 入射的}
\ref{homology-definition-injective} を参照

\noindent
{\it 入射的}
\ref{more-algebra-definition-projective} を参照

\noindent
{\it 非分離次数}
\ref{fields-definition-insep-degree} を参照

\noindent
{\it $\mathcal{O}_X$ の整閉包 （$\mathcal{A}$ におけるもの）}
\ref{morphisms-definition-integral-closure} を参照

\noindent
{\it $\mathcal{O}_X$ の整閉包 （$\mathcal{A}$ におけるもの）}
\ref{spaces-morphisms-definition-integral-closure} を参照

\noindent
{\it 整閉包}
\ref{algebra-definition-integral-closure} を参照

\noindent
{\it 整域}
\ref{fields-definition-domain} を参照

\noindent
{\it $I$ 上整}
\ref{algebra-definition-integral-over-ideal} を参照

\noindent
{\it $R$ 上整}
\ref{algebra-definition-integral-ring-map} を参照

\noindent
{\it 整閉}
\ref{algebra-definition-integral-closure} を参照

\noindent
{\it 整}
\ref{algebra-definition-integral-ring-map} を参照

\noindent
{\it 整}
\ref{properties-definition-integral} を参照

\noindent
{\it 整}
\ref{morphisms-definition-integral} を参照

\noindent
{\it 整}
\ref{spaces-morphisms-definition-integral} を参照

\noindent
{\it 整}
\ref{spaces-over-fields-definition-integral-algebraic-space} を参照

\noindent
{\it 整}
\ref{stacks-morphisms-definition-integral} を参照

\noindent
{\it 整}
\ref{stacks-morphisms-definition-integral-algebraic-stack} を参照

\noindent
{\it 整}
\ref{exercises-definition-integral} を参照

\noindent
{\it 内部}
\ref{topology-definition-nowhere-dense} を参照

\noindent
{\it 正しく交わる}
\ref{intersection-definition-proper-intersection} を参照

\noindent
{\it 正しく交わる}
\ref{intersection-definition-proper-intersection} を参照

\noindent
{\it 交点数}
\ref{varieties-definition-intersection-number} を参照

\noindent
{\it 交点数}
\ref{spaces-over-fields-definition-intersection-number} を参照

\noindent
{\it $\mathcal{E}$ の第 $j$ Chern 類との交叉}
\ref{chow-definition-cap-chern-classes} を参照

\noindent
{\it $\mathcal{L}$ の第一 Chern 類との交叉積}
\ref{chow-definition-cap-c1} を参照

\noindent
{\it $\mathcal{L}$ の第一 Chern 類との交叉積}
\ref{spaces-chow-definition-cap-c1} を参照

\noindent
{\it 閉部分スキーム $Z$ の逆像 $f^{-1}(Z)$}
\ref{schemes-definition-inverse-image-closed-subscheme} を参照

\noindent
{\it 閉部分空間 $Z$ の逆像 $f^{-1}(Z)$}
\ref{spaces-morphisms-definition-inverse-image-closed-subspace} を参照

\noindent
{\it 逆像}
\ref{etale-cohomology-definition-inverse-image} を参照

\noindent
{\it $\mathcal{C}$ における $I$ 上の逆系}
\ref{categories-definition-system-over-poset} を参照

\noindent
{\it 可逆 $\mathcal{O}_X$-加群の層}
\ref{modules-definition-invertible} を参照

\noindent
{\it 可逆 ${\mathcal O}_X$-加群}
\ref{exercises-definition-invertible-sheaf} を参照

\noindent
{\it 可逆加群 $M$}
\ref{exercises-definition-invertible-module} を参照

\noindent
{\it 可逆加群}
\ref{exercises-definition-finite-locally-free} を参照

\noindent
{\it $D$ に付随する可逆層 $\mathcal{O}_S(D)$}
\ref{divisors-definition-invertible-sheaf-effective-Cartier-divisor} を参照

\noindent
{\it 可逆層 $\mathcal{O}_X(D)$（$D$ に伴うもの）}
\ref{spaces-divisors-definition-invertible-sheaf-effective-Cartier-divisor} を参照

\noindent
{\it 可逆}
\ref{categories-definition-invertible} を参照

\noindent
{\it 可逆}
\ref{more-algebra-definition-invertible} を参照

\noindent
{\it 可逆}
\ref{sites-modules-definition-invertible-sheaf} を参照

\noindent
{\it 既約成分}
\ref{topology-definition-irreducible-components} を参照

\noindent
{\it 既約成分}
\ref{exercises-definition-irreducible-component} を参照

\noindent
{\it 既約}
\ref{topology-definition-irreducible-components} を参照

\noindent
{\it 既約}
\ref{algebra-definition-irreducible-prime-element} を参照

\noindent
{\it 既約}
\ref{exercises-definition-irreducible} を参照

\noindent
{\it 既約}
\ref{exercises-definition-irreducible} を参照

\noindent
{\it 孤立点}
\ref{topology-definition-isolated-point} を参照

\noindent
{\it 同型射}
\ref{categories-definition-isomorphism} を参照

\noindent
{\it J-0}
\ref{more-algebra-definition-J} を参照

\noindent
{\it J-1}
\ref{more-algebra-definition-J} を参照

\noindent
{\it J-2}
\ref{more-algebra-definition-J} を参照

\noindent
{\it J-2}
\ref{morphisms-definition-J} を参照

\noindent
{\it ジャコブソン環}
\ref{algebra-definition-ring-jacobson} を参照

\noindent
{\it Jacobson}
\ref{topology-definition-space-jacobson} を参照

\noindent
{\it Jacobson}
\ref{properties-definition-jacobson} を参照

\noindent
{\it 日本スキーム}
\ref{algebra-definition-N} を参照

\noindent
{\it 日本スキーム}
\ref{properties-definition-nagata} を参照

\noindent
{\it K-平坦}
\ref{more-algebra-definition-K-flat} を参照

\noindent
{\it K-平坦}
\ref{cohomology-definition-K-flat} を参照

\noindent
{\it K-平坦}
\ref{sites-cohomology-definition-K-flat} を参照

\noindent
{\it K-入射的}
\ref{derived-definition-K-injective} を参照

\noindent
{\it K-入射的}
\ref{sdga-definition-K-injective} を参照

\noindent
{\it K\"ahler の差イデアル}
\ref{discriminant-definition-kahler-different} を参照

\noindent
{\it Kan 複体}
\ref{simplicial-definition-kan} を参照

\noindent
{\it Kan ファイブレーション}
\ref{simplicial-definition-kan} を参照

\noindent
{\it Kaplansky デヴィサージュ}
\ref{algebra-definition-devissage} を参照

\noindent
{\it カルービ圏}
\ref{homology-definition-karoubian} を参照

\noindent
{\it $F$ の核}
\ref{derived-definition-kernel-category} を参照

\noindent
{\it $H$ の核}
\ref{derived-definition-kernel-category} を参照

\noindent
{\it 関手 $F$ の核}
\ref{homology-definition-kernel-category} を参照

\noindent
{\it 核}
\ref{homology-definition-kernel} を参照

\noindent
{\it コルモゴロフ}
\ref{topology-definition-generic-point} を参照

\noindent
{\it $x$ で Koszul}
\ref{more-morphisms-definition-lci} を参照

\noindent
{\it $x$ で Koszul}
\ref{spaces-more-morphisms-definition-lci} を参照

\noindent
{\it $f_1, \ldots, f_n$ 上の Koszul 複体}
\ref{modules-definition-koszul-complex} を参照

\noindent
{\it $f_1, \ldots, f_r$ 上の Koszul 複体}
\ref{more-algebra-definition-koszul-complex} を参照

\noindent
{\it Koszul 複体}
\ref{more-algebra-definition-koszul} を参照

\noindent
{\it Koszul 複体}
\ref{modules-definition-koszul} を参照

\noindent
{\it Koszul 射}
\ref{more-morphisms-definition-lci} を参照

\noindent
{\it Koszul 射}
\ref{spaces-more-morphisms-definition-lci} を参照

\noindent
{\it Koszul-正則イデアル}
\ref{more-algebra-definition-regular-ideal} を参照

\noindent
{\it Koszul 正則埋め込み}
\ref{divisors-definition-regular-immersion} を参照

\noindent
{\it Koszul 正則埋め込み}
\ref{spaces-more-morphisms-definition-regular-immersion} を参照

\noindent
{\it Koszul-正則}
\ref{more-algebra-definition-koszul-regular-sequence} を参照

\noindent
{\it Koszul-正則}
\ref{divisors-definition-regular-ideal-sheaf} を参照

\noindent
{\it Koszul}
\ref{stacks-morphisms-definition-lci} を参照

\noindent
{\it $x$ における $X$ のクルル次元}
\ref{topology-definition-Krull} を参照

\noindent
{\it クルル次元}
\ref{topology-definition-Krull} を参照

\noindent
{\it クルル次元}
\ref{algebra-definition-Krull} を参照

\noindent
{\it $V$ の格子}
\ref{algebra-definition-lattice} を参照

\noindent
{\it $F$ に関して左非輪状}
\ref{derived-definition-derived-functor} を参照

\noindent
{\it 左随伴}
\ref{categories-definition-adjoint} を参照

\noindent
{\it 左許容}
\ref{derived-definition-admissible} を参照

\noindent
{\it 左導来可能}
\ref{derived-definition-everywhere-defined} を参照

\noindent
{\it 左導来関手 $LF$ が定義される}
\ref{derived-definition-right-derived-functor-defined} を参照

\noindent
{\it $F$ の左導来関手}
\ref{derived-definition-derived-functor} を参照

\noindent
{\it 左双対}
\ref{categories-definition-dual} を参照

\noindent
{\it 左完全}
\ref{categories-definition-exact} を参照

\noindent
{\it 左乗法系}
\ref{categories-definition-multiplicative-system} を参照

\noindent
{\it 左直交}
\ref{derived-definition-orthogonal} を参照

\noindent
{\it Leibniz 則}
\ref{algebra-definition-derivation} を参照

\noindent
{\it Leibniz 則}
\ref{modules-definition-derivation} を参照

\noindent
{\it Leibniz 則}
\ref{sites-modules-definition-derivation} を参照

\noindent
{\it 長さ}
\ref{topology-definition-Krull} を参照

\noindent
{\it 長さ}
\ref{algebra-definition-length} を参照

\noindent
{\it 長さ}
\ref{algebra-definition-chain-primes} を参照

\noindent
{\it 長さ}
\ref{exercises-definition-length} を参照

\noindent
{\it 上にある}
\ref{etale-cohomology-definition-geometric-point} を参照

\noindent
{\it 上にある}
\ref{artin-definition-formal-objects} を参照

\noindent
{\it $f$ に沿う $x$ の持ち上げ}
\ref{formal-defos-definition-lifts} を参照

\noindent
{\it 持ち上げ}
\ref{categories-definition-fibre-category} を参照

\noindent
{\it 持ち上げ}
\ref{categories-definition-fibre-category} を参照

\noindent
{\it 極限を保つ}
\ref{spaces-limits-definition-locally-finite-presentation} を参照

\noindent
{\it 極限を保つ}
\ref{spaces-limits-definition-locally-finite-presentation} を参照

\noindent
{\it 極限を保つ}
\ref{artin-definition-limit-preserving} を参照

\noindent
{\it 極限を保つ}
\ref{stacks-limits-definition-limit-preserving} を参照

\noindent
{\it 極限}
\ref{categories-definition-limit} を参照

\noindent
{\it 極限}
\ref{homology-definition-limit-spectral-sequence} を参照

\noindent
{\it 次数 $d$、次元 $r$ の線形系}
\ref{curves-definition-linear-series} を参照

\noindent
{\it 線形適切}
\ref{adequate-definition-adequate-functor} を参照

\noindent
{\it $\Omega$ において $k$ 上線形無関連}
\ref{fields-definition-linearly-disjoint} を参照

\noindent
{\it 線形位相をもつ}
\ref{more-algebra-definition-topological-ring} を参照

\noindent
{\it 線形位相をもつ}
\ref{more-algebra-definition-topological-ring} を参照

\noindent
{\it 滑らかエタール・サイト}
\ref{stacks-cohomology-definition-lisse-etale} を参照

\noindent
{\it 滑らか}
\ref{proetale-definition-adic} を参照

\noindent
{\it 滑らか}
\ref{trace-definition-l-adic-sheaf} を参照

\noindent
{\it 局所完全交叉射}
\ref{more-morphisms-definition-lci} を参照

\noindent
{\it 局所完全交叉射}
\ref{spaces-more-morphisms-definition-lci} を参照

\noindent
{\it 局所完全交叉射}
\ref{stacks-morphisms-definition-lci} を参照

\noindent
{\it $k$ 上の局所完全交叉}
\ref{algebra-definition-lci-field} を参照

\noindent
{\it $k$ 上の局所完全交叉}
\ref{morphisms-definition-syntomic} を参照

\noindent
{\it 局所完全交叉}
\ref{more-algebra-definition-local-complete-intersection} を参照

\noindent
{\it 局所完全交叉}
\ref{dpa-definition-lci} を参照

\noindent
{\it 局所環の局所準同型}
\ref{algebra-definition-local-ring} を参照

\noindent
{\it $\tau$-位相について局所的}
\ref{descent-definition-property-local} を参照

\noindent
{\it 局所同型}
\ref{proetale-definition-local-isomorphism} を参照

\noindent
{\it 局所 Lefschetz 数}
\ref{trace-definition-local-lefschetz-number} を参照

\noindent
{\it $\tau$-位相について基底上局所的}
\ref{descent-definition-property-morphisms-local} を参照

\noindent
{\it $\tau$-位相について基底上局所的}
\ref{spaces-descent-definition-property-morphisms-local} を参照

\noindent
{\it $\tau$-位相について始域上局所的}
\ref{descent-definition-property-morphisms-local-source} を参照

\noindent
{\it $\tau$-位相について始域上局所的}
\ref{spaces-descent-definition-property-morphisms-local-source} を参照

\noindent
{\it 局所環準同型 $\varphi : R \to S$}
\ref{algebra-definition-local-ring} を参照

\noindent
{\it $X$ の $x$ における局所環}
\ref{schemes-definition-locally-ringed-space} を参照

\noindent
{\it $\mathfrak q$ におけるファイバーの局所環}
\ref{algebra-definition-fibre} を参照

\noindent
{\it 局所環}
\ref{algebra-definition-local-ring} を参照

\noindent
{\it 局所化射}
\ref{sites-definition-localize} を参照

\noindent
{\it 局所化射}
\ref{sites-definition-localize-topos} を参照

\noindent
{\it 局所化射}
\ref{sites-modules-definition-localize-ringed-site} を参照

\noindent
{\it 局所化射}
\ref{sites-modules-definition-localize-ringed-topos} を参照

\noindent
{\it $S$ に関する $A$ の局所化}
\ref{algebra-definition-localization} を参照

\noindent
{\it 環付きサイト $(\mathcal{C}, \mathcal{O})$ の対象 $U$ における 局所化}
\ref{sites-modules-definition-localize-ringed-site} を参照

\noindent
{\it 環付きトポス $(\Sh(\mathcal{C}), \mathcal{O})$ の $\mathcal{F}$ における局所化}
\ref{sites-modules-definition-localize-ringed-topos} を参照

\noindent
{\it 対象 $U$ におけるサイト $\mathcal{C}$ の局所化}
\ref{sites-definition-localize} を参照

\noindent
{\it $\mathcal{F}$ におけるトポス $\Sh(\mathcal{C})$ の局所化}
\ref{sites-definition-localize-topos} を参照

\noindent
{\it 局所化}
\ref{algebra-definition-localization-module} を参照

\noindent
{\it 局所化第 $p$ Chern 類}
\ref{chow-definition-localized-chern} を参照

\noindent
{\it 局所化 Chern 指標}
\ref{chow-definition-localized-chern} を参照

\noindent
{\it 局所的に $P$}
\ref{properties-definition-locally-P} を参照

\noindent
{\it $\overline{x}$ において $K$ に関して局所非輪状}
\ref{etale-cohomology-definition-locally-acyclic} を参照

\noindent
{\it $K$ に関して局所非輪状}
\ref{etale-cohomology-definition-locally-acyclic} を参照

\noindent
{\it 局所非輪状}
\ref{etale-cohomology-definition-locally-acyclic} を参照

\noindent
{\it 局所アディック*}
\ref{formal-spaces-definition-types-formal-algebraic-spaces} を参照

\noindent
{\it 局所代数的 $k$-スキーム}
\ref{varieties-definition-algebraic-scheme} を参照

\noindent
{\it 局所閉埋め込み}
\ref{schemes-definition-immersion} を参照

\noindent
{\it 局所閉部分空間}
\ref{spaces-definition-immersion} を参照

\noindent
{\it 局所閉部分スタック}
\ref{stacks-properties-definition-substacks} を参照

\noindent
{\it 局所連結}
\ref{topology-definition-locally-connected} を参照

\noindent
{\it 局所定数}
\ref{sites-modules-definition-locally-constant} を参照

\noindent
{\it 局所定数}
\ref{etale-cohomology-definition-finite-locally-constant} を参照

\noindent
{\it 局所定数}
\ref{etale-cohomology-definition-finite-locally-constant} を参照

\noindent
{\it 局所定数}
\ref{etale-cohomology-definition-finite-locally-constant} を参照

\noindent
{\it 局所構成可能}
\ref{topology-definition-constructible} を参照

\noindent
{\it 局所可算添字付きかつ古典的}
\ref{formal-spaces-definition-types-formal-algebraic-spaces} を参照

\noindent
{\it 局所可算添字付き}
\ref{formal-spaces-definition-types-formal-algebraic-spaces} を参照

\noindent
{\it 局所有限}
\ref{topology-definition-locally-finite} を参照

\noindent
{\it 局所有限}
\ref{cohomology-definition-covering-locally-finite} を参照

\noindent
{\it 局所有限}
\ref{exercises-definition-graded-module} を参照

\noindent
{\it 局所自由}
\ref{algebra-definition-locally-free} を参照

\noindent
{\it 局所自由}
\ref{modules-definition-locally-free} を参照

\noindent
{\it 局所自由}
\ref{sites-modules-definition-site-local} を参照

\noindent
{\it 局所的に $r$ 個の切断で生成される}
\ref{sites-modules-definition-site-local} を参照

\noindent
{\it 局所的に切断で生成される}
\ref{modules-definition-locally-generated} を参照

\noindent
{\it 局所的に切断で生成される}
\ref{sites-modules-definition-site-local} を参照

\noindent
{\it 局所的に有限 Tor 次元をもつ}
\ref{cohomology-definition-tor-amplitude} を参照

\noindent
{\it 局所的に有限 Tor 次元をもつ}
\ref{sites-cohomology-definition-tor-amplitude} を参照

\noindent
{\it 局所冪零}
\ref{algebra-definition-locally-nilpotent-ideal} を参照

\noindent
{\it 局所 Noether}
\ref{topology-definition-noetherian} を参照

\noindent
{\it 局所 Noether}
\ref{properties-definition-noetherian} を参照

\noindent
{\it 局所 Noether}
\ref{formal-spaces-definition-types-formal-algebraic-spaces} を参照

\noindent
{\it 局所 Noether}
\ref{exercises-definition-Noetherian-scheme} を参照

\noindent
{\it $S$ 上局所有限表示}
\ref{spaces-limits-definition-locally-finite-presentation} を参照

\noindent
{\it 局所有限表示}
\ref{morphisms-definition-finite-presentation} を参照

\noindent
{\it 局所有限表示}
\ref{spaces-morphisms-definition-locally-finite-presentation} を参照

\noindent
{\it 局所有限表示}
\ref{spaces-limits-definition-locally-finite-presentation} を参照

\noindent
{\it 局所有限表示}
\ref{spaces-limits-definition-locally-finite-presentation} を参照

\noindent
{\it 局所有限表示}
\ref{stacks-morphisms-definition-locally-finite-presentation} を参照

\noindent
{\it 局所有限型}
\ref{morphisms-definition-finite-type} を参照

\noindent
{\it 局所有限型}
\ref{spaces-morphisms-definition-locally-finite-type} を参照

\noindent
{\it 局所有限型}
\ref{formal-spaces-definition-finite-type} を参照

\noindent
{\it 局所有限型}
\ref{stacks-morphisms-definition-locally-finite-type} を参照

\noindent
{\it 局所的に $P$ 型}
\ref{morphisms-definition-locally-P} を参照

\noindent
{\it 局所主閉部分スキーム}
\ref{divisors-definition-effective-Cartier-divisor} を参照

\noindent
{\it 局所主閉部分空間}
\ref{spaces-divisors-definition-effective-Cartier-divisor} を参照

\noindent
{\it 局所射影的}
\ref{properties-definition-locally-projective} を参照

\noindent
{\it 局所射影的}
\ref{morphisms-definition-projective} を参照

\noindent
{\it 局所射影的}
\ref{spaces-properties-definition-locally-projective} を参照

\noindent
{\it 局所準連接}
\ref{crystalline-definition-modules} を参照

\noindent
{\it 局所準連接}
\ref{stacks-sheaves-definition-locally-quasi-coherent} を参照

\noindent
{\it 局所準コンパクト}
\ref{topology-definition-locally-quasi-compact} を参照

\noindent
{\it 局所準有限}
\ref{morphisms-definition-quasi-finite} を参照

\noindent
{\it 局所準有限}
\ref{spaces-morphisms-definition-locally-quasi-finite} を参照

\noindent
{\it 局所準有限}
\ref{stacks-morphisms-definition-locally-quasi-finite} を参照

\noindent
{\it 局所準射影的}
\ref{morphisms-definition-quasi-projective} を参照

\noindent
{\it 局所環付きサイト}
\ref{sites-modules-definition-locally-ringed} を参照

\noindent
{\it 局所環付き空間 $(X, \mathcal{O}_X)$}
\ref{schemes-definition-locally-ringed-space} を参照

\noindent
{\it 局所環付き}
\ref{sites-modules-definition-locally-ringed-topos} を参照

\noindent
{\it $S$ 上局所分離的}
\ref{spaces-definition-separated} を参照

\noindent
{\it 局所分離的}
\ref{spaces-properties-definition-separated} を参照

\noindent
{\it 局所分離的}
\ref{spaces-properties-definition-separated} を参照

\noindent
{\it 局所分離的}
\ref{spaces-morphisms-definition-separated} を参照

\noindent
{\it 局所自明}
\ref{groupoids-definition-principal-homogeneous-space} を参照

\noindent
{\it 局所自明}
\ref{spaces-groupoids-definition-principal-homogeneous-space} を参照

\noindent
{\it 局所弱アディック}
\ref{formal-spaces-definition-types-formal-algebraic-spaces} を参照

\noindent
{\it 局所的}
\ref{properties-definition-property-local} を参照

\noindent
{\it 局所的}
\ref{morphisms-definition-property-local} を参照

\noindent
{\it 極大 Cohen--Macaulay}
\ref{algebra-definition-maximal-CM} を参照

\noindent
{\it McQuillan}
\ref{formal-spaces-definition-types-affine-formal-algebraic-space} を参照

\noindent
{\it 有理型関数}
\ref{divisors-definition-sheaf-meromorphic-functions} を参照

\noindent
{\it 有理型関数}
\ref{spaces-divisors-definition-sheaf-meromorphic-functions} を参照

\noindent
{\it $\mathcal{F}$ の有理型切断}
\ref{divisors-definition-meromorphic-section} を参照

\noindent
{\it $\mathcal{F}$ の有理型切断}
\ref{spaces-divisors-definition-meromorphic-section} を参照

\noindent
{\it 極小モデル}
\ref{models-definition-minimal-model} を参照

\noindent
{\it 最小多項式}
\ref{fields-definition-minimal-polynomial} を参照

\noindent
{\it 極小}
\ref{models-definition-type-minimal} を参照

\noindent
{\it 極小}
\ref{formal-defos-definition-minimal-versal} を参照

\noindent
{\it 極小}
\ref{formal-defos-definition-minimal-groupoid-in-functors} を参照

\noindent
{\it 最小半普遍的}
\ref{formal-defos-definition-minimal-versal} を参照

\noindent
{\it Mittag-Leffler 条件}
\ref{homology-definition-Mittag-Leffler} を参照

\noindent
{\it 加群の Mittag-Leffler 有向系}
\ref{algebra-definition-ML-inductive-system} を参照

\noindent
{\it Mittag-Leffler}
\ref{algebra-definition-ML-system} を参照

\noindent
{\it Mittag-Leffler}
\ref{algebra-definition-mittag-leffler-module} を参照

\noindent
{\it 混標数}
\ref{more-algebra-definition-mixed} を参照

\noindent
{\it ML}
\ref{homology-definition-Mittag-Leffler} を参照

\noindent
{\it $X$ の修正}
\ref{morphisms-definition-modification} を参照

\noindent
{\it $X$ の修正}
\ref{spaces-over-fields-definition-modification} を参照

\noindent
{\it 微分加群}
\ref{algebra-definition-differentials} を参照

\noindent
{\it 微分加群}
\ref{modules-definition-module-differentials} を参照

\noindent
{\it 微分加群}
\ref{sites-modules-definition-module-differentials} を参照

\noindent
{\it K\"ahler 微分加群}
\ref{algebra-definition-differentials} を参照

\noindent
{\it $k$ 階主部加群の層}
\ref{algebra-definition-module-principal-parts} を参照

\noindent
{\it $k$ 階主部加群の層}
\ref{modules-definition-module-principal-parts} を参照

\noindent
{\it $k$ 階主部加群の層}
\ref{sites-modules-definition-module-principal-parts} を参照

\noindent
{\it 加群値函手}
\ref{adequate-definition-module-valued-functor} を参照

\noindent
{\it 種数 $g$ の滑らかな固有曲線のモジュライ・スタック}
\ref{moduli-curves-definition-deligne-mumford-smooth} を参照

\noindent
{\it 滑らかな固有曲線のモジュライ・スタック}
\ref{moduli-curves-definition-deligne-mumford-smooth} を参照

\noindent
{\it 種数 $g$ の安定曲線のモジュライスタック}
\ref{moduli-curves-definition-deligne-mumford} を参照

\noindent
{\it 安定曲線のモジュライスタック}
\ref{moduli-curves-definition-deligne-mumford} を参照

\noindent
{\it モノイド圏}
\ref{categories-definition-monoidal-category} を参照

\noindent
{\it モノ射}
\ref{categories-definition-mono-epi} を参照

\noindent
{\it モノ射}
\ref{schemes-definition-monomorphism} を参照

\noindent
{\it モノ射}
\ref{spaces-morphisms-definition-monomorphism} を参照

\noindent
{\it モノ射}
\ref{formal-spaces-definition-monomorphism} を参照

\noindent
{\it モノ射}
\ref{stacks-properties-definition-monomorphism} を参照

\noindent
{\it フィルトレーション付き対象の射 $(A, F) \to (B, F)$}
\ref{homology-definition-filtered} を参照

\noindent
{\it 降下データの射 $(N, \varphi) \to (N', \varphi')$}
\ref{descent-definition-descent-datum-modules} を参照

\noindent
{\it $\mathcal{C}$ 上の関手における亜群の射 $(U, R, s, t, c) \to (U', R', s', t', c')$}
\ref{formal-defos-definition-groupoid-in-functors} を参照

\noindent
{\it 降下データの射 $\psi : (\mathcal{F}_i, \varphi_{ij}) \to (\mathcal{F}'_i, \varphi'_{ij})$}
\ref{descent-definition-descent-datum-quasi-coherent} を参照

\noindent
{\it 降下データの射 $\psi : (\mathcal{F}_i, \varphi_{ij}) \to (\mathcal{F}'_i, \varphi'_{ij})$}
\ref{spaces-descent-definition-descent-datum-quasi-coherent} を参照

\noindent
{\it $B$ 上の群代数空間の射 $\psi : (G, m) \to (G', m')$}
\ref{spaces-groupoids-definition-group-space} を参照

\noindent
{\it $S$ 上の群スキームの射 $\psi : (G, m) \to (G', m')$}
\ref{groupoids-definition-group-scheme} を参照

\noindent
{\it 降下データの射 $\psi : (V_i, \varphi_{ij}) \to (V'_i, \varphi'_{ij})$}
\ref{descent-definition-descent-datum-for-family-of-morphisms} を参照

\noindent
{\it 降下データの射 $\psi : (V_i, \varphi_{ij}) \to (V'_i, \varphi'_{ij})$}
\ref{spaces-descent-definition-descent-datum-for-family-of-morphisms} を参照

\noindent
{\it 降下データの射 $\psi : (X_i, \varphi_{ij}) \to (X'_i, \varphi'_{ij})$}
\ref{stacks-definition-descent-data} を参照

\noindent
{\it $\mathcal{B}$ 上の $\mathcal{O}$-加群の前層の射 $\varphi : \mathcal{F} \to \mathcal{G}$}
\ref{sheaves-definition-sheaf-modules-basis} を参照

\noindent
{\it $\mathcal{O}$-加群の前層の射 $\varphi : \mathcal{F} \to \mathcal{G}$}
\ref{sheaves-definition-presheaf-modules} を参照

\noindent
{\it $\mathcal{O}$-加群の前層の射 $\varphi : \mathcal{F} \to \mathcal{G}$}
\ref{sites-modules-definition-presheaf-modules} を参照

\noindent
{\it $\mathcal{B}$ 上の集合の前層の射 $\varphi : \mathcal{F} \to \mathcal{G}$}
\ref{sheaves-definition-presheaf-basis} を参照

\noindent
{\it 集合の前層の射 $\varphi : \mathcal{F} \to \mathcal{G}$ （$X$ 上）}
\ref{sheaves-definition-presheaf} を参照

\noindent
{\it $\mathcal{C}$ に値をとる前層の射 $\varphi : \mathcal{F} \to \mathcal{G}$}
\ref{sheaves-definition-presheaf-values-in-category} を参照

\noindent
{\it $\mathcal{B}$ 上の $\mathcal{C}$ に値をもつ前層の射 $\varphi : \mathcal{F} \to \mathcal{G}$}
\ref{sheaves-definition-sheaf-structures-basis} を参照

\noindent
{\it 形式対象の射 $a : \xi \to \eta$}
\ref{formal-defos-definition-formal-objects} を参照

\noindent
{\it $S$ 上の亜群スキームの射 $f : (U, R, s, t, c) \to (U', R', s', t', c')$}
\ref{groupoids-definition-groupoid} を参照

\noindent
{\it $B$ 上の代数空間における亜群の射 $f : (U, R, s, t, c) \to (U', R', s', t', c')$}
\ref{spaces-groupoids-definition-groupoid} を参照

\noindent
{\it $Y \to X$ に関する降下データの射 $f : (V/Y, \varphi) \to (V'/Y, \varphi')$}
\ref{spaces-descent-definition-descent-datum} を参照

\noindent
{\it $S$ 上の代数空間の射 $f : F \to F'$}
\ref{spaces-definition-morphism-algebraic-spaces} を参照

\noindent
{\it 射 $f : p \to p'$}
\ref{sites-definition-morphism-points} を参照

\noindent
{\it スキームの射 $f : X \to Y$（$S$ 上のもの）}
\ref{schemes-definition-base-change} を参照

\noindent
{\it $X \to S$ に関する降下データの射 $g : (V/X, \varphi) \to (V'/X, \varphi')$}
\ref{descent-definition-descent-datum} を参照

\noindent
{\it $\mathcal{U}$ から $\mathcal{V}$ への射}
\ref{sites-definition-morphism-coverings} を参照

\noindent
{\it $F$ から $G$ への $\delta$ 関手の射}
\ref{homology-definition-morphism-delta-functors} を参照

\noindent
{\it $\mathcal{G}$-トーサーの射}
\ref{cohomology-definition-torsor} を参照

\noindent
{\it $\mathcal{G}$-トーサーの射}
\ref{sites-cohomology-definition-torsor} を参照

\noindent
{\it $G$-加群の射}
\ref{etale-cohomology-definition-G-module-continuous} を参照

\noindent
{\it $G$-集合の射}
\ref{pione-definition-G-set-continuous} を参照

\noindent
{\it $n$-切断単体的対象の射}
\ref{simplicial-definition-truncated-simplicial-object} を参照

\noindent
{\it $R\text{-}G$-加群の射}
\ref{etale-cohomology-definition-G-module-continuous} を参照

\noindent
{\it エタール近傍の射}
\ref{etale-cohomology-definition-geometric-point} を参照

\noindent
{\it エタール近傍の射}
\ref{spaces-properties-definition-etale-neighbourhood} を参照

\noindent
{\it エタール近傍の射}
\ref{more-morphisms-definition-etale-neighbourhood} を参照

\noindent
{\it $X$ 上のアーベル前層の射}
\ref{sheaves-definition-abelian-presheaves} を参照

\noindent
{\it アフィン形式代数空間の射}
\ref{formal-spaces-definition-affine-formal-algebraic-space} を参照

\noindent
{\it アフィンスキームの射}
\ref{schemes-definition-affine-scheme} を参照

\noindent
{\it 錐の射}
\ref{constructions-definition-abstract-cone} を参照

\noindent
{\it 射 $U \to U'$}
\ref{simplicial-definition-cosimplicial-object} を参照

\noindent
{\it 微分対象の射}
\ref{homology-definition-differential-object} を参照

\noindent
{\it 分割冪スキームの射}
\ref{crystalline-definition-divided-power-scheme} を参照

\noindent
{\it $(S, \mathcal{I}, \gamma)$ に相対的な $X$ の分割冪肥厚化の射}
\ref{crystalline-definition-divided-power-thickening-X} を参照

\noindent
{\it 点線矢印の射}
\ref{categories-definition-dotted-arrows} を参照

\noindent
{\it 点線矢印の射}
\ref{stacks-morphisms-definition-fill-in-diagram} を参照

\noindent
{\it 基本エタール近傍の射}
\ref{decent-spaces-definition-elemenary-etale-neighbourhood} を参照

\noindent
{\it 完全対の射}
\ref{homology-definition-exact-couple} を参照

\noindent
{\it 拡大の射}
\ref{homology-definition-extension} を参照

\noindent
{\it $\mathcal{C}$ の終域を固定した射の族の $\mathcal{U}$ から $\mathcal{V}$ への射}
\ref{sites-definition-morphism-coverings} を参照

\noindent
{\it 形式代数空間の射}
\ref{formal-spaces-definition-formal-algebraic-space} を参照

\noindent
{\it 形式対象の射}
\ref{artin-definition-formal-objects} を参照

\noindent
{\it 関手の射}
\ref{categories-definition-transformation-functors} を参照

\noindent
{\it 芽の射}
\ref{descent-definition-germs} を参照

\noindent
{\it $(U, R, s, t, c)$ 上カルテジアンな亜群スキームの射}
\ref{groupoids-definition-cartesian-morphism} を参照

\noindent
{\it 持ち上げの射}
\ref{formal-defos-definition-lifts} を参照

\noindent
{\it 局所環付きサイトの射}
\ref{sites-modules-definition-morphism-locally-ringed-topoi} を参照

\noindent
{\it 局所環付き空間の射}
\ref{schemes-definition-locally-ringed-space} を参照

\noindent
{\it 局所環付きトポスの射}
\ref{sites-modules-definition-morphism-locally-ringed-topoi} を参照

\noindent
{\it 加群値函手の射}
\ref{adequate-definition-module-valued-functor} を参照

\noindent
{\it Postnikov 系の射}
\ref{derived-definition-postnikov-system} を参照

\noindent
{\it 前変形圏の射}
\ref{formal-defos-definition-predeformation-category} を参照

\noindent
{\it $\mathcal{X}$ 上の前層の射}
\ref{stacks-sheaves-definition-presheaves} を参照

\noindent
{\it 擬 $\mathcal{G}$-トーサーの射}
\ref{cohomology-definition-torsor} を参照

\noindent
{\it 擬 $\mathcal{G}$-トーサーの射}
\ref{sites-cohomology-definition-torsor} を参照

\noindent
{\it 環付きサイトの射}
\ref{sites-modules-definition-ringed-site} を参照

\noindent
{\it 環付き空間の射}
\ref{sheaves-definition-ringed-space} を参照

\noindent
{\it 環付きトポスの射}
\ref{sites-modules-definition-ringed-topos} を参照

\noindent
{\it スキームの射}
\ref{schemes-definition-scheme} を参照

\noindent
{\it $\mathcal{O}$-加群の層の射}
\ref{sheaves-definition-sheaf-modules} を参照

\noindent
{\it $\mathcal{O}$-加群の層の射}
\ref{sites-modules-definition-sheaf-modules} を参照

\noindent
{\it $\mathcal{B}$ 上の集合の層の射}
\ref{sheaves-definition-sheaf-basis} を参照

\noindent
{\it 集合の層の射}
\ref{sheaves-definition-sheaf} を参照

\noindent
{\it 射 $U \to U'$}
\ref{simplicial-definition-simplicial-object} を参照

\noindent
{\it 加群の導来圏の単体的系の射}
\ref{spaces-simplicial-definition-cartesian-derived-modules} を参照

\noindent
{\it 導来圏の単体的系の射}
\ref{spaces-simplicial-definition-cartesian-derived} を参照

\noindent
{\it サイトの射}
\ref{sites-definition-morphism-sites} を参照

\noindent
{\it スペクトル系列の射}
\ref{homology-definition-spectral-sequence} を参照

\noindent
{\it 肥厚化の射}
\ref{more-morphisms-definition-thickening} を参照

\noindent
{\it 肥厚化の射}
\ref{spaces-more-morphisms-definition-thickening} を参照

\noindent
{\it 肥厚化の射}
\ref{stacks-more-morphisms-definition-thickening} を参照

\noindent
{\it トポスの射}
\ref{sites-definition-topos} を参照

\noindent
{\it 三角の射}
\ref{derived-definition-triangle} を参照

\noindent
{\it $S$ 上のベクトル束の射}
\ref{constructions-definition-abstract-vector-bundle} を参照

\noindent
{\it 前層の射}
\ref{sites-definition-presheaves-sets} を参照

\noindent
{\it $\mathcal{Z}$ 上の厚化の射}
\ref{stacks-more-morphisms-definition-thickening} を参照

\noindent
{\it $B$ 上の肥厚化の射}
\ref{spaces-more-morphisms-definition-thickening} を参照

\noindent
{\it $S$ 上の肥厚化の射}
\ref{more-morphisms-definition-thickening} を参照

\noindent
{\it 型 $\mathcal{P}$ の射は $\tau$-被覆に対する降下を満たす}
\ref{descent-definition-descending-types-morphisms} を参照

\noindent
{\it 射}
\ref{sites-definition-presheaf} を参照

\noindent
{\it 射}
\ref{trace-definition-l-adic-sheaf} を参照

\noindent
{\it 多重交叉特異点}
\ref{curves-definition-multicross} を参照

\noindent
{\it $R$ の乗法的集合}
\ref{algebra-definition-multiplicative-subset} を参照

\noindent
{\it 乗法系}
\ref{categories-definition-multiplicative-system} を参照

\noindent
{\it 定義イデアル $I$ に関する $M$ の重複度}
\ref{intersection-definition-multiplicity} を参照

\noindent
{\it $\mathcal{F}$ における $Z$ の重複度}
\ref{spaces-chow-definition-cycle-associated-to-coherent-sheaf} を参照

\noindent
{\it $Y$ における $Z$ の重複度}
\ref{spaces-chow-definition-cycle-associated-to-closed-subscheme} を参照

\noindent
{\it $\mathcal{F}$ における $Z'$ の重複度}
\ref{chow-definition-cycle-associated-to-coherent-sheaf} を参照

\noindent
{\it $Z$ における $Z'$ の重複度}
\ref{chow-definition-cycle-associated-to-closed-subscheme} を参照

\noindent
{\it $\mathcal{X}$ の $x_0$ を通る形式的分枝の重複度}
\ref{stacks-geometry-definition-multiplicity-formal-branches} を参照

\noindent
{\it 重複度}
\ref{chow-definition-periodic-length} を参照

\noindent
{\it 重複度}
\ref{stacks-geometry-definition-multiplicity} を参照

\noindent
{\it N-1}
\ref{algebra-definition-N} を参照

\noindent
{\it N-2}
\ref{algebra-definition-N} を参照

\noindent
{\it Nagata 環}
\ref{algebra-definition-nagata} を参照

\noindent
{\it 永田スキーム}
\ref{properties-definition-nagata} を参照

\noindent
{\it $f$ の素朴な余接複体}
\ref{more-morphisms-definition-netherlander} を参照

\noindent
{\it $f$ の素朴な余接複体}
\ref{spaces-more-morphisms-definition-netherlander} を参照

\noindent
{\it 素朴な余接複体}
\ref{algebra-definition-naive-cotangent-complex} を参照

\noindent
{\it 素朴な余接複体}
\ref{modules-definition-naive-cotangent-complex} を参照

\noindent
{\it 素朴な余接複体}
\ref{modules-definition-cotangent-complex-morphism-ringed-topoi} を参照

\noindent
{\it 素朴な余接複体}
\ref{sites-modules-definition-naive-cotangent-complex} を参照

\noindent
{\it 素朴な余接複体}
\ref{sites-modules-definition-cotangent-complex-morphism-ringed-topoi} を参照

\noindent
{\it 素朴な障害理論}
\ref{artin-definition-naive-obstruction-theory} を参照

\noindent
{\it 素朴な rig-平坦}
\ref{restricted-definition-naively-rig-flat} を参照

\noindent
{\it 自然変換}
\ref{categories-definition-transformation-functors} を参照

\noindent
{\it 冪零}
\ref{algebra-definition-locally-nilpotent-ideal} を参照

\noindent
{\it 節点}
\ref{curves-definition-multicross} を参照

\noindent
{\it 節点}
\ref{curves-definition-nodal} を参照

\noindent
{\it Noether}
\ref{topology-definition-noetherian} を参照

\noindent
{\it Noether}
\ref{homology-definition-Noetherian} を参照

\noindent
{\it Noether}
\ref{homology-definition-Noetherian} を参照

\noindent
{\it Noether}
\ref{properties-definition-noetherian} を参照

\noindent
{\it Noether}
\ref{spaces-properties-definition-noetherian} を参照

\noindent
{\it Noether}
\ref{formal-spaces-definition-types-affine-formal-algebraic-space} を参照

\noindent
{\it Noether}
\ref{stacks-morphisms-definition-noetherian} を参照

\noindent
{\it Noether}
\ref{exercises-definition-Noetherian-space} を参照

\noindent
{\it Noether}
\ref{exercises-definition-Noetherian-scheme} を参照

\noindent
{\it 非退化}
\ref{crystalline-definition-F-crystal} を参照

\noindent
{\it $X$ の非特異射影モデル}
\ref{curves-definition-normal-projective-model} を参照

\noindent
{\it 非特異}
\ref{properties-definition-regular} を参照

\noindent
{\it 非自明解}
\ref{etale-cohomology-definition-Cr} を参照

\noindent
{\it $x$ で正規}
\ref{more-morphisms-definition-normal} を参照

\noindent
{\it 法束}
\ref{divisors-definition-normal-cone} を参照

\noindent
{\it 法束}
\ref{spaces-more-morphisms-definition-normal-cone} を参照

\noindent
{\it $F$ 上の $E$ の正規閉包}
\ref{fields-definition-normal-closure} を参照

\noindent
{\it 法錐 $C_ZX$}
\ref{divisors-definition-normal-cone} を参照

\noindent
{\it 法錐 $C_ZX$}
\ref{spaces-more-morphisms-definition-normal-cone} を参照

\noindent
{\it 正規交差因子}
\ref{etale-definition-normal-crossings} を参照

\noindent
{\it 正規射}
\ref{more-morphisms-definition-normal} を参照

\noindent
{\it $X$ の正規化 （$Y$ におけるもの）}
\ref{morphisms-definition-normalization-X-in-Y} を参照

\noindent
{\it $X$ の正規化 （$Y$ におけるもの）}
\ref{spaces-morphisms-definition-normalization-X-in-Y} を参照

\noindent
{\it 正規化}
\ref{morphisms-definition-normalization} を参照

\noindent
{\it 正規化}
\ref{spaces-morphisms-definition-normalization} を参照

\noindent
{\it 正規化}
\ref{stacks-morphisms-definition-normalization} を参照

\noindent
{\it $x$ における $X$ の正規化ブローアップ}
\ref{resolve-definition-normalized-blowup} を参照

\noindent
{\it $x$ における $X$ の正規化ブローアップ}
\ref{spaces-resolve-definition-normalized-blowup} を参照

\noindent
{\it 正規化されている}
\ref{formal-defos-definition-minimal-groupoid-in-functors} を参照

\noindent
{\it 正規}
\ref{fields-definition-normal} を参照

\noindent
{\it 正規}
\ref{fields-definition-separable-algebraic} を参照

\noindent
{\it 正規}
\ref{algebra-definition-domain-normal} を参照

\noindent
{\it 正規}
\ref{algebra-definition-ring-normal} を参照

\noindent
{\it 正規}
\ref{properties-definition-normal} を参照

\noindent
{\it ノルム}
\ref{fields-definition-trace-norm} を参照

\noindent
{\it 至る所疎}
\ref{topology-definition-nowhere-dense} を参照

\noindent
{\it 数体}
\ref{fields-definition-number-field} を参照

\noindent
{\it $A$ の枝数}
\ref{more-algebra-definition-number-of-branches} を参照

\noindent
{\it $X$ の $x$ における枝数}
\ref{properties-definition-number-of-branches} を参照

\noindent
{\it $X$ の $x$ における枝数}
\ref{decent-spaces-definition-number-of-branches} を参照

\noindent
{\it $x$ における $\mathcal{X}$ の幾何的枝数}
\ref{stacks-properties-definition-number-of-geometric-branches} を参照

\noindent
{\it $A$ の幾何的枝数}
\ref{more-algebra-definition-number-of-branches} を参照

\noindent
{\it $x$ における $X$ の幾何的枝数}
\ref{properties-definition-number-of-branches} を参照

\noindent
{\it $x$ における $X$ の幾何的枝数}
\ref{spaces-properties-definition-number-of-geometric-branches} を参照

\noindent
{\it 数値多項式}
\ref{algebra-definition-numerical-polynomial} を参照

\noindent
{\it 数値多項式}
\ref{exercises-definition-numerical-polynomial} を参照

\noindent
{\it $X$ に付随する数値型}
\ref{models-definition-numerical-type-model} を参照

\noindent
{\it 種数 $g$ の数値型}
\ref{models-definition-genus} を参照

\noindent
{\it 数値型}
\ref{models-definition-type} を参照

\noindent
{\it 障害加群}
\ref{artin-definition-obstruction-theory} を参照

\noindent
{\it 障害理論}
\ref{artin-definition-obstruction-theory} を参照

\noindent
{\it 障害}
\ref{artin-definition-obstruction-theory} を参照

\noindent
{\it $S$ に相対的有限表示をもつ}
\ref{more-morphisms-definition-relatively-finitely-presented-sheaf} を参照

\noindent
{\it 有限表示}
\ref{sites-modules-definition-site-local} を参照

\noindent
{\it 有限表示}
\ref{spaces-morphisms-definition-locally-finite-presentation} を参照

\noindent
{\it 有限表示}
\ref{stacks-morphisms-definition-locally-finite-presentation} を参照

\noindent
{\it 有限型}
\ref{sites-modules-definition-site-local} を参照

\noindent
{\it 有限型}
\ref{spaces-morphisms-definition-locally-finite-type} を参照

\noindent
{\it 有限型}
\ref{stacks-morphisms-definition-locally-finite-type} を参照

\noindent
{\it 岡族}
\ref{algebra-definition-oka-family} を参照

\noindent
{\it $x$ における $\mathcal{F}/X/S$ の一段階デヴィサージュ}
\ref{flat-definition-one-step-devissage-at-x} を参照

\noindent
{\it $s$ 上の $\mathcal{F}/X/S$ の一段階デヴィサージュ}
\ref{flat-definition-one-step-devissage} を参照

\noindent
{\it 開埋め込み}
\ref{sites-definition-immersion-topoi} を参照

\noindent
{\it 開埋め込み}
\ref{schemes-definition-immersion-locally-ringed-spaces} を参照

\noindent
{\it 開埋め込み}
\ref{schemes-definition-immersion} を参照

\noindent
{\it 開埋め込み}
\ref{spaces-definition-immersion} を参照

\noindent
{\it 開埋め込み}
\ref{stacks-properties-definition-immersion} を参照

\noindent
{\it 開部分群スキーム}
\ref{groupoids-definition-closed-subgroup-scheme} を参照

\noindent
{\it 開部分スキーム}
\ref{schemes-definition-immersion} を参照

\noindent
{\it $U$ に付随する $(X, \mathcal{O})$ の開部分空間}
\ref{sheaves-definition-restriction} を参照

\noindent
{\it $X$ の、$U$ に付随する開部分空間}
\ref{schemes-definition-open-subspace} を参照

\noindent
{\it 開部分空間}
\ref{spaces-definition-immersion} を参照

\noindent
{\it 開部分スタック}
\ref{stacks-properties-definition-substacks} を参照

\noindent
{\it 開部分トポス}
\ref{sites-definition-open-subtopos} を参照

\noindent
{\it 半普遍性の開性}
\ref{artin-definition-openness-versality} を参照

\noindent
{\it 半普遍性の開性}
\ref{artin-definition-openness-versality} を参照

\noindent
{\it 開}
\ref{morphisms-definition-open} を参照

\noindent
{\it 開}
\ref{trace-definition-open} を参照

\noindent
{\it 開}
\ref{spaces-morphisms-definition-open} を参照

\noindent
{\it 開}
\ref{stacks-morphisms-definition-open} を参照

\noindent
{\it 反対代数}
\ref{brauer-definition-opposite} を参照

\noindent
{\it 反対圏}
\ref{categories-definition-opposite} を参照

\noindent
{\it 反対微分次数付き代数}
\ref{dga-definition-opposite-dga} を参照

\noindent
{\it $R$ の軌道空間}
\ref{groupoids-quotients-definition-orbit-space} を参照

\noindent
{\it 軌道}
\ref{groupoids-quotients-definition-orbit} を参照

\noindent
{\it 軌道}
\ref{groupoids-quotients-definition-geometric-orbits} を参照

\noindent
{\it $R$ に沿う消滅位数}
\ref{algebra-definition-ord} を参照

\noindent
{\it $f$ の $Z$ に沿う 消滅次数}
\ref{divisors-definition-order-vanishing} を参照

\noindent
{\it $f$ の $Z$ に沿う 消滅次数}
\ref{spaces-over-fields-definition-order-vanishing} を参照

\noindent
{\it $s$ の $Z$ に沿う消滅次数}
\ref{divisors-definition-order-vanishing-meromorphic} を参照

\noindent
{\it $s$ の $Z$ に沿う消滅次数}
\ref{spaces-over-fields-definition-order-vanishing-meromorphic} を参照

\noindent
{\it 順序付き {\v C}ech 複体}
\ref{cohomology-definition-ordered-cech-complex} を参照

\noindent
{\it 通常二重点}
\ref{curves-definition-multicross} を参照

\noindent
{\it 通常二重点}
\ref{curves-definition-nodal} を参照

\noindent
{\it $k$ 上の $K$ の p-基底}
\ref{more-algebra-definition-p-basis} を参照

\noindent
{\it $k$ 上 p-独立}
\ref{more-algebra-definition-p-basis} を参照

\noindent
{\it $\tau$-位相に関して寄生的}
\ref{descent-definition-parasitic} を参照

\noindent
{\it 寄生的}
\ref{descent-definition-parasitic} を参照

\noindent
{\it 寄生的}
\ref{stacks-cohomology-definition-parasitic} を参照

\noindent
{\it 半順序}
\ref{categories-definition-directed-set} を参照

\noindent
{\it 半順序集合}
\ref{categories-definition-directed-set} を参照

\noindent
{\it 分割}
\ref{topology-definition-paritition} を参照

\noindent
{\it 部分}
\ref{topology-definition-paritition} を参照

\noindent
{\it $x$ で完全}
\ref{spaces-more-morphisms-definition-perfect} を参照

\noindent
{\it 完全閉包}
\ref{algebra-definition-perfection} を参照

\noindent
{\it $R$ に関して完全}
\ref{more-algebra-definition-relatively-perfect} を参照

\noindent
{\it $S$ に関して完全}
\ref{perfect-definition-relatively-perfect} を参照

\noindent
{\it $Y$ に関して完全}
\ref{spaces-more-morphisms-definition-relatively-perfect} を参照

\noindent
{\it 完全環写像}
\ref{more-algebra-definition-pseudo-coherent-perfect} を参照

\noindent
{\it 完全}
\ref{algebra-definition-perfect} を参照

\noindent
{\it 完全}
\ref{more-algebra-definition-perfect} を参照

\noindent
{\it 完全}
\ref{more-algebra-definition-perfect} を参照

\noindent
{\it 完全}
\ref{cohomology-definition-perfect} を参照

\noindent
{\it 完全}
\ref{cohomology-definition-perfect} を参照

\noindent
{\it 完全}
\ref{sites-cohomology-definition-perfect} を参照

\noindent
{\it 完全}
\ref{sites-cohomology-definition-perfect} を参照

\noindent
{\it 完全}
\ref{more-morphisms-definition-perfect} を参照

\noindent
{\it 完全}
\ref{trace-definition-perfect} を参照

\noindent
{\it 完全}
\ref{spaces-more-morphisms-definition-perfect} を参照

\noindent
{\it $T$ の ph 被覆}
\ref{topologies-definition-ph-covering} を参照

\noindent
{\it $X$ の ph 被覆}
\ref{spaces-topologies-definition-ph-covering} を参照

\noindent
{\it Picard 函手}
\ref{pic-definition-picard-functor} を参照

\noindent
{\it $A$ の Picard 群}
\ref{exercises-definition-class-group} を参照

\noindent
{\it $T$ の Picard 群}
\ref{models-definition-picard-group} を参照

\noindent
{\it $X$ の Picard 群}
\ref{exercises-definition-picard-group} を参照

\noindent
{\it Picard 群}
\ref{modules-definition-pic} を参照

\noindent
{\it Picard 群}
\ref{sites-modules-definition-pic} を参照

\noindent
{\it PID}
\ref{algebra-definition-PID} を参照

\noindent
{\it サイト $\mathcal{C}$ の点 $p$}
\ref{sites-definition-point} を参照

\noindent
{\it 点 $p$}
\ref{sites-definition-point-topology} を参照

\noindent
{\it トポス $\Sh(\mathcal{C})$ の点}
\ref{sites-definition-point-topos} を参照

\noindent
{\it 点}
\ref{spaces-properties-definition-points} を参照

\noindent
{\it 点}
\ref{stacks-properties-definition-points} を参照

\noindent
{\it 重み付け}
\ref{more-morphisms-definition-weighting} を参照

\noindent
{\it Postnikov 系}
\ref{derived-definition-postnikov-system} を参照

\noindent
{\it 前アディック}
\ref{more-algebra-definition-topological-ring} を参照

\noindent
{\it 前許容}
\ref{more-algebra-definition-topological-ring} を参照

\noindent
{\it 前同値関係}
\ref{groupoids-definition-equivalence-relation} を参照

\noindent
{\it 前同値関係}
\ref{spaces-groupoids-definition-equivalence-relation} を参照

\noindent
{\it 前関係}
\ref{groupoids-definition-equivalence-relation} を参照

\noindent
{\it 前関係}
\ref{spaces-groupoids-definition-equivalence-relation} を参照

\noindent
{\it 前三角圏}
\ref{derived-definition-triangulated-category} を参照

\noindent
{\it 前三角部分圏}
\ref{derived-definition-triangulated-subcategory} を参照

\noindent
{\it 前加法圏}
\ref{homology-definition-preadditive} を参照

\noindent
{\it 前変形圏}
\ref{formal-defos-definition-predeformation-category} を参照

\noindent
{\it 前順序集合}
\ref{categories-definition-directed-set} を参照

\noindent
{\it 前順序}
\ref{categories-definition-directed-set} を参照

\noindent
{\it $(U, R, s, t, c)$ による $\mathcal{F}$ の表示}
\ref{formal-defos-definition-presentation} を参照

\noindent
{\it 表示}
\ref{spaces-definition-presentation} を参照

\noindent
{\it 表示}
\ref{algebraic-definition-presentation} を参照

\noindent
{\it 任意の基底変換で保たれる}
\ref{schemes-definition-preserved-by-base-change} を参照

\noindent
{\it 任意の基底変換で保たれる}
\ref{schemes-definition-preserved-by-base-change} を参照

\noindent
{\it 基底変換で保たれる}
\ref{schemes-definition-preserved-by-base-change} を参照

\noindent
{\it 基底変換で保たれる}
\ref{schemes-definition-preserved-by-base-change} を参照

\noindent
{\it $\mathcal{B}$ 上の集合の前層 $\mathcal{F}$}
\ref{sheaves-definition-presheaf-basis} を参照

\noindent
{\it 集合の前層 $\mathcal{F}$（$X$ 上）}
\ref{sheaves-definition-presheaf} を参照

\noindent
{\it $\mathcal{C}$ に値をとる $X$ 上の前層 $\mathcal{F}$}
\ref{sheaves-definition-presheaf-values-in-category} を参照

\noindent
{\it $\mathcal{B}$ 上の $\mathcal{C}$ に値をもつ前層 $\mathcal{F}$}
\ref{sheaves-definition-sheaf-structures-basis} を参照

\noindent
{\it $\mathcal{B}$ 上の $\mathcal{O}$-加群の前層 $\mathcal{F}$}
\ref{sheaves-definition-sheaf-modules-basis} を参照

\noindent
{\it $\mathcal{O}$-加群の前層}
\ref{sheaves-definition-presheaf-modules} を参照

\noindent
{\it $\mathcal{O}$-加群の前層}
\ref{sites-modules-definition-presheaf-modules} を参照

\noindent
{\it $X$ 上のアーベル群の前層}
\ref{sheaves-definition-abelian-presheaves} を参照

\noindent
{\it $x$ から $y$ への同型の前層}
\ref{stacks-definition-mor-presheaf} を参照

\noindent
{\it $\mathcal{X}$ 上の加群の前層}
\ref{stacks-sheaves-definition-modules} を参照

\noindent
{\it $x$ から $y$ への射の前層}
\ref{stacks-definition-mor-presheaf} を参照

\noindent
{\it $\mathcal{C}$ 上の集合の前層}
\ref{categories-definition-presheaf} を参照

\noindent
{\it 集合の前層}
\ref{sites-definition-presheaves-sets} を参照

\noindent
{\it 集合の前層}
\ref{etale-cohomology-definition-presheaf} を参照

\noindent
{\it $\mathcal{X}$ 上の前層}
\ref{stacks-sheaves-definition-presheaves} を参照

\noindent
{\it 前層}
\ref{categories-definition-presheaf} を参照

\noindent
{\it 前層}
\ref{sites-definition-presheaf} を参照

\noindent
{\it 前安定曲線族}
\ref{moduli-curves-definition-prestable} を参照

\noindent
{\it 素因子}
\ref{divisors-definition-Weil-divisor} を参照

\noindent
{\it 素因子}
\ref{spaces-over-fields-definition-Weil-divisor} を参照

\noindent
{\it 素因子}
\ref{exercises-definition-divisor} を参照

\noindent
{\it $F$ の素体}
\ref{fields-definition-characteristic} を参照

\noindent
{\it 素元}
\ref{algebra-definition-irreducible-prime-element} を参照

\noindent
{\it $f$ に付随する主因子}
\ref{chow-definition-principal-divisor} を参照

\noindent
{\it $f$ に付随する主因子}
\ref{spaces-chow-definition-principal-divisor} を参照

\noindent
{\it 主等質 $G$-空間で $B$ 上にあるもの}
\ref{spaces-groupoids-definition-principal-homogeneous-space} を参照

\noindent
{\it 主等質空間}
\ref{groupoids-definition-principal-homogeneous-space} を参照

\noindent
{\it 主等質空間}
\ref{spaces-groupoids-definition-principal-homogeneous-space} を参照

\noindent
{\it 単項イデアル整域}
\ref{algebra-definition-PID} を参照

\noindent
{\it $f$ に付随する主 Weil 因子}
\ref{divisors-definition-principal-divisor} を参照

\noindent
{\it $f$ に付随する主 Weil 因子}
\ref{spaces-over-fields-definition-principal-divisor} を参照

\noindent
{\it $T$ の pro-エタール被覆}
\ref{proetale-definition-fpqc-covering} を参照

\noindent
{\it 積 $U \times V$ が存在する}
\ref{simplicial-definition-product-with-simplicial-set} を参照

\noindent
{\it $U$ と $V$ の積 $U \times V$}
\ref{simplicial-definition-product-with-simplicial-set} を参照

\noindent
{\it 直積圏}
\ref{categories-definition-product-category} を参照

\noindent
{\it $U$ と $V$ の積}
\ref{simplicial-definition-product} を参照

\noindent
{\it $U$ と $V$ の積}
\ref{simplicial-definition-product-cosimplicial-objects} を参照

\noindent
{\it 積}
\ref{categories-definition-products} を参照

\noindent
{\it 積}
\ref{categories-definition-product} を参照

\noindent
{\it 副有限群}
\ref{topology-definition-profinite-group} を参照

\noindent
{\it 副有限}
\ref{topology-definition-profinite} を参照

\noindent
{\it 射影 $n$ 次元空間（$\mathbf{Z}$ 上）}
\ref{constructions-definition-projective-space} を参照

\noindent
{\it 射影 $n$ 次元空間（$R$ 上）}
\ref{constructions-definition-projective-space} を参照

\noindent
{\it 射影 $n$ 次元空間（$S$ 上）}
\ref{constructions-definition-projective-space} を参照

\noindent
{\it $\mathcal{E}$ に付随する射影束}
\ref{constructions-definition-projective-bundle} を参照

\noindent
{\it 射影被覆}
\ref{dualizing-definition-projective-cover} を参照

\noindent
{\it 射影次元}
\ref{algebra-definition-finite-proj-dim} を参照

\noindent
{\it 射影包絡}
\ref{dualizing-definition-projective-cover} を参照

\noindent
{\it $A$ の射影分解}
\ref{derived-definition-projective-resolution} を参照

\noindent
{\it $K^\bullet$ の射影分解}
\ref{derived-definition-projective-resolution} を参照

\noindent
{\it $\mathcal{C}$ における $I$ 上の射影系}
\ref{categories-definition-system-over-poset} を参照

\noindent
{\it 射影多様体}
\ref{varieties-definition-variety-type} を参照

\noindent
{\it $[a, b]$ における射影振幅}
\ref{more-algebra-definition-projective-dimension} を参照

\noindent
{\it 射影的}
\ref{algebra-definition-projective} を参照

\noindent
{\it 射影的}
\ref{homology-definition-projective} を参照

\noindent
{\it 射影的}
\ref{morphisms-definition-projective} を参照

\noindent
{\it 射影的}
\ref{trace-definition-filtered} を参照

\noindent
{\it 固有相対サイクル}
\ref{relative-cycles-definition-proper} を参照

\noindent
{\it 固有多様体}
\ref{varieties-definition-variety-type} を参照

\noindent
{\it 性質 $\mathcal{P}$}
\ref{spaces-definition-relative-representable-property} を参照

\noindent
{\it 性質 $\mathcal{P}$}
\ref{bootstrap-definition-property-transformation} を参照

\noindent
{\it 性質 $\mathcal{P}$}
\ref{algebraic-definition-relative-representable-property} を参照

\noindent
{\it 固有}
\ref{topology-definition-proper-map} を参照

\noindent
{\it 固有}
\ref{morphisms-definition-proper} を参照

\noindent
{\it 固有}
\ref{spaces-morphisms-definition-proper} を参照

\noindent
{\it 固有}
\ref{formal-spaces-definition-proper} を参照

\noindent
{\it 固有}
\ref{stacks-morphisms-definition-proper} を参照

\noindent
{\it プロ表現可能}
\ref{formal-defos-definition-prorepresentable} を参照

\noindent
{\it プロ表現可能}
\ref{formal-defos-definition-prorepresentable-groupoid-in-functors} を参照

\noindent
{\it 擬 $\mathcal{G}$-トーサー}
\ref{cohomology-definition-torsor} を参照

\noindent
{\it 擬 $\mathcal{G}$-トーサー}
\ref{sites-cohomology-definition-torsor} を参照

\noindent
{\it 擬 $G$-トーサー}
\ref{groupoids-definition-pseudo-torsor} を参照

\noindent
{\it 擬 $G$-トーサー}
\ref{spaces-groupoids-definition-pseudo-torsor} を参照

\noindent
{\it 擬関手}
\ref{categories-definition-functor-into-2-category} を参照

\noindent
{\it 擬トーサー}
\ref{cohomology-definition-torsor} を参照

\noindent
{\it 擬トーサー}
\ref{sites-cohomology-definition-torsor} を参照

\noindent
{\it 擬鎖状}
\ref{stacks-geometry-definition-pseudo-catenary} を参照

\noindent
{\it $x$ で擬連接}
\ref{spaces-more-morphisms-definition-pseudo-coherent} を参照

\noindent
{\it $R$ に関して擬コヒーレント}
\ref{more-algebra-definition-relatively-pseudo-coherent} を参照

\noindent
{\it $R$ に関して擬コヒーレント}
\ref{more-algebra-definition-relatively-pseudo-coherent} を参照

\noindent
{\it $S$ に相対的に擬連接}
\ref{more-morphisms-definition-relative-pseudo-coherence} を参照

\noindent
{\it $S$ に相対的に擬連接}
\ref{more-morphisms-definition-relative-pseudo-coherence} を参照

\noindent
{\it $Y$ に関して擬連接}
\ref{spaces-more-morphisms-definition-relative-pseudo-coherence} を参照

\noindent
{\it $Y$ に関して擬連接}
\ref{spaces-more-morphisms-definition-relative-pseudo-coherence} を参照

\noindent
{\it 擬コヒーレント環写像}
\ref{more-algebra-definition-pseudo-coherent-perfect} を参照

\noindent
{\it 擬連接}
\ref{more-algebra-definition-pseudo-coherent} を参照

\noindent
{\it 擬連接}
\ref{more-algebra-definition-pseudo-coherent} を参照

\noindent
{\it 擬連接}
\ref{cohomology-definition-pseudo-coherent} を参照

\noindent
{\it 擬連接}
\ref{cohomology-definition-pseudo-coherent} を参照

\noindent
{\it 擬連接}
\ref{sites-cohomology-definition-pseudo-coherent} を参照

\noindent
{\it 擬連接}
\ref{sites-cohomology-definition-pseudo-coherent} を参照

\noindent
{\it 擬連接}
\ref{more-morphisms-definition-pseudo-coherent} を参照

\noindent
{\it 擬連接}
\ref{spaces-more-morphisms-definition-pseudo-coherent} を参照

\noindent
{\it $\mathcal{F}$ の引き戻し $x^{-1}\mathcal{F}$}
\ref{stacks-sheaves-definition-pullback} を参照

\noindent
{\it 引き戻し関手}
\ref{categories-definition-pullback-functor-fibred-category} を参照

\noindent
{\it 引き戻し関手}
\ref{stacks-definition-pullback-functor} を参照

\noindent
{\it 引き戻し関手}
\ref{descent-definition-pullback-functor} を参照

\noindent
{\it 引き戻し関手}
\ref{descent-definition-pullback-functor-family} を参照

\noindent
{\it 引き戻し関手}
\ref{spaces-descent-definition-pullback-functor} を参照

\noindent
{\it 引き戻し関手}
\ref{spaces-descent-definition-pullback-functor-family} を参照

\noindent
{\it $f$ に沿う $\mathcal{S}$ の引き戻し}
\ref{stacks-definition-pullback-stack} を参照

\noindent
{\it $D$ の $f$ による引き戻しは定義される}
\ref{divisors-definition-pullback-effective-Cartier-divisor} を参照

\noindent
{\it $D$ の $f$ による引き戻しは定義される}
\ref{spaces-divisors-definition-pullback-effective-Cartier-divisor} を参照

\noindent
{\it $f$ による $S$ の引き戻し}
\ref{sites-definition-pullback-sieve} を参照

\noindent
{\it 有効 Cartier 因子の引き戻し}
\ref{divisors-definition-pullback-effective-Cartier-divisor} を参照

\noindent
{\it 有効 Cartier 因子の引き戻し}
\ref{spaces-divisors-definition-pullback-effective-Cartier-divisor} を参照

\noindent
{\it $f$ に対して有理型関数の引き戻しが定義される}
\ref{divisors-definition-pullback-meromorphic-sections} を参照

\noindent
{\it $f$ に対して有理型関数の引き戻しが定義される}
\ref{spaces-divisors-definition-pullback-meromorphic-sections} を参照

\noindent
{\it 引き戻し}
\ref{sheaves-definition-pushforward} を参照

\noindent
{\it 引き戻し}
\ref{sites-modules-definition-pushforward} を参照

\noindent
{\it 引き戻し}
\ref{groupoids-definition-restrict-relation} を参照

\noindent
{\it 引き戻し}
\ref{etale-cohomology-definition-inverse-image} を参照

\noindent
{\it 引き戻し}
\ref{spaces-groupoids-definition-restrict-relation} を参照

\noindent
{\it $y$ 上で純粋}
\ref{spaces-flat-definition-pure} を参照

\noindent
{\it $y$ 上で純粋}
\ref{spaces-flat-definition-pure} を参照

\noindent
{\it $X_s$ に沿って純粋}
\ref{flat-definition-pure} を参照

\noindent
{\it $X_s$ に沿って純粋}
\ref{flat-definition-pure} を参照

\noindent
{\it 純拡大加群}
\ref{adequate-definition-pure-ext} を参照

\noindent
{\it 純入射分解}
\ref{adequate-definition-pure-resolution} を参照

\noindent
{\it 純入射}
\ref{adequate-definition-pure} を参照

\noindent
{\it 純射影分解}
\ref{adequate-definition-pure-resolution} を参照

\noindent
{\it 純射影}
\ref{adequate-definition-pure} を参照

\noindent
{\it $S$ に関して純粋}
\ref{flat-definition-pure} を参照

\noindent
{\it $S$ に関して純粋}
\ref{flat-definition-pure} を参照

\noindent
{\it $Y$ に関して純粋}
\ref{spaces-flat-definition-pure} を参照

\noindent
{\it $Y$ に関して純粋}
\ref{spaces-flat-definition-pure} を参照

\noindent
{\it 純非分離的}
\ref{fields-definition-purely-inseparable} を参照

\noindent
{\it 純非分離的}
\ref{fields-definition-purely-inseparable} を参照

\noindent
{\it 純非分離的}
\ref{fields-definition-purely-inseparable} を参照

\noindent
{\it 純非分離的}
\ref{fields-definition-separable-algebraic} を参照

\noindent
{\it 純超越拡大}
\ref{fields-definition-transcendence} を参照

\noindent
{\it 純}
\ref{algebra-definition-pure-ideal} を参照

\noindent
{\it $f$ に沿う $\mathcal{S}$ の順像}
\ref{stacks-definition-pushforward-stack} を参照

\noindent
{\it 順像}
\ref{sheaves-definition-pushforward} を参照

\noindent
{\it 順像}
\ref{sites-definition-pushforward-algebraic-structures} を参照

\noindent
{\it 順像}
\ref{sites-modules-definition-pushforward} を参照

\noindent
{\it 順像}
\ref{chow-definition-proper-pushforward} を参照

\noindent
{\it 順像}
\ref{etale-cohomology-definition-direct-image-presheaf} を参照

\noindent
{\it 順像}
\ref{etale-cohomology-definition-direct-image-sheaf} を参照

\noindent
{\it 順像}
\ref{spaces-chow-definition-proper-pushforward} を参照

\noindent
{\it $U$ 上の $V$ と $W$ の押し出し}
\ref{simplicial-definition-push-out} を参照

\noindent
{\it 押し出し}
\ref{categories-definition-pushouts} を参照

\noindent
{\it qc 被覆}
\ref{sites-cohomology-definition-covering-LC} を参照

\noindent
{\it 準アフィン}
\ref{properties-definition-quasi-affine} を参照

\noindent
{\it 準アフィン}
\ref{morphisms-definition-quasi-affine} を参照

\noindent
{\it 準アフィン}
\ref{spaces-morphisms-definition-quasi-affine} を参照

\noindent
{\it 準連接 $\mathcal{O}_\mathcal{X}$-加群}
\ref{stacks-sheaves-definition-quasi-coherent} を参照

\noindent
{\it $(U, R, s, t, c)$ 上の準連接加群}
\ref{groupoids-definition-groupoid-module} を参照

\noindent
{\it $(U, R, s, t, c)$ 上の準連接加群}
\ref{spaces-groupoids-definition-groupoid-module} を参照

\noindent
{\it $\mathcal{X}$ 上の準連接加群}
\ref{stacks-sheaves-definition-quasi-coherent} を参照

\noindent
{\it 準連接 $\mathcal{O}_X$-加群の層}
\ref{modules-definition-quasi-coherent} を参照

\noindent
{\it 準連接}
\ref{sites-modules-definition-site-local} を参照

\noindent
{\it 準連接}
\ref{etale-cohomology-definition-ringed-site} を参照

\noindent
{\it 準連接}
\ref{crystalline-definition-modules} を参照

\noindent
{\it 準連接}
\ref{spaces-properties-definition-quasi-coherent} を参照

\noindent
{\it 準連接}
\ref{exercises-definition-quasi-coherent} を参照

\noindent
{\it 準コンパクト}
\ref{topology-definition-quasi-compact} を参照

\noindent
{\it 準コンパクト}
\ref{topology-definition-quasi-compact} を参照

\noindent
{\it 準コンパクト}
\ref{sites-definition-quasi-compact} を参照

\noindent
{\it 準コンパクト}
\ref{sites-definition-quasi-compact-topos} を参照

\noindent
{\it 準コンパクト}
\ref{sites-definition-quasi-compact-topos} を参照

\noindent
{\it 準コンパクト}
\ref{schemes-definition-quasi-compact} を参照

\noindent
{\it 準コンパクト}
\ref{spaces-properties-definition-quasi-compact} を参照

\noindent
{\it 準コンパクト}
\ref{spaces-morphisms-definition-quasi-compact} を参照

\noindent
{\it 準コンパクト}
\ref{formal-spaces-definition-quasi-compact} を参照

\noindent
{\it 準コンパクト}
\ref{formal-spaces-definition-quasi-compact-morphism} を参照

\noindent
{\it 準コンパクト}
\ref{stacks-properties-definition-quasi-compact} を参照

\noindent
{\it 準コンパクト}
\ref{stacks-morphisms-definition-quasi-compact} を参照

\noindent
{\it 準コンパクト}
\ref{exercises-definition-quasi-compact} を参照

\noindent
{\it $S$ 上準 DM}
\ref{stacks-morphisms-definition-absolute-separated} を参照

\noindent
{\it 準 DM}
\ref{stacks-morphisms-definition-separated} を参照

\noindent
{\it 準 DM}
\ref{stacks-morphisms-definition-absolute-separated} を参照

\noindent
{\it 準優秀}
\ref{more-algebra-definition-excellent} を参照

\noindent
{\it 準優秀}
\ref{morphisms-definition-excellent} を参照

\noindent
{\it $\mathfrak q$ で準有限}
\ref{algebra-definition-quasi-finite} を参照

\noindent
{\it $x$ において準有限}
\ref{spaces-morphisms-definition-locally-quasi-finite} を参照

\noindent
{\it 点 $x \in X$ において準有限}
\ref{morphisms-definition-quasi-finite} を参照

\noindent
{\it 準有限}
\ref{algebra-definition-quasi-finite} を参照

\noindent
{\it 準有限}
\ref{morphisms-definition-quasi-finite} を参照

\noindent
{\it 準有限}
\ref{spaces-morphisms-definition-locally-quasi-finite} を参照

\noindent
{\it 準有限}
\ref{stacks-morphisms-definition-quasi-finite} を参照

\noindent
{\it 擬逆関手}
\ref{categories-definition-equivalence-categories} を参照

\noindent
{\it 擬同型}
\ref{homology-definition-quasi-isomorphism} を参照

\noindent
{\it 擬同型}
\ref{homology-definition-quasi-isomorphism-cochain} を参照

\noindent
{\it 準等自明}
\ref{groupoids-definition-principal-homogeneous-space} を参照

\noindent
{\it 準等自明}
\ref{spaces-groupoids-definition-principal-homogeneous-space} を参照

\noindent
{\it 準射影多様体}
\ref{varieties-definition-variety-type} を参照

\noindent
{\it 準射影的}
\ref{morphisms-definition-quasi-projective} を参照

\noindent
{\it 準固有}
\ref{topology-definition-proper-map} を参照

\noindent
{\it 準正則イデアル}
\ref{more-algebra-definition-regular-ideal} を参照

\noindent
{\it 準正則埋め込み}
\ref{divisors-definition-regular-immersion} を参照

\noindent
{\it 準正則埋め込み}
\ref{spaces-more-morphisms-definition-regular-immersion} を参照

\noindent
{\it 準正則列}
\ref{algebra-definition-quasi-regular-sequence} を参照

\noindent
{\it 準正則}
\ref{divisors-definition-regular-ideal-sheaf} を参照

\noindent
{\it $S$ 上準分離}
\ref{spaces-definition-separated} を参照

\noindent
{\it $S$ 上準分離}
\ref{stacks-morphisms-definition-absolute-separated} を参照

\noindent
{\it 準分離}
\ref{schemes-definition-separated} を参照

\noindent
{\it 準分離}
\ref{schemes-definition-separated} を参照

\noindent
{\it 準分離}
\ref{spaces-properties-definition-separated} を参照

\noindent
{\it 準分離}
\ref{spaces-properties-definition-separated} を参照

\noindent
{\it 準分離}
\ref{spaces-morphisms-definition-separated} を参照

\noindent
{\it 準分離}
\ref{formal-spaces-definition-separated} を参照

\noindent
{\it 準分離}
\ref{formal-spaces-definition-separated-morphism} を参照

\noindent
{\it 準分離}
\ref{stacks-morphisms-definition-separated} を参照

\noindent
{\it 準分離}
\ref{stacks-morphisms-definition-absolute-separated} を参照

\noindent
{\it 準ソバー}
\ref{topology-definition-generic-point} を参照

\noindent
{\it $u$ 上準分裂}
\ref{spaces-more-groupoids-definition-split-at-point} を参照

\noindent
{\it $u$ 上の $R$ の準分裂}
\ref{spaces-more-groupoids-definition-split-at-point} を参照

\noindent
{\it 商圏 $\mathcal{D}/\mathcal{B}$}
\ref{derived-definition-quotient-category} を参照

\noindent
{\it 商となる，群亜群をファイバーとする余ファイバー圏 $[U/R] \to \mathcal{C}$}
\ref{formal-defos-definition-quotient} を参照

\noindent
{\it 商フィルトレーション}
\ref{homology-definition-filtered} を参照

\noindent
{\it 商関手}
\ref{derived-definition-quotient-category} を参照

\noindent
{\it 商射 $U \to [U/R]$}
\ref{formal-defos-definition-quotient} を参照

\noindent
{\it $G$ による $U$ の商}
\ref{spaces-definition-quotient} を参照

\noindent
{\it 代数空間で表現可能な商}
\ref{spaces-groupoids-definition-representable-quotient} を参照

\noindent
{\it 代数空間で表現可能な商}
\ref{spaces-groupoids-definition-representable-quotient} を参照

\noindent
{\it 商層 $U/R$}
\ref{groupoids-definition-quotient-sheaf} を参照

\noindent
{\it 商層 $U/R$}
\ref{spaces-groupoids-definition-quotient-sheaf} を参照

\noindent
{\it 商スタック}
\ref{spaces-groupoids-definition-quotient-stack} を参照

\noindent
{\it 商スタック}
\ref{spaces-groupoids-definition-quotient-stack} を参照

\noindent
{\it 商}
\ref{homology-definition-injective-surjective} を参照

\noindent
{\it 純非分離}
\ref{morphisms-definition-universally-injective} を参照

\noindent
{\it 純非分離}
\ref{spaces-more-morphisms-definition-radicial} を参照

\noindent
{\it 分岐指数}
\ref{more-algebra-definition-extension-discrete-valuation-rings} を参照

\noindent
{\it 階数 $r$}
\ref{sites-modules-definition-invertible-sheaf} を参照

\noindent
{\it 階数}
\ref{algebra-definition-rank} を参照

\noindent
{\it 階数}
\ref{more-algebra-definition-rank} を参照

\noindent
{\it 階数}
\ref{morphisms-definition-finite-locally-free} を参照

\noindent
{\it 階数}
\ref{divisors-definition-rank} を参照

\noindent
{\it 階数}
\ref{spaces-morphisms-definition-finite-locally-free} を参照

\noindent
{\it $X$ 上の有理関数}
\ref{morphisms-definition-rational-function} を参照

\noindent
{\it $X$ 上の有理関数}
\ref{spaces-morphisms-definition-rational-function} を参照

\noindent
{\it $X$ から $Y$ への有理写像}
\ref{morphisms-definition-rational-map} を参照

\noindent
{\it $X$ から $Y$ への有理写像}
\ref{spaces-morphisms-definition-rational-map} を参照

\noindent
{\it 零と有理同値}
\ref{chow-definition-rational-equivalence} を参照

\noindent
{\it 零と有理同値}
\ref{spaces-chow-definition-rational-equivalence} を参照

\noindent
{\it 有理同値}
\ref{chow-definition-rational-equivalence} を参照

\noindent
{\it 有理同値}
\ref{spaces-chow-definition-rational-equivalence} を参照

\noindent
{\it 妥当}
\ref{decent-spaces-definition-very-reasonable} を参照

\noindent
{\it 妥当}
\ref{decent-spaces-definition-relative-conditions} を参照

\noindent
{\it 被約誘導代数空間構造}
\ref{spaces-properties-definition-reduced-induced-space} を参照

\noindent
{\it 誘導被約代数スタック構造}
\ref{stacks-properties-definition-reduced-induced-stack} を参照

\noindent
{\it 被約誘導スキーム構造}
\ref{schemes-definition-reduced-induced-scheme} を参照

\noindent
{\it 被約}
\ref{schemes-definition-reduced} を参照

\noindent
{\it $\mathcal{X}$ の被約化 $\mathcal{X}_{red}$}
\ref{stacks-properties-definition-reduced-induced-stack} を参照

\noindent
{\it $X$ の被約化 $X_{red}$}
\ref{schemes-definition-reduced-induced-scheme} を参照

\noindent
{\it $X$ の被約化 $X_{red}$}
\ref{spaces-properties-definition-reduced-induced-space} を参照

\noindent
{\it $A$ について有理特異点への還元が可能である}
\ref{resolve-definition-reduce-to-rational} を参照

\noindent
{\it Rees 代数}
\ref{algebra-definition-blow-up} を参照

\noindent
{\it 細分}
\ref{sites-definition-morphism-coverings} を参照

\noindent
{\it 細分する}
\ref{topology-definition-paritition} を参照

\noindent
{\it 反射包}
\ref{more-algebra-definition-reflexive-hull} を参照

\noindent
{\it 反射包}
\ref{divisors-definition-reflexive} を参照

\noindent
{\it 反射的}
\ref{more-algebra-definition-reflexive} を参照

\noindent
{\it 反射的}
\ref{divisors-definition-reflexive} を参照

\noindent
{\it 点 $x$ で正則}
\ref{more-morphisms-definition-regular} を参照

\noindent
{\it 正則イデアル}
\ref{more-algebra-definition-regular-ideal} を参照

\noindent
{\it 正則埋め込み}
\ref{divisors-definition-regular-immersion} を参照

\noindent
{\it 余次元 $\leq k$ で正則}
\ref{algebra-definition-conditions} を参照

\noindent
{\it 余次元 $k$ で正則}
\ref{properties-definition-Rk} を参照

\noindent
{\it 正則局所環}
\ref{algebra-definition-regular-local} を参照

\noindent
{\it 正則点集合}
\ref{properties-definition-singular-locus} を参照

\noindent
{\it 正則射}
\ref{more-morphisms-definition-regular} を参照

\noindent
{\it 正則切断}
\ref{divisors-definition-regular-section} を参照

\noindent
{\it 正則切断}
\ref{spaces-divisors-definition-regular-section} を参照

\noindent
{\it 正則列}
\ref{algebra-definition-regular-sequence} を参照

\noindent
{\it 正則パラメータ系}
\ref{algebra-definition-regular-local} を参照

\noindent
{\it 正則}
\ref{algebra-definition-regular} を参照

\noindent
{\it 正則}
\ref{homology-definition-bounded-ss} を参照

\noindent
{\it 正則}
\ref{more-algebra-definition-regular} を参照

\noindent
{\it 正則}
\ref{properties-definition-regular} を参照

\noindent
{\it 正則}
\ref{divisors-definition-regular-ideal-sheaf} を参照

\noindent
{\it 正則}
\ref{divisors-definition-regular-meromorphic-section} を参照

\noindent
{\it 正則}
\ref{spaces-divisors-definition-regular-meromorphic-section} を参照

\noindent
{\it 関係}
\ref{algebra-definition-relation} を参照

\noindent
{\it 関係}
\ref{groupoids-definition-equivalence-relation} を参照

\noindent
{\it 関係}
\ref{spaces-groupoids-definition-equivalence-relation} を参照

\noindent
{\it 相対 $H_1$-正則埋め込み}
\ref{divisors-definition-relative-H1-regular-immersion} を参照

\noindent
{\it $X/S$ 上の相対 $r$-サイクル}
\ref{relative-cycles-definition-relative-cycles} を参照

\noindent
{\it $S$ 上の $X$ における $\mathcal{F}$ の相対随伴点集合}
\ref{divisors-definition-relative-assassin} を参照

\noindent
{\it $S/R$ 上の $N$ の相対随伴素イデアル集合}
\ref{algebra-definition-relative-assassin} を参照

\noindent
{\it 相対余接空間}
\ref{formal-defos-definition-tangent-space-ring} を参照

\noindent
{\it 相対次元が $\leq d$ である（点 $x$ において）}
\ref{morphisms-definition-relative-dimension-d} を参照

\noindent
{\it 相対次元 $\leq d$}
\ref{morphisms-definition-relative-dimension-d} を参照

\noindent
{\it 相対次元 $\leq d$}
\ref{spaces-morphisms-definition-relative-dimension} を参照

\noindent
{\it 相対次元 $d$}
\ref{morphisms-definition-relative-dimension-d} を参照

\noindent
{\it 相対次元 $d$}
\ref{spaces-morphisms-definition-relative-dimension} を参照

\noindent
{\it $\mathfrak q$ における $S/R$ の相対次元}
\ref{algebra-definition-relative-dimension} を参照

\noindent
{\it 相対次元}
\ref{algebra-definition-relative-dimension} を参照

\noindent
{\it 相対次元}
\ref{stacks-geometry-definition-relative-dimension-for-stacks} を参照

\noindent
{\it 相対双対化複体}
\ref{dualizing-definition-relative-dualizing-complex} を参照

\noindent
{\it 相対双対化複体}
\ref{duality-definition-relative-dualizing-complex} を参照

\noindent
{\it 相対双対化複体}
\ref{spaces-duality-definition-relative-dualizing-proper-flat} を参照

\noindent
{\it 相対双対化層}
\ref{moduli-curves-definition-relative-dualizing-sheaf} を参照

\noindent
{\it 相対有効 Cartier 因子}
\ref{divisors-definition-relative-effective-Cartier-divisor} を参照

\noindent
{\it 相対有効 Cartier 因子}
\ref{spaces-divisors-definition-relative-effective-Cartier-divisor} を参照

\noindent
{\it $X/S$ の相対 Frobenius 射}
\ref{varieties-definition-relative-frobenius} を参照

\noindent
{\it 相対的大域完全交叉}
\ref{algebra-definition-relative-global-complete-intersection} を参照

\noindent
{\it $\mathcal{A}$ の $S$ 上の相対斉次スペクトル}
\ref{constructions-definition-relative-proj} を参照

\noindent
{\it $X$ 上の $\mathcal{A}$ の相対斉次スペクトル}
\ref{spaces-divisors-definition-relative-proj} を参照

\noindent
{\it $\mathcal{S}'$ 上の $\mathcal{S}$ の相対惰性}
\ref{categories-definition-inertia-fibred-category} を参照

\noindent
{\it $\mathcal{A}$ の $S$ 上の相対 Proj}
\ref{constructions-definition-relative-proj} を参照

\noindent
{\it $X$ 上の $\mathcal{A}$ の相対 Proj}
\ref{spaces-divisors-definition-relative-proj} を参照

\noindent
{\it 相対準正則埋め込み}
\ref{divisors-definition-relative-H1-regular-immersion} を参照

\noindent
{\it $x$ の相対自己同型の層}
\ref{stacks-morphisms-definition-isom} を参照

\noindent
{\it $x_1$ から $x_2$ への相対同型の層}
\ref{stacks-morphisms-definition-isom} を参照

\noindent
{\it $\mathcal{A}$ の $S$ 上の相対スペクトル}
\ref{constructions-definition-relative-spec} を参照

\noindent
{\it $X$ 上の $\mathcal{A}$ の相対スペクトル}
\ref{spaces-morphisms-definition-relative-spec} を参照

\noindent
{\it $S$ 上の $X$ における $\mathcal{F}$ の相対弱随伴点集合}
\ref{divisors-definition-relative-weak-assassin} を参照

\noindent
{\it $\mathcal{F}$ の $X$ における $Y$ 上の相対弱随伴点集合}
\ref{spaces-divisors-definition-relative-weak-assassin} を参照

\noindent
{\it 相対的に豊富}
\ref{morphisms-definition-relatively-ample} を参照

\noindent
{\it 相対的に豊富}
\ref{spaces-divisors-definition-relatively-ample} を参照

\noindent
{\it 相対的に極限を保つ}
\ref{spaces-limits-definition-locally-finite-presentation} を参照

\noindent
{\it 互いに素}
\ref{fields-definition-relatively-prime} を参照

\noindent
{\it 相対的に非常に豊富}
\ref{morphisms-definition-very-ample} を参照

\noindent
{\it スキームにより表現可能}
\ref{schemes-definition-representable-functor} を参照

\noindent
{\it 代数空間によって表現可能}
\ref{bootstrap-definition-morphism-representable-by-spaces} を参照

\noindent
{\it 代数空間によって表現可能}
\ref{algebraic-definition-representable-by-algebraic-spaces} を参照

\noindent
{\it $S$ 上の代数空間によって表現可能}
\ref{algebraic-definition-representable-by-algebraic-space} を参照

\noindent
{\it 開埋め込みにより表現可能}
\ref{schemes-definition-representable-by-open-immersions} を参照

\noindent
{\it 表現可能な商}
\ref{groupoids-definition-representable-quotient} を参照

\noindent
{\it 表現可能な商}
\ref{groupoids-definition-representable-quotient} を参照

\noindent
{\it 表現可能な商}
\ref{spaces-groupoids-definition-representable-quotient} を参照

\noindent
{\it 表現可能な商}
\ref{spaces-groupoids-definition-representable-quotient} を参照

\noindent
{\it 表現可能層}
\ref{sites-definition-representable-sheaf} を参照

\noindent
{\it 表現可能}
\ref{categories-definition-representable-functor} を参照

\noindent
{\it 表現可能}
\ref{categories-definition-representable-morphism} を参照

\noindent
{\it 表現可能}
\ref{categories-definition-representable-map-presheaves} を参照

\noindent
{\it 表現可能}
\ref{categories-definition-representable-fibred-category} を参照

\noindent
{\it 表現可能}
\ref{categories-definition-representable-map-categories-fibred-in-groupoids} を参照

\noindent
{\it 表現可能}
\ref{schemes-definition-representable-functor} を参照

\noindent
{\it 表現可能}
\ref{formal-defos-definition-representable} を参照

\noindent
{\it 剰余次数}
\ref{more-algebra-definition-extension-discrete-valuation-rings} を参照

\noindent
{\it 剰余次数}
\ref{more-algebra-definition-extension-valuation-rings} を参照

\noindent
{\it $x$ における $\mathcal{X}$ の剰余ガーベが存在する}
\ref{stacks-properties-definition-residual-gerbe} を参照

\noindent
{\it $x$ における $\mathcal{X}$ の剰余ガーベ}
\ref{stacks-properties-definition-residual-gerbe} を参照

\noindent
{\it $x$ における $X$ の残余空間}
\ref{decent-spaces-definition-residual-space} を参照

\noindent
{\it 剰余体次数}
\ref{more-algebra-definition-extension-discrete-valuation-rings} を参照

\noindent
{\it 剰余体次数}
\ref{more-algebra-definition-extension-valuation-rings} を参照

\noindent
{\it $x$ における $X$ の剰余体}
\ref{schemes-definition-locally-ringed-space} を参照

\noindent
{\it $x$ における $X$ の剰余体}
\ref{decent-spaces-definition-residue-field} を参照

\noindent
{\it 剰余体}
\ref{algebra-definition-local-ring} を参照

\noindent
{\it 分解関手}
\ref{derived-definition-localization-functor} を参照

\noindent
{\it 有限自由 $R$-加群による $M$ の分解}
\ref{algebra-definition-finite-free-resolution} を参照

\noindent
{\it 自由 $R$-加群による $M$ の分解}
\ref{algebra-definition-finite-free-resolution} を参照

\noindent
{\it 正規化ブローアップによる特異点解消}
\ref{resolve-definition-resolution-surface} を参照

\noindent
{\it 正規化ブローアップによる特異点解消}
\ref{spaces-resolve-definition-resolution-surface} を参照

\noindent
{\it 特異点解消}
\ref{resolve-definition-resolution} を参照

\noindent
{\it 特異点解消}
\ref{spaces-resolve-definition-resolution} を参照

\noindent
{\it 分解性質}
\ref{perfect-definition-resolution-property} を参照

\noindent
{\it 分解性質}
\ref{spaces-perfect-definition-resolution-property} を参照

\noindent
{\it 分解}
\ref{algebra-definition-finite-free-resolution} を参照

\noindent
{\it $(U, R, s, t, c)$ の $\mathcal{C}'$ への制限 $(U, R, s, t, c)|_{\mathcal{C}'}$}
\ref{formal-defos-definition-restricting-groupoids-in-functors} を参照

\noindent
{\it $(U, R, s, t, c)$ の $U'$ への制限}
\ref{groupoids-definition-restrict-groupoid} を参照

\noindent
{\it $(U, R, s, t, c)$ の $U'$ への制限}
\ref{spaces-groupoids-definition-restrict-groupoid} を参照

\noindent
{\it $\mathcal{F}$ の $\mathcal{C}/U$ への制限}
\ref{sites-definition-localize} を参照

\noindent
{\it $\mathcal{F}$ の $\mathcal{C}/U$ への制限}
\ref{sites-modules-definition-localize-ringed-site} を参照

\noindent
{\it $\mathcal{F}$ の $U_\etale$ への制限}
\ref{stacks-sheaves-definition-pullback} を参照

\noindent
{\it $\mathcal{G}$ の $U$ への制限}
\ref{sheaves-definition-restriction} を参照

\noindent
{\it $\mathcal{G}$ の $U$ への制限}
\ref{sheaves-definition-restriction} を参照

\noindent
{\it $\mathcal{G}$ の $U$ への制限}
\ref{sheaves-definition-restriction} を参照

\noindent
{\it 小エタール・サイトへの制限}
\ref{topologies-definition-restriction-small-etale} を参照

\noindent
{\it 小エタール・サイトへの制限}
\ref{spaces-topologies-definition-restriction-small-etale} を参照

\noindent
{\it 小 pro-エタール・サイトへの制限}
\ref{proetale-definition-restriction-small-proetale} を参照

\noindent
{\it 小 Zariski サイトへの制限}
\ref{topologies-definition-restriction-small-zariski} を参照

\noindent
{\it 制限}
\ref{groupoids-definition-restrict-relation} を参照

\noindent
{\it 制限}
\ref{spaces-groupoids-definition-restrict-relation} を参照

\noindent
{\it レトロコンパクト}
\ref{topology-definition-quasi-compact} を参照

\noindent
{\it $(A, I)$ 上 rig-エタール}
\ref{restricted-definition-rig-etale-homomorphism} を参照

\noindent
{\it rig-エタール}
\ref{restricted-definition-rig-etale} を参照

\noindent
{\it rig-閉}
\ref{restricted-definition-rig-closed} を参照

\noindent
{\it rig-エタール}
\ref{restricted-definition-rig-etale-continuous-homomorphism} を参照

\noindent
{\it rig-平坦}
\ref{restricted-definition-rig-flat-continuous-homomorphism} を参照

\noindent
{\it rig-平坦}
\ref{restricted-definition-rig-flat} を参照

\noindent
{\it $(A, I)$ 上 rig-滑らか}
\ref{restricted-definition-rig-smooth-homomorphism} を参照

\noindent
{\it rig-滑らか}
\ref{restricted-definition-rig-smooth-continuous-homomorphism} を参照

\noindent
{\it rig-滑らか}
\ref{restricted-definition-rig-smooth} を参照

\noindent
{\it rig-全射}
\ref{restricted-definition-rig-surjective} を参照

\noindent
{\it $F$ に関して右非輪状}
\ref{derived-definition-derived-functor} を参照

\noindent
{\it 右随伴}
\ref{categories-definition-adjoint} を参照

\noindent
{\it 右許容}
\ref{derived-definition-admissible} を参照

\noindent
{\it 右導来可能}
\ref{derived-definition-everywhere-defined} を参照

\noindent
{\it 右導来関手 $RF$ が定義される}
\ref{derived-definition-right-derived-functor-defined} を参照

\noindent
{\it $F$ の右導来関手}
\ref{derived-definition-derived-functor} を参照

\noindent
{\it 右双対}
\ref{categories-definition-dual} を参照

\noindent
{\it 右完全}
\ref{categories-definition-exact} を参照

\noindent
{\it 右乗法系}
\ref{categories-definition-multiplicative-system} を参照

\noindent
{\it 右直交}
\ref{derived-definition-orthogonal} を参照

\noindent
{\it $X$ 上の有理関数環}
\ref{morphisms-definition-ring-of-rational-functions} を参照

\noindent
{\it $X$ 上の有理関数環}
\ref{spaces-morphisms-definition-ring-of-rational-functions} を参照

\noindent
{\it 環付きサイト}
\ref{sites-modules-definition-ringed-site} を参照

\noindent
{\it 環付きサイト}
\ref{etale-cohomology-definition-ringed-site} を参照

\noindent
{\it 環付き空間}
\ref{sheaves-definition-ringed-space} を参照

\noindent
{\it 環付きトポス}
\ref{sites-modules-definition-ringed-topos} を参照

\noindent
{\it 付値判定法の存在部分を満たす}
\ref{schemes-definition-valuative-criterion} を参照

\noindent
{\it 付値判定法の存在部分を満たす}
\ref{spaces-morphisms-definition-valuative-criterion} を参照

\noindent
{\it fpqc 位相に対する層条件を満たす}
\ref{topologies-definition-sheaf-property-fpqc} を参照

\noindent
{\it fpqc 位相に対する層条件を満たす}
\ref{etale-cohomology-definition-sheaf-property-fpqc} を参照

\noindent
{\it 与えられた族に対する層条件を満たす}
\ref{topologies-definition-sheaf-property-fpqc} を参照

\noindent
{\it V 位相に対する層性質を満たす}
\ref{topologies-definition-sheaf-property-V} を参照

\noindent
{\it Zariski 位相に関する層条件を満たす}
\ref{schemes-definition-representable-by-open-immersions} を参照

\noindent
{\it 付値判定法の一意性部分を満たす}
\ref{schemes-definition-valuative-criterion} を参照

\noindent
{\it 付値判定法の一意性部分を満たす}
\ref{spaces-morphisms-definition-valuative-criterion} を参照

\noindent
{\it 付値判定法を満たす}
\ref{spaces-morphisms-definition-valuative-criterion} を参照

\noindent
{\it 飽和している}
\ref{categories-definition-saturated-multiplicative-system} を参照

\noindent
{\it 飽和している}
\ref{derived-definition-saturated} を参照

\noindent
{\it $R$ 上のスキーム}
\ref{schemes-definition-base-change} を参照

\noindent
{\it $S$ 上のスキーム}
\ref{schemes-definition-base-change} を参照

\noindent
{\it $Z$ 上のスキーム構造}
\ref{schemes-definition-reduced-induced-scheme} を参照

\noindent
{\it $X$ における $U$ のスキーム論的閉包}
\ref{morphisms-definition-scheme-theoretically-dense} を参照

\noindent
{\it $X$ における $U$ のスキーム論的閉包}
\ref{spaces-morphisms-definition-scheme-theoretically-dense} を参照

\noindent
{\it スキーム論的ファイバー $X_s$（$f$ の $s$ 上のもの）}
\ref{schemes-definition-fibre} を参照

\noindent
{\it スキーム論的像}
\ref{morphisms-definition-scheme-theoretic-image} を参照

\noindent
{\it スキーム論的像}
\ref{spaces-morphisms-definition-scheme-theoretic-image} を参照

\noindent
{\it スキーム論的像}
\ref{stacks-morphisms-definition-scheme-theoretic-image} を参照

\noindent
{\it スキーム論的交叉}
\ref{morphisms-definition-scheme-theoretic-intersection-union} を参照

\noindent
{\it スキーム論的交叉}
\ref{spaces-morphisms-definition-scheme-theoretic-intersection-union} を参照

\noindent
{\it $\mathcal{F}$ のスキーム論的台}
\ref{morphisms-definition-scheme-theoretic-support} を参照

\noindent
{\it $\mathcal{F}$ のスキーム論的台}
\ref{spaces-morphisms-definition-scheme-theoretic-support} を参照

\noindent
{\it スキーム論的和}
\ref{morphisms-definition-scheme-theoretic-intersection-union} を参照

\noindent
{\it スキーム論的和}
\ref{spaces-morphisms-definition-scheme-theoretic-intersection-union} を参照

\noindent
{\it $X$ においてスキーム論的に稠密}
\ref{morphisms-definition-scheme-theoretically-dense} を参照

\noindent
{\it $X$ においてスキーム論的に稠密}
\ref{spaces-morphisms-definition-scheme-theoretically-dense} を参照

\noindent
{\it スキーム}
\ref{schemes-definition-scheme} を参照

\noindent
{\it コンパクト台付き切断}
\ref{more-etale-definition-compact-support} を参照

\noindent
{\it $X$ 上の半表現可能対象}
\ref{hypercovering-definition-SR} を参照

\noindent
{\it 半表現可能対象}
\ref{hypercovering-definition-SR} を参照

\noindent
{\it $X$ の $Y$ における半正規化}
\ref{morphisms-definition-relative-seminormalization} を参照

\noindent
{\it 半正規化}
\ref{morphisms-definition-seminormalization} を参照

\noindent
{\it 半正規}
\ref{morphisms-definition-seminormal-ring} を参照

\noindent
{\it 半正規}
\ref{morphisms-definition-seminormal} を参照

\noindent
{\it 半安定曲線族}
\ref{moduli-curves-definition-semistable} を参照

\noindent
{\it 半安定還元}
\ref{models-definition-semistable} を参照

\noindent
{\it 分離次数}
\ref{fields-definition-separable-degree} を参照

\noindent
{\it 分離次数}
\ref{fields-definition-insep-degree} を参照

\noindent
{\it $k$ 上分離的}
\ref{algebra-definition-separable-field-extension} を参照

\noindent
{\it 分離解}
\ref{more-algebra-definition-solution} を参照

\noindent
{\it 分離的}
\ref{fields-definition-separable} を参照

\noindent
{\it 分離的}
\ref{fields-definition-separable} を参照

\noindent
{\it 分離的}
\ref{fields-definition-separable} を参照

\noindent
{\it 分離的}
\ref{fields-definition-separable-algebraic} を参照

\noindent
{\it $k$ 上分離生成}
\ref{algebra-definition-separable-field-extension} を参照

\noindent
{\it 分離的な群スキーム}
\ref{groupoids-definition-smooth-group-scheme} を参照

\noindent
{\it $S$ 上分離}
\ref{spaces-definition-separated} を参照

\noindent
{\it $S$ 上分離}
\ref{stacks-morphisms-definition-absolute-separated} を参照

\noindent
{\it 分離前層}
\ref{etale-cohomology-definition-sheaf} を参照

\noindent
{\it 分離的}
\ref{topology-definition-separated} を参照

\noindent
{\it 分離的}
\ref{sheaves-definition-separated} を参照

\noindent
{\it 分離的}
\ref{sites-definition-separated} を参照

\noindent
{\it 分離的}
\ref{sites-definition-presheaf-separated-topology} を参照

\noindent
{\it 分離的}
\ref{homology-definition-filtered} を参照

\noindent
{\it 分離的}
\ref{schemes-definition-separated} を参照

\noindent
{\it 分離的}
\ref{schemes-definition-separated} を参照

\noindent
{\it 分離的}
\ref{spaces-properties-definition-separated} を参照

\noindent
{\it 分離的}
\ref{spaces-properties-definition-separated} を参照

\noindent
{\it 分離的}
\ref{spaces-morphisms-definition-separated} を参照

\noindent
{\it 分離的}
\ref{formal-spaces-definition-separated} を参照

\noindent
{\it 分離的}
\ref{formal-spaces-definition-separated-morphism} を参照

\noindent
{\it 分離的}
\ref{stacks-morphisms-definition-separated} を参照

\noindent
{\it 分離的}
\ref{stacks-morphisms-definition-absolute-separated} を参照

\noindent
{\it $R$-軌道を分離する}
\ref{groupoids-quotients-definition-set-theoretically-invariant} を参照

\noindent
{\it 軌道を分離する}
\ref{groupoids-quotients-definition-set-theoretically-invariant} を参照

\noindent
{\it Serre 関手}
\ref{equiv-definition-Serre-functor} を参照

\noindent
{\it セール部分圏}
\ref{homology-definition-serre-subcategory} を参照

\noindent
{\it 集合論的同値関係}
\ref{groupoids-quotients-definition-set-theoretic-equivalence} を参照

\noindent
{\it 集合論的前同値関係}
\ref{groupoids-quotients-definition-set-theoretic-equivalence} を参照

\noindent
{\it 集合論的に $R$-不変}
\ref{groupoids-definition-invariant-open} を参照

\noindent
{\it 集合論的に $R$-不変}
\ref{groupoids-quotients-definition-set-theoretically-invariant} を参照

\noindent
{\it セトイド}
\ref{categories-definition-setoid} を参照

\noindent
{\it $\mathcal{B}$ 上の $\mathcal{O}$-加群の層 $\mathcal{F}$}
\ref{sheaves-definition-sheaf-modules-basis} を参照

\noindent
{\it $\mathcal{B}$ 上の集合の層 $\mathcal{F}$}
\ref{sheaves-definition-sheaf-basis} を参照

\noindent
{\it 集合の層 $\mathcal{F}$（$X$ 上）}
\ref{sheaves-definition-sheaf} を参照

\noindent
{\it $\mathcal{B}$ 上の $\mathcal{C}$ に値をもつ層 $\mathcal{F}$}
\ref{sheaves-definition-sheaf-structures-basis} を参照

\noindent
{\it $\mathcal{F}$ に付随する層}
\ref{sites-definition-associated-sheaf} を参照

\noindent
{\it $\mathcal{F}$ に付随する層}
\ref{sites-definition-associated-sheaf-topology} を参照

\noindent
{\it 加群 $M$ と環準同型 $\alpha$ に付随する層}
\ref{modules-definition-sheaf-associated} を参照

\noindent
{\it 加群 $M$ に付随する層}
\ref{modules-definition-sheaf-associated} を参照

\noindent
{\it エタール位相に関する層}
\ref{stacks-sheaves-definition-sheaves} を参照

\noindent
{\it fppf 位相に関する層}
\ref{stacks-sheaves-definition-sheaves} を参照

\noindent
{\it 滑らかな位相に関する層}
\ref{stacks-sheaves-definition-sheaves} を参照

\noindent
{\it シントミック位相に関する層}
\ref{stacks-sheaves-definition-sheaves} を参照

\noindent
{\it Zariski 位相に関する層}
\ref{stacks-sheaves-definition-sheaves} を参照

\noindent
{\it $\mathcal{F}$ に付随する $\mathcal{O}$-加群の層}
\ref{descent-definition-structure-sheaf} を参照

\noindent
{\it $\mathcal{F}$ に付随する $\mathcal{O}$-加群の層}
\ref{descent-definition-structure-sheaf} を参照

\noindent
{\it $\mathcal{O}$-加群の層}
\ref{sheaves-definition-sheaf-modules} を参照

\noindent
{\it $\mathcal{O}$-加群の層}
\ref{sites-modules-definition-sheaf-modules} を参照

\noindent
{\it $\mathcal{O}_\mathcal{X}$-加群の層}
\ref{stacks-sheaves-definition-modules} を参照

\noindent
{\it $X$ 上の $R$-不変関数の層}
\ref{groupoids-quotients-definition-functions} を参照

\noindent
{\it $X$ 上のアーベル群の層}
\ref{sheaves-definition-abelian-sheaf} を参照

\noindent
{\it $x$ の自己同型の層}
\ref{stacks-morphisms-definition-isom} を参照

\noindent
{\it 微分次数付き $\mathcal{O}$-代数の層}
\ref{sdga-definition-dga} を参照

\noindent
{\it 微分次数付き代数の層}
\ref{sdga-definition-dga} を参照

\noindent
{\it $S$ 上の $X$ の微分層 $\Omega_{X/S}$}
\ref{modules-definition-differentials} を参照

\noindent
{\it $S$ 上の $X$ の微分層 $\Omega_{X/S}$}
\ref{morphisms-definition-sheaf-differentials} を参照

\noindent
{\it $Y$ 上の $X$ の微分層 $\Omega_{X/Y}$}
\ref{sites-modules-definition-sheaf-differentials} を参照

\noindent
{\it $Y$ 上の $X$ の微分層 $\Omega_{X/Y}$}
\ref{spaces-more-morphisms-definition-sheaf-differentials} を参照

\noindent
{\it 次数付き $\mathcal{O}$-代数の層}
\ref{sdga-definition-ga} を参照

\noindent
{\it 次数付き代数の層}
\ref{sdga-definition-ga} を参照

\noindent
{\it $x_1$ から $x_2$ への同型の層}
\ref{stacks-morphisms-definition-isom} を参照

\noindent
{\it $X$ 上の有理型関数の層}
\ref{divisors-definition-sheaf-meromorphic-functions} を参照

\noindent
{\it $X$ 上の有理型関数の層}
\ref{spaces-divisors-definition-sheaf-meromorphic-functions} を参照

\noindent
{\it 全商環の層 ${\mathcal K}_S$}
\ref{exercises-definition-divisor} を参照

\noindent
{\it 層論的に空}
\ref{sites-definition-empty} を参照

\noindent
{\it 層}
\ref{sheaves-definition-sheaf-values-in-category} を参照

\noindent
{\it 層}
\ref{sites-definition-sheaf-sets} を参照

\noindent
{\it 層}
\ref{sites-definition-sheaf} を参照

\noindent
{\it 層}
\ref{sites-definition-sheaf-sets-topology} を参照

\noindent
{\it 層}
\ref{etale-cohomology-definition-sheaf} を参照

\noindent
{\it 層}
\ref{stacks-sheaves-definition-sheaves} を参照

\noindent
{\it シフト}
\ref{homology-definition-graded-shift} を参照

\noindent
{\it 短完全列}
\ref{homology-definition-exact} を参照

\noindent
{\it 同胞関手}
\ref{equiv-definition-siblings} を参照

\noindent
{\it 同胞関手}
\ref{equiv-definition-siblings-geometric} を参照

\noindent
{\it 同胞関手}
\ref{equiv-definition-siblings} を参照

\noindent
{\it 同胞関手}
\ref{equiv-definition-siblings-geometric} を参照

\noindent
{\it 射 $f_i$ が生成する $U$ 上のふるい}
\ref{sites-definition-sieve-generated} を参照

\noindent
{\it $U$ 上のふるい $S$}
\ref{sites-definition-sieve} を参照

\noindent
{\it 相似}
\ref{etale-cohomology-definition-brauer-equivalent} を参照

\noindent
{\it 単純}
\ref{algebra-definition-simple-module} を参照

\noindent
{\it 単純}
\ref{brauer-definition-simple} を参照

\noindent
{\it 単純}
\ref{brauer-definition-simple} を参照

\noindent
{\it 単純}
\ref{homology-definition-simple} を参照

\noindent
{\it 単体的 $\mathcal{A}_\bullet$-加群}
\ref{sites-cohomology-definition-simplicial-module} を参照

\noindent
{\it 単体アーベル群}
\ref{simplicial-definition-simplicial-object} を参照

\noindent
{\it $\mathcal{C}$ の単体的対象 $U$}
\ref{simplicial-definition-simplicial-object} を参照

\noindent
{\it $f$ に付随する単体的スキーム}
\ref{spaces-simplicial-definition-fibre-products-simplicial-scheme} を参照

\noindent
{\it 単体集合}
\ref{simplicial-definition-simplicial-object} を参照

\noindent
{\it $\mathcal{A}_\bullet$-加群の単体的層}
\ref{sites-cohomology-definition-simplicial-module} を参照

\noindent
{\it 加群の導来圏の単体的系}
\ref{spaces-simplicial-definition-cartesian-derived-modules} を参照

\noindent
{\it 導来圏の単体的系}
\ref{spaces-simplicial-definition-cartesian-derived} を参照

\noindent
{\it $R$ 上の $A$ の特異イデアル}
\ref{smoothing-definition-singular-ideal} を参照

\noindent
{\it 特異点集合}
\ref{properties-definition-singular-locus} を参照

\noindent
{\it $X$ の特異点は高々節点である}
\ref{curves-definition-nodal} を参照

\noindent
{\it サイト}
\ref{sites-definition-site} を参照

\noindent
{\it サイト}
\ref{etale-cohomology-definition-site} を参照

\noindent
{\it 大きさ}
\ref{injectives-definition-size} を参照

\noindent
{\it 斜体}
\ref{brauer-definition-skew-field} を参照

\noindent
{\it 値 $A$ をもつ $x$ における摩天楼層}
\ref{sheaves-definition-skyscraper-sheaf} を参照

\noindent
{\it 摩天楼層}
\ref{sheaves-definition-skyscraper-sheaf} を参照

\noindent
{\it 摩天楼層}
\ref{sheaves-definition-skyscraper-sheaf} を参照

\noindent
{\it 摩天楼層}
\ref{sheaves-definition-skyscraper-sheaf} を参照

\noindent
{\it 摩天楼層}
\ref{sheaves-definition-skyscraper-sheaf} を参照

\noindent
{\it 摩天楼層}
\ref{sites-definition-pushforward-point} を参照

\noindent
{\it $S$ の小 $\tau$-サイト}
\ref{etale-cohomology-definition-tau-site} を参照

\noindent
{\it 小エタールサイト $X_\etale$}
\ref{spaces-properties-definition-etale-site} を参照

\noindent
{\it $S$ の小エタール・サイト}
\ref{topologies-definition-big-small-etale} を参照

\noindent
{\it $S$ 上の小エタール・サイト}
\ref{etale-cohomology-definition-big-etale-site} を参照

\noindent
{\it 小エタール・サイト}
\ref{formal-spaces-definition-etale-sites} を参照

\noindent
{\it 小エタールトポス}
\ref{etale-cohomology-definition-etale-topos} を参照

\noindent
{\it 小エタールトポス}
\ref{spaces-properties-definition-etale-topos} を参照

\noindent
{\it $S$ の小アフィンエタール・サイト}
\ref{topologies-definition-big-small-etale} を参照

\noindent
{\it $S$ の小アフィン Zariski サイト}
\ref{topologies-definition-big-small-Zariski} を参照

\noindent
{\it 小拡大}
\ref{algebra-definition-small-extension} を参照

\noindent
{\it 小拡大}
\ref{formal-defos-definition-small-extension} を参照

\noindent
{\it $S$ の小 pro-エタール・サイト}
\ref{proetale-definition-big-small-proetale} を参照

\noindent
{\it 小 Zariski サイト $F_{Zar}$}
\ref{spaces-definition-small-Zariski-site} を参照

\noindent
{\it $S$ の小 Zariski サイト}
\ref{topologies-definition-big-small-Zariski} を参照

\noindent
{\it 小 Zariski サイト}
\ref{etale-cohomology-definition-big-etale-site} を参照

\noindent
{\it 小 Zariski トポス}
\ref{etale-cohomology-definition-etale-topos} を参照

\noindent
{\it $\mathfrak q$ で滑らか}
\ref{algebra-definition-smooth-at-prime} を参照

\noindent
{\it 点 $x \in X$ において滑らか}
\ref{morphisms-definition-smooth} を参照

\noindent
{\it $x$ において滑らか}
\ref{spaces-morphisms-definition-smooth} を参照

\noindent
{\it $T$ の滑らか被覆}
\ref{topologies-definition-smooth-covering} を参照

\noindent
{\it $X$ の滑らかな被覆}
\ref{spaces-topologies-definition-smooth-covering} を参照

\noindent
{\it 滑らかな群スキーム}
\ref{groupoids-definition-smooth-group-scheme} を参照

\noindent
{\it 滑らかなグルーポイド}
\ref{algebraic-definition-smooth-groupoid} を参照

\noindent
{\it 始域と終域に関して滑らか局所的}
\ref{spaces-descent-definition-local-source-target} を参照

\noindent
{\it 滑らか局所的}
\ref{descent-definition-local-at-point} を参照

\noindent
{\it 相対次元 $d$ の滑らかな射}
\ref{morphisms-definition-smooth-relative-dimension} を参照

\noindent
{\it 滑らかな層}
\ref{stacks-sheaves-definition-sheaves} を参照

\noindent
{\it 滑らかな多様体}
\ref{varieties-definition-variety-type} を参照

\noindent
{\it 滑らか}
\ref{algebra-definition-smooth} を参照

\noindent
{\it 滑らか}
\ref{morphisms-definition-smooth} を参照

\noindent
{\it 滑らか}
\ref{descent-definition-etale-morphism-germs} を参照

\noindent
{\it 滑らか}
\ref{spaces-morphisms-definition-smooth} を参照

\noindent
{\it 滑らか}
\ref{formal-defos-definition-smooth-morphism} を参照

\noindent
{\it 滑らか}
\ref{formal-defos-definition-cofibered-groupoid-projection-smooth} を参照

\noindent
{\it 滑らか}
\ref{formal-defos-definition-smooth-groupoid-in-functors} を参照

\noindent
{\it 滑らか}
\ref{stacks-morphisms-definition-smooth} を参照

\noindent
{\it 滑らか}
\ref{stacks-introduction-definition-smooth} を参照

\noindent
{\it ソバー}
\ref{topology-definition-generic-point} を参照

\noindent
{\it $A \subset B$ に対する解}
\ref{more-algebra-definition-solution} を参照

\noindent
{\it $\mathcal{C}$ から $\mathcal{D}$ への特殊余連続関手 $u$}
\ref{sites-definition-special-cocontinuous-functor} を参照

\noindent
{\it 特殊化は $f$ に沿って 持ち上がる}
\ref{topology-definition-lift-specializations} を参照

\noindent
{\it 特殊化}
\ref{topology-definition-specialization} を参照

\noindent
{\it 特殊化}
\ref{exercises-definition-closed} を参照

\noindent
{\it 特殊化}
\ref{exercises-definition-specialization} を参照

\noindent
{\it 特殊化的}
\ref{topology-definition-lift-specializations} を参照

\noindent
{\it $(A, d, \alpha)$ に付随するスペクトル系列}
\ref{homology-definition-differential-object-selfmap} を参照

\noindent
{\it 完全対に付随するスペクトル系列}
\ref{homology-definition-spectral-sequence-associated-exact-couple} を参照

\noindent
{\it $\mathcal{A}$ におけるスペクトル系列}
\ref{homology-definition-spectral-sequence} を参照

\noindent
{\it スペクトル}
\ref{topology-definition-spectral-space} を参照

\noindent
{\it スペクトル}
\ref{topology-definition-spectral-space} を参照

\noindent
{\it $\mathcal{A}$ の $S$ 上のスペクトル}
\ref{constructions-definition-relative-spec} を参照

\noindent
{\it $X$ 上の $\mathcal{A}$ のスペクトル}
\ref{spaces-morphisms-definition-relative-spec} を参照

\noindent
{\it スペクトル}
\ref{algebra-definition-spectrum-ring} を参照

\noindent
{\it スペクトル}
\ref{schemes-definition-structure-sheaf} を参照

\noindent
{\it 分裂亜群ファイバー圏}
\ref{categories-definition-split-category-fibred-in-groupoids} を参照

\noindent
{\it 分裂等化子}
\ref{descent-definition-split-equalizer} を参照

\noindent
{\it 分裂ファイバー圏}
\ref{categories-definition-split-fibred-category} を参照

\noindent
{\it 分裂節点}
\ref{curves-definition-split-node} を参照

\noindent
{\it $u$ 上分裂}
\ref{spaces-more-groupoids-definition-split-at-point} を参照

\noindent
{\it 分解する}
\ref{brauer-definition-splitting} を参照

\noindent
{\it $F$ 上の $P$ の分解体}
\ref{fields-definition-splitting-field} を参照

\noindent
{\it 分解体}
\ref{brauer-definition-splitting} を参照

\noindent
{\it $u$ 上の $R$ の分裂}
\ref{spaces-more-groupoids-definition-split-at-point} を参照

\noindent
{\it 分裂する}
\ref{homology-definition-ses-split} を参照

\noindent
{\it 分裂する}
\ref{simplicial-definition-split} を参照

\noindent
{\it 分裂する}
\ref{simplicial-definition-split} を参照

\noindent
{\it 代数空間における亜群 $(U, R, s, t, c)$ の固定化群}
\ref{spaces-groupoids-definition-stabilizer-groupoid} を参照

\noindent
{\it 亜群スキーム $(U, R, s, t, c)$ の安定化群}
\ref{groupoids-definition-stabilizer-groupoid} を参照

\noindent
{\it 安定曲線族}
\ref{moduli-curves-definition-stable} を参照

\noindent
{\it 基底変換で安定}
\ref{morphisms-definition-property-local} を参照

\noindent
{\it 合成で安定}
\ref{morphisms-definition-property-local} を参照

\noindent
{\it 一般化の下で安定}
\ref{topology-definition-specialization} を参照

\noindent
{\it 特殊化の下で安定}
\ref{topology-definition-specialization} を参照

\noindent
{\it 安定自由}
\ref{more-algebra-definition-stably-free} を参照

\noindent
{\it 安定同型}
\ref{more-algebra-definition-stably-free} を参照

\noindent
{\it 離散圏のスタック}
\ref{stacks-definition-stack-in-sets} を参照

\noindent
{\it グルーポイドのスタック}
\ref{stacks-definition-stack-in-groupoids} を参照

\noindent
{\it セトイドのスタック}
\ref{stacks-definition-stack-in-sets} を参照

\noindent
{\it 集合のスタック}
\ref{stacks-definition-stack-in-sets} を参照

\noindent
{\it スタック}
\ref{stacks-definition-stack} を参照

\noindent
{\it 茎}
\ref{etale-cohomology-definition-stalk} を参照

\noindent
{\it 茎}
\ref{trace-definition-torsion-l-adic-sheaf} を参照

\noindent
{\it 茎}
\ref{spaces-properties-definition-stalk} を参照

\noindent
{\it 標準 $\tau$-被覆}
\ref{etale-cohomology-definition-standard-tau} を参照

\noindent
{\it 標準エタール被覆}
\ref{topologies-definition-standard-etale} を参照

\noindent
{\it 標準エタール}
\ref{algebra-definition-standard-etale} を参照

\noindent
{\it 標準エタール}
\ref{morphisms-definition-etale} を参照

\noindent
{\it 標準エタール}
\ref{etale-cohomology-definition-standard-etale} を参照

\noindent
{\it 標準 fppf 被覆}
\ref{topologies-definition-standard-fppf} を参照

\noindent
{\it 標準 fpqc 被覆}
\ref{topologies-definition-standard-fpqc} を参照

\noindent
{\it 標準 h 被覆}
\ref{flat-definition-standard-h} を参照

\noindent
{\it 標準開被覆}
\ref{schemes-definition-standard-covering} を参照

\noindent
{\it 標準開被覆}
\ref{schemes-definition-standard-covering} を参照

\noindent
{\it 標準開被覆}
\ref{constructions-definition-standard-covering} を参照

\noindent
{\it 標準開集合}
\ref{algebra-definition-Zariski-topology} を参照

\noindent
{\it 標準 ph 被覆}
\ref{topologies-definition-standard-ph-covering} を参照

\noindent
{\it 標準 pro-エタール被覆}
\ref{proetale-definition-standard-proetale} を参照

\noindent
{\it $\mathcal{A}$ 上の $\mathcal{B}$ の標準分解}
\ref{cotangent-definition-standard-resolution-sheaves-rings} を参照

\noindent
{\it $A$ 上の $B$ の標準分解}
\ref{cotangent-definition-standard-resolution} を参照

\noindent
{\it 標準縮小}
\ref{flat-definition-shrink} を参照

\noindent
{\it 標準縮小}
\ref{flat-definition-shrink-complete} を参照

\noindent
{\it $R$ 上の標準滑らかな代数}
\ref{algebra-definition-standard-smooth} を参照

\noindent
{\it 標準平滑被覆}
\ref{topologies-definition-standard-smooth} を参照

\noindent
{\it 標準滑らか}
\ref{algebra-definition-standard-smooth} を参照

\noindent
{\it 標準滑らか}
\ref{algebra-definition-standard-smooth} を参照

\noindent
{\it 標準滑らか}
\ref{morphisms-definition-smooth} を参照

\noindent
{\it 標準シントミック被覆}
\ref{topologies-definition-standard-syntomic} を参照

\noindent
{\it 標準シントミック}
\ref{morphisms-definition-syntomic} を参照

\noindent
{\it 標準 V 被覆}
\ref{topologies-definition-standard-V-covering} を参照

\noindent
{\it 標準 Zariski 被覆}
\ref{topologies-definition-standard-Zariski} を参照

\noindent
{\it 層}
\ref{topology-definition-stratification} を参照

\noindent
{\it 層化}
\ref{topology-definition-stratification} を参照

\noindent
{\it $\mathcal{O}_{S, s}$ の強ヘンゼル化}
\ref{etale-cohomology-definition-etale-local-rings} を参照

\noindent
{\it $\kappa \subset \kappa^{sep}$ に関する $R$ の強ヘンゼル化}
\ref{algebra-definition-henselization} を参照

\noindent
{\it $\overline{s}$ における $S$ の強ヘンゼル化}
\ref{etale-cohomology-definition-etale-local-rings} を参照

\noindent
{\it $\overline{x}$ における $X$ の強ヘンゼル化}
\ref{spaces-properties-definition-etale-local-rings} を参照

\noindent
{\it 強ヘンゼル化}
\ref{algebra-definition-henselization} を参照

\noindent
{\it 位相空間の厳密写像}
\ref{topology-definition-submersive} を参照

\noindent
{\it 厳密な厚化の射}
\ref{defos-definition-strict-morphism-thickenings} を参照

\noindent
{\it 厳密な厚化の射}
\ref{defos-definition-strict-morphism-thickenings-ringed-topoi} を参照

\noindent
{\it 単純正規交差因子}
\ref{etale-definition-strict-normal-crossings} を参照

\noindent
{\it $R \to R[\frac{I}{a}]$ に沿う $M$ の狭義変換}
\ref{more-algebra-definition-strict-transform} を参照

\noindent
{\it 狭義変換}
\ref{divisors-definition-strict-transform} を参照

\noindent
{\it 狭義変換}
\ref{divisors-definition-strict-transform} を参照

\noindent
{\it 狭義変換}
\ref{spaces-divisors-definition-strict-transform} を参照

\noindent
{\it 狭義変換}
\ref{spaces-divisors-definition-strict-transform} を参照

\noindent
{\it 厳密可換}
\ref{dga-definition-cdga} を参照

\noindent
{\it 厳密充満}
\ref{categories-definition-subcategory} を参照

\noindent
{\it 強ヘンゼル的}
\ref{algebra-definition-henselian} を参照

\noindent
{\it 強ヘンゼル的}
\ref{etale-cohomology-definition-strictly-henselian} を参照

\noindent
{\it 狭義完全}
\ref{cohomology-definition-strictly-perfect} を参照

\noindent
{\it 狭義完全}
\ref{sites-cohomology-definition-strictly-perfect} を参照

\noindent
{\it $R$ 上の $A$ における厳密標準元}
\ref{smoothing-definition-strictly-standard} を参照

\noindent
{\it 厳密}
\ref{homology-definition-strict} を参照

\noindent
{\it 強生成元}
\ref{derived-definition-generators} を参照

\noindent
{\it $u$ 上の $R$ の強分裂}
\ref{spaces-more-groupoids-definition-split-at-point} を参照

\noindent
{\it より強い}
\ref{sites-definition-finer} を参照

\noindent
{\it 強 $\mathcal{C}$-カルテシアン射}
\ref{categories-definition-cartesian-over-C} を参照

\noindent
{\it 強カルテシアン射}
\ref{categories-definition-cartesian-over-C} を参照

\noindent
{\it $u$ 上強分裂}
\ref{spaces-more-groupoids-definition-split-at-point} を参照

\noindent
{\it $R$ 上強超越的}
\ref{algebra-definition-strongly-transcendental} を参照

\noindent
{\it 構造射}
\ref{schemes-definition-base-change} を参照

\noindent
{\it $\mathcal{C}$ から継承される $\mathcal{S}$ 上のサイト構造}
\ref{stacks-definition-topology-inherited} を参照

\noindent
{\it $R$ のスペクトルの構造層 $\mathcal{O}_{\Spec(R)}$}
\ref{schemes-definition-structure-sheaf} を参照

\noindent
{\it 斉次スペクトルの構造層 $\mathcal{O}_{\text{Proj}(S)}$（$S$ の）}
\ref{constructions-definition-structure-sheaf} を参照

\noindent
{\it $\mathcal{X}$ の構造層}
\ref{stacks-sheaves-definition-structure-sheaf} を参照

\noindent
{\it 大サイト $(\Sch/S)_\tau$ の構造層}
\ref{descent-definition-structure-sheaf} を参照

\noindent
{\it 構造層}
\ref{descent-definition-structure-sheaf} を参照

\noindent
{\it 構造層}
\ref{sites-modules-definition-ringed-site} を参照

\noindent
{\it 構造層}
\ref{sites-modules-definition-ringed-topos} を参照

\noindent
{\it 構造層}
\ref{etale-cohomology-definition-structure-sheaf} を参照

\noindent
{\it 構造層}
\ref{spaces-properties-definition-structure-sheaf} を参照

\noindent
{\it 部分 $2$-圏}
\ref{categories-definition-sub-2-category} を参照

\noindent
{\it $X$ の位相の準開基}
\ref{topology-definition-subbase} を参照

\noindent
{\it $X$ の位相の部分基底}
\ref{topology-definition-subbase} を参照

\noindent
{\it 準標準的}
\ref{sites-definition-weaker-than-canonical} を参照

\noindent
{\it 部分圏}
\ref{categories-definition-subcategory} を参照

\noindent
{\it 部分体}
\ref{fields-definition-field} を参照

\noindent
{\it 部分関手 $H \subset F$}
\ref{schemes-definition-representable-by-open-immersions} を参照

\noindent
{\it 商位相的}
\ref{topology-definition-submersive} を参照

\noindent
{\it 商位相的}
\ref{morphisms-definition-submersive} を参照

\noindent
{\it 商位相的}
\ref{spaces-morphisms-definition-submersive} を参照

\noindent
{\it 商位相的}
\ref{stacks-morphisms-definition-submersive} を参照

\noindent
{\it 部分対象}
\ref{homology-definition-injective-surjective} を参照

\noindent
{\it 部分前層}
\ref{sheaves-definition-injective-surjective} を参照

\noindent
{\it 部分前層}
\ref{sites-definition-sub-presheaf} を参照

\noindent
{\it $s_i$ によって生成される部分層}
\ref{modules-definition-generated-by-local-sections} を参照

\noindent
{\it $\mathcal{I}$ により零化される切断の部分層}
\ref{properties-definition-subsheaf-sections-annihilated-by-ideal} を参照

\noindent
{\it $\mathcal{I}$ により零化される切断の部分層}
\ref{spaces-limits-definition-subsheaf-sections-annihilated-by-ideal} を参照

\noindent
{\it $T$ 上に台を持つ切断の部分層}
\ref{properties-definition-subsheaf-sections-supported-on-closed} を参照

\noindent
{\it $T$ 上に台を持つ切断の部分層}
\ref{spaces-limits-definition-subsheaf-sections-supported-on-closed} を参照

\noindent
{\it 部分層}
\ref{sheaves-definition-injective-surjective} を参照

\noindent
{\it 部分トポス}
\ref{sites-definition-subtopos} を参照

\noindent
{\it 有効 Cartier 因子 $D_1$ と $D_2$ の和}
\ref{divisors-definition-sum-effective-Cartier-divisors} を参照

\noindent
{\it 有効 Cartier 因子 $D_1$ と $D_2$ の和}
\ref{spaces-divisors-definition-sum-effective-Cartier-divisors} を参照

\noindent
{\it 有効 Cartier 因子の和}
\ref{obsolete-definition-locally-finite-sum-effective-Cartier-divisors} を参照

\noindent
{\it $\mathcal{F}$ の台}
\ref{modules-definition-support} を参照

\noindent
{\it $\mathcal{F}$ の台}
\ref{etale-cohomology-definition-support} を参照

\noindent
{\it $\mathcal{F}$ の台}
\ref{spaces-properties-definition-support} を参照

\noindent
{\it $\sigma$ の台}
\ref{etale-cohomology-definition-support} を参照

\noindent
{\it $\sigma$ の台}
\ref{spaces-properties-definition-support} を参照

\noindent
{\it $M$ の台}
\ref{algebra-definition-support-module} を参照

\noindent
{\it $s$ の台}
\ref{modules-definition-support} を参照

\noindent
{\it $T$ に台をもつ}
\ref{perfect-definition-supported-on} を参照

\noindent
{\it $T$ に台をもつ}
\ref{spaces-perfect-definition-supported-on} を参照

\noindent
{\it 台}
\ref{chow-definition-support-cycle} を参照

\noindent
{\it 全射}
\ref{sheaves-definition-injective-surjective} を参照

\noindent
{\it 全射}
\ref{sheaves-definition-injective-surjective} を参照

\noindent
{\it 全射}
\ref{sites-definition-presheaves-injective-surjective} を参照

\noindent
{\it 全射}
\ref{sites-definition-sheaves-injective-surjective} を参照

\noindent
{\it 全射}
\ref{homology-definition-injective-surjective} を参照

\noindent
{\it 全射}
\ref{morphisms-definition-surjective} を参照

\noindent
{\it 全射}
\ref{spaces-morphisms-definition-surjective} を参照

\noindent
{\it 全射}
\ref{formal-spaces-definition-surjective} を参照

\noindent
{\it 全射}
\ref{stacks-properties-definition-surjective} を参照

\noindent
{\it $M, a, b$ に付随する記号}
\ref{chow-definition-symbol-M} を参照

\noindent
{\it 記号冪}
\ref{algebra-definition-symbolic-power} を参照

\noindent
{\it 記号}
\ref{chow-definition-determinant} を参照

\noindent
{\it 対称モノイド圏}
\ref{categories-definition-symmetric-monoidal-category} を参照

\noindent
{\it 点 $x \in X$ においてシントミック}
\ref{morphisms-definition-syntomic} を参照

\noindent
{\it $x$ においてシントミック}
\ref{spaces-morphisms-definition-syntomic} を参照

\noindent
{\it $T$ のシントミック被覆}
\ref{topologies-definition-syntomic-covering} を参照

\noindent
{\it $X$ のシントミック被覆}
\ref{spaces-topologies-definition-syntomic-covering} を参照

\noindent
{\it 相対次元 $d$ のシントミック射}
\ref{morphisms-definition-syntomic-relative-dimension} を参照

\noindent
{\it シントミック層}
\ref{stacks-sheaves-definition-sheaves} を参照

\noindent
{\it シントミック}
\ref{algebra-definition-lci} を参照

\noindent
{\it シントミック}
\ref{morphisms-definition-syntomic} を参照

\noindent
{\it シントミック}
\ref{spaces-morphisms-definition-syntomic} を参照

\noindent
{\it $(X_i, f_{i'i})$ 上の層の系 $(\mathcal{F}_i, \varphi_{i'i})$}
\ref{etale-cohomology-definition-inverse-system-sheaves} を参照

\noindent
{\it $I$ 上の $R$-加群の系 $(M_i, \mu_{ij})$}
\ref{algebra-definition-directed-system} を参照

\noindent
{\it $R$ のパラメータ系}
\ref{algebra-definition-regular-local} を参照

\noindent
{\it 環の系}
\ref{exercises-definition-directed-poset} を参照

\noindent
{\it $\mathcal{C}$ における $I$ 上の系}
\ref{categories-definition-system-over-poset} を参照

\noindent
{\it $\mathfrak m$ の馴性惰性群}
\ref{more-algebra-definition-wild-inertia} を参照

\noindent
{\it テイム記号}
\ref{chow-definition-tame-symbol} を参照

\noindent
{\it $A$ に関して馴分岐}
\ref{more-algebra-definition-types-of-extensions} を参照

\noindent
{\it $\mathcal{F}$ の接空間 $T \mathcal{F}$}
\ref{formal-defos-definition-tangent-space} を参照

\noindent
{\it $F$ の接空間 $TF$}
\ref{formal-defos-definition-tangent-space-over-R} を参照

\noindent
{\it $S$ 上の $X$ の $x$ における接空間}
\ref{varieties-definition-tangent-space} を参照

\noindent
{\it $S$ 上の $X$ の接空間}
\ref{exercises-definition-tangent-space} を参照

\noindent
{\it 接ベクトル}
\ref{varieties-definition-tangent-space} を参照

\noindent
{\it 接ベクトル}
\ref{exercises-definition-tangent-space} を参照

\noindent
{\it 自明に同値}
\ref{sites-definition-combinatorial-tautological} を参照

\noindent
{\it タウト}
\ref{formal-spaces-definition-taut} を参照

\noindent
{\it テンソル冪}
\ref{modules-definition-powers} を参照

\noindent
{\it テンソル積微分次数付き代数}
\ref{dga-definition-tensor-product} を参照

\noindent
{\it テンソル積}
\ref{formal-spaces-definition-toplogy-tensor-product} を参照

\noindent
{\it $\mathcal{A}$ の複体の項別分裂完全列}
\ref{derived-definition-split-ses} を参照

\noindent
{\it 項ごとに分裂する単射 $\alpha : A^\bullet \to B^\bullet$}
\ref{derived-definition-termwise-split-map} を参照

\noindent
{\it 項ごとに分裂する全射 $\beta : B^\bullet \to C^\bullet$}
\ref{derived-definition-termwise-split-map} を参照

\noindent
{\it $y$ における $f : X \to Y$ のファイバーは幾何学的に被約である}
\ref{spaces-more-morphisms-definition-geometrically-reduced-fibre} を参照

\noindent
{\it $f$ の $y$ 上のファイバーは局所 Noether である}
\ref{spaces-divisors-definition-locally-Noetherian-fibre} を参照

\noindent
{\it $X$ の $z$ 上のファイバーは、$x$ で、$Y$ の $z$ 上の ファイバー上平坦である}
\ref{spaces-more-morphisms-definition-module-flat-on-fibre} を参照

\noindent
{\it $X$ の $z$ 上のファイバーは $Y$ の $z$ 上のファイバー上平坦である}
\ref{spaces-more-morphisms-definition-module-flat-on-fibre} を参照

\noindent
{\it $f$ のファイバーは局所 Noether である}
\ref{spaces-divisors-definition-locally-Noetherian-fibre} を参照

\noindent
{\it $X$ から $Y$ への $S$ 上の相対同値の Fourier--Mukai 核}
\ref{equiv-definition-relative-equivalence-kernel} を参照

\noindent
{\it $X$ 上の関数は $U$ 上の $R$-不変関数である}
\ref{groupoids-quotients-definition-functions} を参照

\noindent
{\it $f$ に対する Gysin 写像が存在する}
\ref{chow-definition-lci-gysin} を参照

\noindent
{\it 相対次元}
\ref{stacks-geometry-definition-relative-dimension} を参照

\noindent
{\it $\mathcal{F}$ の $z$ 上のファイバーへの制限は、$x$ で、$Y$ の $z$ 上の ファイバー上平坦である}
\ref{spaces-more-morphisms-definition-module-flat-on-fibre} を参照

\noindent
{\it $\mathcal{Z}$ 上の厚化}
\ref{stacks-more-morphisms-definition-thickening} を参照

\noindent
{\it $B$ 上の肥厚化}
\ref{spaces-more-morphisms-definition-thickening} を参照

\noindent
{\it $S$ 上の肥厚化}
\ref{more-morphisms-definition-thickening} を参照

\noindent
{\it 肥厚化}
\ref{more-morphisms-definition-thickening} を参照

\noindent
{\it 肥厚化}
\ref{spaces-more-morphisms-definition-thickening} を参照

\noindent
{\it 肥厚化}
\ref{stacks-more-morphisms-definition-thickening} を参照

\noindent
{\it $T$ の位相種数}
\ref{models-definition-top-genus} を参照

\noindent
{\it 位相群}
\ref{topology-definition-topological-group} を参照

\noindent
{\it 位相加群}
\ref{topology-definition-topological-module} を参照

\noindent
{\it 位相加群}
\ref{more-algebra-definition-topological-ring} を参照

\noindent
{\it 位相環}
\ref{topology-definition-topological-ring} を参照

\noindent
{\it 位相環}
\ref{more-algebra-definition-topological-ring} を参照

\noindent
{\it 位相空間}
\ref{spaces-properties-definition-topological-space} を参照

\noindent
{\it 位相空間}
\ref{stacks-properties-definition-topological-space} を参照

\noindent
{\it 位相的冪零}
\ref{formal-spaces-definition-weakly-admissible} を参照

\noindent
{\it 位相的有限型}
\ref{formal-spaces-definition-topologically-finite-type} を参照

\noindent
{\it $\mathcal{C}$ に付随する位相}
\ref{sites-definition-topology-associated-site} を参照

\noindent
{\it $\mathcal{C}$ 上の位相}
\ref{sites-definition-topology} を参照

\noindent
{\it トポス}
\ref{sites-definition-topos} を参照

\noindent
{\it Tor 次元 $\leq d$}
\ref{more-algebra-definition-tor-amplitude} を参照

\noindent
{\it Tor 次元 $\leq d$}
\ref{cohomology-definition-tor-amplitude} を参照

\noindent
{\it Tor 次元 $\leq d$}
\ref{sites-cohomology-definition-tor-amplitude} を参照

\noindent
{\it $B$ 上 Tor 独立}
\ref{spaces-perfect-definition-tor-independent} を参照

\noindent
{\it $R$ 上 Tor 独立}
\ref{more-algebra-definition-tor-independent} を参照

\noindent
{\it $S$ 上 Tor 独立}
\ref{perfect-definition-tor-independent} を参照

\noindent
{\it $[a, b]$ における Tor 振幅}
\ref{more-algebra-definition-tor-amplitude} を参照

\noindent
{\it $[a, b]$ における Tor 振幅}
\ref{cohomology-definition-tor-amplitude} を参照

\noindent
{\it $[a, b]$ における Tor 振幅}
\ref{sites-cohomology-definition-tor-amplitude} を参照

\noindent
{\it ねじれがない}
\ref{more-algebra-definition-torsion} を参照

\noindent
{\it ねじれがない}
\ref{divisors-definition-torsion} を参照

\noindent
{\it ねじれ加群}
\ref{more-algebra-definition-torsion} を参照

\noindent
{\it ねじれ}
\ref{more-algebra-definition-torsion} を参照

\noindent
{\it ねじれ}
\ref{divisors-definition-torsion} を参照

\noindent
{\it ねじれ}
\ref{trace-definition-torsion-l-adic-sheaf} を参照

\noindent
{\it トーサー}
\ref{cohomology-definition-torsor} を参照

\noindent
{\it トーサー}
\ref{sites-cohomology-definition-torsor} を参照

\noindent
{\it Tor}
\ref{cohomology-definition-tor} を参照

\noindent
{\it Tor}
\ref{sites-cohomology-definition-tor} を参照

\noindent
{\it $X$ 上の $\mathcal{E}$ の全 Chern 類}
\ref{chow-definition-chern-classes} を参照

\noindent
{\it $\mathcal{E}$ の全 Chern 類}
\ref{spaces-chow-definition-chern-classes} を参照

\noindent
{\it 全 Chern 類}
\ref{chow-definition-chern-classes-final} を参照

\noindent
{\it $G$ の全左導来関手}
\ref{trace-definition-derived-functor} を参照

\noindent
{\it $F$ の全右導来関手}
\ref{trace-definition-derived-functor} を参照

\noindent
{\it 完全非輪状}
\ref{sites-cohomology-definition-limp} を参照

\noindent
{\it 完全不連結}
\ref{topology-definition-totally-disconnected} を参照

\noindent
{\it $A$ に関して完全分岐}
\ref{more-algebra-definition-types-of-extensions} を参照

\noindent
{\it 塔}
\ref{fields-definition-tower} を参照

\noindent
{\it トレース元}
\ref{discriminant-definition-trace-element} を参照

\noindent
{\it トレース形式}
\ref{fields-definition-trace-pairing} を参照

\noindent
{\it トレース}
\ref{fields-definition-trace-norm} を参照

\noindent
{\it トレース}
\ref{etale-cohomology-definition-trace-map} を参照

\noindent
{\it トレース}
\ref{trace-definition-trace} を参照

\noindent
{\it 超越基底}
\ref{fields-definition-transcendence} を参照

\noindent
{\it $x/f(x)$ の超越次数}
\ref{spaces-morphisms-definition-dimension-fibre} を参照

\noindent
{\it 超越次数}
\ref{fields-definition-transcendence-degree} を参照

\noindent
{\it 遷移写像}
\ref{categories-definition-system-over-poset} を参照

\noindent
{\it $0 \to K \to L \to M \to 0$ に付随する三角形}
\ref{dga-definition-distinguished-triangle} を参照

\noindent
{\it 項ごとに分裂する複体の列に付随する三角}
\ref{derived-definition-split-ses} を参照

\noindent
{\it 三角}
\ref{derived-definition-triangle} を参照

\noindent
{\it 導来圏における準連接対象の三角圏}
\ref{stacks-sheaves-definition-QC} を参照

\noindent
{\it 三角圏}
\ref{derived-definition-triangulated-category} を参照

\noindent
{\it 三角関手}
\ref{derived-definition-exact-functor-triangulated-categories} を参照

\noindent
{\it 三角部分圏}
\ref{derived-definition-triangulated-subcategory} を参照

\noindent
{\it 自明な $\mathcal{G}$-トーサー}
\ref{cohomology-definition-torsor} を参照

\noindent
{\it 自明な $\mathcal{G}$-トーサー}
\ref{sites-cohomology-definition-torsor} を参照

\noindent
{\it 自明な降下データ}
\ref{stacks-definition-effective-descent-datum} を参照

\noindent
{\it 自明な降下データ}
\ref{descent-definition-descent-datum-effective-quasi-coherent} を参照

\noindent
{\it 自明な降下データ}
\ref{descent-definition-effective} を参照

\noindent
{\it 自明な降下データ}
\ref{spaces-descent-definition-descent-datum-effective-quasi-coherent} を参照

\noindent
{\it 自明な降下データ}
\ref{spaces-descent-definition-effective} を参照

\noindent
{\it 自明 Kan ファイブレーション}
\ref{simplicial-definition-trivial-kan} を参照

\noindent
{\it 自明}
\ref{more-algebra-definition-invertible} を参照

\noindent
{\it 自明}
\ref{modules-definition-invertible} を参照

\noindent
{\it 自明}
\ref{groupoids-definition-pseudo-torsor} を参照

\noindent
{\it 自明}
\ref{spaces-groupoids-definition-pseudo-torsor} を参照

\noindent
{\it 自明}
\ref{exercises-definition-invertible-module} を参照

\noindent
{\it $\text{Proj}(S)$ の構造層の捩り}
\ref{constructions-definition-twist} を参照

\noindent
{\it 構造層のねじり}
\ref{constructions-definition-projective-bundle} を参照

\noindent
{\it 両側許容}
\ref{derived-definition-admissible} を参照

\noindent
{\it 代数構造の型}
\ref{sheaves-definition-algebraic-structure} を参照

\noindent
{\it UFD}
\ref{algebra-definition-UFD} を参照

\noindent
{\it $\mathcal{F}$ の台集合の前層}
\ref{sheaves-definition-underlying-presheaf-sets} を参照

\noindent
{\it $x$ において単枝}
\ref{properties-definition-unibranch} を参照

\noindent
{\it $x$ において単枝}
\ref{decent-spaces-definition-unibranch} を参照

\noindent
{\it 単枝}
\ref{more-algebra-definition-unibranch} を参照

\noindent
{\it 単枝}
\ref{properties-definition-unibranch} を参照

\noindent
{\it 単枝}
\ref{decent-spaces-definition-unibranch} を参照

\noindent
{\it $\mathcal{C}$ における一様圏論的モジュライ空間}
\ref{stacks-more-morphisms-definition-categorical-quotient} を参照

\noindent
{\it 一様圏論的モジュライ空間}
\ref{stacks-more-morphisms-definition-categorical-quotient} を参照

\noindent
{\it 一様圏論的商}
\ref{groupoids-quotients-definition-universal-categorical} を参照

\noindent
{\it 一意化元}
\ref{algebra-definition-uniformizer} を参照

\noindent
{\it 一様に}
\ref{groupoids-quotients-definition-topological} を参照

\noindent
{\it 一意分解整域}
\ref{algebra-definition-UFD} を参照

\noindent
{\it 付値判定法の一意性部分}
\ref{stacks-morphisms-definition-uniqueness} を参照

\noindent
{\it 普遍 $\delta$ 関手}
\ref{homology-definition-universal-delta-functor} を参照

\noindent
{\it 普遍 $\varphi$-導分}
\ref{modules-definition-module-differentials} を参照

\noindent
{\it 普遍 $\varphi$-導分}
\ref{sites-modules-definition-module-differentials} を参照

\noindent
{\it 普遍 $S$-導分}
\ref{morphisms-definition-sheaf-differentials} を参照

\noindent
{\it 普遍 $Y$-導分}
\ref{sites-modules-definition-sheaf-differentials} を参照

\noindent
{\it 普遍 $Y$-導分}
\ref{spaces-more-morphisms-definition-sheaf-differentials} を参照

\noindent
{\it 普遍圏論的商}
\ref{groupoids-quotients-definition-universal-categorical} を参照

\noindent
{\it 普遍有効エピ族}
\ref{sites-definition-universal-effective-epimorphisms} を参照

\noindent
{\it 普遍一次肥厚化}
\ref{algebra-definition-universal-thickening} を参照

\noindent
{\it 普遍一次肥厚化}
\ref{more-morphisms-definition-universal-thickening} を参照

\noindent
{\it 普遍一次肥厚化}
\ref{spaces-more-morphisms-definition-universal-thickening} を参照

\noindent
{\it $\mathcal{F}$ の普遍平坦化が存在する}
\ref{flat-definition-flattening} を参照

\noindent
{\it $\mathcal{F}$ の普遍平坦化が存在する}
\ref{spaces-flat-definition-flattening} を参照

\noindent
{\it $X$ の普遍平坦化が存在する}
\ref{flat-definition-flattening} を参照

\noindent
{\it $X$ の普遍平坦化が存在する}
\ref{spaces-flat-definition-flattening} を参照

\noindent
{\it 普遍同相射}
\ref{morphisms-definition-universal-homeomorphism} を参照

\noindent
{\it 普遍同相射}
\ref{spaces-morphisms-definition-universal-homeomorphism} を参照

\noindent
{\it 普遍同相射}
\ref{stacks-morphisms-definition-universal-homeomorphism} を参照

\noindent
{\it 普遍的 $S$-純粋}
\ref{flat-definition-pure} を参照

\noindent
{\it 普遍的に $Y$-純粋}
\ref{spaces-flat-definition-pure} を参照

\noindent
{\it 普遍鎖状}
\ref{algebra-definition-universally-catenary} を参照

\noindent
{\it 普遍鎖状}
\ref{morphisms-definition-universally-catenary} を参照

\noindent
{\it 普遍鎖状}
\ref{decent-spaces-definition-universally-catenary} を参照

\noindent
{\it 普遍閉}
\ref{topology-definition-proper-map} を参照

\noindent
{\it 普遍閉}
\ref{schemes-definition-universally-closed} を参照

\noindent
{\it 普遍閉}
\ref{spaces-morphisms-definition-closed} を参照

\noindent
{\it 普遍閉}
\ref{stacks-morphisms-definition-closed} を参照

\noindent
{\it 純完全}
\ref{algebra-definition-universally-injective} を参照

\noindent
{\it 普遍単射}
\ref{algebra-definition-universally-injective} を参照

\noindent
{\it 普遍単射}
\ref{morphisms-definition-universally-injective} を参照

\noindent
{\it 普遍単射}
\ref{descent-definition-universally-injective} を参照

\noindent
{\it 普遍単射}
\ref{spaces-morphisms-definition-universally-injective} を参照

\noindent
{\it 普遍単射}
\ref{stacks-morphisms-definition-universally-injective} を参照

\noindent
{\it 普遍的日本スキーム}
\ref{algebra-definition-nagata} を参照

\noindent
{\it 普遍的日本スキーム}
\ref{properties-definition-nagata} を参照

\noindent
{\it $K$ に関して普遍局所非輪状}
\ref{etale-cohomology-definition-locally-acyclic} を参照

\noindent
{\it 普遍局所非輪状}
\ref{etale-cohomology-definition-locally-acyclic} を参照

\noindent
{\it 普遍開}
\ref{morphisms-definition-open} を参照

\noindent
{\it 普遍開}
\ref{spaces-morphisms-definition-open} を参照

\noindent
{\it 普遍開}
\ref{stacks-morphisms-definition-open} を参照

\noindent
{\it $y$ 上で普遍的に純粋}
\ref{spaces-flat-definition-pure} を参照

\noindent
{\it $X_s$ に沿って普遍的に純粋}
\ref{flat-definition-pure} を参照

\noindent
{\it $S$ に関して普遍的に純粋}
\ref{flat-definition-pure} を参照

\noindent
{\it $Y$ に関して普遍的に純粋}
\ref{spaces-flat-definition-pure} を参照

\noindent
{\it 普遍的に商写像的}
\ref{morphisms-definition-submersive} を参照

\noindent
{\it 普遍的に商写像的}
\ref{spaces-morphisms-definition-submersive} を参照

\noindent
{\it 普遍的に商写像的}
\ref{stacks-morphisms-definition-submersive} を参照

\noindent
{\it 普遍的に}
\ref{groupoids-quotients-definition-topological} を参照

\noindent
{\it 非障害的}
\ref{formal-defos-definition-cofibered-groupoid-projection-smooth} を参照

\noindent
{\it $\mathfrak q$ で不分岐}
\ref{algebra-definition-unramified} を参照

\noindent
{\it $x \in X$ において不分岐}
\ref{morphisms-definition-unramified} を参照

\noindent
{\it $x$ において不分岐}
\ref{etale-definition-unramified-schemes} を参照

\noindent
{\it $x$ において不分岐}
\ref{spaces-morphisms-definition-unramified} を参照

\noindent
{\it $\text{GL}_2(\mathbf{A})$ 上の $\Lambda$ に値を取る 不分岐尖点形式}
\ref{trace-definition-unramified} を参照

\noindent
{\it 局所環の不分岐準同型}
\ref{etale-definition-unramified-rings} を参照

\noindent
{\it $A$ に関して不分岐}
\ref{more-algebra-definition-types-of-extensions} を参照

\noindent
{\it 不分岐}
\ref{algebra-definition-unramified} を参照

\noindent
{\it 不分岐}
\ref{morphisms-definition-unramified} を参照

\noindent
{\it 不分岐}
\ref{etale-definition-unramified-schemes} を参照

\noindent
{\it 不分岐}
\ref{spaces-morphisms-definition-unramified} を参照

\noindent
{\it 不分岐}
\ref{stacks-morphisms-definition-unramified} を参照

\noindent
{\it $T$ の V 被覆}
\ref{topologies-definition-V-covering} を参照

\noindent
{\it 付値環}
\ref{algebra-definition-valuation-ring} を参照

\noindent
{\it 付値}
\ref{algebra-definition-value-group} を参照

\noindent
{\it 値群}
\ref{algebra-definition-value-group} を参照

\noindent
{\it $X$ における $LF$ の値}
\ref{derived-definition-right-derived-functor-defined} を参照

\noindent
{\it $X$ における $RF$ の値}
\ref{derived-definition-right-derived-functor-defined} を参照

\noindent
{\it 本質的に定値}
\ref{categories-definition-essentially-constant-diagram} を参照

\noindent
{\it 本質的に定値}
\ref{categories-definition-essentially-constant-diagram} を参照

\noindent
{\it 多様体}
\ref{varieties-definition-variety} を参照

\noindent
{\it 多様体}
\ref{etale-cohomology-definition-variety} を参照

\noindent
{\it ベクトル束 $\pi : V \to S$（$S$ 上）}
\ref{constructions-definition-abstract-vector-bundle} を参照

\noindent
{\it $\mathcal{E}$ に付随するベクトル束}
\ref{constructions-definition-vector-bundle} を参照

\noindent
{\it $x_0$ における $\mathcal{X}$ の半普遍変形環}
\ref{stacks-geometry-definition-versal-ring-at-x} を参照

\noindent
{\it 半普遍的}
\ref{formal-defos-definition-versal} を参照

\noindent
{\it 半普遍的}
\ref{artin-definition-versal-formal-object} を参照

\noindent
{\it 半普遍的}
\ref{artin-definition-versal} を参照

\noindent
{\it 垂直}
\ref{categories-definition-2-category} を参照

\noindent
{\it $X/S$ 上非常に豊富}
\ref{morphisms-definition-very-ample} を参照

\noindent
{\it 非常に妥当}
\ref{decent-spaces-definition-very-reasonable} を参照

\noindent
{\it 非常に妥当}
\ref{decent-spaces-definition-relative-conditions} を参照

\noindent
{\it $S'$ 上の代数空間とみなしたもの}
\ref{spaces-definition-base-change} を参照

\noindent
{\it $S'$ 上の代数スタックとみなしたもの}
\ref{algebraic-definition-viewed-as} を参照

\noindent
{\it w-可縮}
\ref{proetale-definition-w-contractible} を参照

\noindent
{\it w-局所的}
\ref{proetale-definition-w-local} を参照

\noindent
{\it w-局所的}
\ref{proetale-definition-w-local} を参照

\noindent
{\it 弱 $R$-軌道}
\ref{groupoids-quotients-definition-geometric-orbits} を参照

\noindent
{\it 弱次元は $\leq d$}
\ref{more-algebra-definition-weak-dimension} を参照

\noindent
{\it 弱関手}
\ref{categories-definition-functor-into-2-category} を参照

\noindent
{\it 弱生成元}
\ref{derived-definition-generators} を参照

\noindent
{\it 弱定義イデアル}
\ref{formal-spaces-definition-weakly-admissible} を参照

\noindent
{\it $X$ の $Y$ における弱正規化}
\ref{morphisms-definition-relative-seminormalization} を参照

\noindent
{\it 弱正規化}
\ref{morphisms-definition-weak-normalization} を参照

\noindent
{\it 弱軌道}
\ref{groupoids-quotients-definition-geometric-orbits} を参照

\noindent
{\it 弱セール部分圏}
\ref{homology-definition-serre-subcategory} を参照

\noindent
{\it $A \subset B$ に対する弱解}
\ref{more-algebra-definition-solution} を参照

\noindent
{\it 標準位相より弱い}
\ref{sites-definition-weaker-than-canonical} を参照

\noindent
{\it より弱い}
\ref{sites-definition-finer} を参照

\noindent
{\it 弱 $R$-同値}
\ref{groupoids-quotients-definition-geometric-orbits} を参照

\noindent
{\it 弱エタール}
\ref{more-algebra-definition-weakly-etale} を参照

\noindent
{\it 弱エタール}
\ref{more-morphisms-definition-weakly-etale} を参照

\noindent
{\it 弱アディック}
\ref{formal-spaces-definition-weakly-adic} を参照

\noindent
{\it 弱アディック}
\ref{formal-spaces-definition-types-affine-formal-algebraic-space} を参照

\noindent
{\it 弱許容}
\ref{formal-spaces-definition-weakly-admissible} を参照

\noindent
{\it $X$ の弱随伴点}
\ref{divisors-definition-weakly-associated} を参照

\noindent
{\it $X$ の弱随伴点}
\ref{spaces-divisors-definition-weakly-associated} を参照

\noindent
{\it 弱随伴する}
\ref{algebra-definition-weakly-associated} を参照

\noindent
{\it 弱随伴する}
\ref{divisors-definition-weakly-associated} を参照

\noindent
{\it 弱随伴する}
\ref{spaces-divisors-definition-weakly-associated} を参照

\noindent
{\it 弱可縮}
\ref{sites-definition-w-contractible} を参照

\noindent
{\it $H(K)$ に弱収束する}
\ref{homology-definition-filtered-differential-ss-converges} を参照

\noindent
{\it $H^*(K^\bullet)$ に弱収束する}
\ref{homology-definition-filtered-complex-ss-converges} を参照

\noindent
{\it $H^n(\text{Tot}(K^{\bullet, \bullet}))$ に弱収束する}
\ref{homology-definition-ss-double-complex-converge} を参照

\noindent
{\it $H^n(\text{Tot}(K^{\bullet, \bullet}))$ に弱収束する}
\ref{homology-definition-ss-double-complex-converge} を参照

\noindent
{\it 弱正規}
\ref{morphisms-definition-weakly-normal} を参照

\noindent
{\it 弱前アディック}
\ref{formal-spaces-definition-weakly-adic} を参照

\noindent
{\it 弱前許容}
\ref{formal-spaces-definition-weakly-admissible} を参照

\noindent
{\it 弱不分岐}
\ref{more-algebra-definition-extension-discrete-valuation-rings} を参照

\noindent
{\it 弱不分岐}
\ref{more-algebra-definition-extension-valuation-rings} を参照

\noindent
{\it 重み付け}
\ref{more-morphisms-definition-weighting} を参照

\noindent
{\it Weil コホモロジー理論}
\ref{weil-definition-weil-cohomology-theory} を参照

\noindent
{\it 有効 Cartier 因子 $D \subset X$ に付随する Weil 因子 $[D]$}
\ref{exercises-definition-divisor} を参照

\noindent
{\it $\mathcal{L}$ に付随する Weil 因子}
\ref{chow-definition-divisor-invertible-sheaf} を参照

\noindent
{\it $\mathcal{L}$ に付随する Weil 因子}
\ref{spaces-chow-definition-divisor-invertible-sheaf} を参照

\noindent
{\it $s$ に付随する Weil 因子}
\ref{divisors-definition-divisor-invertible-sheaf} を参照

\noindent
{\it $s$ に付随する Weil 因子}
\ref{chow-definition-divisor-invertible-sheaf} を参照

\noindent
{\it $s$ に付随する Weil 因子}
\ref{spaces-over-fields-definition-divisor-invertible-sheaf} を参照

\noindent
{\it $s$ に付随する Weil 因子}
\ref{spaces-chow-definition-divisor-invertible-sheaf} を参照

\noindent
{\it 付随する Weil 因子}
\ref{exercises-definition-divisor} を参照

\noindent
{\it 有理関数 $f \in K(X)^\ast$ に付随する Weil 因子}
\ref{exercises-definition-divisor} を参照

\noindent
{\it $\mathcal{L}$ に付随する Weil 因子類}
\ref{divisors-definition-divisor-invertible-sheaf} を参照

\noindent
{\it $\mathcal{L}$ に付随する Weil 因子類}
\ref{spaces-over-fields-definition-divisor-invertible-sheaf} を参照

\noindent
{\it Weil 因子類群}
\ref{divisors-definition-class-group} を参照

\noindent
{\it Weil 因子類群}
\ref{spaces-over-fields-definition-class-group} を参照

\noindent
{\it Weil 因子}
\ref{divisors-definition-Weil-divisor} を参照

\noindent
{\it Weil 因子}
\ref{spaces-over-fields-definition-Weil-divisor} を参照

\noindent
{\it Weil 因子}
\ref{exercises-definition-divisor} を参照

\noindent
{\it ほとんどアフィン}
\ref{stacks-more-morphisms-definition-well-nigh-affine} を参照

\noindent
{\it 半表現可能対象に前層を対応させる関手}
\ref{hypercovering-definition-SR-F} を参照

\noindent
{\it $\mathfrak m$ の野性惰性群}
\ref{more-algebra-definition-wild-inertia} を参照

\noindent
{\it Yoneda 拡大}
\ref{derived-definition-yoneda-extension} を参照

\noindent
{\it $T$ の Zariski 被覆}
\ref{topologies-definition-zariski-covering} を参照

\noindent
{\it $X$ の Zariski 被覆}
\ref{spaces-topologies-definition-zariski-covering} を参照

\noindent
{\it Zariski 被覆}
\ref{spaces-definition-Zariski-open-covering} を参照

\noindent
{\it $S$ 上 Zariski 局所準分離的}
\ref{spaces-definition-separated} を参照

\noindent
{\it Zariski 局所準分離的}
\ref{spaces-properties-definition-separated} を参照

\noindent
{\it Zariski 局所準分離的}
\ref{spaces-properties-definition-separated} を参照

\noindent
{\it Zariski 対}
\ref{more-algebra-definition-zariski-pair} を参照

\noindent
{\it Zariski 層}
\ref{stacks-sheaves-definition-sheaves} を参照

\noindent
{\it Zariski トポス}
\ref{etale-cohomology-definition-etale-topos} を参照

\noindent
{\it Zariski、エタール、滑らか、シントミック、または fppf 被覆}
\ref{crystalline-definition-big-crystalline-site} を参照

\noindent
{\it ザリスキー}
\ref{algebra-definition-Zariski-topology} を参照

\noindent
{\it 零対象}
\ref{homology-definition-zero-object} を参照

\noindent
{\it 零点スキーム}
\ref{divisors-definition-zero-scheme-s} を参照

\noindent
{\it 零点スキーム}
\ref{spaces-divisors-definition-zero-scheme-s} を参照

\noindent
{\it $\mathcal{A}$ の第零 $K$ 群}
\ref{homology-definition-K-zero} を参照

\noindent
{\it $\mathcal{D}$ の第零 $K$ 群}
\ref{derived-definition-K-zero} を参照

\noindent
{\it 第零 {\v C}ech コホモロジー群}
\ref{etale-cohomology-definition-0-cech} を参照

\noindent
{\it {\v C}ech コホモロジー群}
\ref{cohomology-definition-cech-complex} を参照

\noindent
{\it {\v C}ech コホモロジー群}
\ref{sites-cohomology-definition-cech-complex} を参照

\noindent
{\it {\v C}ech コホモロジー群}
\ref{etale-cohomology-definition-cech-complex} を参照

\noindent
{\it {\v C}ech 複体}
\ref{cohomology-definition-cech-complex} を参照

\noindent
{\it {\v C}ech 複体}
\ref{sites-cohomology-definition-cech-complex} を参照

\noindent
{\it {\v C}ech 複体}
\ref{etale-cohomology-definition-cech-complex} を参照

\end{multicols}

\begin{multicols}{2}[\section{章別の定義}\label{index-section-per-chapter}]

\medskip\noindent
{\bf はじめに}

\medskip

\medskip\noindent
{\bf 規約}

\medskip

\medskip\noindent
{\bf 集合}

\medskip

\medskip\noindent
{\bf 圏}

\medskip

\noindent
\ref{categories-definition-category} を参照：
{\it 圏}


\noindent
\ref{categories-definition-isomorphism} を参照：
{\it 同型射}


\noindent
\ref{categories-definition-groupoid} を参照：
{\it グルーポイド}


\noindent
\ref{categories-definition-functor} を参照：
{\it 関手}


\noindent
\ref{categories-definition-faithful} を参照：
{\it 忠実},
{\it 充満忠実},
{\it 本質的全射}


\noindent
\ref{categories-definition-subcategory} を参照：
{\it 部分圏},
{\it 充満部分圏},
{\it 厳密充満}


\noindent
\ref{categories-definition-transformation-functors} を参照：
{\it 自然変換},
{\it 関手の射}


\noindent
\ref{categories-definition-equivalence-categories} を参照：
{\it 圏同値},
{\it 擬逆関手}


\noindent
\ref{categories-definition-product-category} を参照：
{\it 直積圏}


\noindent
\ref{categories-definition-opposite} を参照：
{\it 反対圏}


\noindent
\ref{categories-definition-contravariant} を参照：
{\it 反変}


\noindent
\ref{categories-definition-presheaf} を参照：
{\it $\mathcal{C}$ 上の集合の前層},
{\it 前層}


\noindent
\ref{categories-definition-representable-functor} を参照：
{\it 表現可能}


\noindent
\ref{categories-definition-products} を参照：
{\it 積}


\noindent
\ref{categories-definition-has-products-of-pairs} を参照：
{\it 二対象の積をもつ}


\noindent
\ref{categories-definition-coproducts} を参照：
{\it 余積},
{\it 融合和}


\noindent
\ref{categories-definition-has-coproducts-of-pairs} を参照：
{\it 二対象の余積をもつ}


\noindent
\ref{categories-definition-fibre-products} を参照：
{\it ファイバー積}


\noindent
\ref{categories-definition-cartesian} を参照：
{\it カルテジアン}


\noindent
\ref{categories-definition-has-fibre-products} を参照：
{\it ファイバー積をもつ}


\noindent
\ref{categories-definition-representable-morphism} を参照：
{\it 表現可能}


\noindent
\ref{categories-definition-representable-map-presheaves} を参照：
{\it 表現可能},
{\it $F$ は $G$ 上相対的に表現可能}


\noindent
\ref{categories-definition-pushouts} を参照：
{\it 押し出し}


\noindent
\ref{categories-definition-cocartesian} を参照：
{\it 余カルテジアン}


\noindent
\ref{categories-definition-equalizers} を参照：
{\it 等化子}


\noindent
\ref{categories-definition-coequalizers} を参照：
{\it 余等化子}


\noindent
\ref{categories-definition-initial-final} を参照：
{\it 始対象},
{\it 終対象}


\noindent
\ref{categories-definition-mono-epi} を参照：
{\it モノ射},
{\it エピ射}


\noindent
\ref{categories-definition-limit} を参照：
{\it 極限}


\noindent
\ref{categories-definition-colimit} を参照：
{\it 余極限}


\noindent
\ref{categories-definition-product} を参照：
{\it 積}


\noindent
\ref{categories-definition-coproduct} を参照：
{\it 余積}


\noindent
\ref{categories-definition-category-connected} を参照：
{\it 連結}


\noindent
\ref{categories-definition-cofinal} を参照：
{\it $\mathcal{I}$ は $\mathcal{J}$ において共終},
{\it 共終}


\noindent
\ref{categories-definition-initial} を参照：
{\it $\mathcal{I}$ は $\mathcal{J}$ において始的},
{\it 始対象}


\noindent
\ref{categories-definition-directed} を参照：
{\it 有向},
{\it フィルター付き},
{\it 有向},
{\it フィルター付き}


\noindent
\ref{categories-definition-codirected} を参照：
{\it 余有向},
{\it 余フィルター付き},
{\it 余有向},
{\it 余フィルター付き}


\noindent
\ref{categories-definition-directed-set} を参照：
{\it 前順序},
{\it 前順序集合},
{\it 有向集合},
{\it 半順序},
{\it 半順序集合},
{\it 有向半順序集合}


\noindent
\ref{categories-definition-system-over-poset} を参照：
{\it $\mathcal{C}$ における $I$ 上の系},
{\it $\mathcal{C}$ における $I$ 上の帰納系},
{\it $\mathcal{C}$ における $I$ 上の逆系},
{\it $\mathcal{C}$ における $I$ 上の射影系},
{\it 遷移写像}


\noindent
\ref{categories-definition-directed-system} を参照：
{\it 有向系},
{\it 有向逆系}


\noindent
\ref{categories-definition-essentially-constant-diagram} を参照：
{\it 値},
{\it 本質的に定値},
{\it 値},
{\it 本質的に定値}


\noindent
\ref{categories-definition-essentially-constant-system} を参照：
{\it 本質的に定値な系},
{\it 本質的に定値な逆系}


\noindent
\ref{categories-definition-exact} を参照：
{\it 左完全},
{\it 右完全},
{\it 完全}


\noindent
\ref{categories-definition-adjoint} を参照：
{\it 左随伴},
{\it 右随伴}


\noindent
\ref{categories-definition-compact-object} を参照：
{\it 圏論的コンパクト}


\noindent
\ref{categories-definition-multiplicative-system} を参照：
{\it 左乗法系},
{\it 右乗法系},
{\it 乗法系}


\noindent
\ref{categories-definition-left-localization-as-fraction} を参照：
{\it $s^{-1} f$}


\noindent
\ref{categories-definition-right-localization-as-fraction} を参照：
{\it $f s^{-1}$}


\noindent
\ref{categories-definition-saturated-multiplicative-system} を参照：
{\it 飽和している}


\noindent
\ref{categories-definition-horizontal-composition} を参照：
{\it 水平}


\noindent
\ref{categories-definition-2-category} を参照：
{\it $2$-圏},
{\it $1$-射},
{\it $2$-射},
{\it 垂直},
{\it 合成},
{\it 水平}


\noindent
\ref{categories-definition-sub-2-category} を参照：
{\it 部分 $2$-圏}


\noindent
\ref{categories-definition-equivalence} を参照：
{\it 同値}


\noindent
\ref{categories-definition-functor-into-2-category} を参照：
{\it 関手},
{\it 弱関手},
{\it 擬関手}


\noindent
\ref{categories-definition-2-1-category} を参照：
{\it $(2, 1)$-圏}


\noindent
\ref{categories-definition-final-object-2-category} を参照：
{\it 終対象}


\noindent
\ref{categories-definition-2-fibre-products} を参照：
{\it $f$ と $g$ の 2-ファイバー積}


\noindent
\ref{categories-definition-categories-over-C} を参照：
{\it $\mathcal{C}$ 上の圏の $2$-圏}


\noindent
\ref{categories-definition-fibre-category} を参照：
{\it ファイバー圏},
{\it 持ち上げ},
{\it $x$ は $U$ 上にある},
{\it 持ち上げ},
{\it $\phi$ は $f$ 上にある}


\noindent
\ref{categories-definition-cartesian-over-C} を参照：
{\it 強カルテシアン射},
{\it 強 $\mathcal{C}$-カルテシアン射}


\noindent
\ref{categories-definition-fibred-category} を参照：
{\it $\mathcal{C}$ 上のファイバー圏}


\noindent
\ref{categories-definition-pullback-functor-fibred-category} を参照：
{\it 引き戻しの選択},
{\it 引き戻し関手}


\noindent
\ref{categories-definition-fibred-categories-over-C} を参照：
{\it $\mathcal{C}$ 上のファイバー圏の $2$-圏}


\noindent
\ref{categories-definition-inertia-fibred-category} を参照：
{\it $\mathcal{S}'$ 上の $\mathcal{S}$ の相対惰性},
{\it $\mathcal{S}$ の惰性ファイバー圏 $\mathcal{I}_\mathcal{S}$}


\noindent
\ref{categories-definition-fibred-groupoids} を参照：
{\it グルーポイドにファイバー化されている}


\noindent
\ref{categories-definition-categories-fibred-in-groupoids-over-C} を参照：
{\it $\mathcal{C}$ 上グルーポイドに ファイバー化された圏の $2$-圏}


\noindent
\ref{categories-definition-split-fibred-category} を参照：
{\it 分裂ファイバー圏},
{\it $\mathcal{S}_F$}


\noindent
\ref{categories-definition-split-category-fibred-in-groupoids} を参照：
{\it 分裂亜群ファイバー圏},
{\it $\mathcal{S}_F$}


\noindent
\ref{categories-definition-discrete} を参照：
{\it 離散的}


\noindent
\ref{categories-definition-category-fibred-sets} を参照：
{\it 集合にファイバー化された圏},
{\it 離散圏にファイバー化された圏}


\noindent
\ref{categories-definition-categories-fibred-in-sets-over-C} を参照：
{\it $\mathcal{C}$ 上集合にファイバー化された圏の $2$-圏}


\noindent
\ref{categories-definition-setoid} を参照：
{\it セトイド}


\noindent
\ref{categories-definition-category-fibred-setoids} を参照：
{\it セトイドにファイバー化された圏}


\noindent
\ref{categories-definition-categories-fibred-in-setoids-over-C} を参照：
{\it $\mathcal{C}$ 上セトイドにファイバー化された圏の $2$-圏}


\noindent
\ref{categories-definition-representable-fibred-category} を参照：
{\it 表現可能}


\noindent
\ref{categories-definition-representable-map-categories-fibred-in-groupoids} を参照：
{\it 表現可能},
{\it $\mathcal{X}$ が $\mathcal{Y}$ 上で相対的に表現可能}


\noindent
\ref{categories-definition-monoidal-category} を参照：
{\it モノイド圏}


\noindent
\ref{categories-definition-functor-monoidal-categories} を参照：
{\it モノイド圏の関手}


\noindent
\ref{categories-definition-invertible} を参照：
{\it 可逆}


\noindent
\ref{categories-definition-dual} を参照：
{\it 左双対},
{\it 右双対}


\noindent
\ref{categories-definition-symmetric-monoidal-category} を参照：
{\it 対称モノイド圏}


\noindent
\ref{categories-definition-functor-symmetric-monoidal-categories} を参照：
{\it 対称モノイド圏の関手}


\noindent
\ref{categories-definition-dotted-arrows} を参照：
{\it 点線矢印の射}


\medskip\noindent
{\bf 位相空間}

\medskip

\noindent
\ref{topology-definition-separated} を参照：
{\it 分離的}


\noindent
\ref{topology-definition-base} を参照：
{\it $X$ の位相の開基},
{\it $X$ の位相の基底}


\noindent
\ref{topology-definition-subbase} を参照：
{\it $X$ の位相の準開基},
{\it $X$ の位相の部分基底}


\noindent
\ref{topology-definition-submersive} を参照：
{\it 位相空間の厳密写像},
{\it 商位相的}


\noindent
\ref{topology-definition-connected-components} を参照：
{\it 連結},
{\it 連結成分}


\noindent
\ref{topology-definition-totally-disconnected} を参照：
{\it 完全不連結}


\noindent
\ref{topology-definition-locally-connected} を参照：
{\it 局所連結}


\noindent
\ref{topology-definition-irreducible-components} を参照：
{\it 既約},
{\it 既約成分}


\noindent
\ref{topology-definition-generic-point} を参照：
{\it 生成点},
{\it コルモゴロフ},
{\it 準ソバー},
{\it ソバー}


\noindent
\ref{topology-definition-noetherian} を参照：
{\it Noether},
{\it 局所 Noether}


\noindent
\ref{topology-definition-Krull} を参照：
{\it 既約閉部分集合の鎖},
{\it 長さ},
{\it 次元},
{\it クルル次元},
{\it $x$ における $X$ のクルル次元}


\noindent
\ref{topology-definition-equidimensional} を参照：
{\it 等次元}


\noindent
\ref{topology-definition-codimension} を参照：
{\it 余次元}


\noindent
\ref{topology-definition-catenary} を参照：
{\it 鎖状}


\noindent
\ref{topology-definition-quasi-compact} を参照：
{\it 準コンパクト},
{\it 準コンパクト},
{\it レトロコンパクト}


\noindent
\ref{topology-definition-locally-quasi-compact} を参照：
{\it 局所準コンパクト}


\noindent
\ref{topology-definition-constructible} を参照：
{\it 構成可能},
{\it 局所構成可能}


\noindent
\ref{topology-definition-proper-map} を参照：
{\it 閉},
{\it Bourbaki 固有},
{\it 準固有},
{\it 普遍閉},
{\it 固有}


\noindent
\ref{topology-definition-space-jacobson} を参照：
{\it Jacobson}


\noindent
\ref{topology-definition-specialization} を参照：
{\it 特殊化},
{\it 一般化},
{\it 特殊化の下で安定},
{\it 一般化の下で安定}


\noindent
\ref{topology-definition-lift-specializations} を参照：
{\it 特殊化は $f$ に沿って 持ち上がる},
{\it 特殊化的},
{\it 一般化は $f$ に沿って 持ち上がる},
{\it 一般化的}


\noindent
\ref{topology-definition-dimension-function} を参照：
{\it 直近特殊化},
{\it 次元関数}


\noindent
\ref{topology-definition-nowhere-dense} を参照：
{\it 内部},
{\it 至る所疎}


\noindent
\ref{topology-definition-profinite} を参照：
{\it 副有限}


\noindent
\ref{topology-definition-spectral-space} を参照：
{\it スペクトル},
{\it スペクトル}


\noindent
\ref{topology-definition-extremally-disconnected} を参照：
{\it 極値非連結}


\noindent
\ref{topology-definition-isolated-point} を参照：
{\it 孤立点}


\noindent
\ref{topology-definition-paritition} を参照：
{\it 分割},
{\it 部分},
{\it 細分する}


\noindent
\ref{topology-definition-good-stratification} を参照：
{\it 良い層化}


\noindent
\ref{topology-definition-stratification} を参照：
{\it 層化},
{\it 層}


\noindent
\ref{topology-definition-locally-finite} を参照：
{\it 局所有限}


\noindent
\ref{topology-definition-topological-group} を参照：
{\it 位相群},
{\it 位相群の準同型}


\noindent
\ref{topology-definition-profinite-group} を参照：
{\it 副有限群}


\noindent
\ref{topology-definition-topological-ring} を参照：
{\it 位相環},
{\it 位相環の準同型}


\noindent
\ref{topology-definition-topological-module} を参照：
{\it 位相加群},
{\it 位相加群の準同型}


\medskip\noindent
{\bf 空間上の層}

\medskip

\noindent
\ref{sheaves-definition-presheaf} を参照：
{\it 集合の前層 $\mathcal{F}$（$X$ 上）},
{\it 集合の前層の射 $\varphi : \mathcal{F} \to \mathcal{G}$ （$X$ 上）}


\noindent
\ref{sheaves-definition-constant-presheaf} を参照：
{\it 値 $A$ の定数前層}


\noindent
\ref{sheaves-definition-abelian-presheaves} を参照：
{\it $X$ 上のアーベル群の前層},
{\it $X$ 上のアーベル前層},
{\it $X$ 上のアーベル前層の射}


\noindent
\ref{sheaves-definition-presheaf-values-in-category} を参照：
{\it $\mathcal{C}$ に値をとる $X$ 上の前層 $\mathcal{F}$},
{\it $\mathcal{C}$ に値をとる前層の射 $\varphi : \mathcal{F} \to \mathcal{G}$}


\noindent
\ref{sheaves-definition-underlying-presheaf-sets} を参照：
{\it $\mathcal{F}$ の台集合の前層}


\noindent
\ref{sheaves-definition-presheaf-modules} を参照：
{\it $\mathcal{O}$-加群の前層},
{\it $\mathcal{O}$-加群の前層の射 $\varphi : \mathcal{F} \to \mathcal{G}$}


\noindent
\ref{sheaves-definition-sheaf} を参照：
{\it 集合の層 $\mathcal{F}$（$X$ 上）},
{\it 集合の層の射}


\noindent
\ref{sheaves-definition-constant-sheaf} を参照：
{\it 値 $A$ の定数層}


\noindent
\ref{sheaves-definition-abelian-sheaf} を参照：
{\it $X$ 上のアーベル層},
{\it $X$ 上のアーベル群の層}


\noindent
\ref{sheaves-definition-sheaf-values-in-category} を参照：
{\it 層}


\noindent
\ref{sheaves-definition-sheaf-modules} を参照：
{\it $\mathcal{O}$-加群の層},
{\it $\mathcal{O}$-加群の層の射}


\noindent
\ref{sheaves-definition-separated} を参照：
{\it 分離的}


\noindent
\ref{sheaves-definition-algebraic-structure} を参照：
{\it 代数構造の型}


\noindent
\ref{sheaves-definition-injective-surjective} を参照：
{\it 部分前層},
{\it 部分層},
{\it 入射的},
{\it 全射},
{\it 入射的},
{\it 全射}


\noindent
\ref{sheaves-definition-f-map} を参照：
{\it $f$-写像 $\xi : \mathcal{G} \to \mathcal{F}$}


\noindent
\ref{sheaves-definition-composition-f-maps} を参照：
{\it $\varphi$ と $\psi$ の合成}


\noindent
\ref{sheaves-definition-ringed-space} を参照：
{\it 環付き空間},
{\it 環付き空間の射}


\noindent
\ref{sheaves-definition-composition-maps-ringed-spaces} を参照：
{\it 環付き空間の射の合成}


\noindent
\ref{sheaves-definition-pushforward} を参照：
{\it 順像},
{\it 引き戻し}


\noindent
\ref{sheaves-definition-skyscraper-sheaf} を参照：
{\it 値 $A$ をもつ $x$ における摩天楼層},
{\it 摩天楼層},
{\it 摩天楼層},
{\it 摩天楼層},
{\it 摩天楼層}


\noindent
\ref{sheaves-definition-presheaf-basis} を参照：
{\it $\mathcal{B}$ 上の集合の前層 $\mathcal{F}$},
{\it $\mathcal{B}$ 上の集合の前層の射 $\varphi : \mathcal{F} \to \mathcal{G}$}


\noindent
\ref{sheaves-definition-sheaf-basis} を参照：
{\it $\mathcal{B}$ 上の集合の層 $\mathcal{F}$},
{\it $\mathcal{B}$ 上の集合の層の射}


\noindent
\ref{sheaves-definition-sheaf-structures-basis} を参照：
{\it $\mathcal{B}$ 上の $\mathcal{C}$ に値をもつ前層 $\mathcal{F}$},
{\it $\mathcal{B}$ 上の $\mathcal{C}$ に値をもつ前層の射 $\varphi : \mathcal{F} \to \mathcal{G}$},
{\it $\mathcal{B}$ 上の $\mathcal{C}$ に値をもつ層 $\mathcal{F}$}


\noindent
\ref{sheaves-definition-sheaf-modules-basis} を参照：
{\it $\mathcal{B}$ 上の $\mathcal{O}$-加群の前層 $\mathcal{F}$},
{\it $\mathcal{B}$ 上の $\mathcal{O}$-加群の前層の射 $\varphi : \mathcal{F} \to \mathcal{G}$},
{\it $\mathcal{B}$ 上の $\mathcal{O}$-加群の層 $\mathcal{F}$}


\noindent
\ref{sheaves-definition-restriction} を参照：
{\it $\mathcal{G}$ の $U$ への制限},
{\it $\mathcal{G}$ の $U$ への制限},
{\it $U$ に付随する $(X, \mathcal{O})$ の開部分空間},
{\it $\mathcal{G}$ の $U$ への制限}


\noindent
\ref{sheaves-definition-j-shriek} を参照：
{\it $\mathcal{F}$ の空集合による延長 $j_{p!}\mathcal{F}$},
{\it $\mathcal{F}$ の空集合による延長 $j_!\mathcal{F}$}


\noindent
\ref{sheaves-definition-j-shriek-structures} を参照：
{\it $\mathcal{F}$ の $0$ による延長 $j_{p!}\mathcal{F}$},
{\it $\mathcal{F}$ の $0$ による延長 $j_!\mathcal{F}$},
{\it $\mathcal{F}$ の $e$ による延長 $j_{p!}\mathcal{F}$},
{\it $\mathcal{F}$ の $e$ による延長 $j_!\mathcal{F}$},
{\it $0$ による延長},
{\it $0$ による延長}


\medskip\noindent
{\bf サイトと層}

\medskip

\noindent
\ref{sites-definition-presheaves-sets} を参照：
{\it 集合の前層},
{\it 前層の射}


\noindent
\ref{sites-definition-presheaf} を参照：
{\it 前層},
{\it 射}


\noindent
\ref{sites-definition-presheaves-injective-surjective} を参照：
{\it 入射的},
{\it 全射}


\noindent
\ref{sites-definition-sub-presheaf} を参照：
{\it 部分前層}


\noindent
\ref{sites-definition-image} を参照：
{\it $\varphi$ の像}


\noindent
\ref{sites-definition-family-morphisms-fixed-target} を参照：
{\it 終域を固定した射の族}


\noindent
\ref{sites-definition-site} を参照：
{\it サイト},
{\it $\mathcal{C}$ の被覆}


\noindent
\ref{sites-definition-sheaf-sets} を参照：
{\it 層}


\noindent
\ref{sites-definition-category-sheaves-sets} を参照：
{\it $\Sh(\mathcal{C})$}


\noindent
\ref{sites-definition-sheaf} を参照：
{\it 層}


\noindent
\ref{sites-definition-morphism-coverings} を参照：
{\it $\mathcal{C}$ の終域を固定した射の族の $\mathcal{U}$ から $\mathcal{V}$ への射},
{\it $\mathcal{U}$ から $\mathcal{V}$ への射},
{\it 細分}


\noindent
\ref{sites-definition-combinatorial-tautological} を参照：
{\it 組合せ的に同値},
{\it 自明に同値}


\noindent
\ref{sites-definition-separated} を参照：
{\it 分離的}


\noindent
\ref{sites-definition-associated-sheaf} を参照：
{\it $\mathcal{F}$ に付随する層}


\noindent
\ref{sites-definition-sheaves-injective-surjective} を参照：
{\it 入射的},
{\it 全射}


\noindent
\ref{sites-definition-universal-effective-epimorphisms} を参照：
{\it 有効エピ族},
{\it 普遍有効エピ族}


\noindent
\ref{sites-definition-weaker-than-canonical} を参照：
{\it 標準位相より弱い},
{\it 準標準的}


\noindent
\ref{sites-definition-representable-sheaf} を参照：
{\it 表現可能層},
{\it $\underline{U}$}


\noindent
\ref{sites-definition-continuous} を参照：
{\it 連続}


\noindent
\ref{sites-definition-morphism-sites} を参照：
{\it サイトの射}


\noindent
\ref{sites-definition-composition-morphisms-sites} を参照：
{\it 合成}


\noindent
\ref{sites-definition-topos} を参照：
{\it トポス},
{\it トポスの射},
{\it 合成 $f\circ g$}


\noindent
\ref{sites-definition-quasi-compact} を参照：
{\it 準コンパクト}


\noindent
\ref{sites-definition-quasi-compact-topos} を参照：
{\it 準コンパクト},
{\it 準コンパクト}


\noindent
\ref{sites-definition-cocontinuous} を参照：
{\it 余連続}


\noindent
\ref{sites-definition-localize} を参照：
{\it 対象 $U$ におけるサイト $\mathcal{C}$ の局所化},
{\it 局所化射},
{\it 順像関手},
{\it $\mathcal{F}$ の $\mathcal{C}/U$ への制限},
{\it 空集合による $\mathcal{G}$ の延長}


\noindent
\ref{sites-definition-special-cocontinuous-functor} を参照：
{\it $\mathcal{C}$ から $\mathcal{D}$ への特殊余連続関手 $u$}


\noindent
\ref{sites-definition-localize-topos} を参照：
{\it $\mathcal{F}$ におけるトポス $\Sh(\mathcal{C})$ の局所化},
{\it 局所化射}


\noindent
\ref{sites-definition-point-topos} を参照：
{\it トポス $\Sh(\mathcal{C})$ の点}


\noindent
\ref{sites-definition-point} を参照：
{\it サイト $\mathcal{C}$ の点 $p$}


\noindent
\ref{sites-definition-pushforward-point} を参照：
{\it 摩天楼層}


\noindent
\ref{sites-definition-2-morphism-topoi} を参照：
{\it $f$ から $g$ への 2-射}


\noindent
\ref{sites-definition-morphism-points} を参照：
{\it 射 $f : p \to p'$}


\noindent
\ref{sites-definition-enough-points} を参照：
{\it 保守的},
{\it 十分多くの点をもつ}


\noindent
\ref{sites-definition-w-contractible} を参照：
{\it 弱可縮},
{\it 十分多くの弱可縮対象をもつ},
{\it 十分多くの $P$ 対象をもつ}


\noindent
\ref{sites-definition-empty} を参照：
{\it 層論的に空}


\noindent
\ref{sites-definition-almost-cocontinuous} を参照：
{\it ほとんど余連続}


\noindent
\ref{sites-definition-embedding} を参照：
{\it 埋め込み}


\noindent
\ref{sites-definition-subtopos} を参照：
{\it 部分トポス}


\noindent
\ref{sites-definition-open-subtopos} を参照：
{\it 開部分トポス}


\noindent
\ref{sites-definition-closed-subtopos} を参照：
{\it 閉部分トポス}


\noindent
\ref{sites-definition-immersion-topoi} を参照：
{\it 開埋め込み},
{\it 閉埋め込み}


\noindent
\ref{sites-definition-pushforward-algebraic-structures} を参照：
{\it 順像}


\noindent
\ref{sites-definition-global-sections} を参照：
{\it 大域切断}


\noindent
\ref{sites-definition-sieve} を参照：
{\it $U$ 上のふるい $S$}


\noindent
\ref{sites-definition-sieve-generated} を参照：
{\it 射 $f_i$ が生成する $U$ 上のふるい}


\noindent
\ref{sites-definition-pullback-sieve} を参照：
{\it $f$ による $S$ の引き戻し}


\noindent
\ref{sites-definition-topology} を参照：
{\it $\mathcal{C}$ 上の位相}


\noindent
\ref{sites-definition-finer} を参照：
{\it より細かい},
{\it より強い},
{\it より粗い},
{\it より弱い}


\noindent
\ref{sites-definition-sheaf-sets-topology} を参照：
{\it 層}


\noindent
\ref{sites-definition-canonical-topology} を参照：
{\it 標準位相}


\noindent
\ref{sites-definition-topology-associated-site} を参照：
{\it $\mathcal{C}$ に付随する位相}


\noindent
\ref{sites-definition-presheaf-separated-topology} を参照：
{\it 分離的}


\noindent
\ref{sites-definition-associated-sheaf-topology} を参照：
{\it $\mathcal{F}$ に付随する層}


\noindent
\ref{sites-definition-point-topology} を参照：
{\it 点 $p$}


\medskip\noindent
{\bf スタック}

\medskip

\noindent
\ref{stacks-definition-mor-presheaf} を参照：
{\it $x$ から $y$ への射の前層},
{\it $x$ から $y$ への同型の前層}


\noindent
\ref{stacks-definition-descent-data} を参照：
{\it 族 $\{f_i : U_i \to U\}$ に関する $\mathcal{S}$ の降下データ $(X_i, \varphi_{ij})$},
{\it コサイクル条件},
{\it 降下データの射 $\psi : (X_i, \varphi_{ij}) \to (X'_i, \varphi'_{ij})$}


\noindent
\ref{stacks-definition-pullback-functor} を参照：
{\it 引き戻し関手}


\noindent
\ref{stacks-definition-effective-descent-datum} を参照：
{\it 自明な降下データ},
{\it 標準降下データ},
{\it 有効}


\noindent
\ref{stacks-definition-stack} を参照：
{\it スタック}


\noindent
\ref{stacks-definition-stacks-over-C} を参照：
{\it $\mathcal{C}$ 上のスタックの $2$-圏}


\noindent
\ref{stacks-definition-stack-in-groupoids} を参照：
{\it グルーポイドのスタック}


\noindent
\ref{stacks-definition-stacks-in-groupoids-over-C} を参照：
{\it $\mathcal{C}$ 上のグルーポイドのスタックの $2$-圏}


\noindent
\ref{stacks-definition-stack-in-sets} を参照：
{\it セトイドのスタック},
{\it 集合のスタック},
{\it 離散圏のスタック}


\noindent
\ref{stacks-definition-stacks-in-setoids-over-C} を参照：
{\it $\mathcal{C}$ 上のセトイドのスタックの $2$-圏}


\noindent
\ref{stacks-definition-topology-inherited} を参照：
{\it $\mathcal{C}$ から継承される $\mathcal{S}$ 上のサイト構造},
{\it $\mathcal{S}$ は $\mathcal{C}$ から誘導される位相を備える}


\noindent
\ref{stacks-definition-gerbe} を参照：
{\it ガーベ}


\noindent
\ref{stacks-definition-gerbe-over-stack-in-groupoids} を参照：
{\it 上のガーベ}


\noindent
\ref{stacks-definition-pushforward-stack} を参照：
{\it $f_*\mathcal{S}$},
{\it $f$ に沿う $\mathcal{S}$ の順像}


\noindent
\ref{stacks-definition-pullback-stack} を参照：
{\it $f^{-1}\mathcal{S}$},
{\it $f$ に沿う $\mathcal{S}$ の引き戻し}


\medskip\noindent
{\bf 体}

\medskip

\noindent
\ref{fields-definition-field} を参照：
{\it 体},
{\it 部分体}


\noindent
\ref{fields-definition-domain} を参照：
{\it 整域},
{\it 整域}


\noindent
\ref{fields-definition-characteristic} を参照：
{\it 標数},
{\it $F$ の素体}


\noindent
\ref{fields-definition-extension} を参照：
{\it 拡大体}


\noindent
\ref{fields-definition-tower} を参照：
{\it 塔}


\noindent
\ref{fields-definition-generated-by} を参照：
{\it 生成する},
{\it 有限生成体拡大}


\noindent
\ref{fields-definition-degree} を参照：
{\it 次数},
{\it 有限}


\noindent
\ref{fields-definition-number-field} を参照：
{\it 数体}


\noindent
\ref{fields-definition-algebraic} を参照：
{\it 代数的},
{\it 代数拡大}


\noindent
\ref{fields-definition-minimal-polynomial} を参照：
{\it 最小多項式}


\noindent
\ref{fields-definition-algebraically-closed} を参照：
{\it 代数閉}


\noindent
\ref{fields-definition-algebraic-closure} を参照：
{\it 代数閉包}


\noindent
\ref{fields-definition-relatively-prime} を参照：
{\it 互いに素}


\noindent
\ref{fields-definition-separable} を参照：
{\it 分離的},
{\it 分離的},
{\it 分離的}


\noindent
\ref{fields-definition-separable-degree} を参照：
{\it 分離次数}


\noindent
\ref{fields-definition-purely-inseparable} を参照：
{\it 純非分離的},
{\it 純非分離的},
{\it 純非分離的}


\noindent
\ref{fields-definition-insep-degree} を参照：
{\it 分離次数},
{\it 非分離次数},
{\it 非分離度}


\noindent
\ref{fields-definition-normal} を参照：
{\it 正規}


\noindent
\ref{fields-definition-automorphisms} を参照：
{\it $F$ 上の $E$ の自己同型},
{\it $E/F$ の自己同型}


\noindent
\ref{fields-definition-splitting-field} を参照：
{\it $F$ 上の $P$ の分解体}


\noindent
\ref{fields-definition-normal-closure} を参照：
{\it $F$ 上の $E$ の正規閉包}


\noindent
\ref{fields-definition-trace-norm} を参照：
{\it トレース},
{\it ノルム}


\noindent
\ref{fields-definition-trace-pairing} を参照：
{\it トレース形式}


\noindent
\ref{fields-definition-discriminant} を参照：
{\it $L/K$ の判別式}


\noindent
\ref{fields-definition-galois} を参照：
{\it ガロア}


\noindent
\ref{fields-definition-galois-group} を参照：
{\it ガロア群}


\noindent
\ref{fields-definition-transcendence} を参照：
{\it 代数的独立},
{\it 純超越拡大},
{\it 超越基底}


\noindent
\ref{fields-definition-transcendence-degree} を参照：
{\it 超越次数}


\noindent
\ref{fields-definition-algebraically-closed-in} を参照：
{\it $K$ における $k$ の代数閉包},
{\it $K$ において代数的に閉じている}


\noindent
\ref{fields-definition-compositum} を参照：
{\it $\Omega$ における $K$ と $L$ の合成体}


\noindent
\ref{fields-definition-linearly-disjoint} を参照：
{\it $\Omega$ において $k$ 上線形無関連}


\noindent
\ref{fields-definition-separable-algebraic} を参照：
{\it 代数的},
{\it 分離的},
{\it 純非分離的},
{\it 正規},
{\it ガロア}


\medskip\noindent
{\bf 可換環論}

\medskip

\noindent
\ref{algebra-definition-module-finite-type} を参照：
{\it 有限 $R$-加群},
{\it 有限生成 $R$-加群},
{\it 有限表示 $R$-加群},
{\it 有限表示の $R$-加群}


\noindent
\ref{algebra-definition-finite-type} を参照：
{\it 有限型},
{\it $S$ は有限型 $R$-代数である},
{\it 有限表示}


\noindent
\ref{algebra-definition-finite-ring-map} を参照：
{\it 有限}


\noindent
\ref{algebra-definition-directed-system} を参照：
{\it $I$ 上の $R$-加群の系 $(M_i, \mu_{ij})$},
{\it 有向系}


\noindent
\ref{algebra-definition-homomorphism-directed-systems} を参照：
{\it 系の準同型}


\noindent
\ref{algebra-definition-multiplicative-subset} を参照：
{\it $R$ の乗法的集合}


\noindent
\ref{algebra-definition-localization} を参照：
{\it $S$ に関する $A$ の局所化}


\noindent
\ref{algebra-definition-localization-module} を参照：
{\it 局所化}


\noindent
\ref{algebra-definition-relation} を参照：
{\it 関係}


\noindent
\ref{algebra-definition-bilinear} を参照：
{\it $R$-双線形}


\noindent
\ref{algebra-definition-bimodule} を参照：
{\it $(A, B)$-両側加群}


\noindent
\ref{algebra-definition-base-change} を参照：
{\it 基底変換},
{\it 基底変換}


\noindent
\ref{algebra-definition-spectrum-ring} を参照：
{\it スペクトル}


\noindent
\ref{algebra-definition-Zariski-topology} を参照：
{\it ザリスキー},
{\it 標準開集合}


\noindent
\ref{algebra-definition-local-ring} を参照：
{\it 局所環},
{\it 剰余体},
{\it 局所環の局所準同型},
{\it 局所環準同型 $\varphi : R \to S$}


\noindent
\ref{algebra-definition-oka-family} を参照：
{\it 岡族}


\noindent
\ref{algebra-definition-locally-nilpotent-ideal} を参照：
{\it 局所冪零},
{\it 冪零}


\noindent
\ref{algebra-definition-ring-jacobson} を参照：
{\it ジャコブソン環}


\noindent
\ref{algebra-definition-integral-ring-map} を参照：
{\it $R$ 上整},
{\it 整}


\noindent
\ref{algebra-definition-integral-closure} を参照：
{\it 整閉包},
{\it 整閉}


\noindent
\ref{algebra-definition-domain-normal} を参照：
{\it 正規}


\noindent
\ref{algebra-definition-almost-integral} を参照：
{\it $R$ 上概整},
{\it 完整閉}


\noindent
\ref{algebra-definition-ring-normal} を参照：
{\it 正規}


\noindent
\ref{algebra-definition-integral-over-ideal} を参照：
{\it $I$ 上整}


\noindent
\ref{algebra-definition-flat} を参照：
{\it 平坦},
{\it 忠実平坦},
{\it 平坦},
{\it 忠実平坦}


\noindent
\ref{algebra-definition-support-module} を参照：
{\it $M$ の台}


\noindent
\ref{algebra-definition-annihilator} を参照：
{\it $m$ の零化イデアル},
{\it $M$ の零化イデアル}


\noindent
\ref{algebra-definition-going-up-down} を参照：
{\it 上昇性},
{\it 下降性}


\noindent
\ref{algebra-definition-separable-field-extension} を参照：
{\it $k$ 上分離生成},
{\it $k$ 上分離的}


\noindent
\ref{algebra-definition-geometrically-reduced} を参照：
{\it $k$ 上幾何学的に被約}


\noindent
\ref{algebra-definition-perfect} を参照：
{\it 完全}


\noindent
\ref{algebra-definition-perfection} を参照：
{\it 完全閉包}


\noindent
\ref{algebra-definition-geometrically-irreducible} を参照：
{\it $k$ 上幾何学的既約}


\noindent
\ref{algebra-definition-geometrically-connected} を参照：
{\it $k$ 上幾何学的連結}


\noindent
\ref{algebra-definition-geometrically-integral} を参照：
{\it $k$ 上幾何学的整}


\noindent
\ref{algebra-definition-valuation-ring} を参照：
{\it 支配する},
{\it 付値環},
{\it 中心をもつ}


\noindent
\ref{algebra-definition-value-group} を参照：
{\it 値群},
{\it 付値},
{\it 離散付値環}


\noindent
\ref{algebra-definition-length} を参照：
{\it 長さ}


\noindent
\ref{algebra-definition-simple-module} を参照：
{\it 単純}


\noindent
\ref{algebra-definition-artinian} を参照：
{\it アルティン的}


\noindent
\ref{algebra-definition-essentially-finite-p-t} を参照：
{\it 本質的に有限型},
{\it 本質的に有限表示}


\noindent
\ref{algebra-definition-proj} を参照：
{\it 斉次スペクトル}


\noindent
\ref{algebra-definition-numerical-polynomial} を参照：
{\it 数値多項式}


\noindent
\ref{algebra-definition-ideal-definition} を参照：
{\it $R$ の定義イデアル}


\noindent
\ref{algebra-definition-hilbert-polynomial} を参照：
{\it ヒルベルト多項式}


\noindent
\ref{algebra-definition-d} を参照：
{\it $d(M)$}


\noindent
\ref{algebra-definition-chain-primes} を参照：
{\it 素イデアルの鎖},
{\it 長さ}


\noindent
\ref{algebra-definition-Krull} を参照：
{\it クルル次元}


\noindent
\ref{algebra-definition-height} を参照：
{\it 高さ}


\noindent
\ref{algebra-definition-regular-local} を参照：
{\it $R$ のパラメータ系},
{\it 正則局所環},
{\it 正則パラメータ系}


\noindent
\ref{algebra-definition-associated} を参照：
{\it 随伴する}


\noindent
\ref{algebra-definition-symbolic-power} を参照：
{\it 記号冪}


\noindent
\ref{algebra-definition-relative-assassin} を参照：
{\it $S/R$ 上の $N$ の相対随伴素イデアル集合}


\noindent
\ref{algebra-definition-weakly-associated} を参照：
{\it 弱随伴する}


\noindent
\ref{algebra-definition-embedded-primes} を参照：
{\it 埋没随伴素イデアル},
{\it $R$ の埋没素イデアル}


\noindent
\ref{algebra-definition-regular-sequence} を参照：
{\it $M$-正則列},
{\it $I$ 内の $M$-正則列},
{\it 正則列}


\noindent
\ref{algebra-definition-quasi-regular-sequence} を参照：
{\it $M$-準正則},
{\it 準正則列}


\noindent
\ref{algebra-definition-blow-up} を参照：
{\it ブローアップ代数},
{\it Rees 代数},
{\it アフィン・ブローアップ代数}


\noindent
\ref{algebra-definition-finite-free-resolution} を参照：
{\it 分解},
{\it 自由 $R$-加群による $M$ の分解},
{\it 有限自由 $R$-加群による $M$ の分解}


\noindent
\ref{algebra-definition-depth} を参照：
{\it $I$-深さ},
{\it 深さ}


\noindent
\ref{algebra-definition-projective} を参照：
{\it 射影的}


\noindent
\ref{algebra-definition-locally-free} を参照：
{\it 局所自由},
{\it 有限局所自由},
{\it 階数 $r$ の有限局所自由}


\noindent
\ref{algebra-definition-universally-injective} を参照：
{\it 普遍単射},
{\it 純完全}


\noindent
\ref{algebra-definition-devissage} を参照：
{\it 直和デヴィサージュ},
{\it Kaplansky デヴィサージュ}


\noindent
\ref{algebra-definition-ML-system} を参照：
{\it Mittag-Leffler}


\noindent
\ref{algebra-definition-ML-inductive-system} を参照：
{\it 加群の Mittag-Leffler 有向系}


\noindent
\ref{algebra-definition-domination} を参照：
{\it 支配する}


\noindent
\ref{algebra-definition-mittag-leffler-module} を参照：
{\it Mittag-Leffler}


\noindent
\ref{algebra-definition-coherent} を参照：
{\it コヒーレント加群},
{\it コヒーレント環}


\noindent
\ref{algebra-definition-complete} を参照：
{\it $I$-進完備},
{\it $I$-進完備}


\noindent
\ref{algebra-definition-rank} を参照：
{\it 階数}


\noindent
\ref{algebra-definition-CM} を参照：
{\it Cohen--Macaulay}


\noindent
\ref{algebra-definition-maximal-CM} を参照：
{\it 極大 Cohen--Macaulay}


\noindent
\ref{algebra-definition-module-CM} を参照：
{\it Cohen--Macaulay}


\noindent
\ref{algebra-definition-local-ring-CM} を参照：
{\it Cohen--Macaulay}


\noindent
\ref{algebra-definition-ring-CM} を参照：
{\it Cohen--Macaulay}


\noindent
\ref{algebra-definition-catenary} を参照：
{\it 鎖状}


\noindent
\ref{algebra-definition-universally-catenary} を参照：
{\it 普遍鎖状}


\noindent
\ref{algebra-definition-pure-ideal} を参照：
{\it 純}


\noindent
\ref{algebra-definition-finite-proj-dim} を参照：
{\it 有限射影次元},
{\it 射影次元}


\noindent
\ref{algebra-definition-finite-gl-dim} を参照：
{\it 有限大域次元},
{\it 大域次元}


\noindent
\ref{algebra-definition-regular} を参照：
{\it 正則}


\noindent
\ref{algebra-definition-fibre} を参照：
{\it $\mathfrak q$ におけるファイバーの局所環}


\noindent
\ref{algebra-definition-uniformizer} を参照：
{\it 一意化元}


\noindent
\ref{algebra-definition-irreducible-prime-element} を参照：
{\it 同伴},
{\it 既約},
{\it 素元}


\noindent
\ref{algebra-definition-UFD} を参照：
{\it 一意分解整域},
{\it UFD}


\noindent
\ref{algebra-definition-PID} を参照：
{\it 単項イデアル整域},
{\it PID}


\noindent
\ref{algebra-definition-dedekind-domain} を参照：
{\it Dedekind 整域}


\noindent
\ref{algebra-definition-ord} を参照：
{\it $R$ に沿う消滅位数}


\noindent
\ref{algebra-definition-lattice} を参照：
{\it $V$ の格子}


\noindent
\ref{algebra-definition-distance} を参照：
{\it $M$ と $M'$ の距離}


\noindent
\ref{algebra-definition-quasi-finite} を参照：
{\it $\mathfrak q$ で準有限},
{\it 準有限}


\noindent
\ref{algebra-definition-strongly-transcendental} を参照：
{\it $R$ 上強超越的}


\noindent
\ref{algebra-definition-relative-dimension} を参照：
{\it $\mathfrak q$ における $S/R$ の相対次元},
{\it 相対次元}


\noindent
\ref{algebra-definition-derivation} を参照：
{\it 導分},
{\it $R$-導分},
{\it Leibniz 則}


\noindent
\ref{algebra-definition-differentials} を参照：
{\it K\"ahler 微分加群},
{\it 微分加群}


\noindent
\ref{algebra-definition-differential-operators} を参照：
{\it $k$ 階の微分作用素 $D : M \to N$}


\noindent
\ref{algebra-definition-module-principal-parts} を参照：
{\it $k$ 階主部加群の層}


\noindent
\ref{algebra-definition-naive-cotangent-complex} を参照：
{\it 素朴な余接複体}


\noindent
\ref{algebra-definition-lci-field} を参照：
{\it $k$ 上大域的完全交叉},
{\it $k$ 上の局所完全交叉}


\noindent
\ref{algebra-definition-lci-local-ring} を参照：
{\it $k$ 上完全交叉}


\noindent
\ref{algebra-definition-lci} を参照：
{\it シントミック},
{\it $R$ 上の平坦な局所完全交叉}


\noindent
\ref{algebra-definition-relative-global-complete-intersection} を参照：
{\it 相対的大域完全交叉}


\noindent
\ref{algebra-definition-smooth} を参照：
{\it 滑らか}


\noindent
\ref{algebra-definition-standard-smooth} を参照：
{\it $R$ 上の標準滑らかな代数},
{\it 標準滑らか},
{\it 標準滑らか}


\noindent
\ref{algebra-definition-smooth-at-prime} を参照：
{\it $\mathfrak q$ で滑らか}


\noindent
\ref{algebra-definition-formally-smooth} を参照：
{\it $R$ 上形式的に滑らか}


\noindent
\ref{algebra-definition-small-extension} を参照：
{\it 小拡大}


\noindent
\ref{algebra-definition-etale} を参照：
{\it エタール},
{\it $\mathfrak q$ においてエタール}


\noindent
\ref{algebra-definition-standard-etale} を参照：
{\it 標準エタール}


\noindent
\ref{algebra-definition-formally-unramified} を参照：
{\it $R$ 上形式的不分岐}


\noindent
\ref{algebra-definition-universal-thickening} を参照：
{\it 普遍一次肥厚化},
{\it 余法加群},
{\it $C_{S/R}$}


\noindent
\ref{algebra-definition-formally-etale} を参照：
{\it $R$ 上形式的エタール}


\noindent
\ref{algebra-definition-unramified} を参照：
{\it 不分岐},
{\it G-不分岐},
{\it $\mathfrak q$ で不分岐},
{\it $\mathfrak q$ において G-不分岐}


\noindent
\ref{algebra-definition-henselian} を参照：
{\it ヘンゼル的},
{\it 強ヘンゼル的}


\noindent
\ref{algebra-definition-henselization} を参照：
{\it ヘンゼル化},
{\it $\kappa \subset \kappa^{sep}$ に関する $R$ の強ヘンゼル化},
{\it 強ヘンゼル化}


\noindent
\ref{algebra-definition-conditions} を参照：
{\it $(R_k)$},
{\it 余次元 $\leq k$ で正則},
{\it $(S_k)$}


\noindent
\ref{algebra-definition-complete-local-ring} を参照：
{\it 完備局所環}


\noindent
\ref{algebra-definition-coefficient-ring} を参照：
{\it 係数環}


\noindent
\ref{algebra-definition-cohen-ring} を参照：
{\it Cohen 環}


\noindent
\ref{algebra-definition-N} を参照：
{\it N-1},
{\it N-2},
{\it 日本スキーム}


\noindent
\ref{algebra-definition-nagata} を参照：
{\it 普遍的日本スキーム},
{\it Nagata 環}


\noindent
\ref{algebra-definition-analytically-unramified} を参照：
{\it 解析的不分岐},
{\it 解析的不分岐}


\noindent
\ref{algebra-definition-geometrically-normal} を参照：
{\it 幾何学的に正規}


\noindent
\ref{algebra-definition-geometrically-regular} を参照：
{\it 幾何学的に正則}


\medskip\noindent
{\bf Brauer 群}

\medskip

\noindent
\ref{brauer-definition-finite} を参照：
{\it 有限}


\noindent
\ref{brauer-definition-skew-field} を参照：
{\it 斜体}


\noindent
\ref{brauer-definition-simple} を参照：
{\it 単純},
{\it 単純}


\noindent
\ref{brauer-definition-central} を参照：
{\it 中心的}


\noindent
\ref{brauer-definition-opposite} を参照：
{\it 反対代数}


\noindent
\ref{brauer-definition-brauer-group} を参照：
{\it ブラウアー群}


\noindent
\ref{brauer-definition-splitting} を参照：
{\it 分解する},
{\it 分解体}


\medskip\noindent
{\bf ホモロジー代数}

\medskip

\noindent
\ref{homology-definition-preadditive} を参照：
{\it 前加法圏},
{\it 加法的}


\noindent
\ref{homology-definition-zero-object} を参照：
{\it 零対象}


\noindent
\ref{homology-definition-direct-sum} を参照：
{\it 直和}


\noindent
\ref{homology-definition-additive-category} を参照：
{\it 加法的}


\noindent
\ref{homology-definition-kernel} を参照：
{\it 核},
{\it 余核},
{\it $f$ の余像},
{\it $f$ の像}


\noindent
\ref{homology-definition-karoubian} を参照：
{\it カルービ圏}


\noindent
\ref{homology-definition-abelian-category} を参照：
{\it アーベル}


\noindent
\ref{homology-definition-injective-surjective} を参照：
{\it 入射的},
{\it 全射},
{\it 部分対象},
{\it 商}


\noindent
\ref{homology-definition-exact} を参照：
{\it 複体},
{\it $y$ で完全},
{\it $x_i$ で完全},
{\it 完全},
{\it 完全列},
{\it 完全複体},
{\it 短完全列}


\noindent
\ref{homology-definition-ses-split} を参照：
{\it 分裂する}


\noindent
\ref{homology-definition-extension} を参照：
{\it $A$ による $B$ の拡大 $E$},
{\it 拡大の射}


\noindent
\ref{homology-definition-ext-group} を参照：
{\it $\Ext$-群}


\noindent
\ref{homology-definition-simple} を参照：
{\it 単純}


\noindent
\ref{homology-definition-Artinian} を参照：
{\it アルティン的},
{\it アルティン的}


\noindent
\ref{homology-definition-Noetherian} を参照：
{\it Noether},
{\it Noether}


\noindent
\ref{homology-definition-serre-subcategory} を参照：
{\it セール部分圏},
{\it 弱セール部分圏}


\noindent
\ref{homology-definition-kernel-category} を参照：
{\it 関手 $F$ の核}


\noindent
\ref{homology-definition-K-zero} を参照：
{\it $\mathcal{A}$ の第零 $K$ 群}


\noindent
\ref{homology-definition-cohomological-delta-functor} を参照：
{\it コホモロジー的 $\delta$ 関手},
{\it $\delta$-関手}


\noindent
\ref{homology-definition-morphism-delta-functors} を参照：
{\it $F$ から $G$ への $\delta$ 関手の射}


\noindent
\ref{homology-definition-universal-delta-functor} を参照：
{\it 普遍 $\delta$ 関手}


\noindent
\ref{homology-definition-homotopy-equivalent} を参照：
{\it ホモトピー同値},
{\it ホモトピー同値}


\noindent
\ref{homology-definition-quasi-isomorphism} を参照：
{\it 擬同型},
{\it 非輪状}


\noindent
\ref{homology-definition-homotopy-equivalent-cochain} を参照：
{\it ホモトピー同値},
{\it ホモトピー同値}


\noindent
\ref{homology-definition-quasi-isomorphism-cochain} を参照：
{\it 擬同型},
{\it 非輪状}


\noindent
\ref{homology-definition-shift} を参照：
{\it $k$-シフト鎖複体 $A[k]_\bullet$}


\noindent
\ref{homology-definition-homology-shift} を参照：
{\it $H_{i + k}(A_\bullet) \rightarrow H_i(A[k]_\bullet)$}


\noindent
\ref{homology-definition-shift-cochain} を参照：
{\it $k$-シフト余鎖複体 $A[k]^\bullet$}


\noindent
\ref{homology-definition-cohomology-shift} を参照：
{\it $H^{i + k}(A^\bullet) \longrightarrow H^i(A[k]^\bullet)$}


\noindent
\ref{homology-definition-graded} を参照：
{\it $\mathcal{A}$ の次数付き対象の圏}


\noindent
\ref{homology-definition-graded-shift} を参照：
{\it シフト}


\noindent
\ref{homology-definition-additive-monoidal} を参照：
{\it 加法的モノイド圏}


\noindent
\ref{homology-definition-double-complex} を参照：
{\it 二重複体}


\noindent
\ref{homology-definition-associated-simple-complex} を参照：
{\it 付随単複体},
{\it 付随全複体}


\noindent
\ref{homology-definition-filtered} を参照：
{\it 減少フィルトレーション},
{\it $\mathcal{A}$ のフィルター付き対象},
{\it フィルトレーション付き対象の射 $(A, F) \to (B, F)$},
{\it 誘導フィルトレーション},
{\it 商フィルトレーション},
{\it 有限},
{\it 分離的},
{\it 網羅的}


\noindent
\ref{homology-definition-strict} を参照：
{\it 厳密}


\noindent
\ref{homology-definition-spectral-sequence} を参照：
{\it $\mathcal{A}$ におけるスペクトル系列},
{\it スペクトル系列の射}


\noindent
\ref{homology-definition-limit-spectral-sequence} を参照：
{\it 極限},
{\it $E_r$ で退化する}


\noindent
\ref{homology-definition-exact-couple} を参照：
{\it 完全対},
{\it 完全対の射}


\noindent
\ref{homology-definition-spectral-sequence-associated-exact-couple} を参照：
{\it 完全対に付随するスペクトル系列}


\noindent
\ref{homology-definition-differential-object} を参照：
{\it 微分対象},
{\it 微分対象の射}


\noindent
\ref{homology-definition-differential-object-homology} を参照：
{\it ホモロジー}


\noindent
\ref{homology-definition-differential-object-selfmap} を参照：
{\it $(A, d, \alpha)$ に付随するスペクトル系列}


\noindent
\ref{homology-definition-filtered-differential} を参照：
{\it フィルトレーション付き微分対象}


\noindent
\ref{homology-definition-filtration-cohomology-filtered-differential} を参照：
{\it 誘導フィルトレーション}


\noindent
\ref{homology-definition-filtered-differential-ss-converges} を参照：
{\it $H(K)$ に弱収束する},
{\it $H(K)$ に収束する}


\noindent
\ref{homology-definition-filtered-complex} を参照：
{\it $\mathcal{A}$ のフィルター付き複体 $K^\bullet$}


\noindent
\ref{homology-definition-filtration-cohomology-filtered-complex} を参照：
{\it 誘導フィルトレーション}


\noindent
\ref{homology-definition-bounded-ss} を参照：
{\it 正則},
{\it 余正則},
{\it 有界},
{\it 下に有界},
{\it 上に有界}


\noindent
\ref{homology-definition-filtered-complex-ss-converges} を参照：
{\it $H^*(K^\bullet)$ に弱収束する},
{\it $H^*(K^\bullet)$ に収束する},
{\it $H^*(K^\bullet)$ に収束する}


\noindent
\ref{homology-definition-ss-double-complex-converge} を参照：
{\it $H^n(\text{Tot}(K^{\bullet, \bullet}))$ に弱収束する},
{\it $H^n(\text{Tot}(K^{\bullet, \bullet}))$ に収束する},
{\it $H^n(\text{Tot}(K^{\bullet, \bullet}))$ に収束する},
{\it $H^n(\text{Tot}(K^{\bullet, \bullet}))$ に弱収束する},
{\it $H^n(\text{Tot}(K^{\bullet, \bullet}))$ に収束する},
{\it $H^n(\text{Tot}(K^{\bullet, \bullet}))$ に収束する}


\noindent
\ref{homology-definition-injective} を参照：
{\it 入射的}


\noindent
\ref{homology-definition-enough-injectives} を参照：
{\it 十分な入射対象をもつ}


\noindent
\ref{homology-definition-functorial-injective-embedding} を参照：
{\it 関手的入射埋め込みをもつ}


\noindent
\ref{homology-definition-projective} を参照：
{\it 射影的}


\noindent
\ref{homology-definition-enough-projectives} を参照：
{\it 十分な射影対象をもつ}


\noindent
\ref{homology-definition-functorial-projective-surjections} を参照：
{\it 関手的射影全射をもつ}


\noindent
\ref{homology-definition-Mittag-Leffler} を参照：
{\it Mittag-Leffler 条件},
{\it ML}


\medskip\noindent
{\bf 導来圏}

\medskip

\noindent
\ref{derived-definition-triangle} を参照：
{\it 三角},
{\it 三角の射}


\noindent
\ref{derived-definition-triangulated-category} を参照：
{\it 三角圏},
{\it 完全三角},
{\it 前三角圏}


\noindent
\ref{derived-definition-exact-functor-triangulated-categories} を参照：
{\it 完全関手},
{\it 三角関手}


\noindent
\ref{derived-definition-triangulated-subcategory} を参照：
{\it 前三角部分圏},
{\it 三角部分圏}


\noindent
\ref{derived-definition-homological} を参照：
{\it ホモロジー的},
{\it コホモロジー的}


\noindent
\ref{derived-definition-delta-functor} を参照：
{\it $\mathcal{A}$ から $\mathcal{D}$ への $\delta$-関手},
{\it 与えられた $\delta$-関手による短完全列の像}


\noindent
\ref{derived-definition-localization} を参照：
{\it 三角構造と両立する}


\noindent
\ref{derived-definition-saturated} を参照：
{\it 飽和している}


\noindent
\ref{derived-definition-kernel-category} を参照：
{\it $F$ の核},
{\it $H$ の核}


\noindent
\ref{derived-definition-quotient-category} を参照：
{\it 商圏 $\mathcal{D}/\mathcal{B}$},
{\it 商関手}


\noindent
\ref{derived-definition-complexes-notation} を参照：
{\it （余鎖）複体の圏},
{\it 下に有界},
{\it 上に有界},
{\it 有界}


\noindent
\ref{derived-definition-cone} を参照：
{\it 錐}


\noindent
\ref{derived-definition-termwise-split-map} を参照：
{\it 項ごとに分裂する単射 $\alpha : A^\bullet \to B^\bullet$},
{\it 項ごとに分裂する全射 $\beta : B^\bullet \to C^\bullet$}


\noindent
\ref{derived-definition-split-ses} を参照：
{\it $\mathcal{A}$ の複体の項別分裂完全列},
{\it 項ごとに分裂する複体の列に付随する三角}


\noindent
\ref{derived-definition-distinguished-triangle} を参照：
{\it $K(\mathcal{A})$ の完全三角}


\noindent
\ref{derived-definition-unbounded-derived-category} を参照：
{\it $\mathcal{A}$ の導来圏},
{\it 有界導来圏}


\noindent
\ref{derived-definition-finite-filtered} を参照：
{\it $\mathcal{A}$ の有限フィルトレーション付き 対象の圏}


\noindent
\ref{derived-definition-filtered-acyclic} を参照：
{\it フィルター付き擬同型},
{\it フィルター付き非輪状}


\noindent
\ref{derived-definition-filtered-derived} を参照：
{\it $\mathcal{A}$ のフィルトレーション付き導来圏}


\noindent
\ref{derived-definition-filtered-derived-bounded} を参照：
{\it 有界フィルトレーション付き導来圏}


\noindent
\ref{derived-definition-right-derived-functor-defined} を参照：
{\it 右導来関手 $RF$ が定義される},
{\it $X$ における $RF$ の値},
{\it 左導来関手 $LF$ が定義される},
{\it $X$ における $LF$ の値}


\noindent
\ref{derived-definition-everywhere-defined} を参照：
{\it 右導来可能},
{\it 至る所で定義される},
{\it 左導来可能},
{\it 至る所で定義される}


\noindent
\ref{derived-definition-computes} を参照：
{\it 計算する},
{\it 計算する}


\noindent
\ref{derived-definition-derived-functor} を参照：
{\it $F$ の右導来関手},
{\it $F$ の左導来関手},
{\it $F$ に関して右非輪状},
{\it $RF$ に関して非輪状},
{\it $F$ に関して左非輪状},
{\it $LF$ に関して非輪状}


\noindent
\ref{derived-definition-higher-derived-functors} を参照：
{\it $F$ の第 $i$ 右導来関手 $R^iF$}


\noindent
\ref{derived-definition-injective-resolution} を参照：
{\it $A$ の入射分解},
{\it $K^\bullet$ の入射分解}


\noindent
\ref{derived-definition-projective-resolution} を参照：
{\it $A$ の射影分解},
{\it $K^\bullet$ の射影分解}


\noindent
\ref{derived-definition-cartan-eilenberg} を参照：
{\it Cartan--Eilenberg 分解}


\noindent
\ref{derived-definition-localization-functor} を参照：
{\it 分解関手}


\noindent
\ref{derived-definition-filtered-complexes-notation} を参照：
{\it フィルター付き入射的}


\noindent
\ref{derived-definition-ext} を参照：
{\it 第 $i$ 拡大群}


\noindent
\ref{derived-definition-yoneda-extension} を参照：
{\it Yoneda 拡大},
{\it 同値}


\noindent
\ref{derived-definition-K-zero} を参照：
{\it $\mathcal{D}$ の第零 $K$ 群}


\noindent
\ref{derived-definition-K-injective} を参照：
{\it K-入射的}


\noindent
\ref{derived-definition-derived-colimit} を参照：
{\it 導来余極限},
{\it ホモトピー余極限}


\noindent
\ref{derived-definition-derived-limit} を参照：
{\it 導来極限},
{\it ホモトピー極限}


\noindent
\ref{derived-definition-generators} を参照：
{\it 古典的生成元},
{\it 強生成元},
{\it 弱生成元},
{\it 生成対象}


\noindent
\ref{derived-definition-compact-object} を参照：
{\it コンパクト対象}


\noindent
\ref{derived-definition-compactly-generated} を参照：
{\it コンパクト生成}


\noindent
\ref{derived-definition-orthogonal} を参照：
{\it 右直交},
{\it 左直交}


\noindent
\ref{derived-definition-admissible} を参照：
{\it 右許容},
{\it 左許容},
{\it 両側許容}


\noindent
\ref{derived-definition-postnikov-system} を参照：
{\it Postnikov 系},
{\it Postnikov 系の射}


\medskip\noindent
{\bf 単体的方法}

\medskip

\noindent
\ref{simplicial-definition-face-degeneracy} を参照：
{\it $\delta^n_j : [n-1] \to [n]$},
{\it $\sigma^n_j : [n + 1]  \to [n]$}


\noindent
\ref{simplicial-definition-simplicial-object} を参照：
{\it $\mathcal{C}$ の単体的対象 $U$},
{\it 単体集合},
{\it 単体アーベル群},
{\it 射 $U \to U'$},
{\it $\mathcal{C}$ の単体的対象の圏}


\noindent
\ref{simplicial-definition-cosimplicial-object} を参照：
{\it $\mathcal{C}$ の余単体的対象 $U$},
{\it 余単体集合},
{\it 余単体アーベル群},
{\it 射 $U \to U'$},
{\it $\mathcal{C}$ の余単体的対象の圏}


\noindent
\ref{simplicial-definition-product} を参照：
{\it $U$ と $V$ の積}


\noindent
\ref{simplicial-definition-fibre-product} を参照：
{\it $U$ 上の $V$ と $W$ のファイバー積}


\noindent
\ref{simplicial-definition-push-out} を参照：
{\it $U$ 上の $V$ と $W$ の押し出し}


\noindent
\ref{simplicial-definition-product-cosimplicial-objects} を参照：
{\it $U$ と $V$ の積}


\noindent
\ref{simplicial-definition-fibre-product-cosimplicial-objects} を参照：
{\it $U$ 上の $V$ と $W$ のファイバー積}


\noindent
\ref{simplicial-definition-terminology-simplicial-sets} を参照：
{\it $U$ の $n$-単体},
{\it $x$ の面},
{\it $x$ の退化},
{\it 退化している}


\noindent
\ref{simplicial-definition-truncated-simplicial-object} を参照：
{\it $\mathcal{C}$ の $n$-切断単体対象},
{\it $n$-切断単体的対象の射}


\noindent
\ref{simplicial-definition-product-with-simplicial-set} を参照：
{\it $U$ と $V$ の積 $U \times V$},
{\it 積 $U \times V$ が存在する}


\noindent
\ref{simplicial-definition-hom-deltak-cosimplicial} を参照：
{\it $\Hom(U, V)$}


\noindent
\ref{simplicial-definition-hom-deltak-simplicial} を参照：
{\it $\Hom(U, V)$}


\noindent
\ref{simplicial-definition-hom-from-simplicial-set} を参照：
{\it $\Hom(U, V)$}


\noindent
\ref{simplicial-definition-split} を参照：
{\it 分裂する},
{\it 分裂する}


\noindent
\ref{simplicial-definition-augmentation} を参照：
{\it $\mathcal{C}$ の対象 $X$ へ向かう $U$ の増大写像 $\epsilon : U \to X$}


\noindent
\ref{simplicial-definition-eilenberg-maclane} を参照：
{\it Eilenberg--Mac Lane 対象 $K(A, k)$}


\noindent
\ref{simplicial-definition-homotopy} を参照：
{\it $a$ から $b$ へのホモトピー},
{\it ホモトピック},
{\it 同じホモトピー類に属する}


\noindent
\ref{simplicial-definition-homotopy-equivalent} を参照：
{\it ホモトピー同値},
{\it ホモトピー同値}


\noindent
\ref{simplicial-definition-homotopy-cosimplicial} を参照：
{\it $a$ から $b$ へのホモトピー},
{\it ホモトピック},
{\it 同じホモトピー類に属する}


\noindent
\ref{simplicial-definition-trivial-kan} を参照：
{\it 自明 Kan ファイブレーション}


\noindent
\ref{simplicial-definition-kan} を参照：
{\it Kan ファイブレーション},
{\it Kan 複体}


\medskip\noindent
{\bf 代数の補足}

\medskip

\noindent
\ref{more-algebra-definition-stably-free} を参照：
{\it 安定同型},
{\it 安定自由}


\noindent
\ref{more-algebra-definition-fitting-ideal} を参照：
{\it 第 $k$ Fitting イデアル}


\noindent
\ref{more-algebra-definition-zariski-pair} を参照：
{\it Zariski 対}


\noindent
\ref{more-algebra-definition-henselian-pair} を参照：
{\it ヘンゼル対}


\noindent
\ref{more-algebra-definition-absolutely-integrally-closed} を参照：
{\it 絶対整閉}


\noindent
\ref{more-algebra-definition-auto-ass} を参照：
{\it 自己随伴}


\noindent
\ref{more-algebra-definition-torsion} を参照：
{\it ねじれ},
{\it ねじれがない},
{\it ねじれ加群}


\noindent
\ref{more-algebra-definition-rank} を参照：
{\it 階数}


\noindent
\ref{more-algebra-definition-reflexive} を参照：
{\it 反射的}


\noindent
\ref{more-algebra-definition-reflexive-hull} を参照：
{\it 反射包}


\noindent
\ref{more-algebra-definition-content-ideal} を参照：
{\it $x$ の内容イデアル}


\noindent
\ref{more-algebra-definition-strict-transform} を参照：
{\it $R \to R[\frac{I}{a}]$ に沿う $M$ の狭義変換}


\noindent
\ref{more-algebra-definition-koszul} を参照：
{\it Koszul 複体}


\noindent
\ref{more-algebra-definition-koszul-complex} を参照：
{\it $f_1, \ldots, f_r$ 上の Koszul 複体}


\noindent
\ref{more-algebra-definition-koszul-regular-sequence} を参照：
{\it $M$-Koszul-正則},
{\it $M$-$H_1$-正則},
{\it Koszul-正則},
{\it $H_1$-正則}


\noindent
\ref{more-algebra-definition-regular-ideal} を参照：
{\it 正則イデアル},
{\it Koszul-正則イデアル},
{\it $H_1$-正則イデアル},
{\it 準正則イデアル}


\noindent
\ref{more-algebra-definition-local-complete-intersection} を参照：
{\it 局所完全交叉}


\noindent
\ref{more-algebra-definition-topological-ring} を参照：
{\it 位相環},
{\it 位相加群},
{\it 位相加群の準同型},
{\it 位相環の準同型},
{\it 線形位相をもつ},
{\it 線形位相をもつ},
{\it 定義イデアル},
{\it 前許容},
{\it 許容},
{\it 前アディック},
{\it アディック}


\noindent
\ref{more-algebra-definition-formally-smooth} を参照：
{\it $R$ 上形式的に滑らか}


\noindent
\ref{more-algebra-definition-formally-smooth-adic} を参照：
{\it $\mathfrak n$-進位相に関して形式的に 滑らか}


\noindent
\ref{more-algebra-definition-regular} を参照：
{\it 正則}


\noindent
\ref{more-algebra-definition-p-basis} を参照：
{\it $k$ 上 p-独立},
{\it $k$ 上の $K$ の p-基底}


\noindent
\ref{more-algebra-definition-J} を参照：
{\it J-0},
{\it J-1},
{\it J-2}


\noindent
\ref{more-algebra-definition-G-ring} を参照：
{\it G-環}


\noindent
\ref{more-algebra-definition-excellent} を参照：
{\it 準優秀},
{\it 優秀}


\noindent
\ref{more-algebra-definition-projective} を参照：
{\it 入射的}


\noindent
\ref{more-algebra-definition-simple-functors} を参照：
{\it $M \mapsto M^\vee$},
{\it 自由加群}


\noindent
\ref{more-algebra-definition-K-flat} を参照：
{\it K-平坦}


\noindent
\ref{more-algebra-definition-derived-tor} を参照：
{\it 導来テンソル積}


\noindent
\ref{more-algebra-definition-tor-independent} を参照：
{\it $R$ 上 Tor 独立}


\noindent
\ref{more-algebra-definition-pseudo-coherent} を参照：
{\it $m$-擬連接},
{\it 擬連接},
{\it $m$-擬連接},
{\it 擬連接}


\noindent
\ref{more-algebra-definition-tor-amplitude} を参照：
{\it $[a, b]$ における Tor 振幅},
{\it 有限 Tor 次元をもつ},
{\it Tor 次元 $\leq d$},
{\it 有限 Tor 次元をもつ}


\noindent
\ref{more-algebra-definition-projective-dimension} を参照：
{\it 有限射影次元},
{\it $[a, b]$ における射影振幅}


\noindent
\ref{more-algebra-definition-injective-dimension} を参照：
{\it 有限入射次元},
{\it $[a, b]$ における入射振幅}


\noindent
\ref{more-algebra-definition-near-projective} を参照：
{\it $I$-射影的}


\noindent
\ref{more-algebra-definition-perfect} を参照：
{\it 完全},
{\it 完全}


\noindent
\ref{more-algebra-definition-relatively-finitely-presented} を参照：
{\it $R$ に関して有限表示}


\noindent
\ref{more-algebra-definition-relatively-pseudo-coherent} を参照：
{\it $R$ に関して $m$-擬連接},
{\it $R$ に関して擬コヒーレント},
{\it $R$ に関して $m$-擬連接},
{\it $R$ に関して擬コヒーレント}


\noindent
\ref{more-algebra-definition-pseudo-coherent-perfect} を参照：
{\it 擬コヒーレント環写像},
{\it 完全環写像}


\noindent
\ref{more-algebra-definition-relatively-perfect} を参照：
{\it $R$-完全},
{\it $R$ に関して完全}


\noindent
\ref{more-algebra-definition-f-power-torsion} を参照：
{\it $I$-冪ねじれ加群},
{\it $f$-冪ねじれ加群}


\noindent
\ref{more-algebra-definition-derived-complete} を参照：
{\it $I$ に関して導来完備},
{\it $I$ に関して導来完備}


\noindent
\ref{more-algebra-definition-weakly-etale} を参照：
{\it 絶対平坦},
{\it 弱エタール},
{\it 絶対平坦}


\noindent
\ref{more-algebra-definition-weak-dimension} を参照：
{\it 弱次元は $\leq d$}


\noindent
\ref{more-algebra-definition-unibranch} を参照：
{\it 単枝},
{\it 幾何的単枝}


\noindent
\ref{more-algebra-definition-number-of-branches} を参照：
{\it $A$ の枝数},
{\it $A$ の幾何的枝数}


\noindent
\ref{more-algebra-definition-formally-catenary} を参照：
{\it 形式鎖状}


\noindent
\ref{more-algebra-definition-extension-discrete-valuation-rings} を参照：
{\it 離散付値環の拡大},
{\it 分岐指数},
{\it 弱不分岐},
{\it 剰余次数},
{\it 剰余体次数}


\noindent
\ref{more-algebra-definition-types-of-extensions} を参照：
{\it $A$ に関して不分岐},
{\it $A$ に関して馴分岐},
{\it $A$ に関して完全分岐}


\noindent
\ref{more-algebra-definition-decomposition-inertia} を参照：
{\it $\mathfrak m$ の分解群},
{\it $\mathfrak m$ の惰性群}


\noindent
\ref{more-algebra-definition-wild-inertia} を参照：
{\it $\mathfrak m$ の野性惰性群},
{\it $\mathfrak m$ の馴性惰性群}


\noindent
\ref{more-algebra-definition-mixed} を参照：
{\it 混標数},
{\it 絶対分岐指数}


\noindent
\ref{more-algebra-definition-solution} を参照：
{\it $A \subset B$ に対する弱解},
{\it $A \subset B$ に対する解},
{\it 分離解}


\noindent
\ref{more-algebra-definition-invertible} を参照：
{\it 可逆},
{\it 自明}


\noindent
\ref{more-algebra-definition-extension-valuation-rings} を参照：
{\it 付値環の拡大},
{\it 弱不分岐},
{\it 剰余次数},
{\it 剰余体次数}


\noindent
\ref{more-algebra-definition-bezout} を参照：
{\it B\'ezout 整域},
{\it 単因子整域}


\medskip\noindent
{\bf 環準同型の平滑化}

\medskip

\noindent
\ref{smoothing-definition-singular-ideal} を参照：
{\it $R$ 上の $A$ の特異イデアル}


\noindent
\ref{smoothing-definition-strictly-standard} を参照：
{\it $R$ 上の $A$ における初等標準元},
{\it $R$ 上の $A$ における厳密標準元}


\medskip\noindent
{\bf 加群の層}

\medskip

\noindent
\ref{modules-definition-globally-generated} を参照：
{\it 大域切断で生成される},
{\it 生成する}


\noindent
\ref{modules-definition-generated-by-local-sections} を参照：
{\it $s_i$ によって生成される部分層}


\noindent
\ref{modules-definition-support} を参照：
{\it $\mathcal{F}$ の台},
{\it $s$ の台}


\noindent
\ref{modules-definition-locally-generated} を参照：
{\it 局所的に切断で生成される}


\noindent
\ref{modules-definition-finite-type} を参照：
{\it 有限型}


\noindent
\ref{modules-definition-quasi-coherent} を参照：
{\it 準連接 $\mathcal{O}_X$-加群の層}


\noindent
\ref{modules-definition-sheaf-associated} を参照：
{\it 加群 $M$ と環準同型 $\alpha$ に付随する層},
{\it 加群 $M$ に付随する層}


\noindent
\ref{modules-definition-finite-presentation} を参照：
{\it 有限表示}


\noindent
\ref{modules-definition-coherent} を参照：
{\it 連接 $\mathcal{O}_X$-加群の層}


\noindent
\ref{modules-definition-closed-immersion} を参照：
{\it 環付き空間の閉埋め込み}


\noindent
\ref{modules-definition-locally-free} を参照：
{\it 局所自由},
{\it 有限局所自由},
{\it 階数 $r$ の有限局所自由}


\noindent
\ref{modules-definition-flat} を参照：
{\it 平坦}


\noindent
\ref{modules-definition-flat-at-point} を参照：
{\it $x$ で平坦}


\noindent
\ref{modules-definition-flat-morphism} を参照：
{\it $x$ で平坦},
{\it 平坦}


\noindent
\ref{modules-definition-flat-module} を参照：
{\it 点 $x \in X$ で $Y$ 上平坦},
{\it $Y$ 上平坦}


\noindent
\ref{modules-definition-annihilator-sheaf} を参照：
{\it 零化イデアル}


\noindent
\ref{modules-definition-koszul} を参照：
{\it Koszul 複体}


\noindent
\ref{modules-definition-koszul-complex} を参照：
{\it $f_1, \ldots, f_n$ 上の Koszul 複体}


\noindent
\ref{modules-definition-invertible} を参照：
{\it 可逆 $\mathcal{O}_X$-加群の層},
{\it 自明}


\noindent
\ref{modules-definition-powers} を参照：
{\it テンソル冪}


\noindent
\ref{modules-definition-gamma-star} を参照：
{\it 付随次数付き環}


\noindent
\ref{modules-definition-pic} を参照：
{\it Picard 群}


\noindent
\ref{modules-definition-derivation} を参照：
{\it $\mathcal{O}_1$-導分},
{\it $\varphi$-導分},
{\it Leibniz 則}


\noindent
\ref{modules-definition-module-differentials} を参照：
{\it 微分加群},
{\it 普遍 $\varphi$-導分}


\noindent
\ref{modules-definition-differentials} を参照：
{\it $S$-導分},
{\it $S$ 上の $X$ の微分層 $\Omega_{X/S}$}


\noindent
\ref{modules-definition-differential-operators} を参照：
{\it $k$ 階微分作用素 $D : \mathcal{F} \to \mathcal{G}$}


\noindent
\ref{modules-definition-module-principal-parts} を参照：
{\it $k$ 階主部加群の層}


\noindent
\ref{modules-definition-relative-differential-operators} を参照：
{\it $X/S$ 上の $k$ 階微分作用素}


\noindent
\ref{modules-definition-de-rham-complex} を参照：
{\it $\mathcal{A}$ 上の $\mathcal{B}$ の de Rham 複体}


\noindent
\ref{modules-definition-de-rham-complex-morphism-ringed-spaces} を参照：
{\it de Rham 複体}


\noindent
\ref{modules-definition-naive-cotangent-complex} を参照：
{\it 素朴な余接複体}


\noindent
\ref{modules-definition-cotangent-complex-morphism-ringed-topoi} を参照：
{\it 素朴な余接複体}


\medskip\noindent
{\bf サイト上の加群}

\medskip

\noindent
\ref{sites-modules-definition-free-abelian-presheaf-on} を参照：
{\it 自由アーベル群の前層}


\noindent
\ref{sites-modules-definition-free-abelian-sheaf-on} を参照：
{\it 自由アーベル群の層}


\noindent
\ref{sites-modules-definition-ringed-site} を参照：
{\it 環付きサイト},
{\it 構造層},
{\it 環付きサイトの射},
{\it 環付きサイトの射の合成}


\noindent
\ref{sites-modules-definition-ringed-topos} を参照：
{\it 環付きトポス},
{\it 構造層},
{\it 環付きトポスの射},
{\it 環付きトポスの射の合成}


\noindent
\ref{sites-modules-definition-2-morphism-ringed-topoi} を参照：
{\it $f$ から $g$ への 2-射}


\noindent
\ref{sites-modules-definition-presheaf-modules} を参照：
{\it $\mathcal{O}$-加群の前層},
{\it $\mathcal{O}$-加群の前層の射 $\varphi : \mathcal{F} \to \mathcal{G}$}


\noindent
\ref{sites-modules-definition-sheaf-modules} を参照：
{\it $\mathcal{O}$-加群の層},
{\it $\mathcal{O}$-加群の層の射}


\noindent
\ref{sites-modules-definition-pushforward} を参照：
{\it 順像},
{\it 引き戻し}


\noindent
\ref{sites-modules-definition-g-shriek} を参照：
{\it $g_{p!}\mathcal{F}$},
{\it $g_!\mathcal{F} = (g_{p!}\mathcal{F})^\#$}


\noindent
\ref{sites-modules-definition-global} を参照：
{\it 自由 $\mathcal{O}$-加群},
{\it 有限自由},
{\it 大域切断で生成される},
{\it $r$ 個の大域切断で生成される},
{\it 有限個の大域切断で生成される},
{\it 大域表示},
{\it 大域有限表示}


\noindent
\ref{sites-modules-definition-localize-ringed-site} を参照：
{\it 環付きサイト $(\mathcal{C}, \mathcal{O})$ の対象 $U$ における 局所化},
{\it 局所化射},
{\it 順像関手},
{\it $\mathcal{F}$ の $\mathcal{C}/U$ への制限},
{\it 零による延長}


\noindent
\ref{sites-modules-definition-localize-ringed-topos} を参照：
{\it 環付きトポス $(\Sh(\mathcal{C}), \mathcal{O})$ の $\mathcal{F}$ における局所化},
{\it 局所化射}


\noindent
\ref{sites-modules-definition-site-local} を参照：
{\it 局所自由},
{\it 有限局所自由},
{\it 局所的に切断で生成される},
{\it 局所的に $r$ 個の切断で生成される},
{\it 有限型},
{\it 準連接},
{\it 有限表示},
{\it 連接}


\noindent
\ref{sites-modules-definition-flat} を参照：
{\it 平坦},
{\it 平坦},
{\it 平坦},
{\it 平坦}


\noindent
\ref{sites-modules-definition-flat-morphism} を参照：
{\it 平坦},
{\it 平坦}


\noindent
\ref{sites-modules-definition-flat-module} を参照：
{\it $(\Sh(\mathcal{D}), \mathcal{O}')$ 上平坦}


\noindent
\ref{sites-modules-definition-invertible-sheaf} を参照：
{\it 階数 $r$},
{\it 可逆},
{\it $\mathcal{O}^*$}


\noindent
\ref{sites-modules-definition-pic} を参照：
{\it Picard 群}


\noindent
\ref{sites-modules-definition-derivation} を参照：
{\it $\mathcal{O}_1$-導分},
{\it $\varphi$-導分},
{\it Leibniz 則}


\noindent
\ref{sites-modules-definition-module-differentials} を参照：
{\it 微分加群},
{\it 普遍 $\varphi$-導分}


\noindent
\ref{sites-modules-definition-sheaf-differentials} を参照：
{\it $Y$-導分},
{\it $Y$ 上の $X$ の微分層 $\Omega_{X/Y}$},
{\it 普遍 $Y$-導分}


\noindent
\ref{sites-modules-definition-differential-operators} を参照：
{\it $k$ 階微分作用素 $D : \mathcal{F} \to \mathcal{G}$}


\noindent
\ref{sites-modules-definition-module-principal-parts} を参照：
{\it $k$ 階主部加群の層}


\noindent
\ref{sites-modules-definition-naive-cotangent-complex} を参照：
{\it 素朴な余接複体}


\noindent
\ref{sites-modules-definition-cotangent-complex-morphism-ringed-topoi} を参照：
{\it 素朴な余接複体}


\noindent
\ref{sites-modules-definition-locally-ringed} を参照：
{\it 局所環付きサイト}


\noindent
\ref{sites-modules-definition-locally-ringed-topos} を参照：
{\it 局所環付き}


\noindent
\ref{sites-modules-definition-morphism-locally-ringed-topoi} を参照：
{\it 局所環付きトポスの射},
{\it 局所環付きサイトの射}


\noindent
\ref{sites-modules-definition-locally-constant} を参照：
{\it 定数層},
{\it 局所定数},
{\it 有限局所定数}


\medskip\noindent
{\bf 入射対象}

\medskip

\noindent
\ref{injectives-definition-small} を参照：
{\it $I$ に関して $\alpha$-小}


\noindent
\ref{injectives-definition-grothendieck-conditions} を参照：
{\it 生成対象},
{\it Grothendieck アーベル圏}


\noindent
\ref{injectives-definition-size} を参照：
{\it 大きさ}


\medskip\noindent
{\bf 層コホモロジー}

\medskip

\noindent
\ref{cohomology-definition-torsor} を参照：
{\it 擬トーサー},
{\it 擬 $\mathcal{G}$-トーサー},
{\it 擬 $\mathcal{G}$-トーサーの射},
{\it トーサー},
{\it $\mathcal{G}$-トーサー},
{\it $\mathcal{G}$-トーサーの射},
{\it 自明な $\mathcal{G}$-トーサー}


\noindent
\ref{cohomology-definition-cech-complex} を参照：
{\it {\v C}ech 複体},
{\it {\v C}ech コホモロジー群}


\noindent
\ref{cohomology-definition-flasque} を参照：
{\it フラスク},
{\it フラビー}


\noindent
\ref{cohomology-definition-alternating-cech-complex} を参照：
{\it 交代 {\v C}ech 複体}


\noindent
\ref{cohomology-definition-ordered-cech-complex} を参照：
{\it 順序付き {\v C}ech 複体}


\noindent
\ref{cohomology-definition-covering-locally-finite} を参照：
{\it 局所有限}


\noindent
\ref{cohomology-definition-K-flat} を参照：
{\it K-平坦}


\noindent
\ref{cohomology-definition-derived-tor} を参照：
{\it 導来テンソル積}


\noindent
\ref{cohomology-definition-tor} を参照：
{\it Tor}


\noindent
\ref{cohomology-definition-strictly-perfect} を参照：
{\it 狭義完全}


\noindent
\ref{cohomology-definition-pseudo-coherent} を参照：
{\it $m$-擬連接},
{\it 擬連接},
{\it $m$-擬連接},
{\it 擬連接}


\noindent
\ref{cohomology-definition-tor-amplitude} を参照：
{\it $[a, b]$ における Tor 振幅},
{\it 有限 Tor 次元をもつ},
{\it 局所的に有限 Tor 次元をもつ},
{\it Tor 次元 $\leq d$}


\noindent
\ref{cohomology-definition-perfect} を参照：
{\it 完全},
{\it 完全}


\medskip\noindent
{\bf サイト上のコホモロジー}

\medskip

\noindent
\ref{sites-cohomology-definition-torsor} を参照：
{\it 擬トーサー},
{\it 擬 $\mathcal{G}$-トーサー},
{\it 擬 $\mathcal{G}$-トーサーの射},
{\it トーサー},
{\it $\mathcal{G}$-トーサー},
{\it $\mathcal{G}$-トーサーの射},
{\it 自明な $\mathcal{G}$-トーサー}


\noindent
\ref{sites-cohomology-definition-cech-complex} を参照：
{\it {\v C}ech 複体},
{\it {\v C}ech コホモロジー群}


\noindent
\ref{sites-cohomology-definition-limp} を参照：
{\it 完全非輪状}


\noindent
\ref{sites-cohomology-definition-K-flat} を参照：
{\it K-平坦}


\noindent
\ref{sites-cohomology-definition-derived-tor} を参照：
{\it 導来テンソル積}


\noindent
\ref{sites-cohomology-definition-tor} を参照：
{\it Tor}


\noindent
\ref{sites-cohomology-definition-covering-LC} を参照：
{\it qc 被覆}


\noindent
\ref{sites-cohomology-definition-simplicial-module} を参照：
{\it 単体的 $\mathcal{A}_\bullet$-加群},
{\it $\mathcal{A}_\bullet$-加群の単体的層}


\noindent
\ref{sites-cohomology-definition-cartesian} を参照：
{\it $\mathit{QC}(\mathcal{O})$}


\noindent
\ref{sites-cohomology-definition-strictly-perfect} を参照：
{\it 狭義完全}


\noindent
\ref{sites-cohomology-definition-pseudo-coherent} を参照：
{\it $m$-擬連接},
{\it 擬連接},
{\it $m$-擬連接},
{\it 擬連接}


\noindent
\ref{sites-cohomology-definition-tor-amplitude} を参照：
{\it $[a, b]$ における Tor 振幅},
{\it 有限 Tor 次元をもつ},
{\it 局所的に有限 Tor 次元をもつ},
{\it Tor 次元 $\leq d$}


\noindent
\ref{sites-cohomology-definition-perfect} を参照：
{\it 完全},
{\it 完全}


\medskip\noindent
{\bf 微分付き次数代数}

\medskip

\noindent
\ref{dga-definition-dga} を参照：
{\it $R$ 上の微分次数付き代数}


\noindent
\ref{dga-definition-homomorphism-dga} を参照：
{\it 微分次数付き代数の準同型}


\noindent
\ref{dga-definition-cdga} を参照：
{\it 可換},
{\it 厳密可換}


\noindent
\ref{dga-definition-tensor-product} を参照：
{\it テンソル積微分次数付き代数}


\noindent
\ref{dga-definition-dgm} を参照：
{\it 微分次数付き加群},
{\it 微分次数付き加群の準同型}


\noindent
\ref{dga-definition-shift} を参照：
{\it $k$-シフト加群}


\noindent
\ref{dga-definition-homotopy} を参照：
{\it $f$ と $g$ の間のホモトピー},
{\it ホモトピック}


\noindent
\ref{dga-definition-complexes-notation} を参照：
{\it ホモトピー圏}


\noindent
\ref{dga-definition-cone} を参照：
{\it 錐}


\noindent
\ref{dga-definition-admissible-ses} を参照：
{\it 許容単射},
{\it 許容全射},
{\it 許容短完全列}


\noindent
\ref{dga-definition-distinguished-triangle} を参照：
{\it $0 \to K \to L \to M \to 0$ に付随する三角形},
{\it 識別三角形}


\noindent
\ref{dga-definition-opposite-dga} を参照：
{\it 反対微分次数付き代数}


\noindent
\ref{dga-definition-shift-graded-module} を参照：
{\it 第 $k$ シフト $A$-加群},
{\it 第 $k$ シフト $A$-加群}


\noindent
\ref{dga-definition-unbounded-derived-category} を参照：
{\it $(A, \text{d})$ の導来圏}


\noindent
\ref{dga-definition-linear-category} を参照：
{\it $R$-線形圏 $\mathcal{A}$}


\noindent
\ref{dga-definition-functor-linear-categories} を参照：
{\it $R$-線形圏の間の関手},
{\it $R$-線形関手}


\noindent
\ref{dga-definition-graded-category} を参照：
{\it $R$ 上の次数付き圏 $\mathcal{A}$}


\noindent
\ref{dga-definition-functor-graded-categories} を参照：
{\it $R$ 上の次数付き圏の間の関手},
{\it 次数付き関手}


\noindent
\ref{dga-definition-H0-of-graded-category} を参照：
{\it $\mathcal{A}^0$}


\noindent
\ref{dga-definition-graded-direct-sum} を参照：
{\it 次数付き直和}


\noindent
\ref{dga-definition-dga-category} を参照：
{\it $R$ 上の微分次数付き圏 $\mathcal{A}$}


\noindent
\ref{dga-definition-functor-dga-categories} を参照：
{\it $R$ 上の微分次数付き圏の関手}


\noindent
\ref{dga-definition-homotopy-category-of-dga-category} を参照：
{\it $\mathcal{A}$ の複体の圏},
{\it $\mathcal{A}$ のホモトピー圏}


\noindent
\ref{dga-definition-dg-direct-sum} を参照：
{\it 微分次数付き直和}


\noindent
\ref{dga-definition-bimodule} を参照：
{\it $(A, B)$-両側加群},
{\it 次数付き $(A, B)$-両側加群},
{\it 微分次数付き $(A, B)$-両側加群}


\medskip\noindent
{\bf 分割冪代数}

\medskip

\noindent
\ref{dpa-definition-divided-powers} を参照：
{\it 分割冪構造}


\noindent
\ref{dpa-definition-divided-power-ring} を参照：
{\it 分割冪環},
{\it 分割冪環の準同型}


\noindent
\ref{dpa-definition-extends} を参照：
{\it 延長される}


\noindent
\ref{dpa-definition-divided-powers-graded} を参照：
{\it 分割冪構造}


\noindent
\ref{dpa-definition-divided-powers-dga} を参照：
{\it 微分次数付き構造と両立する}


\noindent
\ref{dpa-definition-lci} を参照：
{\it 完全交叉},
{\it 局所完全交叉}


\medskip\noindent
{\bf 微分付き次数層}

\medskip

\noindent
\ref{sdga-definition-ga} を参照：
{\it 次数付き $\mathcal{O}$-代数の層},
{\it 次数付き代数の層},
{\it 次数付き $\mathcal{O}$-代数の準同型}


\noindent
\ref{sdga-definition-gm} を参照：
{\it 次数付き $\mathcal{A}$-加群},
{\it 次数付き加群},
{\it 次数付き $\mathcal{A}$-加群の準同型}


\noindent
\ref{sdga-definition-bimodule} を参照：
{\it 次数付き $(\mathcal{A}, \mathcal{B})$-両側加群},
{\it 次数付き $(\mathcal{A}, \mathcal{B})$-両側加群の準同型}


\noindent
\ref{sdga-definition-dga} を参照：
{\it 微分次数付き $\mathcal{O}$-代数の層},
{\it 微分次数付き代数の層},
{\it 微分次数付き $\mathcal{O}$-代数の準同型}


\noindent
\ref{sdga-definition-dgm} を参照：
{\it 微分次数付き $\mathcal{A}$-加群},
{\it 微分次数付き加群},
{\it 微分次数付き $\mathcal{A}$-加群の準同型}


\noindent
\ref{sdga-definition-dg-bimodule} を参照：
{\it 微分次数付き $(\mathcal{A}, \mathcal{B})$-双加群},
{\it 微分次数付き $(\mathcal{A}, \mathcal{B})$-両側加群の準同型}


\noindent
\ref{sdga-definition-homotopy} を参照：
{\it $f$ と $g$ の間のホモトピー},
{\it ホモトピック}


\noindent
\ref{sdga-definition-complexes-notation} を参照：
{\it ホモトピー圏}


\noindent
\ref{sdga-definition-cone} を参照：
{\it 錐}


\noindent
\ref{sdga-definition-graded-injective} を参照：
{\it 次数付き入射的}


\noindent
\ref{sdga-definition-K-injective} を参照：
{\it K-入射的}


\noindent
\ref{sdga-definition-derived-category} を参照：
{\it $(\mathcal{A}, \text{d})$ の導来圏}


\noindent
\ref{sdga-definition-pullback} を参照：
{\it 導来テンソル積},
{\it 導来引き戻し}


\noindent
\ref{sdga-definition-pushforward} を参照：
{\it 導来内部 Hom},
{\it 導来順像}


\noindent
\ref{sdga-definition-cartesian} を参照：
{\it $\mathit{QC}(\mathcal{A}, \text{d})$}


\medskip\noindent
{\bf 超被覆}

\medskip

\noindent
\ref{hypercovering-definition-SR} を参照：
{\it 半表現可能対象},
{\it $X$ 上の半表現可能対象}


\noindent
\ref{hypercovering-definition-SR-F} を参照：
{\it 半表現可能対象に前層を対応させる関手}


\noindent
\ref{hypercovering-definition-covering-SR} を参照：
{\it 被覆},
{\it 被覆}


\noindent
\ref{hypercovering-definition-hypercovering} を参照：
{\it $X$ の超被覆}


\noindent
\ref{hypercovering-definition-homology} を参照：
{\it $K$ のホモロジー}


\noindent
\ref{hypercovering-definition-hypercovering-variant} を参照：
{\it $\mathcal{G}$ の超被覆},
{\it 超被覆}


\medskip\noindent
{\bf スキーム}

\medskip

\noindent
\ref{schemes-definition-locally-ringed-space} を参照：
{\it 局所環付き空間 $(X, \mathcal{O}_X)$},
{\it $X$ の $x$ における局所環},
{\it $x$ における $X$ の剰余体},
{\it 局所環付き空間の射}


\noindent
\ref{schemes-definition-immersion-locally-ringed-spaces} を参照：
{\it 開埋め込み}


\noindent
\ref{schemes-definition-open-subspace} を参照：
{\it $X$ の、$U$ に付随する開部分空間}


\noindent
\ref{schemes-definition-closed-immersion-locally-ringed-spaces} を参照：
{\it 閉埋め込み}


\noindent
\ref{schemes-definition-closed-subspace} を参照：
{\it $X$ の、イデアルの層 $\mathcal{I}$ に付随する 閉部分空間}


\noindent
\ref{schemes-definition-standard-covering} を参照：
{\it 標準開被覆},
{\it 標準開被覆}


\noindent
\ref{schemes-definition-structure-sheaf} を参照：
{\it $R$ のスペクトルの構造層 $\mathcal{O}_{\Spec(R)}$},
{\it スペクトル}


\noindent
\ref{schemes-definition-affine-scheme} を参照：
{\it アフィンスキーム},
{\it アフィンスキームの射}


\noindent
\ref{schemes-definition-scheme} を参照：
{\it スキーム},
{\it スキームの射}


\noindent
\ref{schemes-definition-immersion} を参照：
{\it 開埋め込み},
{\it 開部分スキーム},
{\it 閉埋め込み},
{\it 閉部分スキーム},
{\it 埋め込み},
{\it 局所閉埋め込み}


\noindent
\ref{schemes-definition-reduced} を参照：
{\it 被約}


\noindent
\ref{schemes-definition-reduced-induced-scheme} を参照：
{\it $Z$ 上のスキーム構造},
{\it 被約誘導スキーム構造},
{\it $X$ の被約化 $X_{red}$}


\noindent
\ref{schemes-definition-representable-functor} を参照：
{\it スキームにより表現可能},
{\it 表現可能}


\noindent
\ref{schemes-definition-representable-by-open-immersions} を参照：
{\it Zariski 位相に関する層条件を満たす},
{\it 部分関手 $H \subset F$},
{\it 開埋め込みにより表現可能},
{\it $F$ を被覆する}


\noindent
\ref{schemes-definition-fibre-product} を参照：
{\it ファイバー積}


\noindent
\ref{schemes-definition-inverse-image-closed-subscheme} を参照：
{\it 閉部分スキーム $Z$ の逆像 $f^{-1}(Z)$}


\noindent
\ref{schemes-definition-base-change} を参照：
{\it $S$ 上のスキーム},
{\it 構造射},
{\it $R$ 上のスキーム},
{\it スキームの射 $f : X \to Y$（$S$ 上のもの）},
{\it 基底変換},
{\it 基底変換},
{\it 基底変換}


\noindent
\ref{schemes-definition-preserved-by-base-change} を参照：
{\it 任意の基底変換で保たれる},
{\it 基底変換で保たれる},
{\it 任意の基底変換で保たれる},
{\it 基底変換で保たれる}


\noindent
\ref{schemes-definition-fibre} を参照：
{\it スキーム論的ファイバー $X_s$（$f$ の $s$ 上のもの）},
{\it ファイバー（$f$ の $s$ 上のもの）}


\noindent
\ref{schemes-definition-quasi-compact} を参照：
{\it 準コンパクト}


\noindent
\ref{schemes-definition-universally-closed} を参照：
{\it 普遍閉}


\noindent
\ref{schemes-definition-valuative-criterion} を参照：
{\it 付値判定法の存在部分を満たす},
{\it 付値判定法の一意性部分を満たす}


\noindent
\ref{schemes-definition-separated} を参照：
{\it 分離的},
{\it 準分離},
{\it 分離的},
{\it 準分離}


\noindent
\ref{schemes-definition-monomorphism} を参照：
{\it モノ射}


\medskip\noindent
{\bf スキームの構成}

\medskip

\noindent
\ref{constructions-definition-relative-spec} を参照：
{\it $\mathcal{A}$ の $S$ 上の相対スペクトル},
{\it $\mathcal{A}$ の $S$ 上のスペクトル}


\noindent
\ref{constructions-definition-affine-n-space} を参照：
{\it アフィン $n$-空間（$S$ 上）},
{\it アフィン $n$-空間（$R$ 上）}


\noindent
\ref{constructions-definition-vector-bundle} を参照：
{\it $\mathcal{E}$ に付随するベクトル束}


\noindent
\ref{constructions-definition-abstract-vector-bundle} を参照：
{\it ベクトル束 $\pi : V \to S$（$S$ 上）},
{\it $S$ 上のベクトル束の射}


\noindent
\ref{constructions-definition-cone} を参照：
{\it $\mathcal{A}$ に付随する錐},
{\it $\mathcal{A}$ に付随するアフィン錐}


\noindent
\ref{constructions-definition-abstract-cone} を参照：
{\it 錐 $\pi : C \to S$（$S$ 上）},
{\it 錐の射}


\noindent
\ref{constructions-definition-standard-covering} を参照：
{\it 標準開被覆}


\noindent
\ref{constructions-definition-structure-sheaf} を参照：
{\it 斉次スペクトルの構造層 $\mathcal{O}_{\text{Proj}(S)}$（$S$ の）},
{\it 斉次スペクトル}


\noindent
\ref{constructions-definition-twist} を参照：
{\it $\text{Proj}(S)$ の構造層の捩り}


\noindent
\ref{constructions-definition-projective-space} を参照：
{\it 射影 $n$ 次元空間（$\mathbf{Z}$ 上）},
{\it 射影 $n$ 次元空間（$S$ 上）},
{\it 射影 $n$ 次元空間（$R$ 上）}


\noindent
\ref{constructions-definition-relative-proj} を参照：
{\it $\mathcal{A}$ の $S$ 上の相対斉次スペクトル},
{\it $\mathcal{A}$ の $S$ 上の斉次スペクトル},
{\it $\mathcal{A}$ の $S$ 上の相対 Proj}


\noindent
\ref{constructions-definition-projective-bundle} を参照：
{\it $\mathcal{E}$ に付随する射影束},
{\it 構造層のねじり}


\noindent
\ref{constructions-definition-grassmannian} を参照：
{\it $\mathbf{Z}$ 上のグラスマン多様体},
{\it $S$ 上のグラスマン多様体},
{\it $R$ 上のグラスマン多様体}


\medskip\noindent
{\bf スキームの性質}

\medskip

\noindent
\ref{properties-definition-integral} を参照：
{\it 整}


\noindent
\ref{properties-definition-property-local} を参照：
{\it 局所的}


\noindent
\ref{properties-definition-locally-P} を参照：
{\it 局所的に $P$}


\noindent
\ref{properties-definition-noetherian} を参照：
{\it 局所 Noether},
{\it Noether}


\noindent
\ref{properties-definition-jacobson} を参照：
{\it Jacobson}


\noindent
\ref{properties-definition-normal} を参照：
{\it 正規}


\noindent
\ref{properties-definition-Cohen-Macaulay} を参照：
{\it Cohen--Macaulay}


\noindent
\ref{properties-definition-regular} を参照：
{\it 正則},
{\it 非特異}


\noindent
\ref{properties-definition-dimension} を参照：
{\it 次元},
{\it $x$ における $X$ の次元}


\noindent
\ref{properties-definition-catenary} を参照：
{\it 鎖状}


\noindent
\ref{properties-definition-Rk} を参照：
{\it 余次元 $k$ で正則},
{\it $(R_k)$},
{\it $(S_k)$}


\noindent
\ref{properties-definition-nagata} を参照：
{\it 日本スキーム},
{\it 普遍的日本スキーム},
{\it 永田スキーム}


\noindent
\ref{properties-definition-G-scheme} を参照：
{\it G-スキーム}


\noindent
\ref{properties-definition-singular-locus} を参照：
{\it 正則点集合},
{\it 特異点集合}


\noindent
\ref{properties-definition-unibranch} を参照：
{\it $x$ において単枝},
{\it $x$ において幾何的単枝},
{\it 単枝},
{\it 幾何的単枝}


\noindent
\ref{properties-definition-number-of-branches} を参照：
{\it $X$ の $x$ における枝数},
{\it $x$ における $X$ の幾何的枝数}


\noindent
\ref{properties-definition-quasi-affine} を参照：
{\it 準アフィン}


\noindent
\ref{properties-definition-locally-projective} を参照：
{\it 局所射影的}


\noindent
\ref{properties-definition-kappa-generated} を参照：
{\it $\kappa$-生成}


\noindent
\ref{properties-definition-subsheaf-sections-annihilated-by-ideal} を参照：
{\it $\mathcal{I}$ により零化される切断の部分層}


\noindent
\ref{properties-definition-subsheaf-sections-supported-on-closed} を参照：
{\it $T$ 上に台を持つ切断の部分層}


\noindent
\ref{properties-definition-ample} を参照：
{\it 豊富}


\medskip\noindent
{\bf スキームの射}

\medskip

\noindent
\ref{morphisms-definition-scheme-theoretic-intersection-union} を参照：
{\it スキーム論的交叉},
{\it スキーム論的和}


\noindent
\ref{morphisms-definition-scheme-theoretic-support} を参照：
{\it $\mathcal{F}$ のスキーム論的台}


\noindent
\ref{morphisms-definition-scheme-theoretic-image} を参照：
{\it スキーム論的像}


\noindent
\ref{morphisms-definition-scheme-theoretically-dense} を参照：
{\it $X$ における $U$ のスキーム論的閉包},
{\it $X$ においてスキーム論的に稠密}


\noindent
\ref{morphisms-definition-dominant} を参照：
{\it 支配的}


\noindent
\ref{morphisms-definition-surjective} を参照：
{\it 全射}


\noindent
\ref{morphisms-definition-universally-injective} を参照：
{\it 普遍単射},
{\it 純非分離}


\noindent
\ref{morphisms-definition-affine} を参照：
{\it アフィン}


\noindent
\ref{morphisms-definition-family-ample-invertible-modules} を参照：
{\it $X$ 上の豊富な可逆層の族}


\noindent
\ref{morphisms-definition-quasi-affine} を参照：
{\it 準アフィン}


\noindent
\ref{morphisms-definition-property-local} を参照：
{\it 局所的},
{\it 基底変換で安定},
{\it 合成で安定}


\noindent
\ref{morphisms-definition-locally-P} を参照：
{\it 局所的に $P$ 型}


\noindent
\ref{morphisms-definition-finite-type} を参照：
{\it $x \in X$ において有限型},
{\it 局所有限型},
{\it 有限型}


\noindent
\ref{morphisms-definition-finite-type-point} を参照：
{\it 有限型点}


\noindent
\ref{morphisms-definition-universally-catenary} を参照：
{\it 普遍鎖状}


\noindent
\ref{morphisms-definition-J} を参照：
{\it J-2}


\noindent
\ref{morphisms-definition-excellent} を参照：
{\it 準優秀},
{\it 優秀}


\noindent
\ref{morphisms-definition-quasi-finite} を参照：
{\it 点 $x \in X$ において準有限},
{\it 局所準有限},
{\it 準有限}


\noindent
\ref{morphisms-definition-finite-presentation} を参照：
{\it $x \in X$ において有限表示},
{\it 局所有限表示},
{\it 有限表示}


\noindent
\ref{morphisms-definition-open} を参照：
{\it 開},
{\it 普遍開}


\noindent
\ref{morphisms-definition-submersive} を参照：
{\it 商位相的},
{\it 普遍的に商写像的}


\noindent
\ref{morphisms-definition-flat} を参照：
{\it 点 $x \in X$ において平坦},
{\it $S$ 上、点 $x \in X$ において平坦},
{\it 平坦},
{\it $S$ 上平坦}


\noindent
\ref{morphisms-definition-scheme-structure-connected-component} を参照：
{\it $T$ 上の標準的なスキーム構造}


\noindent
\ref{morphisms-definition-relative-dimension-d} を参照：
{\it 相対次元が $\leq d$ である（点 $x$ において）},
{\it 相対次元 $\leq d$},
{\it 相対次元 $d$}


\noindent
\ref{morphisms-definition-syntomic} を参照：
{\it 点 $x \in X$ においてシントミック},
{\it シントミック},
{\it $k$ 上の局所完全交叉},
{\it 標準シントミック}


\noindent
\ref{morphisms-definition-syntomic-relative-dimension} を参照：
{\it 相対次元 $d$ のシントミック射}


\noindent
\ref{morphisms-definition-conormal-sheaf} を参照：
{\it $Z$ の $X$ における余法層 $\mathcal{C}_{Z/X}$},
{\it $i$ の余法層}


\noindent
\ref{morphisms-definition-sheaf-differentials} を参照：
{\it $S$ 上の $X$ の微分層 $\Omega_{X/S}$},
{\it 普遍 $S$-導分}


\noindent
\ref{morphisms-definition-smooth} を参照：
{\it 点 $x \in X$ において滑らか},
{\it 滑らか},
{\it 標準滑らか}


\noindent
\ref{morphisms-definition-smooth-relative-dimension} を参照：
{\it 相対次元 $d$ の滑らかな射}


\noindent
\ref{morphisms-definition-unramified} を参照：
{\it $x \in X$ において不分岐},
{\it $x \in X$ において G-不分岐},
{\it 不分岐},
{\it G-不分岐}


\noindent
\ref{morphisms-definition-etale} を参照：
{\it $x \in X$ においてエタール},
{\it エタール},
{\it 標準エタール}


\noindent
\ref{morphisms-definition-relatively-ample} を参照：
{\it 相対的に豊富},
{\it $f$-相対的に豊富},
{\it $X/S$ 上豊富},
{\it $f$-豊富}


\noindent
\ref{morphisms-definition-very-ample} を参照：
{\it 相対的に非常に豊富},
{\it $f$-相対的に非常に豊富},
{\it $X/S$ 上非常に豊富},
{\it $f$-非常に豊富}


\noindent
\ref{morphisms-definition-quasi-projective} を参照：
{\it 準射影的},
{\it H-準射影的},
{\it 局所準射影的}


\noindent
\ref{morphisms-definition-proper} を参照：
{\it 固有}


\noindent
\ref{morphisms-definition-projective} を参照：
{\it 射影的},
{\it H-射影的},
{\it 局所射影的}


\noindent
\ref{morphisms-definition-integral} を参照：
{\it 整},
{\it 有限}


\noindent
\ref{morphisms-definition-universal-homeomorphism} を参照：
{\it 普遍同相射}


\noindent
\ref{morphisms-definition-seminormal-ring} を参照：
{\it 半正規},
{\it 絶対弱正規}


\noindent
\ref{morphisms-definition-seminormal} を参照：
{\it 半正規},
{\it 絶対弱正規}


\noindent
\ref{morphisms-definition-seminormalization} を参照：
{\it 半正規化},
{\it 絶対弱正規化}


\noindent
\ref{morphisms-definition-finite-locally-free} を参照：
{\it 有限局所自由},
{\it 階数},
{\it 次数}


\noindent
\ref{morphisms-definition-rational-map} を参照：
{\it 同値},
{\it $X$ から $Y$ への有理写像},
{\it $S$-有理写像（$X$ から $Y$ への）}


\noindent
\ref{morphisms-definition-rational-function} を参照：
{\it $X$ 上の有理関数}


\noindent
\ref{morphisms-definition-ring-of-rational-functions} を参照：
{\it $X$ 上の有理関数環}


\noindent
\ref{morphisms-definition-function-field} を参照：
{\it 関数体},
{\it 有理関数体}


\noindent
\ref{morphisms-definition-domain-of-definition} を参照：
{\it 点 $x \in X$ で定義される},
{\it 定義域}


\noindent
\ref{morphisms-definition-dominant-rational} を参照：
{\it 支配的}


\noindent
\ref{morphisms-definition-birational-schemes} を参照：
{\it 双有理},
{\it $S$-双有理}


\noindent
\ref{morphisms-definition-birational} を参照：
{\it 双有理}


\noindent
\ref{morphisms-definition-degree} を参照：
{\it $X$ の $Y$ 上の次数}


\noindent
\ref{morphisms-definition-modification} を参照：
{\it $X$ の修正}


\noindent
\ref{morphisms-definition-alteration} を参照：
{\it $X$ のオルタレーション}


\noindent
\ref{morphisms-definition-integral-closure} を参照：
{\it $\mathcal{O}_X$ の整閉包 （$\mathcal{A}$ におけるもの）}


\noindent
\ref{morphisms-definition-normalization-X-in-Y} を参照：
{\it $X$ の正規化 （$Y$ におけるもの）}


\noindent
\ref{morphisms-definition-normalization} を参照：
{\it 正規化}


\noindent
\ref{morphisms-definition-relative-seminormalization} を参照：
{\it $X$ の $Y$ における半正規化},
{\it $X$ の $Y$ における弱正規化}


\noindent
\ref{morphisms-definition-weak-normalization} を参照：
{\it 弱正規化}


\noindent
\ref{morphisms-definition-weakly-normal} を参照：
{\it 弱正規}


\noindent
\ref{morphisms-definition-universally-bounded} を参照：
{\it $f$ のファイバーの次数を有界にする},
{\it $f$ のファイバーは普遍的に有界である}


\medskip\noindent
{\bf スキームのコホモロジー}

\medskip

\noindent
\ref{coherent-definition-depth} を参照：
{\it 点における深さ $k$},
{\it 点における深さ $k$},
{\it $(S_k)$},
{\it $(S_k)$}


\noindent
\ref{coherent-definition-Cohen-Macaulay} を参照：
{\it Cohen--Macaulay}


\noindent
\ref{coherent-definition-proper-over-base} を参照：
{\it $Z$ は $S$ 上固有}


\medskip\noindent
{\bf 除数}

\medskip

\noindent
\ref{divisors-definition-associated} を参照：
{\it 随伴する},
{\it $X$ の随伴点}


\noindent
\ref{divisors-definition-embedded} を参照：
{\it 埋め込まれた随伴点},
{\it 埋め込み点},
{\it 埋め込み成分}


\noindent
\ref{divisors-definition-weakly-associated} を参照：
{\it 弱随伴する},
{\it $X$ の弱随伴点}


\noindent
\ref{divisors-definition-relative-assassin} を参照：
{\it $S$ 上の $X$ における $\mathcal{F}$ の相対随伴点集合}


\noindent
\ref{divisors-definition-relative-weak-assassin} を参照：
{\it $S$ 上の $X$ における $\mathcal{F}$ の相対弱随伴点集合}


\noindent
\ref{divisors-definition-torsion} を参照：
{\it ねじれ},
{\it ねじれがない}


\noindent
\ref{divisors-definition-rank} を参照：
{\it 階数}


\noindent
\ref{divisors-definition-reflexive} を参照：
{\it 反射包},
{\it 反射的}


\noindent
\ref{divisors-definition-effective-Cartier-divisor} を参照：
{\it 局所主閉部分スキーム},
{\it 有効 Cartier 因子}


\noindent
\ref{divisors-definition-sum-effective-Cartier-divisors} を参照：
{\it 有効 Cartier 因子 $D_1$ と $D_2$ の和}


\noindent
\ref{divisors-definition-pullback-effective-Cartier-divisor} を参照：
{\it $D$ の $f$ による引き戻しは定義される},
{\it 有効 Cartier 因子の引き戻し}


\noindent
\ref{divisors-definition-invertible-sheaf-effective-Cartier-divisor} を参照：
{\it $D$ に付随する可逆層 $\mathcal{O}_S(D)$},
{\it 標準切断}


\noindent
\ref{divisors-definition-regular-section} を参照：
{\it 正則切断}


\noindent
\ref{divisors-definition-zero-scheme-s} を参照：
{\it 零点スキーム}


\noindent
\ref{divisors-definition-relative-effective-Cartier-divisor} を参照：
{\it 相対有効 Cartier 因子}


\noindent
\ref{divisors-definition-conormal-sheaf} を参照：
{\it 余法代数 $\mathcal{C}_{Z/X, *}$、すなわち $Z$ の $X$ における余法代数},
{\it $f$ の余法代数}


\noindent
\ref{divisors-definition-normal-cone} を参照：
{\it 法錐 $C_ZX$},
{\it 法束}


\noindent
\ref{divisors-definition-regular-ideal-sheaf} を参照：
{\it 正則},
{\it Koszul-正則},
{\it $H_1$-正則},
{\it 準正則}


\noindent
\ref{divisors-definition-regular-immersion} を参照：
{\it 正則埋め込み},
{\it Koszul 正則埋め込み},
{\it $H_1$-正則埋め込み},
{\it 準正則埋め込み}


\noindent
\ref{divisors-definition-relative-H1-regular-immersion} を参照：
{\it 相対準正則埋め込み},
{\it 相対 $H_1$-正則埋め込み}


\noindent
\ref{divisors-definition-sheaf-meromorphic-functions} を参照：
{\it $X$ 上の有理型関数の層},
{\it $\mathcal{K}_X$},
{\it 有理型関数}


\noindent
\ref{divisors-definition-meromorphic-section} を参照：
{\it $\mathcal{F}$ の有理型切断}


\noindent
\ref{divisors-definition-pullback-meromorphic-sections} を参照：
{\it $f$ に対して有理型関数の引き戻しが定義される}


\noindent
\ref{divisors-definition-regular-meromorphic-section} を参照：
{\it 正則}


\noindent
\ref{divisors-definition-regular-meromorphic-ideal-denominators} を参照：
{\it $s$ の分母イデアル層}


\noindent
\ref{divisors-definition-Weil-divisor} を参照：
{\it 素因子},
{\it Weil 因子}


\noindent
\ref{divisors-definition-order-vanishing} を参照：
{\it $f$ の $Z$ に沿う 消滅次数}


\noindent
\ref{divisors-definition-principal-divisor} を参照：
{\it $f$ に付随する主 Weil 因子}


\noindent
\ref{divisors-definition-class-group} を参照：
{\it Weil 因子類群}


\noindent
\ref{divisors-definition-order-vanishing-meromorphic} を参照：
{\it $s$ の $Z$ に沿う消滅次数}


\noindent
\ref{divisors-definition-divisor-invertible-sheaf} を参照：
{\it $s$ に付随する Weil 因子},
{\it $\mathcal{L}$ に付随する Weil 因子類}


\noindent
\ref{divisors-definition-blow-up} を参照：
{\it $X$ の $Z$ に沿ったブローアップ},
{\it $X$ のイデアル層 $\mathcal{I}$ におけるブローアップ},
{\it 例外因子},
{\it 中心}


\noindent
\ref{divisors-definition-strict-transform} を参照：
{\it 狭義変換},
{\it 狭義変換}


\noindent
\ref{divisors-definition-admissible-blowup} を参照：
{\it $U$-許容ブローアップ}


\medskip\noindent
{\bf スキームの極限}

\medskip

\medskip\noindent
{\bf 代数多様体}

\medskip

\noindent
\ref{varieties-definition-variety} を参照：
{\it 多様体}


\noindent
\ref{varieties-definition-geometrically-reduced} を参照：
{\it $x$ で幾何学的被約},
{\it 幾何学的被約}


\noindent
\ref{varieties-definition-geometrically-connected} を参照：
{\it 幾何学的連結}


\noindent
\ref{varieties-definition-geometrically-irreducible} を参照：
{\it 幾何学的既約}


\noindent
\ref{varieties-definition-geometrically-integral} を参照：
{\it $x$ において幾何学的に点ごとに整},
{\it 幾何学的に点ごとに整},
{\it 幾何学的整}


\noindent
\ref{varieties-definition-geometrically-normal} を参照：
{\it $x$ において幾何学的に正規},
{\it 幾何学的に正規}


\noindent
\ref{varieties-definition-geometrically-regular} を参照：
{\it $x$ において幾何学的に正則},
{\it $k$ 上幾何学的に正則}


\noindent
\ref{varieties-definition-dual-numbers} を参照：
{\it 双対数}


\noindent
\ref{varieties-definition-tangent-space} を参照：
{\it $S$ 上の $X$ の $x$ における接空間},
{\it 接ベクトル}


\noindent
\ref{varieties-definition-algebraic-scheme} を参照：
{\it 代数的 $k$-スキーム},
{\it 局所代数的 $k$-スキーム}


\noindent
\ref{varieties-definition-variety-type} を参照：
{\it アフィン多様体},
{\it 射影多様体},
{\it 準射影多様体},
{\it 固有多様体},
{\it 滑らかな多様体}


\noindent
\ref{varieties-definition-euler-characteristic} を参照：
{\it $\mathcal{F}$ の Euler 標数}


\noindent
\ref{varieties-definition-regularity} を参照：
{\it $m$-正則}


\noindent
\ref{varieties-definition-hilbert-polynomial} を参照：
{\it ヒルベルト多項式}


\noindent
\ref{varieties-definition-absolute-frobenius} を参照：
{\it $X$ の絶対 Frobenius}


\noindent
\ref{varieties-definition-relative-frobenius} を参照：
{\it $X/S$ の相対 Frobenius 射}


\noindent
\ref{varieties-definition-delta-invariant-algebra} を参照：
{\it $A$ の $\delta$-不変量}


\noindent
\ref{varieties-definition-delta-invariant} を参照：
{\it $x$ における $X$ の $\delta$-不変量}


\noindent
\ref{varieties-definition-wedge} を参照：
{\it $A$ は $A_1, \ldots, A_n$ のウェッジである}


\noindent
\ref{varieties-definition-curve} を参照：
{\it 曲線}


\noindent
\ref{varieties-definition-degree-invertible-sheaf} を参照：
{\it 次数},
{\it 次数}


\noindent
\ref{varieties-definition-intersection-number} を参照：
{\it 交点数}


\noindent
\ref{varieties-definition-degree} を参照：
{\it $\mathcal{L}$ に関する $Z$ の次数}


\noindent
\ref{varieties-definition-embed-dim} を参照：
{\it $x$ における $X$ の埋め込み次元}


\noindent
\ref{varieties-definition-embedding-dimension} を参照：
{\it $x$ における $X/k$ の埋め込み次元}


\medskip\noindent
{\bf スキーム上の位相}

\medskip

\noindent
\ref{topologies-definition-zariski-covering} を参照：
{\it $T$ の Zariski 被覆}


\noindent
\ref{topologies-definition-standard-Zariski} を参照：
{\it 標準 Zariski 被覆}


\noindent
\ref{topologies-definition-big-zariski-site} を参照：
{\it 大 Zariski サイト}


\noindent
\ref{topologies-definition-big-small-Zariski} を参照：
{\it $S$ の大 Zariski サイト},
{\it $S$ の小 Zariski サイト},
{\it $S$ の大アフィン Zariski サイト},
{\it $S$ の小アフィン Zariski サイト}


\noindent
\ref{topologies-definition-restriction-small-zariski} を参照：
{\it 小 Zariski サイトへの制限}


\noindent
\ref{topologies-definition-etale-covering} を参照：
{\it $T$ のエタール被覆}


\noindent
\ref{topologies-definition-standard-etale} を参照：
{\it 標準エタール被覆}


\noindent
\ref{topologies-definition-big-etale-site} を参照：
{\it 大エタールサイト}


\noindent
\ref{topologies-definition-big-small-etale} を参照：
{\it $S$ の大エタール・サイト},
{\it $S$ の小エタール・サイト},
{\it $S$ の大アフィンエタール・サイト},
{\it $S$ の小アフィンエタール・サイト}


\noindent
\ref{topologies-definition-restriction-small-etale} を参照：
{\it 小エタール・サイトへの制限}


\noindent
\ref{topologies-definition-smooth-covering} を参照：
{\it $T$ の滑らか被覆}


\noindent
\ref{topologies-definition-standard-smooth} を参照：
{\it 標準平滑被覆}


\noindent
\ref{topologies-definition-big-smooth-site} を参照：
{\it 大平滑サイト}


\noindent
\ref{topologies-definition-big-small-smooth} を参照：
{\it $S$ の大滑らかサイト},
{\it $S$ の大アフィン滑らかサイト}


\noindent
\ref{topologies-definition-syntomic-covering} を参照：
{\it $T$ のシントミック被覆}


\noindent
\ref{topologies-definition-standard-syntomic} を参照：
{\it 標準シントミック被覆}


\noindent
\ref{topologies-definition-big-syntomic-site} を参照：
{\it 大シントミック・サイト}


\noindent
\ref{topologies-definition-big-small-syntomic} を参照：
{\it $S$ の大シントミック・サイト},
{\it $S$ の大アフィンシントミック・サイト}


\noindent
\ref{topologies-definition-fppf-covering} を参照：
{\it $T$ の fppf 被覆}


\noindent
\ref{topologies-definition-standard-fppf} を参照：
{\it 標準 fppf 被覆}


\noindent
\ref{topologies-definition-big-fppf-site} を参照：
{\it 大 fppf サイト}


\noindent
\ref{topologies-definition-big-small-fppf} を参照：
{\it $S$ の大 fppf サイト},
{\it $S$ の大アフィン fppf サイト}


\noindent
\ref{topologies-definition-standard-ph-covering} を参照：
{\it 標準 ph 被覆}


\noindent
\ref{topologies-definition-ph-covering} を参照：
{\it $T$ の ph 被覆}


\noindent
\ref{topologies-definition-big-ph-site} を参照：
{\it 大 ph サイト}


\noindent
\ref{topologies-definition-big-small-ph} を参照：
{\it $S$ の大 ph サイト},
{\it $S$ の大アフィン ph サイト}


\noindent
\ref{topologies-definition-fpqc-covering} を参照：
{\it $T$ の fpqc 被覆}


\noindent
\ref{topologies-definition-standard-fpqc} を参照：
{\it 標準 fpqc 被覆}


\noindent
\ref{topologies-definition-sheaf-property-fpqc} を参照：
{\it 与えられた族に対する層条件を満たす},
{\it fpqc 位相に対する層条件を満たす}


\noindent
\ref{topologies-definition-standard-V-covering} を参照：
{\it 標準 V 被覆}


\noindent
\ref{topologies-definition-V-covering} を参照：
{\it $T$ の V 被覆}


\noindent
\ref{topologies-definition-sheaf-property-V} を参照：
{\it V 位相に対する層性質を満たす}


\medskip\noindent
{\bf 降下}

\medskip

\noindent
\ref{descent-definition-descent-datum-quasi-coherent} を参照：
{\it 準連接層の降下データ $(\mathcal{F}_i, \varphi_{ij})$},
{\it コサイクル条件},
{\it 降下データの射 $\psi : (\mathcal{F}_i, \varphi_{ij}) \to (\mathcal{F}'_i, \varphi'_{ij})$}


\noindent
\ref{descent-definition-descent-datum-effective-quasi-coherent} を参照：
{\it 自明な降下データ},
{\it 標準降下データ},
{\it 有効}


\noindent
\ref{descent-definition-descent-datum-modules} を参照：
{\it $R \to A$ に関する加群の降下データ $(N, \varphi)$},
{\it コサイクル条件},
{\it 降下データの射 $(N, \varphi) \to (N', \varphi')$}


\noindent
\ref{descent-definition-descent-datum-effective-module} を参照：
{\it 有効}


\noindent
\ref{descent-definition-split-equalizer} を参照：
{\it 分裂等化子}


\noindent
\ref{descent-definition-universally-injective} を参照：
{\it 普遍単射}


\noindent
\ref{descent-definition-C} を参照：
{\it $C$}


\noindent
\ref{descent-definition-effective-descent} を参照：
{\it $f$ に沿う基底拡大},
{\it 加群の降下射},
{\it 加群の有効降下射}


\noindent
\ref{descent-definition-pushforward} を参照：
{\it $f_*$}


\noindent
\ref{descent-definition-structure-sheaf} を参照：
{\it 大サイト $(\Sch/S)_\tau$ の構造層},
{\it 構造層},
{\it $\mathcal{F}$ に付随する $\mathcal{O}$-加群の層},
{\it $\mathcal{F}$ に付随する $\mathcal{O}$-加群の層}


\noindent
\ref{descent-definition-parasitic} を参照：
{\it 寄生的},
{\it $\tau$-位相に関して寄生的}


\noindent
\ref{descent-definition-property-local} を参照：
{\it $\tau$-位相について局所的}


\noindent
\ref{descent-definition-germs} を参照：
{\it $x$ における $X$ の芽},
{\it 芽の射},
{\it 芽の射の合成}


\noindent
\ref{descent-definition-etale-morphism-germs} を参照：
{\it エタール},
{\it 滑らか}


\noindent
\ref{descent-definition-local-at-point} を参照：
{\it エタール局所的},
{\it 滑らか局所的}


\noindent
\ref{descent-definition-property-morphisms-local} を参照：
{\it 基底上 $\tau$-局所的},
{\it 標的上 $\tau$-局所的},
{\it $\tau$-位相について基底上局所的}


\noindent
\ref{descent-definition-property-morphisms-local-source} を参照：
{\it 始域上 $\tau$-局所的},
{\it $\tau$-位相について始域上局所的}


\noindent
\ref{descent-definition-local-source-target} を参照：
{\it 始域と標的上エタール局所的}


\noindent
\ref{descent-definition-local-source-target-at-point} を参照：
{\it 始域と標的上エタール局所的}


\noindent
\ref{descent-definition-descent-datum} を参照：
{\it $V/X/S$ の降下データ},
{\it コサイクル条件},
{\it $X \to S$ に関する降下データ},
{\it $X \to S$ に関する降下データの射 $g : (V/X, \varphi) \to (V'/X, \varphi')$}


\noindent
\ref{descent-definition-descent-datum-for-family-of-morphisms} を参照：
{\it 族 $\{X_i \to S\}$ に関する降下データ $(V_i, \varphi_{ij})$},
{\it 降下データの射 $\psi : (V_i, \varphi_{ij}) \to (V'_i, \varphi'_{ij})$}


\noindent
\ref{descent-definition-pullback-functor} を参照：
{\it 引き戻し関手}


\noindent
\ref{descent-definition-pullback-functor-family} を参照：
{\it 引き戻し関手}


\noindent
\ref{descent-definition-effective} を参照：
{\it 自明な降下データ},
{\it 標準降下データ},
{\it 有効}


\noindent
\ref{descent-definition-effective-family} を参照：
{\it 標準降下データ},
{\it 有効}


\noindent
\ref{descent-definition-descending-types-morphisms} を参照：
{\it 型 $\mathcal{P}$ の射は $\tau$-被覆に対する降下を満たす}


\medskip\noindent
{\bf スキームの導来圏}

\medskip

\noindent
\ref{perfect-definition-supported-on} を参照：
{\it $T$ に台をもつ}


\noindent
\ref{perfect-definition-approximation-holds} を参照：
{\it 三つ組に対して近似が成立する}


\noindent
\ref{perfect-definition-approximation} を参照：
{\it 完全複体による近似が 成立する}


\noindent
\ref{perfect-definition-tor-independent} を参照：
{\it $S$ 上 Tor 独立}


\noindent
\ref{perfect-definition-relatively-perfect} を参照：
{\it $S$ に関して完全},
{\it $S$-完全}


\noindent
\ref{perfect-definition-resolution-property} を参照：
{\it 分解性質}


\noindent
\ref{perfect-definition-K-group} を参照：
{\it $X$ のグロタンディーク群},
{\it $X$ 上のコヒーレント層のグロタンディーク群}


\medskip\noindent
{\bf 射の補足}

\medskip

\noindent
\ref{more-morphisms-definition-thickening} を参照：
{\it 肥厚化},
{\it 一次肥厚化},
{\it 有限次肥厚化},
{\it $n$ 次肥厚化},
{\it 肥厚化の射},
{\it $S$ 上の肥厚化},
{\it $S$ 上の肥厚化の射}


\noindent
\ref{more-morphisms-definition-first-order-infinitesimal-neighbourhood} を参照：
{\it 一次無限小近傍},
{\it $n$ 次無限小近傍}


\noindent
\ref{more-morphisms-definition-formally-unramified} を参照：
{\it 形式的不分岐}


\noindent
\ref{more-morphisms-definition-universal-thickening} を参照：
{\it 普遍一次肥厚化},
{\it $Z$ の $X$ 上の余法層}


\noindent
\ref{more-morphisms-definition-formally-etale} を参照：
{\it 形式的エタール}


\noindent
\ref{more-morphisms-definition-formally-smooth} を参照：
{\it 形式的に滑らか}


\noindent
\ref{more-morphisms-definition-netherlander} を参照：
{\it $f$ の素朴な余接複体}


\noindent
\ref{more-morphisms-definition-normal} を参照：
{\it $x$ で正規},
{\it 正規射}


\noindent
\ref{more-morphisms-definition-regular} を参照：
{\it 点 $x$ で正則},
{\it 正則射}


\noindent
\ref{more-morphisms-definition-CM} を参照：
{\it 点 $x$ で Cohen--Macaulay},
{\it Cohen--Macaulay 射}


\noindent
\ref{more-morphisms-definition-etale-neighbourhood} を参照：
{\it $(S, s)$ のエタール近傍},
{\it エタール近傍の射},
{\it 初等エタール近傍}


\noindent
\ref{more-morphisms-definition-relatively-finitely-presented-sheaf} を参照：
{\it $S$ に相対的に有限表示される},
{\it $S$ に相対的有限表示をもつ}


\noindent
\ref{more-morphisms-definition-relative-pseudo-coherence} を参照：
{\it $S$ に相対的に $m$-擬連接},
{\it $S$ に相対的に擬連接},
{\it $S$ に相対的に $m$-擬連接},
{\it $S$ に相対的に擬連接}


\noindent
\ref{more-morphisms-definition-pseudo-coherent} を参照：
{\it 擬連接}


\noindent
\ref{more-morphisms-definition-perfect} を参照：
{\it 完全}


\noindent
\ref{more-morphisms-definition-lci} を参照：
{\it $x$ で Koszul},
{\it Koszul 射},
{\it 局所完全交叉射}


\noindent
\ref{more-morphisms-definition-weakly-etale} を参照：
{\it 弱エタール},
{\it 絶対平坦}


\noindent
\ref{more-morphisms-definition-ind-quasi-affine} を参照：
{\it ind-準アフィン},
{\it ind-準アフィン}


\noindent
\ref{more-morphisms-definition-affine-stratification} を参照：
{\it アフィン層別}


\noindent
\ref{more-morphisms-definition-affine-stratification-number} を参照：
{\it アフィン層別数}


\noindent
\ref{more-morphisms-definition-weighting} を参照：
{\it 重み付け},
{\it 重み付け}


\noindent
\ref{more-morphisms-definition-cd-morphism} を参照：
{\it 完全分解},
{\it 完全分解}


\medskip\noindent
{\bf 平坦性の補足}

\medskip

\noindent
\ref{flat-definition-one-step-devissage} を参照：
{\it $s$ 上の $\mathcal{F}/X/S$ の一段階デヴィサージュ}


\noindent
\ref{flat-definition-one-step-devissage-at-x} を参照：
{\it $x$ における $\mathcal{F}/X/S$ の一段階デヴィサージュ}


\noindent
\ref{flat-definition-shrink} を参照：
{\it 標準縮小}


\noindent
\ref{flat-definition-complete-devissage} を参照：
{\it $s$ 上の $\mathcal{F}/X/S$ の完全デヴィサージュ}


\noindent
\ref{flat-definition-complete-devissage-at-x} を参照：
{\it $x$ における $\mathcal{F}/X/S$ の完全デヴィサージュ}


\noindent
\ref{flat-definition-shrink-complete} を参照：
{\it 標準縮小}


\noindent
\ref{flat-definition-elementary-etale-neighbourhood} を参照：
{\it $\mathfrak q$ における環準同型 $R \to S$ の基本エタール局所化}


\noindent
\ref{flat-definition-complete-devissage-algebra} を参照：
{\it $\mathfrak r$ 上の $N/S/R$ の完全デヴィサージュ}


\noindent
\ref{flat-definition-complete-devissage-at-x-algebra} を参照：
{\it $\mathfrak q$ における $N/S/R$ の完全デヴィサージュ}


\noindent
\ref{flat-definition-impurity} を参照：
{\it $s$ 上の $\mathcal{F}$ の非純性}


\noindent
\ref{flat-definition-pure} を参照：
{\it $X_s$ に沿って純粋},
{\it $X_s$ に沿って普遍的に純粋},
{\it $X_s$ に沿って純粋},
{\it 普遍的 $S$-純粋},
{\it $S$ に関して普遍的に純粋},
{\it $S$-純粋},
{\it $S$ に関して純粋},
{\it $S$-純粋},
{\it $S$ に関して純粋}


\noindent
\ref{flat-definition-flat-dimension-n} を参照：
{\it $\mathcal{F}$ は次元 $\geq n$ で $S$ 上平坦である}


\noindent
\ref{flat-definition-flattening} を参照：
{\it $\mathcal{F}$ の普遍平坦化が存在する},
{\it $X$ の普遍平坦化が存在する}


\noindent
\ref{flat-definition-flattening-stratification} を参照：
{\it 平坦化層別},
{\it 平坦化層別}


\noindent
\ref{flat-definition-h-covering} を参照：
{\it $T$ の h 被覆}


\noindent
\ref{flat-definition-big-h-site} を参照：
{\it 大 h サイト}


\noindent
\ref{flat-definition-standard-h} を参照：
{\it 標準 h 被覆}


\noindent
\ref{flat-definition-big-small-h} を参照：
{\it $S$ の大 h サイト},
{\it $S$ の大アフィン h サイト}


\medskip\noindent
{\bf 亜群スキーム}

\medskip

\noindent
\ref{groupoids-definition-equivalence-relation} を参照：
{\it 前関係},
{\it 関係},
{\it 前同値関係},
{\it $S$ 上の $U$ の同値関係}


\noindent
\ref{groupoids-definition-restrict-relation} を参照：
{\it 制限},
{\it 引き戻し}


\noindent
\ref{groupoids-definition-group-scheme} を参照：
{\it $S$ 上の群スキーム},
{\it $S$ 上の群スキームの射 $\psi : (G, m) \to (G', m')$}


\noindent
\ref{groupoids-definition-closed-subgroup-scheme} を参照：
{\it 閉部分群スキーム},
{\it 開部分群スキーム}


\noindent
\ref{groupoids-definition-smooth-group-scheme} を参照：
{\it 滑らかな群スキーム},
{\it 平坦な群スキーム},
{\it 分離的な群スキーム}


\noindent
\ref{groupoids-definition-abelian-variety} を参照：
{\it アーベル多様体}


\noindent
\ref{groupoids-definition-action-group-scheme} を参照：
{\it $G$ のスキーム $X/S$ への作用},
{\it 同変},
{\it $G$-同変}


\noindent
\ref{groupoids-definition-free-action} を参照：
{\it 自由}


\noindent
\ref{groupoids-definition-pseudo-torsor} を参照：
{\it 擬 $G$-トーサー},
{\it $G$ の下で形式的に主等質},
{\it 自明}


\noindent
\ref{groupoids-definition-principal-homogeneous-space} を参照：
{\it 主等質空間},
{\it $G$-トーサー},
{\it $\tau$ 位相における $G$-トーサー},
{\it $\tau$ $G$-トーサー},
{\it $\tau$ トーサー},
{\it 準等自明},
{\it 局所自明}


\noindent
\ref{groupoids-definition-equivariant-module} を参照：
{\it $G$-同変準連接 $\mathcal{O}_X$-加群},
{\it 同変準連接 $\mathcal{O}_X$-加群}


\noindent
\ref{groupoids-definition-groupoid} を参照：
{\it $S$ 上の亜群スキーム},
{\it $S$ 上の亜群},
{\it $S$ 上の亜群スキームの射 $f : (U, R, s, t, c) \to (U', R', s', t', c')$}


\noindent
\ref{groupoids-definition-groupoid-module} を参照：
{\it $(U, R, s, t, c)$ 上の準連接加群}


\noindent
\ref{groupoids-definition-stabilizer-groupoid} を参照：
{\it 亜群スキーム $(U, R, s, t, c)$ の安定化群}


\noindent
\ref{groupoids-definition-restrict-groupoid} を参照：
{\it $(U, R, s, t, c)$ の $U'$ への制限}


\noindent
\ref{groupoids-definition-invariant-open} を参照：
{\it 集合論的に $R$-不変},
{\it $R$-不変},
{\it $R$-不変},
{\it $R$-不変}


\noindent
\ref{groupoids-definition-quotient-sheaf} を参照：
{\it 商層 $U/R$}


\noindent
\ref{groupoids-definition-representable-quotient} を参照：
{\it 表現可能な商},
{\it 表現可能な商}


\noindent
\ref{groupoids-definition-cartesian-morphism} を参照：
{\it カルテジアン},
{\it $(U', R', s', t', c')$ は $(U, R, s, t, c)$ 上カルテジアンである},
{\it $(U, R, s, t, c)$ 上カルテジアンな亜群スキームの射}


\medskip\noindent
{\bf 亜群スキームの補足}

\medskip

\medskip\noindent
{\bf スキームのエタール射}

\medskip

\noindent
\ref{etale-definition-unramified-rings} を参照：
{\it 局所環の不分岐準同型}


\noindent
\ref{etale-definition-unramified-schemes} を参照：
{\it $x$ において不分岐},
{\it 不分岐}


\noindent
\ref{etale-definition-flat-rings} を参照：
{\it 平坦},
{\it 忠実平坦},
{\it 平坦（それぞれ忠実平坦）}


\noindent
\ref{etale-definition-flat-schemes} を参照：
{\it $x \in X$ において $Y$ 上平坦},
{\it $x \in X$ において平坦},
{\it 平坦},
{\it 忠実平坦}


\noindent
\ref{etale-definition-etale-ring} を参照：
{\it 局所環のエタール準同型}


\noindent
\ref{etale-definition-etale-schemes-1} を参照：
{\it $x \in X$ においてエタール},
{\it エタール}


\noindent
\ref{etale-definition-strict-normal-crossings} を参照：
{\it 単純正規交差因子}


\noindent
\ref{etale-definition-normal-crossings} を参照：
{\it 正規交差因子}


\medskip\noindent
{\bf Chow ホモロジー}

\medskip

\noindent
\ref{chow-definition-periodic-complex} を参照：
{\it $2$-周期複体},
{\it コホモロジー加群},
{\it 完全},
{\it $(2, 1)$-周期複体},
{\it コホモロジー加群}


\noindent
\ref{chow-definition-periodic-length} を参照：
{\it 重複度},
{\it （加法的）Herbrand 商}


\noindent
\ref{chow-definition-delta-dimension} を参照：
{\it $Z$ の $\delta$-次元}


\noindent
\ref{chow-definition-cycles} を参照：
{\it $X$ 上のサイクル},
{\it $k$-サイクル}


\noindent
\ref{chow-definition-support-cycle} を参照：
{\it 台}


\noindent
\ref{chow-definition-effective-cycle} を参照：
{\it 有効}


\noindent
\ref{chow-definition-cycle-associated-to-closed-subscheme} を参照：
{\it $Z$ における $Z'$ の重複度},
{\it $Z$ に付随する $k$-サイクル}


\noindent
\ref{chow-definition-cycle-associated-to-coherent-sheaf} を参照：
{\it $\mathcal{F}$ における $Z'$ の重複度},
{\it $\mathcal{F}$ に付随する $k$-サイクル}


\noindent
\ref{chow-definition-proper-pushforward} を参照：
{\it 順像}


\noindent
\ref{chow-definition-flat-pullback} を参照：
{\it $f$ による $\alpha$ の平坦引き戻し}


\noindent
\ref{chow-definition-principal-divisor} を参照：
{\it $f$ に付随する主因子}


\noindent
\ref{chow-definition-rational-equivalence} を参照：
{\it 零と有理同値},
{\it 有理同値},
{\it $X$ 上の $k$-サイクルの Chow 群},
{\it $X$ 上の有理同値を法とする $k$-サイクルの Chow 群}


\noindent
\ref{chow-definition-envelope} を参照：
{\it 包絡射}


\noindent
\ref{chow-definition-divisor-invertible-sheaf} を参照：
{\it $s$ に付随する Weil 因子},
{\it $\mathcal{L}$ に付随する Weil 因子}


\noindent
\ref{chow-definition-cap-c1} を参照：
{\it $\mathcal{L}$ の第一 Chern 類との交叉積}


\noindent
\ref{chow-definition-gysin-homomorphism} を参照：
{\it Gysin 準同型}


\noindent
\ref{chow-definition-bivariant-class} を参照：
{\it $f$ に対する次数 $p$ の双変類 $c$}


\noindent
\ref{chow-definition-chow-cohomology} を参照：
{\it Chow コホモロジー}


\noindent
\ref{chow-definition-first-chern-class} を参照：
{\it 第一 Chern 類}


\noindent
\ref{chow-definition-chern-classes} を参照：
{\it $X$ 上の $\mathcal{E}$ の Chern 類},
{\it $X$ 上の $\mathcal{E}$ の全 Chern 類}


\noindent
\ref{chow-definition-cap-chern-classes} を参照：
{\it $\mathcal{E}$ の第 $j$ Chern 類との交叉}


\noindent
\ref{chow-definition-chern-classes-final} を参照：
{\it 第 $i$ Chern 類},
{\it 全 Chern 類}


\noindent
\ref{chow-definition-degree-zero-cycle} を参照：
{\it 零サイクルの次数}


\noindent
\ref{chow-definition-defined-on-perfect} を参照：
{\it $E$ の Chern 類は定義される}


\noindent
\ref{chow-definition-localized-chern} を参照：
{\it 局所化 Chern 指標},
{\it 局所化第 $p$ Chern 類}


\noindent
\ref{chow-definition-lci-gysin} を参照：
{\it $f$ に対する Gysin 写像が存在する},
{\it Gysin 写像}


\noindent
\ref{chow-definition-determinant} を参照：
{\it 許容},
{\it 記号},
{\it 許容関係},
{\it 有限長 $R$-加群 $M$ の行列式}


\noindent
\ref{chow-definition-periodic-determinant} を参照：
{\it $(M, \varphi, \psi)$ の行列式}


\noindent
\ref{chow-definition-symbol-M} を参照：
{\it $M, a, b$ に付随する記号}


\noindent
\ref{chow-definition-tame-symbol} を参照：
{\it テイム記号}


\medskip\noindent
{\bf 交叉理論}

\medskip

\noindent
\ref{intersection-definition-proper-intersection} を参照：
{\it 正しく交わる},
{\it 正しく交わる}


\noindent
\ref{intersection-definition-multiplicity} を参照：
{\it 定義イデアル $I$ に関する $M$ の重複度}


\medskip\noindent
{\bf 曲線の Picard スキーム}

\medskip

\noindent
\ref{pic-definition-picard-functor} を参照：
{\it Picard 函手}


\noindent
\ref{pic-definition-genus} を参照：
{\it 種数}


\medskip\noindent
{\bf Weil コホモロジー理論}

\medskip

\noindent
\ref{weil-definition-chow-group-motives} を参照：
{\it $M$ の第 $i$ Chow 群}


\noindent
\ref{weil-definition-weil-cohomology-theory-classical} を参照：
{\it 古典的 Weil コホモロジー理論}


\noindent
\ref{weil-definition-weil-cohomology-theory} を参照：
{\it Weil コホモロジー理論}


\medskip\noindent
{\bf 適切な加群}

\medskip

\noindent
\ref{adequate-definition-module-valued-functor} を参照：
{\it 加群値函手},
{\it 加群値函手の射}


\noindent
\ref{adequate-definition-adequate-functor} を参照：
{\it 適切},
{\it 線形適切}


\noindent
\ref{adequate-definition-adequate} を参照：
{\it 適切}


\noindent
\ref{adequate-definition-category-adequate-modules} を参照：
{\it $\textit{Adeq}(\mathcal{O})$},
{\it $\textit{Adeq}((\Sch/S)_\tau, \mathcal{O})$},
{\it $\textit{Adeq}(S)$}


\noindent
\ref{adequate-definition-pure} を参照：
{\it 純射影},
{\it 純入射}


\noindent
\ref{adequate-definition-pure-resolution} を参照：
{\it 純射影分解},
{\it 純入射分解}


\noindent
\ref{adequate-definition-pure-ext} を参照：
{\it 純拡大加群}


\medskip\noindent
{\bf 双対化複体}

\medskip

\noindent
\ref{dualizing-definition-essential} を参照：
{\it 本質的},
{\it 本質拡大},
{\it 本質的}


\noindent
\ref{dualizing-definition-projective-cover} を参照：
{\it 射影被覆},
{\it 射影包絡}


\noindent
\ref{dualizing-definition-injective-hull} を参照：
{\it 入射包絡}


\noindent
\ref{dualizing-definition-indecomposable} を参照：
{\it 直既約}


\noindent
\ref{dualizing-definition-dualizing} を参照：
{\it 双対化複体}


\noindent
\ref{dualizing-definition-gorenstein} を参照：
{\it Gorenstein},
{\it Gorenstein}


\noindent
\ref{dualizing-definition-relative-dualizing-complex} を参照：
{\it 相対双対化複体}


\medskip\noindent
{\bf スキームの双対性}

\medskip

\noindent
\ref{duality-definition-dualizing-scheme} を参照：
{\it 双対化複体}


\noindent
\ref{duality-definition-good-dualizing} を参照：
{\it $\omega_S^\bullet$ に関して正規化された双対化複体}


\noindent
\ref{duality-definition-gorenstein} を参照：
{\it Gorenstein}


\noindent
\ref{duality-definition-gorenstein-morphism} を参照：
{\it $x$ において Gorenstein},
{\it Gorenstein 射}


\noindent
\ref{duality-definition-relative-dualizing-complex} を参照：
{\it 相対双対化複体}


\medskip\noindent
{\bf 判別式と不同イデアル}

\medskip

\noindent
\ref{discriminant-definition-trace-element} を参照：
{\it トレース元}


\noindent
\ref{discriminant-definition-kahler-different} を参照：
{\it K\"ahler の差イデアル}


\noindent
\ref{discriminant-definition-different} を参照：
{\it 差イデアル}


\medskip\noindent
{\bf de Rham コホモロジー}

\medskip

\noindent
\ref{derham-definition-hodge-filtration} を参照：
{\it ホッジフィルトレーション}


\noindent
\ref{derham-definition-local-product} を参照：
{\it $Y \subset X$ に対する $S$ 上の対数極ド・ラーム複体が定義される}


\noindent
\ref{derham-definition-log-complex} を参照：
{\it $Y \subset X$ に対する $S$ 上の対数極ド・ラーム複体}


\medskip\noindent
{\bf 局所コホモロジー}

\medskip

\noindent
\ref{local-cohomology-definition-cd} を参照：
{\it $I$ の $A$ におけるコホモロジー次元}


\noindent
\ref{local-cohomology-definition-depth-complex} を参照：
{\it $I$-深さ},
{\it 深さ}


\medskip\noindent
{\bf 代数幾何と形式幾何}

\medskip

\noindent
\ref{algebraization-definition-derived-complete} を参照：
{\it $\mathcal{I}$ に関して導来完備}


\noindent
\ref{algebraization-definition-algebraizable} を参照：
{\it $(\mathcal{F}_n)$ は $X$ へ延長される}


\noindent
\ref{algebraization-definition-canonically-algebraizable} を参照：
{\it $(\mathcal{F}_n)$ は標準的に $X$ へ延長される}


\noindent
\ref{algebraization-definition-s-d-inequalities} を参照：
{\it $(\mathcal{F}_n)$ は $(a, b)$-不等式を満たす},
{\it $(\mathcal{F}_n)$ は狭義 $(a, b)$-不等式を満たす}


\medskip\noindent
{\bf 代数曲線}

\medskip

\noindent
\ref{curves-definition-normal-projective-model} を参照：
{\it $X$ の非特異射影モデル}


\noindent
\ref{curves-definition-linear-series} を参照：
{\it 次数 $d$、次元 $r$ の線形系},
{\it $\mathfrak g^r_d$}


\noindent
\ref{curves-definition-genus} を参照：
{\it 種数}


\noindent
\ref{curves-definition-geometric-genus} を参照：
{\it 幾何種数}


\noindent
\ref{curves-definition-multicross} を参照：
{\it 多重交叉特異点},
{\it 節点},
{\it 通常二重点},
{\it 節点特異点を定める}


\noindent
\ref{curves-definition-nodal} を参照：
{\it 節点},
{\it 通常二重点},
{\it 節点特異点を定める},
{\it $X$ の特異点は高々節点である}


\noindent
\ref{curves-definition-split-node} を参照：
{\it 分裂節点}


\noindent
\ref{curves-definition-nodal-family} を参照：
{\it 相対次元 $1$ の高々節点的}


\medskip\noindent
{\bf 曲面の特異点解消}

\medskip

\noindent
\ref{resolve-definition-normalized-blowup} を参照：
{\it $x$ における $X$ の正規化ブローアップ}


\noindent
\ref{resolve-definition-reduce-to-rational} を参照：
{\it 有理特異点を定める},
{\it $A$ について有理特異点への還元が可能である}


\noindent
\ref{resolve-definition-resolution} を参照：
{\it 特異点解消}


\noindent
\ref{resolve-definition-resolution-surface} を参照：
{\it 正規化ブローアップによる特異点解消}


\medskip\noindent
{\bf 半安定還元}

\medskip

\noindent
\ref{models-definition-type} を参照：
{\it 数値型}


\noindent
\ref{models-definition-type-equivalent} を参照：
{\it 同値な型}


\noindent
\ref{models-definition-genus} を参照：
{\it 種数 $g$ の数値型}


\noindent
\ref{models-definition-type-minus-one} を参照：
{\it $(-1)$-指数}


\noindent
\ref{models-definition-top-genus} を参照：
{\it $T$ の位相種数}


\noindent
\ref{models-definition-type-minimal} を参照：
{\it 極小}


\noindent
\ref{models-definition-type-minus-two} を参照：
{\it $(-2)$-指数}


\noindent
\ref{models-definition-picard-group} を参照：
{\it $T$ の Picard 群}


\noindent
\ref{models-definition-minimal-model} を参照：
{\it 極小モデル}


\noindent
\ref{models-definition-numerical-type-model} を参照：
{\it $X$ に付随する数値型}


\noindent
\ref{models-definition-semistable} を参照：
{\it 半安定還元}


\noindent
\ref{models-definition-good} を参照：
{\it 良い還元}


\medskip\noindent
{\bf 関手と射}

\medskip

\medskip\noindent
{\bf 代数多様体の導来圏}

\medskip

\noindent
\ref{equiv-definition-Serre-functor} を参照：
{\it Serre 関手が存在する},
{\it Serre 関手}


\noindent
\ref{equiv-definition-fourier-mukai-functor} を参照：
{\it Fourier--Mukai 関手},
{\it Fourier--Mukai 核}


\noindent
\ref{equiv-definition-siblings} を参照：
{\it 同胞関手},
{\it 同胞関手}


\noindent
\ref{equiv-definition-siblings-geometric} を参照：
{\it 同胞関手},
{\it 同胞関手}


\noindent
\ref{equiv-definition-relative-equivalence-kernel} を参照：
{\it $X$ から $Y$ への $S$ 上の相対同値の Fourier--Mukai 核}


\noindent
\ref{equiv-definition-derived-equivalent} を参照：
{\it 導来同値}


\medskip\noindent
{\bf スキームの基本群}

\medskip

\noindent
\ref{pione-definition-G-set-continuous} を参照：
{\it $G$-集合},
{\it 離散 $G$-集合},
{\it $G$-集合の射},
{\it $G\textit{-Sets}$}


\noindent
\ref{pione-definition-galois-category} を参照：
{\it ガロア圏}


\noindent
\ref{pione-definition-fundamental-group} を参照：
{\it 基本群},
{\it 基点}


\medskip\noindent
{\bf エタールコホモロジー}

\medskip

\noindent
\ref{etale-cohomology-definition-etale-covering-initial} を参照：
{\it エタール被覆}


\noindent
\ref{etale-cohomology-definition-presheaf} を参照：
{\it 集合の前層},
{\it アーベル群の前層}


\noindent
\ref{etale-cohomology-definition-family-morphisms-fixed-target} を参照：
{\it 終域を固定した射の族}


\noindent
\ref{etale-cohomology-definition-site} を参照：
{\it サイト},
{\it 被覆族}


\noindent
\ref{etale-cohomology-definition-sheaf} を参照：
{\it 分離前層},
{\it 層}


\noindent
\ref{etale-cohomology-definition-category-sheaves} を参照：
{\it 集合の層の圏},
{\it アーベル群の層}


\noindent
\ref{etale-cohomology-definition-0-cech} を参照：
{\it 第零 {\v C}ech コホモロジー群}


\noindent
\ref{etale-cohomology-definition-fpqc-covering} を参照：
{\it fpqc 被覆}


\noindent
\ref{etale-cohomology-definition-sheaf-property-fpqc} を参照：
{\it fpqc 位相に対する層条件を満たす}


\noindent
\ref{etale-cohomology-definition-descent-datum} を参照：
{\it 降下データ},
{\it コサイクル条件},
{\it 有効}


\noindent
\ref{etale-cohomology-definition-descent-datum-modules} を参照：
{\it 降下データ}


\noindent
\ref{etale-cohomology-definition-effective-modules} を参照：
{\it 有効}


\noindent
\ref{etale-cohomology-definition-ringed-site} を参照：
{\it 環付きサイト},
{\it 準連接}


\noindent
\ref{etale-cohomology-definition-cech-complex} を参照：
{\it {\v C}ech 複体},
{\it {\v C}ech コホモロジー群}


\noindent
\ref{etale-cohomology-definition-free-abelian-presheaf} を参照：
{\it $\mathcal{G}$ 上の自由アーベル前層}


\noindent
\ref{etale-cohomology-definition-tau-covering} を参照：
{\it $\tau$-被覆}


\noindent
\ref{etale-cohomology-definition-tau-site} を参照：
{\it $S$ の大 $\tau$-サイト},
{\it $S$ の小 $\tau$-サイト}


\noindent
\ref{etale-cohomology-definition-standard-tau} を参照：
{\it 標準 $\tau$-被覆}


\noindent
\ref{etale-cohomology-definition-etale-topos} を参照：
{\it エタールトポス},
{\it 小エタールトポス},
{\it Zariski トポス},
{\it 小 Zariski トポス},
{\it 大 $\tau$-トポス}


\noindent
\ref{etale-cohomology-definition-additive-sheaf} を参照：
{\it 定数層}


\noindent
\ref{etale-cohomology-definition-structure-sheaf} を参照：
{\it 構造層}


\noindent
\ref{etale-cohomology-definition-etale-morphism} を参照：
{\it エタール}


\noindent
\ref{etale-cohomology-definition-standard-etale} を参照：
{\it 標準エタール}


\noindent
\ref{etale-cohomology-definition-etale-covering} を参照：
{\it エタール被覆}


\noindent
\ref{etale-cohomology-definition-big-etale-site} を参照：
{\it $S$ 上の大エタール・サイト},
{\it $S$ 上の小エタール・サイト},
{\it 大},
{\it 小 Zariski サイト}


\noindent
\ref{etale-cohomology-definition-geometric-point} を参照：
{\it 幾何学的点},
{\it 上にある},
{\it エタール近傍},
{\it エタール近傍の射}


\noindent
\ref{etale-cohomology-definition-stalk} を参照：
{\it 茎}


\noindent
\ref{etale-cohomology-definition-support} を参照：
{\it $\mathcal{F}$ の台},
{\it $\sigma$ の台}


\noindent
\ref{etale-cohomology-definition-henselian} を参照：
{\it ヘンゼル的}


\noindent
\ref{etale-cohomology-definition-strictly-henselian} を参照：
{\it 強ヘンゼル的}


\noindent
\ref{etale-cohomology-definition-etale-local-rings} を参照：
{\it $\overline{s}$ における $S$ のエタール局所環},
{\it $\mathcal{O}_{S, s}$ の強ヘンゼル化},
{\it $\mathcal{O}_{S, s}$ のヘンゼル化},
{\it $\overline{s}$ における $S$ の強ヘンゼル化},
{\it $s$ における $S$ のヘンゼル化}


\noindent
\ref{etale-cohomology-definition-direct-image-presheaf} を参照：
{\it 順像},
{\it 順像}


\noindent
\ref{etale-cohomology-definition-direct-image-sheaf} を参照：
{\it 順像},
{\it 順像}


\noindent
\ref{etale-cohomology-definition-higher-direct-images} を参照：
{\it 高次順像}


\noindent
\ref{etale-cohomology-definition-inverse-image} を参照：
{\it 逆像},
{\it 引き戻し}


\noindent
\ref{etale-cohomology-definition-inverse-system-sheaves} を参照：
{\it $(X_i, f_{i'i})$ 上の層の系 $(\mathcal{F}_i, \varphi_{i'i})$}


\noindent
\ref{etale-cohomology-definition-algebraic-geometric-point} を参照：
{\it 絶対ガロア群},
{\it 代数的}


\noindent
\ref{etale-cohomology-definition-G-module-continuous} を参照：
{\it $G$-加群},
{\it 離散 $G$-加群},
{\it $G$-加群の射},
{\it $R\text{-}G$-加群},
{\it $R\text{-}G$-加群の射}


\noindent
\ref{etale-cohomology-definition-galois-cohomology} を参照：
{\it 連続群コホモロジー群},
{\it 群コホモロジー群},
{\it ガロアコホモロジー群},
{\it 係数 $M$ の $K$ のガロアコホモロジー群}


\noindent
\ref{etale-cohomology-definition-brauer-equivalent} を参照：
{\it 相似},
{\it 同値}


\noindent
\ref{etale-cohomology-definition-brauer-group} を参照：
{\it ブラウアー群}


\noindent
\ref{etale-cohomology-definition-finite-locally-constant} を参照：
{\it 値 $E$ をもつ定数層},
{\it 定数層},
{\it 局所定数},
{\it 有限局所定数},
{\it 値 $A$ の定数層},
{\it 定数層},
{\it 局所定数},
{\it 有限局所定数},
{\it 値 $M$ をもつ定数層},
{\it 定数層},
{\it 局所定数}


\noindent
\ref{etale-cohomology-definition-trace-map} を参照：
{\it トレース}


\noindent
\ref{etale-cohomology-definition-Cr} を参照：
{\it $C_r$},
{\it 非自明解}


\noindent
\ref{etale-cohomology-definition-variety} を参照：
{\it 多様体},
{\it 曲線}


\noindent
\ref{etale-cohomology-definition-extension-zero} を参照：
{\it 零による延長},
{\it 零による延長}


\noindent
\ref{etale-cohomology-definition-constructible} を参照：
{\it 構成可能},
{\it 構成可能},
{\it 構成可能}


\noindent
\ref{etale-cohomology-definition-c} を参照：
{\it $D_c(X_\etale, \Lambda)$}


\noindent
\ref{etale-cohomology-definition-ctf} を参照：
{\it $D_{ctf}(X_\etale, \Lambda)$}


\noindent
\ref{etale-cohomology-definition-locally-acyclic} を参照：
{\it $\overline{x}$ において $K$ に関して局所非輪状},
{\it $K$ に関して局所非輪状},
{\it $K$ に関して普遍局所非輪状},
{\it 局所非輪状},
{\it 普遍局所非輪状}


\noindent
\ref{etale-cohomology-definition-cd} を参照：
{\it $X$ のコホモロジー次元}


\noindent
\ref{etale-cohomology-definition-cd-f} を参照：
{\it $f$ のコホモロジー次元}


\medskip\noindent
{\bf クリスタリンコホモロジー}

\medskip

\noindent
\ref{crystalline-definition-divided-power-envelope} を参照：
{\it $(A, I, \gamma)$ に相対的な $B$ における $J$ の分割冪包絡}


\noindent
\ref{crystalline-definition-compatible} を参照：
{\it $\delta$ は $\gamma$ と両立する}


\noindent
\ref{crystalline-definition-affine-thickening} を参照：
{\it 分割冪肥厚化},
{\it 分割冪肥厚化の準同型}


\noindent
\ref{crystalline-definition-derivation} を参照：
{\it 分割冪 $A$-導分}


\noindent
\ref{crystalline-definition-divided-power-structure} を参照：
{\it 分割冪構造 $\gamma$}


\noindent
\ref{crystalline-definition-divided-power-scheme} を参照：
{\it 分割冪スキーム},
{\it 分割冪スキームの射}


\noindent
\ref{crystalline-definition-divided-power-thickening} を参照：
{\it 分割冪肥厚化}


\noindent
\ref{crystalline-definition-divided-power-thickening-X} を参照：
{\it $(S, \mathcal{I}, \gamma)$ に相対的な $X$ の分割冪肥厚化},
{\it $(S, \mathcal{I}, \gamma)$ に相対的な $X$ の分割冪肥厚化の射}


\noindent
\ref{crystalline-definition-big-crystalline-site} を参照：
{\it Zariski、エタール、滑らか、シントミック、または fppf 被覆},
{\it 大結晶サイト}


\noindent
\ref{crystalline-definition-crystalline-site} を参照：
{\it 結晶サイト}


\noindent
\ref{crystalline-definition-modules} を参照：
{\it 局所準連接},
{\it 準連接},
{\it $\mathcal{O}_{X/S}$-加群のクリスタル}


\noindent
\ref{crystalline-definition-crystal-quasi-coherent-modules} を参照：
{\it 準連接加群のクリスタル},
{\it 有限局所自由加群のクリスタル}


\noindent
\ref{crystalline-definition-global-derivation} を参照：
{\it $S$-導分 $D : \mathcal{O}_{X/S} \to \mathcal{F}$}


\noindent
\ref{crystalline-definition-F-crystal} を参照：
{\it $X/S$ 上の（$\sigma$ に関する）$F$-クリスタル},
{\it 非退化}


\medskip\noindent
{\bf pro-エタールコホモロジー}

\medskip

\noindent
\ref{proetale-definition-w-local} を参照：
{\it w-局所的},
{\it w-局所的}


\noindent
\ref{proetale-definition-local-isomorphism} を参照：
{\it 局所同型},
{\it 局所環を同一視する}


\noindent
\ref{proetale-definition-ind-zariski} を参照：
{\it ind-Zariski}


\noindent
\ref{proetale-definition-ind-etale} を参照：
{\it ind-エタール}


\noindent
\ref{proetale-definition-w-contractible} を参照：
{\it w-可縮}


\noindent
\ref{proetale-definition-fpqc-covering} を参照：
{\it $T$ の pro-エタール被覆}


\noindent
\ref{proetale-definition-standard-proetale} を参照：
{\it 標準 pro-エタール被覆}


\noindent
\ref{proetale-definition-big-proetale-site} を参照：
{\it 大 pro-エタール・サイト}


\noindent
\ref{proetale-definition-big-small-proetale} を参照：
{\it $S$ の大 pro-エタール・サイト},
{\it $S$ の小 pro-エタール・サイト},
{\it $S$ の大アフィン pro-エタール・サイト}


\noindent
\ref{proetale-definition-restriction-small-proetale} を参照：
{\it 小 pro-エタール・サイトへの制限}


\noindent
\ref{proetale-definition-extension-zero} を参照：
{\it 零による延長},
{\it 零による延長}


\noindent
\ref{proetale-definition-constructible} を参照：
{\it 構成可能}


\noindent
\ref{proetale-definition-adic} を参照：
{\it 構成可能 $\Lambda$-層},
{\it 滑らか},
{\it アディック滑らか},
{\it アディック構成可能}


\noindent
\ref{proetale-definition-Dbc} を参照：
{\it 構成可能}


\noindent
\ref{proetale-definition-adic-constructible} を参照：
{\it アディック滑らか},
{\it アディック構成可能}


\medskip\noindent
{\bf 相対サイクル}

\medskip

\noindent
\ref{relative-cycles-definition-relative-cycles} を参照：
{\it $X/S$ 上の相対 $r$-サイクル}


\noindent
\ref{relative-cycles-definition-equidimensional} を参照：
{\it 等次元}


\noindent
\ref{relative-cycles-definition-effective} を参照：
{\it 有効}


\noindent
\ref{relative-cycles-definition-proper} を参照：
{\it 固有相対サイクル}


\medskip\noindent
{\bf エタールコホモロジーの補足}

\medskip

\noindent
\ref{more-etale-definition-f-shriek-separated} を参照：
{\it コンパクト台付き直像}


\noindent
\ref{more-etale-definition-compact-support} を参照：
{\it コンパクト台付き切断}


\noindent
\ref{more-etale-definition-f-shriek-lqf} を参照：
{\it コンパクト台付き直像}


\noindent
\ref{more-etale-definition-cohomology-compact-support} を参照：
{\it コンパクト台をもつ $K$ のコホモロジー},
{\it $K$ のコンパクト台付きコホモロジー}


\medskip\noindent
{\bf トレース公式}

\medskip

\noindent
\ref{trace-definition-geometric-frobenius} を参照：
{\it 幾何学的フロベニウス}


\noindent
\ref{trace-definition-arithmetic-frobenius} を参照：
{\it 算術的フロベニウス}


\noindent
\ref{trace-definition-geometric-frobenius-on-stalk} を参照：
{\it 幾何学的フロベニウス}


\noindent
\ref{trace-definition-trace} を参照：
{\it トレース}


\noindent
\ref{trace-definition-derived-functor} を参照：
{\it $F$ の全右導来関手},
{\it $G$ の全左導来関手}


\noindent
\ref{trace-definition-filtered} を参照：
{\it フィルター付き入射的},
{\it 射影的},
{\it フィルター付き擬同型}


\noindent
\ref{trace-definition-filtered-derived-functors} を参照：
{\it フィルター付き導来関手}


\noindent
\ref{trace-definition-perfect} を参照：
{\it 完全}


\noindent
\ref{trace-definition-finite-tor-dimension} を参照：
{\it 有限 $\text{Tor}$ 次元}


\noindent
\ref{trace-definition-global-lefschetz-number} を参照：
{\it 大域 Lefschetz 数}


\noindent
\ref{trace-definition-local-lefschetz-number} を参照：
{\it 局所 Lefschetz 数}


\noindent
\ref{trace-definition-trace-G} を参照：
{\it $P$ 上の $f$ の $G$-トレース}


\noindent
\ref{trace-definition-l-adic-sheaf} を参照：
{\it $\mathbf{Z}_\ell$-層},
{\it $\ell$-進層},
{\it 滑らか},
{\it 射}


\noindent
\ref{trace-definition-torsion-l-adic-sheaf} を参照：
{\it ねじれ},
{\it 茎}


\noindent
\ref{trace-definition-cohomology-l-adic} を参照：
{\it $\ell$-進コホモロジー}


\noindent
\ref{trace-definition-L-function-finite-ring} を参照：
{\it $\mathcal{F}$ の $L$-関数}


\noindent
\ref{trace-definition-L-function-l-adic} を参照：
{\it $\mathcal{F}$ の $L$-関数}


\noindent
\ref{trace-definition-open} を参照：
{\it 開}


\noindent
\ref{trace-definition-unramified} を参照：
{\it $\text{GL}_2(\mathbf{A})$ 上の $\Lambda$ に値を取る 不分岐尖点形式}


\medskip\noindent
{\bf 代数空間}

\medskip

\noindent
\ref{spaces-definition-relative-representable-property} を参照：
{\it 性質 $\mathcal{P}$}


\noindent
\ref{spaces-definition-algebraic-space} を参照：
{\it $S$ 上の代数空間}


\noindent
\ref{spaces-definition-morphism-algebraic-spaces} を参照：
{\it $S$ 上の代数空間の射 $f : F \to F'$}


\noindent
\ref{spaces-definition-etale-equivalence-relation} を参照：
{\it エタール同値関係}


\noindent
\ref{spaces-definition-presentation} を参照：
{\it 表示}


\noindent
\ref{spaces-definition-immersion} を参照：
{\it 開埋め込み},
{\it 開部分空間},
{\it 閉埋め込み},
{\it 閉部分空間},
{\it 埋め込み},
{\it 局所閉部分空間}


\noindent
\ref{spaces-definition-Zariski-open-covering} を参照：
{\it Zariski 被覆}


\noindent
\ref{spaces-definition-small-Zariski-site} を参照：
{\it 小 Zariski サイト $F_{Zar}$}


\noindent
\ref{spaces-definition-separated} を参照：
{\it $S$ 上分離},
{\it $S$ 上局所分離的},
{\it $S$ 上準分離},
{\it $S$ 上 Zariski 局所準分離的}


\noindent
\ref{spaces-definition-quotient} を参照：
{\it 自由に作用する},
{\it $G$ による $U$ の商}


\noindent
\ref{spaces-definition-base-change} を参照：
{\it $F'$ の $S$ への基底変換},
{\it $S'$ 上の代数空間とみなしたもの}


\medskip\noindent
{\bf 代数空間の性質}

\medskip

\noindent
\ref{spaces-properties-definition-separated} を参照：
{\it 分離的},
{\it 局所分離的},
{\it 準分離},
{\it Zariski 局所準分離的},
{\it 分離的},
{\it 局所分離的},
{\it 準分離},
{\it Zariski 局所準分離的}


\noindent
\ref{spaces-properties-definition-points} を参照：
{\it 点}


\noindent
\ref{spaces-properties-definition-topological-space} を参照：
{\it 位相空間}


\noindent
\ref{spaces-properties-definition-quasi-compact} を参照：
{\it 準コンパクト}


\noindent
\ref{spaces-properties-definition-type-property} を参照：
{\it 性質 $\mathcal{P}$ をもつ}


\noindent
\ref{spaces-properties-definition-property-at-point} を参照：
{\it $x$ において性質 $\mathcal{P}$ をもつ}


\noindent
\ref{spaces-properties-definition-locally-constructible} を参照：
{\it エタール局所構成可能}


\noindent
\ref{spaces-properties-definition-dimension-at-point} を参照：
{\it $x$ における $X$ の次元}


\noindent
\ref{spaces-properties-definition-dimension} を参照：
{\it 次元}


\noindent
\ref{spaces-properties-definition-dimension-local-ring} を参照：
{\it $x$ における $X$ の局所環の次元},
{\it $x$ は $X$ 上の余次元 $d$ の点である}


\noindent
\ref{spaces-properties-definition-reduced-induced-space} を参照：
{\it $Z$ 上の代数空間構造},
{\it 被約誘導代数空間構造},
{\it $X$ の被約化 $X_{red}$}


\noindent
\ref{spaces-properties-definition-etale} を参照：
{\it エタール}


\noindent
\ref{spaces-properties-definition-etale-site} を参照：
{\it 小エタールサイト $X_\etale$}


\noindent
\ref{spaces-properties-definition-spaces-etale-site} を参照：
{\it $X_{spaces, \etale}$}


\noindent
\ref{spaces-properties-definition-affine-etale-site} を参照：
{\it $X_{affine, \etale}$}


\noindent
\ref{spaces-properties-definition-etale-topos} を参照：
{\it エタールトポス},
{\it 小エタールトポス}


\noindent
\ref{spaces-properties-definition-f-map} を参照：
{\it $f$-写像 $\varphi : \mathcal{G} \to \mathcal{F}$}


\noindent
\ref{spaces-properties-definition-geometric-point} を参照：
{\it 幾何学的点},
{\it $x$ の上にある幾何学的点}


\noindent
\ref{spaces-properties-definition-etale-neighbourhood} を参照：
{\it エタール近傍},
{\it エタール近傍の射}


\noindent
\ref{spaces-properties-definition-stalk} を参照：
{\it 茎}


\noindent
\ref{spaces-properties-definition-support} を参照：
{\it $\mathcal{F}$ の台},
{\it $\sigma$ の台}


\noindent
\ref{spaces-properties-definition-structure-sheaf} を参照：
{\it 構造層}


\noindent
\ref{spaces-properties-definition-etale-local-rings} を参照：
{\it $\overline{x}$ における $X$ のエタール局所環},
{\it $\overline{x}$ における $X$ の強ヘンゼル化}


\noindent
\ref{spaces-properties-definition-unibranch} を参照：
{\it $x$ において幾何的単枝},
{\it 幾何的単枝}


\noindent
\ref{spaces-properties-definition-number-of-geometric-branches} を参照：
{\it $x$ における $X$ の幾何的枝数}


\noindent
\ref{spaces-properties-definition-noetherian} を参照：
{\it Noether}


\noindent
\ref{spaces-properties-definition-regular-at-point} を参照：
{\it $X$ は $x$ において正則}


\noindent
\ref{spaces-properties-definition-quasi-coherent} を参照：
{\it 準連接}


\noindent
\ref{spaces-properties-definition-locally-projective} を参照：
{\it 局所射影的}


\medskip\noindent
{\bf 代数空間の射}

\medskip

\noindent
\ref{spaces-morphisms-definition-separated} を参照：
{\it 分離的},
{\it 局所分離的},
{\it 準分離}


\noindent
\ref{spaces-morphisms-definition-surjective} を参照：
{\it 全射}


\noindent
\ref{spaces-morphisms-definition-open} を参照：
{\it 開},
{\it 普遍開}


\noindent
\ref{spaces-morphisms-definition-submersive} を参照：
{\it 商位相的},
{\it 普遍的に商写像的}


\noindent
\ref{spaces-morphisms-definition-quasi-compact} を参照：
{\it 準コンパクト}


\noindent
\ref{spaces-morphisms-definition-closed} を参照：
{\it 閉},
{\it 普遍閉}


\noindent
\ref{spaces-morphisms-definition-monomorphism} を参照：
{\it モノ射}


\noindent
\ref{spaces-morphisms-definition-inverse-image-closed-subspace} を参照：
{\it 閉部分空間 $Z$ の逆像 $f^{-1}(Z)$}


\noindent
\ref{spaces-morphisms-definition-scheme-theoretic-intersection-union} を参照：
{\it スキーム論的交叉},
{\it スキーム論的和}


\noindent
\ref{spaces-morphisms-definition-scheme-theoretic-support} を参照：
{\it $\mathcal{F}$ のスキーム論的台}


\noindent
\ref{spaces-morphisms-definition-scheme-theoretic-image} を参照：
{\it スキーム論的像}


\noindent
\ref{spaces-morphisms-definition-scheme-theoretically-dense} を参照：
{\it $X$ における $U$ のスキーム論的閉包},
{\it $X$ においてスキーム論的に稠密}


\noindent
\ref{spaces-morphisms-definition-dominant} を参照：
{\it 支配的}


\noindent
\ref{spaces-morphisms-definition-universally-injective} を参照：
{\it 普遍単射}


\noindent
\ref{spaces-morphisms-definition-affine} を参照：
{\it アフィン}


\noindent
\ref{spaces-morphisms-definition-relative-spec} を参照：
{\it $X$ 上の $\mathcal{A}$ の相対スペクトル},
{\it $X$ 上の $\mathcal{A}$ のスペクトル}


\noindent
\ref{spaces-morphisms-definition-quasi-affine} を参照：
{\it 準アフィン}


\noindent
\ref{spaces-morphisms-definition-P} を参照：
{\it 性質 $\mathcal{P}$ をもつ}


\noindent
\ref{spaces-morphisms-definition-P-at-point} を参照：
{\it $x$ において性質 $\mathcal{Q}$ をもつ}


\noindent
\ref{spaces-morphisms-definition-locally-finite-type} を参照：
{\it 局所有限型},
{\it $x$ において有限型},
{\it 有限型}


\noindent
\ref{spaces-morphisms-definition-finite-type-point} を参照：
{\it 有限型点}


\noindent
\ref{spaces-morphisms-definition-locally-quasi-finite} を参照：
{\it 局所準有限},
{\it $x$ において準有限},
{\it 準有限}


\noindent
\ref{spaces-morphisms-definition-locally-finite-presentation} を参照：
{\it 局所有限表示},
{\it $x$ において有限表示},
{\it 有限表示}


\noindent
\ref{spaces-morphisms-definition-flat} を参照：
{\it 平坦},
{\it $x$ で平坦}


\noindent
\ref{spaces-morphisms-definition-flat-module} を参照：
{\it $x$ において $Y$ 上平坦},
{\it $Y$ 上平坦}


\noindent
\ref{spaces-morphisms-definition-dimension-fibre} を参照：
{\it $x$ における $f$ のファイバーの局所環の次元},
{\it $x/f(x)$ の超越次数},
{\it $f$ は $x$ において相対次元 $d$ をもつ}


\noindent
\ref{spaces-morphisms-definition-relative-dimension} を参照：
{\it 相対次元 $\leq d$},
{\it 相対次元 $d$}


\noindent
\ref{spaces-morphisms-definition-syntomic} を参照：
{\it シントミック},
{\it $x$ においてシントミック}


\noindent
\ref{spaces-morphisms-definition-smooth} を参照：
{\it 滑らか},
{\it $x$ において滑らか}


\noindent
\ref{spaces-morphisms-definition-unramified} を参照：
{\it 不分岐},
{\it $x$ において不分岐},
{\it G-不分岐},
{\it $x$ において G-非分岐}


\noindent
\ref{spaces-morphisms-definition-etale} を参照：
{\it $x$ においてエタール}


\noindent
\ref{spaces-morphisms-definition-proper} を参照：
{\it 固有}


\noindent
\ref{spaces-morphisms-definition-valuative-criterion} を参照：
{\it 付値判定法の一意性部分を満たす},
{\it 付値判定法の存在部分を満たす},
{\it 付値判定法を満たす}


\noindent
\ref{spaces-morphisms-definition-integral} を参照：
{\it 整},
{\it 有限}


\noindent
\ref{spaces-morphisms-definition-finite-locally-free} を参照：
{\it 有限局所自由},
{\it 階数},
{\it 次数}


\noindent
\ref{spaces-morphisms-definition-rational-map} を参照：
{\it 同値},
{\it $X$ から $Y$ への有理写像},
{\it $X$ から $Y$ への $B$-有理写像}


\noindent
\ref{spaces-morphisms-definition-rational-function} を参照：
{\it $X$ 上の有理関数}


\noindent
\ref{spaces-morphisms-definition-ring-of-rational-functions} を参照：
{\it $X$ 上の有理関数環}


\noindent
\ref{spaces-morphisms-definition-domain-of-definition} を参照：
{\it 点 $x \in |X|$ において定義される},
{\it 定義域}


\noindent
\ref{spaces-morphisms-definition-dominant-rational} を参照：
{\it 支配的}


\noindent
\ref{spaces-morphisms-definition-birational-spaces} を参照：
{\it 双有理}


\noindent
\ref{spaces-morphisms-definition-integral-closure} を参照：
{\it $\mathcal{O}_X$ の整閉包 （$\mathcal{A}$ におけるもの）}


\noindent
\ref{spaces-morphisms-definition-normalization-X-in-Y} を参照：
{\it $X$ の正規化 （$Y$ におけるもの）}


\noindent
\ref{spaces-morphisms-definition-normalization} を参照：
{\it 正規化}


\noindent
\ref{spaces-morphisms-definition-universal-homeomorphism} を参照：
{\it 普遍同相射}


\medskip\noindent
{\bf 適正な代数空間}

\medskip

\noindent
\ref{decent-spaces-definition-universally-bounded} を参照：
{\it $f$ のファイバーは普遍的に有界である}


\noindent
\ref{decent-spaces-definition-very-reasonable} を参照：
{\it 良好},
{\it 妥当},
{\it 非常に妥当}


\noindent
\ref{decent-spaces-definition-residue-field} を参照：
{\it $x$ における $X$ の剰余体}


\noindent
\ref{decent-spaces-definition-elemenary-etale-neighbourhood} を参照：
{\it 初等エタール近傍},
{\it 基本エタール近傍の射}


\noindent
\ref{decent-spaces-definition-henselian-local-ring} を参照：
{\it $x$ における $X$ のヘンゼル局所環}


\noindent
\ref{decent-spaces-definition-residual-space} を参照：
{\it $x$ における $X$ の残余空間}


\noindent
\ref{decent-spaces-definition-relative-conditions} を参照：
{\it 性質 $(\beta)$ をもつ},
{\it 性質 $(\beta)$ をもつ},
{\it 良好},
{\it 妥当},
{\it 非常に妥当}


\noindent
\ref{decent-spaces-definition-birational} を参照：
{\it 双有理}


\noindent
\ref{decent-spaces-definition-unibranch} を参照：
{\it $x$ において単枝},
{\it 単枝}


\noindent
\ref{decent-spaces-definition-number-of-branches} を参照：
{\it $X$ の $x$ における枝数}


\noindent
\ref{decent-spaces-definition-catenary} を参照：
{\it 鎖状}


\noindent
\ref{decent-spaces-definition-universally-catenary} を参照：
{\it 普遍鎖状}


\medskip\noindent
{\bf 代数空間のコホモロジー}

\medskip

\noindent
\ref{spaces-cohomology-definition-alternating-cech-complex} を参照：
{\it 交代 {\v C}ech 複体}


\noindent
\ref{spaces-cohomology-definition-coherent} を参照：
{\it 連接}


\medskip\noindent
{\bf 代数空間の極限}

\medskip

\noindent
\ref{spaces-limits-definition-locally-finite-presentation} を参照：
{\it 極限を保つ},
{\it 局所有限表示},
{\it $S$ 上局所有限表示},
{\it 極限を保つ},
{\it 局所有限表示},
{\it 相対的に極限を保つ}


\noindent
\ref{spaces-limits-definition-subsheaf-sections-annihilated-by-ideal} を参照：
{\it $\mathcal{I}$ により零化される切断の部分層}


\noindent
\ref{spaces-limits-definition-subsheaf-sections-supported-on-closed} を参照：
{\it $T$ 上に台を持つ切断の部分層}


\medskip\noindent
{\bf 代数空間上の除数}

\medskip

\noindent
\ref{spaces-divisors-definition-weakly-associated} を参照：
{\it 弱随伴する},
{\it $X$ の弱随伴点},
{\it $x$ は $\mathcal{F}$ に付随する},
{\it $x$ は $X$ の付随点}


\noindent
\ref{spaces-divisors-definition-locally-Noetherian-fibre} を参照：
{\it $f$ の $y$ 上のファイバーは局所 Noether である},
{\it $f$ のファイバーは局所 Noether である}


\noindent
\ref{spaces-divisors-definition-relative-weak-assassin} を参照：
{\it $\mathcal{F}$ の $X$ における $Y$ 上の相対弱随伴点集合}


\noindent
\ref{spaces-divisors-definition-effective-Cartier-divisor} を参照：
{\it 局所主閉部分空間},
{\it 有効 Cartier 因子}


\noindent
\ref{spaces-divisors-definition-sum-effective-Cartier-divisors} を参照：
{\it 有効 Cartier 因子 $D_1$ と $D_2$ の和}


\noindent
\ref{spaces-divisors-definition-pullback-effective-Cartier-divisor} を参照：
{\it $D$ の $f$ による引き戻しは定義される},
{\it 有効 Cartier 因子の引き戻し}


\noindent
\ref{spaces-divisors-definition-invertible-sheaf-effective-Cartier-divisor} を参照：
{\it 可逆層 $\mathcal{O}_X(D)$（$D$ に伴うもの）}


\noindent
\ref{spaces-divisors-definition-regular-section} を参照：
{\it 正則切断}


\noindent
\ref{spaces-divisors-definition-zero-scheme-s} を参照：
{\it 零点スキーム}


\noindent
\ref{spaces-divisors-definition-relative-effective-Cartier-divisor} を参照：
{\it 相対有効 Cartier 因子}


\noindent
\ref{spaces-divisors-definition-sheaf-meromorphic-functions} を参照：
{\it $X$ 上の有理型関数の層},
{\it $\mathcal{K}_X$},
{\it 有理型関数}


\noindent
\ref{spaces-divisors-definition-meromorphic-section} を参照：
{\it $\mathcal{F}$ の有理型切断}


\noindent
\ref{spaces-divisors-definition-pullback-meromorphic-sections} を参照：
{\it $f$ に対して有理型関数の引き戻しが定義される}


\noindent
\ref{spaces-divisors-definition-regular-meromorphic-section} を参照：
{\it 正則}


\noindent
\ref{spaces-divisors-definition-relative-proj} を参照：
{\it $X$ 上の $\mathcal{A}$ の相対斉次スペクトル},
{\it $X$ 上の $\mathcal{A}$ の斉次スペクトル},
{\it $X$ 上の $\mathcal{A}$ の相対 Proj}


\noindent
\ref{spaces-divisors-definition-relatively-ample} を参照：
{\it 相対的に豊富},
{\it $f$-相対的に豊富},
{\it $X/Y$ 上豊富},
{\it $f$-豊富}


\noindent
\ref{spaces-divisors-definition-blow-up} を参照：
{\it $X$ の $Z$ に沿ったブローアップ},
{\it $X$ のイデアル層 $\mathcal{I}$ におけるブローアップ},
{\it 例外因子},
{\it 中心}


\noindent
\ref{spaces-divisors-definition-strict-transform} を参照：
{\it 狭義変換},
{\it 狭義変換}


\noindent
\ref{spaces-divisors-definition-admissible-blowup} を参照：
{\it $U$-許容ブローアップ}


\medskip\noindent
{\bf 体上の代数空間}

\medskip

\noindent
\ref{spaces-over-fields-definition-integral-algebraic-space} を参照：
{\it 整}


\noindent
\ref{spaces-over-fields-definition-function-field} を参照：
{\it 関数体},
{\it 有理関数体}


\noindent
\ref{spaces-over-fields-definition-degree} を参照：
{\it $X$ の $Y$ 上の次数}


\noindent
\ref{spaces-over-fields-definition-Weil-divisor} を参照：
{\it 素因子},
{\it Weil 因子}


\noindent
\ref{spaces-over-fields-definition-order-vanishing} を参照：
{\it $f$ の $Z$ に沿う 消滅次数}


\noindent
\ref{spaces-over-fields-definition-principal-divisor} を参照：
{\it $f$ に付随する主 Weil 因子}


\noindent
\ref{spaces-over-fields-definition-class-group} を参照：
{\it Weil 因子類群}


\noindent
\ref{spaces-over-fields-definition-order-vanishing-meromorphic} を参照：
{\it $s$ の $Z$ に沿う消滅次数}


\noindent
\ref{spaces-over-fields-definition-divisor-invertible-sheaf} を参照：
{\it $s$ に付随する Weil 因子},
{\it $\mathcal{L}$ に付随する Weil 因子類}


\noindent
\ref{spaces-over-fields-definition-modification} を参照：
{\it $X$ の修正}


\noindent
\ref{spaces-over-fields-definition-alteration} を参照：
{\it $X$ のオルタレーション}


\noindent
\ref{spaces-over-fields-definition-geometrically-reduced} を参照：
{\it $x$ で幾何学的被約},
{\it 幾何学的被約}


\noindent
\ref{spaces-over-fields-definition-geometrically-connected} を参照：
{\it 幾何学的連結}


\noindent
\ref{spaces-over-fields-definition-geometrically-irreducible} を参照：
{\it 幾何学的既約}


\noindent
\ref{spaces-over-fields-definition-geometrically-integral} を参照：
{\it 幾何学的整}


\noindent
\ref{spaces-over-fields-definition-euler-characteristic} を参照：
{\it $\mathcal{F}$ の Euler 標数}


\noindent
\ref{spaces-over-fields-definition-intersection-number} を参照：
{\it 交点数}


\medskip\noindent
{\bf 代数空間上の位相}

\medskip

\noindent
\ref{spaces-topologies-definition-zariski-covering} を参照：
{\it $X$ の Zariski 被覆}


\noindent
\ref{spaces-topologies-definition-etale-covering} を参照：
{\it $X$ のエタール被覆}


\noindent
\ref{spaces-topologies-definition-big-etale-site} を参照：
{\it $(\textit{Spaces}/S)_\etale$}


\noindent
\ref{spaces-topologies-definition-big-small-etale} を参照：
{\it $(\textit{Spaces}/X)_\etale$}


\noindent
\ref{spaces-topologies-definition-restriction-small-etale} を参照：
{\it 小エタール・サイトへの制限}


\noindent
\ref{spaces-topologies-definition-smooth-covering} を参照：
{\it $X$ の滑らかな被覆}


\noindent
\ref{spaces-topologies-definition-syntomic-covering} を参照：
{\it $X$ のシントミック被覆}


\noindent
\ref{spaces-topologies-definition-fppf-covering} を参照：
{\it $X$ の fppf 被覆}


\noindent
\ref{spaces-topologies-definition-big-fppf-site} を参照：
{\it $(\textit{Spaces}/S)_{fppf}$}


\noindent
\ref{spaces-topologies-definition-big-small-fppf} を参照：
{\it $(\textit{Spaces}/X)_{fppf}$}


\noindent
\ref{spaces-topologies-definition-ph-covering} を参照：
{\it $X$ の ph 被覆}


\noindent
\ref{spaces-topologies-definition-big-ph-site} を参照：
{\it $(\textit{Spaces}/S)_{ph}$}


\noindent
\ref{spaces-topologies-definition-big-small-ph} を参照：
{\it $(\textit{Spaces}/X)_{ph}$}


\noindent
\ref{spaces-topologies-definition-fpqc-covering} を参照：
{\it $X$ の fpqc 被覆}


\medskip\noindent
{\bf 降下と代数空間}

\medskip

\noindent
\ref{spaces-descent-definition-descent-datum-quasi-coherent} を参照：
{\it 準連接層の降下データ $(\mathcal{F}_i, \varphi_{ij})$},
{\it コサイクル条件},
{\it 降下データの射 $\psi : (\mathcal{F}_i, \varphi_{ij}) \to (\mathcal{F}'_i, \varphi'_{ij})$}


\noindent
\ref{spaces-descent-definition-descent-datum-effective-quasi-coherent} を参照：
{\it 自明な降下データ},
{\it 標準降下データ},
{\it 有効}


\noindent
\ref{spaces-descent-definition-property-morphisms-local} を参照：
{\it 基底上 $\tau$-局所的},
{\it 標的上 $\tau$-局所的},
{\it $\tau$-位相について基底上局所的}


\noindent
\ref{spaces-descent-definition-property-morphisms-local-source} を参照：
{\it 始域上 $\tau$-局所的},
{\it $\tau$-位相について始域上局所的}


\noindent
\ref{spaces-descent-definition-local-source-target} を参照：
{\it 始域と終域に関して滑らか局所的}


\noindent
\ref{spaces-descent-definition-etale-smooth-local-source-target} を参照：
{\it エタール平滑は始域と終域について局所的}


\noindent
\ref{spaces-descent-definition-descent-datum} を参照：
{\it $V/Y/X$ に対する降下データ},
{\it コサイクル条件},
{\it $Y \to X$ に関する降下データ},
{\it $Y \to X$ に関する降下データの射 $f : (V/Y, \varphi) \to (V'/Y, \varphi')$}


\noindent
\ref{spaces-descent-definition-descent-datum-for-family-of-morphisms} を参照：
{\it 族 $\{X_i \to X\}$ に関する降下データ $(V_i, \varphi_{ij})$},
{\it 降下データの射 $\psi : (V_i, \varphi_{ij}) \to (V'_i, \varphi'_{ij})$}


\noindent
\ref{spaces-descent-definition-pullback-functor} を参照：
{\it 引き戻し関手}


\noindent
\ref{spaces-descent-definition-pullback-functor-family} を参照：
{\it 引き戻し関手}


\noindent
\ref{spaces-descent-definition-effective} を参照：
{\it 自明な降下データ},
{\it 標準降下データ},
{\it 有効}


\noindent
\ref{spaces-descent-definition-effective-family} を参照：
{\it 標準降下データ},
{\it 有効}


\medskip\noindent
{\bf 代数空間の導来圏}

\medskip

\noindent
\ref{spaces-perfect-definition-supported-on} を参照：
{\it $T$ に台をもつ}


\noindent
\ref{spaces-perfect-definition-derived-quasi-coherent} を参照：
{\it 準連接なコホモロジー層をもつ $\mathcal{O}_X$-加群の導来圏}


\noindent
\ref{spaces-perfect-definition-proper-over-base} を参照：
{\it $T$ は $Y$ 上固有}


\noindent
\ref{spaces-perfect-definition-elementary-distinguished-square} を参照：
{\it 初等識別正方形}


\noindent
\ref{spaces-perfect-definition-approximation-holds} を参照：
{\it 三つ組に対して近似が成立する}


\noindent
\ref{spaces-perfect-definition-approximation} を参照：
{\it 完全複体による近似が 成立する}


\noindent
\ref{spaces-perfect-definition-tor-independent} を参照：
{\it $B$ 上 Tor 独立}


\noindent
\ref{spaces-perfect-definition-resolution-property} を参照：
{\it 分解性質}


\medskip\noindent
{\bf 代数空間の射の補足}

\medskip

\noindent
\ref{spaces-more-morphisms-definition-radicial} を参照：
{\it 純非分離}


\noindent
\ref{spaces-more-morphisms-definition-conormal-sheaf} を参照：
{\it $Z$ の $X$ における余法層 $\mathcal{C}_{Z/X}$},
{\it $i$ の余法層}


\noindent
\ref{spaces-more-morphisms-definition-conormal-algebra} を参照：
{\it 余法代数 $\mathcal{C}_{Z/X, *}$、すなわち $Z$ の $X$ における余法代数},
{\it $i$ の余法代数}


\noindent
\ref{spaces-more-morphisms-definition-normal-cone} を参照：
{\it 法錐 $C_ZX$},
{\it 法束}


\noindent
\ref{spaces-more-morphisms-definition-sheaf-differentials} を参照：
{\it $Y$ 上の $X$ の微分層 $\Omega_{X/Y}$},
{\it 普遍 $Y$-導分}


\noindent
\ref{spaces-more-morphisms-definition-thickening} を参照：
{\it 肥厚化},
{\it 一次肥厚化},
{\it 有限次肥厚化},
{\it $n$ 次肥厚化},
{\it 肥厚化の射},
{\it $B$ 上の肥厚化},
{\it $B$ 上の肥厚化の射}


\noindent
\ref{spaces-more-morphisms-definition-first-order-infinitesimal-neighbourhood} を参照：
{\it 一次無限小近傍},
{\it $n$ 次無限小近傍}


\noindent
\ref{spaces-more-morphisms-definition-formally-smooth-etale-unramified} を参照：
{\it 形式的に滑らか},
{\it 形式的エタール},
{\it 形式的不分岐}


\noindent
\ref{spaces-more-morphisms-definition-formally-unramified} を参照：
{\it 形式的不分岐}


\noindent
\ref{spaces-more-morphisms-definition-universal-thickening} を参照：
{\it 普遍一次肥厚化},
{\it $Z$ の $X$ 上の余法層}


\noindent
\ref{spaces-more-morphisms-definition-formally-etale} を参照：
{\it 形式的エタール}


\noindent
\ref{spaces-more-morphisms-definition-formally-smooth} を参照：
{\it 形式的に滑らか}


\noindent
\ref{spaces-more-morphisms-definition-netherlander} を参照：
{\it $f$ の素朴な余接複体}


\noindent
\ref{spaces-more-morphisms-definition-module-flat-on-fibre} を参照：
{\it $\mathcal{F}$ の $z$ 上のファイバーへの制限は、$x$ で、$Y$ の $z$ 上の ファイバー上平坦である},
{\it $X$ の $z$ 上のファイバーは、$x$ で、$Y$ の $z$ 上の ファイバー上平坦である},
{\it $X$ の $z$ 上のファイバーは $Y$ の $z$ 上のファイバー上平坦である}


\noindent
\ref{spaces-more-morphisms-definition-CM} を参照：
{\it 点 $x$ で Cohen--Macaulay},
{\it Cohen--Macaulay 射}


\noindent
\ref{spaces-more-morphisms-definition-gorenstein} を参照：
{\it $x$ において Gorenstein},
{\it Gorenstein 射}


\noindent
\ref{spaces-more-morphisms-definition-geometrically-reduced-fibre} を参照：
{\it $y$ における $f : X \to Y$ のファイバーは幾何学的に被約である}


\noindent
\ref{spaces-more-morphisms-definition-regular-immersion} を参照：
{\it Koszul 正則埋め込み},
{\it $H_1$-正則埋め込み},
{\it 準正則埋め込み}


\noindent
\ref{spaces-more-morphisms-definition-relative-pseudo-coherence} を参照：
{\it $Y$ に関して $m$-擬連接},
{\it $Y$ に関して擬連接},
{\it $Y$ に関して $m$-擬連接},
{\it $Y$ に関して擬連接}


\noindent
\ref{spaces-more-morphisms-definition-pseudo-coherent} を参照：
{\it 擬連接},
{\it $x$ で擬連接}


\noindent
\ref{spaces-more-morphisms-definition-perfect} を参照：
{\it 完全},
{\it $x$ で完全}


\noindent
\ref{spaces-more-morphisms-definition-lci} を参照：
{\it Koszul 射},
{\it 局所完全交叉射},
{\it $x$ で Koszul}


\noindent
\ref{spaces-more-morphisms-definition-relatively-perfect} を参照：
{\it $Y$ に関して完全},
{\it $Y$-完全}


\noindent
\ref{spaces-more-morphisms-definition-nodal-family} を参照：
{\it 相対次元 $1$ の高々節点的}


\medskip\noindent
{\bf 代数空間上の平坦性}

\medskip

\noindent
\ref{spaces-flat-definition-impurity} を参照：
{\it $\mathcal{F}$ の $y$ 上の不純性}


\noindent
\ref{spaces-flat-definition-pure} を参照：
{\it $y$ 上で純粋},
{\it $y$ 上で普遍的に純粋},
{\it $y$ 上で純粋},
{\it 普遍的に $Y$-純粋},
{\it $Y$ に関して普遍的に純粋},
{\it $Y$-純粋},
{\it $Y$ に関して純粋},
{\it $Y$-純粋},
{\it $Y$ に関して純粋}


\noindent
\ref{spaces-flat-definition-flattening} を参照：
{\it $\mathcal{F}$ の普遍平坦化が存在する},
{\it $X$ の普遍平坦化が存在する}


\noindent
\ref{spaces-flat-definition-flat-dimension-n} を参照：
{\it $\mathcal{F}$ は次元 $\geq n$ において $Y$ 上平坦}


\medskip\noindent
{\bf 代数空間における亜群}

\medskip

\noindent
\ref{spaces-groupoids-definition-equivalence-relation} を参照：
{\it 前関係},
{\it 関係},
{\it 前同値関係},
{\it $B$ 上の $U$ の同値関係}


\noindent
\ref{spaces-groupoids-definition-restrict-relation} を参照：
{\it 制限},
{\it 引き戻し}


\noindent
\ref{spaces-groupoids-definition-group-space} を参照：
{\it $B$ 上の群代数空間},
{\it $B$ 上の群代数空間の射 $\psi : (G, m) \to (G', m')$}


\noindent
\ref{spaces-groupoids-definition-action-group-space} を参照：
{\it 代数空間 $X/B$ への $G$ の作用},
{\it 同変},
{\it $G$-同変}


\noindent
\ref{spaces-groupoids-definition-free-action} を参照：
{\it 自由}


\noindent
\ref{spaces-groupoids-definition-pseudo-torsor} を参照：
{\it 擬 $G$-トーサー},
{\it $G$ の下で形式的に主等質},
{\it 自明}


\noindent
\ref{spaces-groupoids-definition-principal-homogeneous-space} を参照：
{\it 主等質空間},
{\it 主等質 $G$-空間で $B$ 上にあるもの},
{\it $\tau$ 位相における $G$-トーサー},
{\it $\tau$ $G$-トーサー},
{\it $\tau$ トーサー},
{\it 準等自明},
{\it 局所自明}


\noindent
\ref{spaces-groupoids-definition-equivariant-module} を参照：
{\it $G$-同変準連接 $\mathcal{O}_X$-加群},
{\it 同変準連接 $\mathcal{O}_X$-加群}


\noindent
\ref{spaces-groupoids-definition-groupoid} を参照：
{\it $B$ 上の代数空間における亜群},
{\it $B$ 上の代数空間における亜群の射 $f : (U, R, s, t, c) \to (U', R', s', t', c')$}


\noindent
\ref{spaces-groupoids-definition-groupoid-module} を参照：
{\it $(U, R, s, t, c)$ 上の準連接加群}


\noindent
\ref{spaces-groupoids-definition-stabilizer-groupoid} を参照：
{\it 代数空間における亜群 $(U, R, s, t, c)$ の固定化群}


\noindent
\ref{spaces-groupoids-definition-restrict-groupoid} を参照：
{\it $(U, R, s, t, c)$ の $U'$ への制限}


\noindent
\ref{spaces-groupoids-definition-invariant-open} を参照：
{\it $R$-不変},
{\it $R$-不変},
{\it $R$-不変}


\noindent
\ref{spaces-groupoids-definition-quotient-sheaf} を参照：
{\it 商層 $U/R$}


\noindent
\ref{spaces-groupoids-definition-representable-quotient} を参照：
{\it 代数空間で表現可能な商},
{\it 表現可能な商},
{\it 表現可能な商},
{\it 代数空間で表現可能な商}


\noindent
\ref{spaces-groupoids-definition-quotient-stack} を参照：
{\it 商スタック},
{\it 商スタック}


\medskip\noindent
{\bf 代数空間における亜群の補足}

\medskip

\noindent
\ref{spaces-more-groupoids-definition-split-at-point} を参照：
{\it $u$ 上強分裂},
{\it $u$ 上の $R$ の強分裂},
{\it $u$ 上分裂},
{\it $u$ 上の $R$ の分裂},
{\it $u$ 上準分裂},
{\it $u$ 上の $R$ の準分裂}


\medskip\noindent
{\bf ブートストラップ}

\medskip

\noindent
\ref{bootstrap-definition-morphism-representable-by-spaces} を参照：
{\it 代数空間によって表現可能}


\noindent
\ref{bootstrap-definition-property-transformation} を参照：
{\it 性質 $\mathcal{P}$}


\medskip\noindent
{\bf 代数空間の押し出し}

\medskip

\medskip\noindent
{\bf 代数空間の Chow 群}

\medskip

\noindent
\ref{spaces-chow-definition-delta-dimension} を参照：
{\it $T$ の $\delta$-次元}


\noindent
\ref{spaces-chow-definition-cycles} を参照：
{\it $X$ 上のサイクル},
{\it $k$-サイクル}


\noindent
\ref{spaces-chow-definition-length-at-x} を参照：
{\it $x$ における $\mathcal{F}$ の長さは $d$ である}


\noindent
\ref{spaces-chow-definition-cycle-associated-to-closed-subscheme} を参照：
{\it $Y$ における $Z$ の重複度},
{\it $Y$ に付随する $k$-サイクル}


\noindent
\ref{spaces-chow-definition-cycle-associated-to-coherent-sheaf} を参照：
{\it $\mathcal{F}$ における $Z$ の重複度},
{\it $\mathcal{F}$ に付随する $k$-サイクル}


\noindent
\ref{spaces-chow-definition-proper-pushforward} を参照：
{\it 順像}


\noindent
\ref{spaces-chow-definition-flat-pullback} を参照：
{\it $f$ による $\alpha$ の平坦引き戻し}


\noindent
\ref{spaces-chow-definition-principal-divisor} を参照：
{\it $f$ に付随する主因子}


\noindent
\ref{spaces-chow-definition-rational-equivalence} を参照：
{\it 零と有理同値},
{\it 有理同値},
{\it $X$ 上の $k$-サイクルの Chow 群},
{\it $X$ 上の有理同値を法とする $k$-サイクルの Chow 群}


\noindent
\ref{spaces-chow-definition-divisor-invertible-sheaf} を参照：
{\it $s$ に付随する Weil 因子},
{\it $\mathcal{L}$ に付随する Weil 因子}


\noindent
\ref{spaces-chow-definition-cap-c1} を参照：
{\it $\mathcal{L}$ の第一 Chern 類との交叉積}


\noindent
\ref{spaces-chow-definition-gysin-homomorphism} を参照：
{\it Gysin 準同型}


\noindent
\ref{spaces-chow-definition-bivariant-class} を参照：
{\it $f$ に対する次数 $p$ の双変類 $c$}


\noindent
\ref{spaces-chow-definition-chow-cohomology} を参照：
{\it Chow コホモロジー}


\noindent
\ref{spaces-chow-definition-chern-classes} を参照：
{\it $\mathcal{E}$ の第 $i$ Chern 類},
{\it $\mathcal{E}$ の全 Chern 類}


\noindent
\ref{spaces-chow-definition-degree-zero-cycle} を参照：
{\it 零サイクルの次数}


\medskip\noindent
{\bf 亜群の商}

\medskip

\noindent
\ref{groupoids-quotients-definition-invariant} を参照：
{\it $R$-不変},
{\it $G$-不変}


\noindent
\ref{groupoids-quotients-definition-base-change} を参照：
{\it 基底変換},
{\it 平坦基底変換}


\noindent
\ref{groupoids-quotients-definition-categorical} を参照：
{\it 圏論的商},
{\it $\mathcal{C}$ における圏論的商},
{\it スキームの圏における圏論的商},
{\it スキームにおける圏論的商}


\noindent
\ref{groupoids-quotients-definition-universal-categorical} を参照：
{\it 普遍圏論的商},
{\it 一様圏論的商}


\noindent
\ref{groupoids-quotients-definition-orbit} を参照：
{\it 軌道},
{\it $R$-軌道}


\noindent
\ref{groupoids-quotients-definition-geometric-orbits} を参照：
{\it 弱 $R$-同値},
{\it $R$-同値},
{\it 弱軌道},
{\it 弱 $R$-軌道},
{\it 軌道},
{\it $R$-軌道}


\noindent
\ref{groupoids-quotients-definition-set-theoretically-invariant} を参照：
{\it 集合論的に $R$-不変},
{\it 軌道を分離する},
{\it $R$-軌道を分離する}


\noindent
\ref{groupoids-quotients-definition-set-theoretic-equivalence} を参照：
{\it 集合論的前同値関係},
{\it 集合論的同値関係}


\noindent
\ref{groupoids-quotients-definition-orbit-space} を参照：
{\it $R$ の軌道空間}


\noindent
\ref{groupoids-quotients-definition-coarse} を参照：
{\it 粗商},
{\it スキームにおける粗商}


\noindent
\ref{groupoids-quotients-definition-topological} を参照：
{\it 一様に},
{\it 普遍的に}


\noindent
\ref{groupoids-quotients-definition-functions} を参照：
{\it $X$ 上の $R$-不変関数の層},
{\it $X$ 上の関数は $U$ 上の $R$-不変関数である}


\noindent
\ref{groupoids-quotients-definition-good} を参照：
{\it 良い商}


\noindent
\ref{groupoids-quotients-definition-geometric} を参照：
{\it 幾何学的商}


\medskip\noindent
{\bf 代数空間のコホモロジーの補足}

\medskip

\medskip\noindent
{\bf 単体的代数空間}

\medskip

\noindent
\ref{spaces-simplicial-definition-cartesian-sheaf} を参照：
{\it カルテジアン},
{\it カルテジアン},
{\it カルテジアン},
{\it カルテジアン}


\noindent
\ref{spaces-simplicial-definition-cartesian-derived} を参照：
{\it 導来圏の単体的系},
{\it カルテジアン},
{\it 導来圏の単体的系の射}


\noindent
\ref{spaces-simplicial-definition-cartesian-derived-modules} を参照：
{\it 加群の導来圏の単体的系},
{\it カルテジアン},
{\it 加群の導来圏の単体的系の射}


\noindent
\ref{spaces-simplicial-definition-cartesian-morphism} を参照：
{\it カルテジアン},
{\it $Y$ は $X$ 上カルテジアンである}


\noindent
\ref{spaces-simplicial-definition-fibre-products-simplicial-scheme} を参照：
{\it $f$ に付随する単体的スキーム}


\medskip\noindent
{\bf 代数空間の双対性}

\medskip

\noindent
\ref{spaces-duality-definition-dualizing-scheme} を参照：
{\it 双対化複体}


\noindent
\ref{spaces-duality-definition-relative-dualizing-proper-flat} を参照：
{\it 相対双対化複体}


\medskip\noindent
{\bf 形式代数空間}

\medskip

\noindent
\ref{formal-spaces-definition-toplogy-tensor-product} を参照：
{\it テンソル積},
{\it 完備化テンソル積}


\noindent
\ref{formal-spaces-definition-weakly-admissible} を参照：
{\it 位相的冪零},
{\it 弱定義イデアル},
{\it 弱前許容},
{\it 弱許容}


\noindent
\ref{formal-spaces-definition-taut} を参照：
{\it タウト}


\noindent
\ref{formal-spaces-definition-adic-homomorphism} を参照：
{\it アディック}


\noindent
\ref{formal-spaces-definition-weakly-adic} を参照：
{\it 弱前アディック},
{\it $c$-進},
{\it 弱アディック}


\noindent
\ref{formal-spaces-definition-affine-formal-algebraic-space} を参照：
{\it アフィン形式代数空間},
{\it アフィン形式代数空間の射}


\noindent
\ref{formal-spaces-definition-types-affine-formal-algebraic-space} を参照：
{\it McQuillan},
{\it 古典的},
{\it 弱アディック},
{\it アディック},
{\it アディック*},
{\it Noether}


\noindent
\ref{formal-spaces-definition-affine-formal-spectrum} を参照：
{\it 形式スペクトル}


\noindent
\ref{formal-spaces-definition-countable} を参照：
{\it 可算添字付き}


\noindent
\ref{formal-spaces-definition-formal-algebraic-space} を参照：
{\it 形式代数空間},
{\it 形式代数空間の射}


\noindent
\ref{formal-spaces-definition-completion} を参照：
{\it $X$ の $T$ に沿う完備化}


\noindent
\ref{formal-spaces-definition-separated} を参照：
{\it 準分離},
{\it 分離的}


\noindent
\ref{formal-spaces-definition-quasi-compact} を参照：
{\it 準コンパクト}


\noindent
\ref{formal-spaces-definition-quasi-compact-morphism} を参照：
{\it 準コンパクト}


\noindent
\ref{formal-spaces-definition-types-formal-algebraic-spaces} を参照：
{\it 局所可算添字付き},
{\it 局所可算添字付きかつ古典的},
{\it 局所弱アディック},
{\it 局所アディック*},
{\it 局所 Noether}


\noindent
\ref{formal-spaces-definition-adic-morphism} を参照：
{\it アディック射}


\noindent
\ref{formal-spaces-definition-finite-type} を参照：
{\it 局所有限型},
{\it 有限型}


\noindent
\ref{formal-spaces-definition-surjective} を参照：
{\it 全射}


\noindent
\ref{formal-spaces-definition-monomorphism} を参照：
{\it モノ射}


\noindent
\ref{formal-spaces-definition-closed-immersion} を参照：
{\it 閉埋め込み}


\noindent
\ref{formal-spaces-definition-topologically-finite-type} を参照：
{\it 位相的有限型}


\noindent
\ref{formal-spaces-definition-separated-morphism} を参照：
{\it 分離的},
{\it 準分離}


\noindent
\ref{formal-spaces-definition-proper} を参照：
{\it 固有}


\noindent
\ref{formal-spaces-definition-etale-sites} を参照：
{\it 小エタール・サイト}


\noindent
\ref{formal-spaces-definition-completion-formal-algebraic-space} を参照：
{\it $X$ の $T$ に沿う完備化}


\noindent
\ref{formal-spaces-definition-completion-subspace} を参照：
{\it $Z$ に沿う $X$ の完備化}


\medskip\noindent
{\bf 形式空間の代数化}

\medskip

\noindent
\ref{restricted-definition-rig-smooth-homomorphism} を参照：
{\it $(A, I)$ 上 rig-滑らか}


\noindent
\ref{restricted-definition-rig-etale-homomorphism} を参照：
{\it $(A, I)$ 上 rig-エタール}


\noindent
\ref{restricted-definition-flat} を参照：
{\it 平坦}


\noindent
\ref{restricted-definition-rig-closed} を参照：
{\it rig-閉}


\noindent
\ref{restricted-definition-completed-principal-localization} を参照：
{\it 完備化された主局所化}


\noindent
\ref{restricted-definition-naively-rig-flat} を参照：
{\it 素朴な rig-平坦}


\noindent
\ref{restricted-definition-rig-flat-continuous-homomorphism} を参照：
{\it rig-平坦}


\noindent
\ref{restricted-definition-rig-flat} を参照：
{\it rig-平坦}


\noindent
\ref{restricted-definition-rig-smooth-continuous-homomorphism} を参照：
{\it rig-滑らか}


\noindent
\ref{restricted-definition-rig-smooth} を参照：
{\it rig-滑らか}


\noindent
\ref{restricted-definition-rig-etale-continuous-homomorphism} を参照：
{\it rig-エタール}


\noindent
\ref{restricted-definition-rig-etale} を参照：
{\it rig-エタール}


\noindent
\ref{restricted-definition-rig-surjective} を参照：
{\it rig-全射}


\noindent
\ref{restricted-definition-formal-modification} を参照：
{\it 形式的修正}


\medskip\noindent
{\bf 曲面の特異点解消再論}

\medskip

\noindent
\ref{spaces-resolve-definition-blowup-at-point} を参照：
{\it $x$ における $X$ のブローアップ $X' \to X$}


\noindent
\ref{spaces-resolve-definition-normalized-blowup} を参照：
{\it $x$ における $X$ の正規化ブローアップ}


\noindent
\ref{spaces-resolve-definition-resolution} を参照：
{\it 特異点解消}


\noindent
\ref{spaces-resolve-definition-resolution-surface} を参照：
{\it 正規化ブローアップによる特異点解消}


\medskip\noindent
{\bf 形式変形理論}

\medskip

\noindent
\ref{formal-defos-definition-CLambda} を参照：
{\it $\mathcal{C}_\Lambda$},
{\it 古典的場合}


\noindent
\ref{formal-defos-definition-small-extension} を参照：
{\it 小拡大}


\noindent
\ref{formal-defos-definition-tangent-space-ring} を参照：
{\it 相対余接空間}


\noindent
\ref{formal-defos-definition-essential-surjection} を参照：
{\it 本質的全射}


\noindent
\ref{formal-defos-definition-completion-CLambda} を参照：
{\it $\widehat{\mathcal{C}}_\Lambda$}


\noindent
\ref{formal-defos-definition-category-cofibred-groupoids} を参照：
{\it $\mathcal{C}$ 上の亜群余ファイバー圏}


\noindent
\ref{formal-defos-definition-prorepresentable} を参照：
{\it プロ表現可能}


\noindent
\ref{formal-defos-definition-predeformation-category} を参照：
{\it 前変形圏},
{\it 前変形圏の射}


\noindent
\ref{formal-defos-definition-formal-objects} を参照：
{\it $\mathcal{F}$ の形式対象の圏 $\widehat{\mathcal{F}}$},
{\it $\mathcal{F}$ の形式対象 $\xi = (R, \xi_n, f_n)$},
{\it 形式対象の射 $a : \xi \to \eta$}


\noindent
\ref{formal-defos-definition-completion} を参照：
{\it $\mathcal{F}$ の完備化}


\noindent
\ref{formal-defos-definition-smooth-morphism} を参照：
{\it 滑らか}


\noindent
\ref{formal-defos-definition-versal} を参照：
{\it 半普遍的}


\noindent
\ref{formal-defos-definition-cofibered-groupoid-projection-smooth} を参照：
{\it 滑らか},
{\it 非障害的}


\noindent
\ref{formal-defos-definition-S1-S2} を参照：
{\it 条件 (S1) および (S2)}


\noindent
\ref{formal-defos-definition-linear} を参照：
{\it $R$-線形}


\noindent
\ref{formal-defos-definition-tangent-space-over-R} を参照：
{\it $F$ の接空間 $TF$}


\noindent
\ref{formal-defos-definition-tangent-space} を参照：
{\it $\mathcal{F}$ の接空間 $T \mathcal{F}$}


\noindent
\ref{formal-defos-definition-differential} を参照：
{\it $\varphi$ の微分 $d \varphi : T \mathcal{F} \to T \mathcal{G}$}


\noindent
\ref{formal-defos-definition-minimal-versal} を参照：
{\it 極小},
{\it 最小半普遍的}


\noindent
\ref{formal-defos-definition-RS} を参照：
{\it 条件 (RS)}


\noindent
\ref{formal-defos-definition-deformation-category} を参照：
{\it 変形圏}


\noindent
\ref{formal-defos-definition-lifts} を参照：
{\it $f$ に沿う $x$ の持ち上げ},
{\it 持ち上げの射}


\noindent
\ref{formal-defos-definition-relative-infinitesimal-auts} を参照：
{\it $x$ 上の $x'$ の無限小自己同型群}


\noindent
\ref{formal-defos-definition-infinitesimal-auts} を参照：
{\it $x_0$ の無限小自己同型群}


\noindent
\ref{formal-defos-definition-automorphism-functor} を参照：
{\it $x$ の自己同型関手}


\noindent
\ref{formal-defos-definition-groupoid-in-functors} を参照：
{\it $\mathcal{C}$ 上の関手における亜群の圏},
{\it $\mathcal{C}$ 上の関手における亜群},
{\it $\mathcal{C}$ 上の関手における亜群の射 $(U, R, s, t, c) \to (U', R', s', t', c')$}


\noindent
\ref{formal-defos-definition-representable} を参照：
{\it 表現可能}


\noindent
\ref{formal-defos-definition-restricting-groupoids-in-functors} を参照：
{\it $(U, R, s, t, c)$ の $\mathcal{C}'$ への制限 $(U, R, s, t, c)|_{\mathcal{C}'}$}


\noindent
\ref{formal-defos-definition-quotient} を参照：
{\it 商となる，群亜群をファイバーとする余ファイバー圏 $[U/R] \to \mathcal{C}$},
{\it 商射 $U \to [U/R]$}


\noindent
\ref{formal-defos-definition-prorepresentable-groupoid-in-functors} を参照：
{\it プロ表現可能}


\noindent
\ref{formal-defos-definition-completion-groupoid-in-functors} を参照：
{\it $(U, R, s, t, c)$ の完備化 $(U, R, s, t, c)^{\wedge}$}


\noindent
\ref{formal-defos-definition-smooth-groupoid-in-functors} を参照：
{\it 滑らか}


\noindent
\ref{formal-defos-definition-presentation} を参照：
{\it $(U, R, s, t, c)$ による $\mathcal{F}$ の表示}


\noindent
\ref{formal-defos-definition-minimal-groupoid-in-functors} を参照：
{\it 正規化されている},
{\it 極小}


\medskip\noindent
{\bf 変形理論}

\medskip

\noindent
\ref{defos-definition-strict-morphism-thickenings} を参照：
{\it 厳密な厚化の射}


\noindent
\ref{defos-definition-strict-morphism-thickenings-ringed-topoi} を参照：
{\it 厳密な厚化の射}


\medskip\noindent
{\bf 余接複体}

\medskip

\noindent
\ref{cotangent-definition-standard-resolution} を参照：
{\it $A$ 上の $B$ の標準分解}


\noindent
\ref{cotangent-definition-cotangent-complex-ring-map} を参照：
{\it 余接複体}


\noindent
\ref{cotangent-definition-biderivation} を参照：
{\it $A$-双導分}


\noindent
\ref{cotangent-definition-atiyah-class} を参照：
{\it Atiyah 類}


\noindent
\ref{cotangent-definition-standard-resolution-sheaves-rings} を参照：
{\it $\mathcal{A}$ 上の $\mathcal{B}$ の標準分解}


\noindent
\ref{cotangent-definition-cotangent-complex-morphism-sheaves-rings} を参照：
{\it 余接複体}


\noindent
\ref{cotangent-definition-atiyah-class-general} を参照：
{\it Atiyah 類}


\noindent
\ref{cotangent-definition-cotangent-complex-morphism-ringed-spaces} を参照：
{\it 余接複体}


\noindent
\ref{cotangent-definition-cotangent-complex-morphism-ringed-topoi} を参照：
{\it 余接複体}


\noindent
\ref{cotangent-definition-cotangent-morphism-schemes} を参照：
{\it $Y$ 上の $X$ の余接複体 $L_{X/Y}$}


\noindent
\ref{cotangent-definition-cotangent-morphism-spaces} を参照：
{\it $Y$ 上の $X$ の余接複体 $L_{X/Y}$}


\medskip\noindent
{\bf 変形問題の例}

\medskip

\medskip\noindent
{\bf 代数スタック}

\medskip

\noindent
\ref{algebraic-definition-representable-by-algebraic-space} を参照：
{\it $S$ 上の代数空間によって表現可能}


\noindent
\ref{algebraic-definition-representable-by-algebraic-spaces} を参照：
{\it 代数空間によって表現可能}


\noindent
\ref{algebraic-definition-relative-representable-property} を参照：
{\it 性質 $\mathcal{P}$}


\noindent
\ref{algebraic-definition-algebraic-stack} を参照：
{\it $S$ 上の代数スタック}


\noindent
\ref{algebraic-definition-deligne-mumford} を参照：
{\it Deligne--Mumford スタック}


\noindent
\ref{algebraic-definition-morphism-algebraic-stacks} を参照：
{\it $S$ 上の代数スタックの $2$-圏}


\noindent
\ref{algebraic-definition-smooth-groupoid} を参照：
{\it 滑らかなグルーポイド}


\noindent
\ref{algebraic-definition-presentation} を参照：
{\it 表示}


\noindent
\ref{algebraic-definition-viewed-as} を参照：
{\it $S'$ 上の代数スタックとみなしたもの}


\noindent
\ref{algebraic-definition-change-of-base} を参照：
{\it $\mathcal{X}'$ の基底変換}


\medskip\noindent
{\bf スタックの例}

\medskip

\noindent
\ref{examples-stacks-definition-hilbert-d-stack} を参照：
{\it $\mathcal{Y}$ 上の $\mathcal{X}$ の次数 $d$ の有限 Hilbert スタック}


\medskip\noindent
{\bf 代数スタック上の層}

\medskip

\noindent
\ref{stacks-sheaves-definition-presheaves} を参照：
{\it $\mathcal{X}$ 上の前層},
{\it $\mathcal{X}$ 上の前層の射}


\noindent
\ref{stacks-sheaves-definition-inherited-topologies} を参照：
{\it 随伴 Zariski サイト},
{\it 随伴エタールサイト},
{\it 随伴滑らかなサイト},
{\it 付随シントミック・サイト},
{\it 随伴 fppf サイト}


\noindent
\ref{stacks-sheaves-definition-sheaves} を参照：
{\it Zariski 層},
{\it Zariski 位相に関する層},
{\it エタール層},
{\it エタール位相に関する層},
{\it 滑らかな層},
{\it 滑らかな位相に関する層},
{\it シントミック層},
{\it シントミック位相に関する層},
{\it fppf 層},
{\it 層},
{\it fppf 位相に関する層}


\noindent
\ref{stacks-sheaves-definition-morphism} を参照：
{\it 随伴する fppf トポスの射}


\noindent
\ref{stacks-sheaves-definition-structure-sheaf} を参照：
{\it $\mathcal{X}$ の構造層}


\noindent
\ref{stacks-sheaves-definition-modules} を参照：
{\it $\mathcal{X}$ 上の加群の前層},
{\it $\mathcal{O}_\mathcal{X}$-加群},
{\it $\mathcal{O}_\mathcal{X}$-加群の層}


\noindent
\ref{stacks-sheaves-definition-pullback} を参照：
{\it $\mathcal{F}$ の引き戻し $x^{-1}\mathcal{F}$},
{\it $\mathcal{F}$ の $U_\etale$ への制限}


\noindent
\ref{stacks-sheaves-definition-quasi-coherent} を参照：
{\it $\mathcal{X}$ 上の準連接加群},
{\it 準連接 $\mathcal{O}_\mathcal{X}$-加群}


\noindent
\ref{stacks-sheaves-definition-locally-quasi-coherent} を参照：
{\it 局所準連接}


\noindent
\ref{stacks-sheaves-definition-alternative} を参照：
{\it 付随するアフィンサイト}


\noindent
\ref{stacks-sheaves-definition-alternative-inherited-topologies} を参照：
{\it 付随するアフィン Zariski サイト},
{\it 付随アフィンエタール・サイト},
{\it 付随アフィン滑らかサイト},
{\it 付随アフィンシントミック・サイト},
{\it 付随するアフィン fppf サイト}


\noindent
\ref{stacks-sheaves-definition-QC} を参照：
{\it 導来圏における準連接対象の三角圏}


\medskip\noindent
{\bf 表現可能性の判定条件}

\medskip

\noindent
\ref{criteria-definition-algebraic} を参照：
{\it 代数的}


\medskip\noindent
{\bf Artin の公理}

\medskip

\noindent
\ref{artin-definition-RS} を参照：
{\it 条件 (RS)}


\noindent
\ref{artin-definition-formal-objects} を参照：
{\it 形式対象},
{\it 形式対象の射},
{\it 上にある}


\noindent
\ref{artin-definition-effective} を参照：
{\it 有効}


\noindent
\ref{artin-definition-limit-preserving} を参照：
{\it 極限を保つ}


\noindent
\ref{artin-definition-versal-formal-object} を参照：
{\it 半普遍的}


\noindent
\ref{artin-definition-versal} を参照：
{\it 半普遍的}


\noindent
\ref{artin-definition-openness-versality} を参照：
{\it 半普遍性の開性},
{\it 半普遍性の開性}


\noindent
\ref{artin-definition-RS-star} を参照：
{\it 条件 (RS*)}


\noindent
\ref{artin-definition-obstruction-theory} を参照：
{\it 障害理論},
{\it 障害加群},
{\it 障害}


\noindent
\ref{artin-definition-naive-obstruction-theory} を参照：
{\it 素朴な障害理論}


\medskip\noindent
{\bf Quot 空間と Hilbert 空間}

\medskip

\medskip\noindent
{\bf 代数スタックの性質}

\medskip

\noindent
\ref{stacks-properties-definition-points} を参照：
{\it 点}


\noindent
\ref{stacks-properties-definition-topological-space} を参照：
{\it 位相空間}


\noindent
\ref{stacks-properties-definition-surjective} を参照：
{\it 全射}


\noindent
\ref{stacks-properties-definition-quasi-compact} を参照：
{\it 準コンパクト}


\noindent
\ref{stacks-properties-definition-type-property} を参照：
{\it 性質 $\mathcal{P}$ をもつ}


\noindent
\ref{stacks-properties-definition-property-at-point} を参照：
{\it $x$ において性質 $\mathcal{P}$ をもつ}


\noindent
\ref{stacks-properties-definition-monomorphism} を参照：
{\it モノ射}


\noindent
\ref{stacks-properties-definition-immersion} を参照：
{\it 開埋め込み},
{\it 閉埋め込み},
{\it 埋め込み}


\noindent
\ref{stacks-properties-definition-substacks} を参照：
{\it 開部分スタック},
{\it 閉部分スタック},
{\it 局所閉部分スタック}


\noindent
\ref{stacks-properties-definition-reduced-induced-stack} を参照：
{\it $Z$ 上の代数スタック構造},
{\it 誘導被約代数スタック構造},
{\it $\mathcal{X}$ の被約化 $\mathcal{X}_{red}$}


\noindent
\ref{stacks-properties-definition-residual-gerbe} を参照：
{\it $x$ における $\mathcal{X}$ の剰余ガーベが存在する},
{\it $x$ における $\mathcal{X}$ の剰余ガーベ}


\noindent
\ref{stacks-properties-definition-dimension-at-point} を参照：
{\it $x$ における $\mathcal{X}$ の次元}


\noindent
\ref{stacks-properties-definition-dimension} を参照：
{\it 次元}


\noindent
\ref{stacks-properties-definition-number-of-geometric-branches} を参照：
{\it $x$ における $\mathcal{X}$ の幾何的枝数},
{\it $x$ において幾何的単枝}


\medskip\noindent
{\bf 代数スタックの射}

\medskip

\noindent
\ref{stacks-morphisms-definition-separated} を参照：
{\it DM},
{\it 準 DM},
{\it 分離的},
{\it 準分離}


\noindent
\ref{stacks-morphisms-definition-absolute-separated} を参照：
{\it $S$ 上 DM},
{\it $S$ 上準 DM},
{\it $S$ 上分離},
{\it $S$ 上準分離},
{\it DM},
{\it 準 DM},
{\it 分離的},
{\it 準分離}


\noindent
\ref{stacks-morphisms-definition-isom} を参照：
{\it $x$ の相対自己同型の層},
{\it $x_1$ から $x_2$ への相対同型の層},
{\it $x$ の自己同型の層},
{\it $x_1$ から $x_2$ への同型の層}


\noindent
\ref{stacks-morphisms-definition-quasi-compact} を参照：
{\it 準コンパクト}


\noindent
\ref{stacks-morphisms-definition-noetherian} を参照：
{\it Noether}


\noindent
\ref{stacks-morphisms-definition-affine} を参照：
{\it アフィン}


\noindent
\ref{stacks-morphisms-definition-integral} を参照：
{\it 整},
{\it 有限}


\noindent
\ref{stacks-morphisms-definition-open} を参照：
{\it 開},
{\it 普遍開}


\noindent
\ref{stacks-morphisms-definition-submersive} を参照：
{\it 商位相的},
{\it 普遍的に商写像的}


\noindent
\ref{stacks-morphisms-definition-closed} を参照：
{\it 閉},
{\it 普遍閉}


\noindent
\ref{stacks-morphisms-definition-universally-injective} を参照：
{\it 普遍単射}


\noindent
\ref{stacks-morphisms-definition-universal-homeomorphism} を参照：
{\it 普遍同相射}


\noindent
\ref{stacks-morphisms-definition-P} を参照：
{\it 性質 $\mathcal{P}$ をもつ}


\noindent
\ref{stacks-morphisms-definition-locally-finite-type} を参照：
{\it 局所有限型},
{\it 有限型}


\noindent
\ref{stacks-morphisms-definition-finite-type-point} を参照：
{\it 有限型点}


\noindent
\ref{stacks-morphisms-definition-locally-quasi-finite} を参照：
{\it 局所準有限}


\noindent
\ref{stacks-morphisms-definition-quasi-finite} を参照：
{\it 準有限}


\noindent
\ref{stacks-morphisms-definition-flat} を参照：
{\it 平坦}


\noindent
\ref{stacks-morphisms-definition-flat-at-point} を参照：
{\it $x$ で平坦}


\noindent
\ref{stacks-morphisms-definition-locally-finite-presentation} を参照：
{\it 局所有限表示},
{\it 有限表示}


\noindent
\ref{stacks-morphisms-definition-gerbe} を参照：
{\it 上のガーベ},
{\it ガーベ}


\noindent
\ref{stacks-morphisms-definition-smooth} を参照：
{\it 滑らか}


\noindent
\ref{stacks-morphisms-definition-etale-smooth-P} を参照：
{\it 性質 $\mathcal{P}$ をもつ}


\noindent
\ref{stacks-morphisms-definition-etale} を参照：
{\it エタール}


\noindent
\ref{stacks-morphisms-definition-unramified} を参照：
{\it 不分岐}


\noindent
\ref{stacks-morphisms-definition-proper} を参照：
{\it 固有}


\noindent
\ref{stacks-morphisms-definition-scheme-theoretic-image} を参照：
{\it スキーム論的像}


\noindent
\ref{stacks-morphisms-definition-fill-in-diagram} を参照：
{\it 点線矢印},
{\it 点線矢印の射}


\noindent
\ref{stacks-morphisms-definition-uniqueness} を参照：
{\it 付値判定法の一意性部分}


\noindent
\ref{stacks-morphisms-definition-existence} を参照：
{\it 付値判定法の存在性部分}


\noindent
\ref{stacks-morphisms-definition-lci} を参照：
{\it 局所完全交叉射},
{\it Koszul}


\noindent
\ref{stacks-morphisms-definition-normalization} を参照：
{\it 正規化}


\noindent
\ref{stacks-morphisms-definition-decent} を参照：
{\it 良好}


\noindent
\ref{stacks-morphisms-definition-integral-algebraic-stack} を参照：
{\it 整}


\medskip\noindent
{\bf 代数スタックの極限}

\medskip

\noindent
\ref{stacks-limits-definition-limit-preserving} を参照：
{\it 極限を保つ}


\medskip\noindent
{\bf 代数スタックのコホモロジー}

\medskip

\noindent
\ref{stacks-cohomology-definition-flat-base-change} を参照：
{\it 平坦基底変換性}


\noindent
\ref{stacks-cohomology-definition-parasitic} を参照：
{\it 寄生的}


\noindent
\ref{stacks-cohomology-definition-lisse-etale} を参照：
{\it 滑らかエタール・サイト},
{\it 平坦 fppf サイト}


\noindent
\ref{stacks-cohomology-definition-coherent} を参照：
{\it 連接}


\medskip\noindent
{\bf スタックの導来圏}

\medskip

\noindent
\ref{stacks-perfect-definition-derived} を参照：
{\it 準連接なコホモロジー層をもつ $\mathcal{O}_\mathcal{X}$-加群の導来圏}


\medskip\noindent
{\bf 代数スタック入門}

\medskip

\noindent
\ref{stacks-introduction-definition-smooth} を参照：
{\it 滑らか}


\noindent
\ref{stacks-introduction-definition-algebraic-stack} を参照：
{\it 代数スタック}


\medskip\noindent
{\bf スタックの射の補足}

\medskip

\noindent
\ref{stacks-more-morphisms-definition-thickening} を参照：
{\it 肥厚化},
{\it 肥厚化の射},
{\it $\mathcal{Z}$ 上の厚化},
{\it $\mathcal{Z}$ 上の厚化の射}


\noindent
\ref{stacks-more-morphisms-definition-first-order-thickening} を参照：
{\it 一次肥厚化}


\noindent
\ref{stacks-more-morphisms-definition-formally-smooth} を参照：
{\it 形式的に滑らか}


\noindent
\ref{stacks-more-morphisms-definition-categorical-quotient} を参照：
{\it 圏論的モジュライ空間},
{\it 一様圏論的モジュライ空間},
{\it $\mathcal{C}$ における圏論的モジュライ空間},
{\it $\mathcal{C}$ における一様圏論的モジュライ空間}


\noindent
\ref{stacks-more-morphisms-definition-well-nigh-affine} を参照：
{\it ほとんどアフィン}


\medskip\noindent
{\bf スタックの幾何}

\medskip

\noindent
\ref{stacks-geometry-definition-versal-ring-at-x} を参照：
{\it $x_0$ における $\mathcal{X}$ の半普遍変形環}


\noindent
\ref{stacks-geometry-definition-multiplicity} を参照：
{\it 重複度}


\noindent
\ref{stacks-geometry-definition-formal-branches} を参照：
{\it $\mathcal{X}$ の $x_0$ を通る形式的分枝}


\noindent
\ref{stacks-geometry-definition-multiplicity-formal-branches} を参照：
{\it $\mathcal{X}$ の $x_0$ を通る形式的分枝の重複度}


\noindent
\ref{stacks-geometry-definition-relative-dimension} を参照：
{\it 相対次元}


\noindent
\ref{stacks-geometry-definition-relative-dimension-for-stacks} を参照：
{\it 相対次元}


\noindent
\ref{stacks-geometry-definition-pseudo-catenary} を参照：
{\it 擬鎖状}


\noindent
\ref{stacks-geometry-definition-dimension-local-ring} を参照：
{\it $x$ における $\mathcal{X}$ の局所環の次元}


\medskip\noindent
{\bf モジュライスタック}

\medskip

\medskip\noindent
{\bf 曲線のモジュライ}

\medskip

\noindent
\ref{moduli-curves-definition-deligne-mumford-smooth} を参照：
{\it 滑らかな固有曲線のモジュライ・スタック},
{\it 種数 $g$ の滑らかな固有曲線のモジュライ・スタック}


\noindent
\ref{moduli-curves-definition-relative-dualizing-sheaf} を参照：
{\it 相対双対化層}


\noindent
\ref{moduli-curves-definition-prestable} を参照：
{\it 前安定曲線族}


\noindent
\ref{moduli-curves-definition-semistable} を参照：
{\it 半安定曲線族}


\noindent
\ref{moduli-curves-definition-stable} を参照：
{\it 安定曲線族}


\noindent
\ref{moduli-curves-definition-deligne-mumford} を参照：
{\it 安定曲線のモジュライスタック},
{\it 種数 $g$ の安定曲線のモジュライスタック}


\medskip\noindent
{\bf 例}

\medskip

\medskip\noindent
{\bf 演習}

\medskip

\noindent
\ref{exercises-definition-directed-poset} を参照：
{\it 有向集合},
{\it 環の系}


\noindent
\ref{exercises-definition-colimit} を参照：
{\it 余極限}


\noindent
\ref{exercises-definition-finite-presentation} を参照：
{\it 有限表示}


\noindent
\ref{exercises-definition-quasi-compact} を参照：
{\it 準コンパクト}


\noindent
\ref{exercises-definition-Hausdorff} を参照：
{\it Hausdorff}


\noindent
\ref{exercises-definition-irreducible} を参照：
{\it 既約},
{\it 既約}


\noindent
\ref{exercises-definition-generic-point} を参照：
{\it 生成点}


\noindent
\ref{exercises-definition-Noetherian-space} を参照：
{\it Noether},
{\it アルティン的}


\noindent
\ref{exercises-definition-irreducible-component} を参照：
{\it 既約成分}


\noindent
\ref{exercises-definition-closed} を参照：
{\it 閉},
{\it 特殊化},
{\it 一般化}


\noindent
\ref{exercises-definition-connected-component} を参照：
{\it 連結},
{\it 連結成分}


\noindent
\ref{exercises-definition-length} を参照：
{\it 長さ}


\noindent
\ref{exercises-definition-catenary} を参照：
{\it 鎖状},
{\it 鎖状}


\noindent
\ref{exercises-definition-finite-locally-free} を参照：
{\it 有限局所自由},
{\it 可逆加群}


\noindent
\ref{exercises-definition-class-group} を参照：
{\it $A$ の類群},
{\it $A$ の Picard 群}


\noindent
\ref{exercises-definition-GU-GD} を参照：
{\it 上昇定理},
{\it 下降定理}


\noindent
\ref{exercises-definition-numerical-polynomial} を参照：
{\it 数値多項式}


\noindent
\ref{exercises-definition-graded-module} を参照：
{\it 次数付き加群},
{\it 局所有限},
{\it Euler--Poincar\'e 函数},
{\it ヒルベルト関数},
{\it ヒルベルト多項式}


\noindent
\ref{exercises-definition-graded-algebra} を参照：
{\it 次数付き $A$-代数},
{\it 次数付き $A$-代数 $B$ 上の次数付き加群 $M$},
{\it 次数付き加群・環の準同型},
{\it 次数付き部分加群},
{\it 次数付きイデアル},
{\it 次数付き加群の完全列}


\noindent
\ref{exercises-definition-homogeneous-ideal} を参照：
{\it 斉次}


\noindent
\ref{exercises-definition-Proj-R} を参照：
{\it 斉次スペクトル $\text{Proj}(R)$}


\noindent
\ref{exercises-definition-Dplus-Vplus} を参照：
{\it $R_{(f)}$}


\noindent
\ref{exercises-definition-CM} を参照：
{\it Cohen--Macaulay}


\noindent
\ref{exercises-definition-injective-filtered} を参照：
{\it フィルター付き入射的}


\noindent
\ref{exercises-definition-finite-filtration-category} を参照：
{\it $\text{Fil}^f(\mathcal{A})$}


\noindent
\ref{exercises-definition-filtered-quasi-isomorphism} を参照：
{\it フィルター付き擬同型}


\noindent
\ref{exercises-definition-filtered-acyclic} を参照：
{\it フィルター付き非輪状}


\noindent
\ref{exercises-definition-integral} を参照：
{\it 整}


\noindent
\ref{exercises-definition-dual-numbers} を参照：
{\it 双対数}


\noindent
\ref{exercises-definition-tangent-space} を参照：
{\it $S$ 上の $X$ の接空間},
{\it 接ベクトル}


\noindent
\ref{exercises-definition-quasi-coherent} を参照：
{\it 準連接}


\noindent
\ref{exercises-definition-specialization} を参照：
{\it 特殊化}


\noindent
\ref{exercises-definition-Noetherian-scheme} を参照：
{\it 局所 Noether},
{\it Noether}


\noindent
\ref{exercises-definition-coherent} を参照：
{\it 連接}


\noindent
\ref{exercises-definition-invertible-sheaf} を参照：
{\it 可逆 ${\mathcal O}_X$-加群}


\noindent
\ref{exercises-definition-invertible-module} を参照：
{\it 可逆加群 $M$},
{\it 自明}


\noindent
\ref{exercises-definition-picard-group} を参照：
{\it $X$ の Picard 群}


\noindent
\ref{exercises-definition-delta} を参照：
{\it $\delta(\tau)$}


\noindent
\ref{exercises-definition-divisor} を参照：
{\it Weil 因子},
{\it 素因子},
{\it 有理関数 $f \in K(X)^\ast$ に付随する Weil 因子},
{\it 有効 Cartier 因子},
{\it 有効 Cartier 因子 $D \subset X$ に付随する Weil 因子 $[D]$},
{\it 全商環の層 ${\mathcal K}_S$},
{\it Cartier 因子},
{\it 付随する Weil 因子}


\medskip\noindent
{\bf 文献案内}

\medskip

\medskip\noindent
{\bf 要望事項}

\medskip

\medskip\noindent
{\bf コーディング規約}

\medskip

\medskip\noindent
{\bf 廃止された内容}

\medskip

\noindent
\ref{obsolete-definition-epsilon} を参照：
{\it $\epsilon$-不変量}


\noindent
\ref{obsolete-definition-locally-finite-sum-effective-Cartier-divisors} を参照：
{\it 有効 Cartier 因子の和}


\medskip\noindent
{\bf GNU 自由文書利用許諾契約書}

\medskip
\end{multicols}
